\documentclass{article}
\usepackage{amsthm}

% if you need to pass options to natbib, use, e.g.:
\PassOptionsToPackage{numbers, compress}{natbib}

% before loading neurips_2026

% The authors should use one of these tracks.
% Before accepting by the NeurIPS conference, select one of the options below.
% 0. "default" for submission
\usepackage{neurips_2026}
% the "default" option is equal to the "main" option, which is used for the Main Track with double-blind reviewing.
% 1. "main" option is used for the Main Track
%  \usepackage[main]{neurips_2026}
% 2. "position" option is used for the Position Paper Track
%  \usepackage[position]{neurips_2026}
% 3. "eandd" option is used for the Evaluations & Datasets Track
 % \usepackage[eandd]{neurips_2026}
 % if you need to opt-in for a single-blind submission in the E&D track:
 %\usepackage[eandd, nonanonymous]{neurips_2026}
% 4. "creativeai" option is used for the Creative AI Track
%  \usepackage[creativeai]{neurips_2026}
% 5. "sglblindworkshop" option is used for the Workshop with single-blind reviewing
 % \usepackage[sglblindworkshop]{neurips_2026}
% 6. "dblblindworkshop" option is used for the Workshop with double-blind reviewing
%  \usepackage[dblblindworkshop]{neurips_2026}

% After being accepted, the authors should add "final" behind the track to compile a camera-ready version.
% 1. Main Track
 % \usepackage[main, final]{neurips_2026}
% 2. Position Paper Track
%  \usepackage[position, final]{neurips_2026}
% 3. Evaluations & Datasets Track
 % \usepackage[eandd, final]{neurips_2026}
% 4. Creative AI Track
%  \usepackage[creativeai, final]{neurips_2026}
% 5. Workshop with single-blind reviewing
%  \usepackage[sglblindworkshop, final]{neurips_2026}
% 6. Workshop with double-blind reviewing
%  \usepackage[dblblindworkshop, final]{neurips_2026}
% Note. For the workshop paper template, both \title{} and \workshoptitle{} are required, with the former indicating the paper title shown in the title and the latter indicating the workshop title displayed in the footnote.
% For workshops (5., 6.), the authors should add the name of the workshop, "\workshoptitle" command is used to set the workshop title.
% \workshoptitle{WORKSHOP TITLE}

% "preprint" option is used for arXiv or other preprint submissions
 % \usepackage[preprint]{neurips_2026}

% to avoid loading the natbib package, add option nonatbib:
%    \usepackage[nonatbib]{neurips_2026}

\usepackage[utf8]{inputenc} % allow utf-8 input
\usepackage[T1]{fontenc}    % use 8-bit T1 fonts
\usepackage{hyperref}       % hyperlinks
\usepackage{url}            % simple URL typesetting
\usepackage{booktabs}       % professional-quality tables
\usepackage{amsfonts}       % blackboard math symbols
\usepackage{nicefrac}       % compact symbols for 1/2, etc.
\usepackage{microtype}      % microtypography
\usepackage{xcolor}         % colors
 \usepackage{amsmath}

 \usepackage{multirow}

\newcommand{\minisection}[1]{\vspace{0.0in} \noindent {\bf #1}}

% Recommended, but optional, packages for figures and better typesetting:
\usepackage{microtype}
\usepackage{graphicx}
\usepackage{subcaption}
\usepackage{booktabs} % for professional tables
% \usepackage[numbers]{natbib}


% Note. For the workshop paper template, both \title{} and \workshoptitle{} are required, with the former indicating the paper title shown in the title and the latter indicating the workshop title displayed in the footnote. 



% The \author macro works with any number of authors. There are two commands
% used to separate the names and addresses of multiple authors: \And and \AND.
%
% Using \And between authors leaves it to LaTeX to determine where to break the
% lines. Using \AND forces a line break at that point. So, if LaTeX puts 3 of 4
% authors names on the first line, and the last on the second line, try using
% \AND instead of \And before the third author name.


% \author{%
%   David S.~Hippocampus\thanks{Use footnote for providing further information
%     about author (webpage, alternative address)---\emph{not} for acknowledging
%     funding agencies.} \\
%   Department of Computer Science\\
%   Cranberry-Lemon University\\
%   Pittsburgh, PA 15213 \\
%   \texttt{hippo@cs.cranberry-lemon.edu} \\
%   % examples of more authors
%   % \And
%   % Coauthor \\
%   % Affiliation \\
%   % Address \\
%   % \texttt{email} \\
%   % \AND
%   % Coauthor \\
%   % Affiliation \\
%   % Address \\
%   % \texttt{email} \\
%   % \And
%   % Coauthor \\
%   % Affiliation \\
%   % Address \\
%   % \texttt{email} \\
%   % \And
%   % Coauthor \\
%   % Affiliation \\
%   % Address \\
%   % \texttt{email} \\
% }


\begin{document}



%%%%%%%%%%%%%%%%%%%%%%%%%%%%%%%%
% THEOREMS
%%%%%%%%%%%%%%%%%%%%%%%%%%%%%%%%
\theoremstyle{plain}
\newtheorem{theorem}{Theorem}[section]
\newtheorem{proposition}[theorem]{Proposition}
\newtheorem{lemma}[theorem]{Lemma}
\newtheorem{corollary}[theorem]{Corollary}
\theoremstyle{definition}
\newtheorem{definition}[theorem]{Definition}
\newtheorem{assumption}[theorem]{Assumption}
\theoremstyle{remark}
\newtheorem{remark}[theorem]{Remark}

\title{
Understanding Sharpness and Forgetting under Constrained Update Geometry in Continual Learning
% Sharpness under Constrained Update Geometry: A PAC-Bayesian Perspective on LoRA-Based Continual Learning
}
\maketitle


\begin{abstract}

% Sharpness-aware optimization is motivated by PAC-Bayesian generalization theory, where optimizing perturbation-smoothed empirical risk encourages flat solutions in standard single-task learning. Although sharpness-based methods are effective in continual learning, their relation to forgetting remains theoretically unclear. Existing explanations are mainly geometric: flat old-task basins are assumed to tolerate subsequent updates and thus undergo smaller increases in old-task risk. This does not explain why perturbation-smoothed empirical risk should relate to forgetting, which is past-task performance degradation rather than single-task generalization error. The gap is sharper in LoRA-based parameter-efficient continual learning, where constrained update geometry restricts future updates to adapter-induced directions and decouples perturbation, update, and interference directions. Thus, full-parameter-space sharpness is not automatically the relevant theoretical object.

% Sharpness-aware optimization is motivated by PAC-Bayesian generalization theory, where perturbation-smoothed empirical risk favors flat minima. However, its role in continual learning remains theoretically under-explained. Existing explanations are largely geometric: flatter old-task basins are expected to tolerate later updates and thus suffer less forgetting. This intuition does not provide a unified theory explaining why the sharpness-aware objective optimized on the current task should be related to the degradation of previous-task performance after subsequent learning. This raises our first question: why should current-task perturbation-smoothed empirical risk be theoretically connected to forgetting? In LoRA-based parameter-efficient continual learning, a second question arises because future adaptation is restricted to adapter-induced update directions under a frozen backbone: which perturbation support defines the sharpness term that enters the continual-learning bound? In this setting, full-parameter sharpness, raw factor-space sharpness, and effective adapter-update sharpness are no longer equivalent.

% Sharpness-aware optimization is well grounded in PAC-Bayesian theory as a mechanism for controlling single-task generalization through perturbation-smoothed empirical risk. Continual learning, however, concerns a posterior trajectory  and forgetting is evaluated by the degradation of previous task performance. This creates a theoretical gap: why should the sharpness-aware objective optimized on the current task be connected to forgetting?

Sharpness-aware optimization is well grounded in PAC-Bayesian theory as a mechanism for controlling single-task generalization through perturbation-smoothed empirical risk. Continual Learning (CL) changes the object of analysis from a single posterior to a history-dependent posterior trajectory, where performance is evaluated along a sequence of updates. This raises a theoretical question: why should the sharpness-aware objective optimized on the current task be connected to trajectory-level CL performance? 
This question becomes sharper in LoRA-based continual learning, where constrained adapter-update geometry makes full-parameter sharpness misaligned with the directions that actually govern sequential adaptation.
% The question is further complicated in LoRA-based continual learning, where a frozen backbone restricts future adaptation to adapter-induced update directions. In this constrained geometry, full-parameter sharpness and adapter-update are not aligned.
We investigate the theoretical mechanism underlying the role of flatness in CL through a PAC-Bayesian analysis. 
First, we derive a pathwise hierarchical decomposition showing that task-conditioned perturbation-smoothed empirical risks accumulate in a bound on trajectory-level sequential generalization, together with within-task adaptation and hyperposterior drift complexities. 
% Second, we specialize the analysis to frozen-backbone LoRA and prove that the bound-relevant smoothed empirical risk and KL complexity reduce to the adapter-update geometry.
Second, we specialize the bound to LoRA based CL and prove that the perturbation-smoothed empirical risk and KL complexity reduce to quantities supported on the adapter-update geometry, identifying adapter-update sharpness as the relevant flatness object.
% Second, we prove a PECL support-reduction result showing that the bound-relevant smoothed empirical risk and KL complexity are supported on the adapter-update geometry, thereby identifying adapter-update sharpness as the relevant flatness object under LoRA-constrained continual learning.
% We investigate the theoretical mechanism underlying the role of flatness in CL through a PAC-Bayesian analysis.
% % We address this gap through a PAC-Bayesian analysis under constrained update geometry. 
% First, we derive a pathwise hierarchical decomposition that relates perturbation-smoothed empirical risk to performance along posterior trajectories. Second, we specialize the result to frozen-backbone LoRA adaptation and prove that the bound-relevant smoothed empirical risk and KL complexity reduce to the admissible effective adapter-update geometry. The analysis identifies adapter sharpness as the relevant flatness object for LoRA-based continual learning. 
Experiments across perturbation types, LoRA variants, and architectures support the theory.

% Sharpness-aware optimization is well grounded in PAC-Bayesian theory as a mechanism for controlling single-task generalization through perturbation-smoothed empirical risk. However, Continual learning concerns a posterior trajectory, and forgetting is evaluated by the degradation of earlier task performance. This creates a theoretical gap: why should the sharpness-aware objective optimized on the current task be connected to forgetting? This question becomes sharper in LoRA-based continual learning, where the frozen backbone restricts future adaptation to adapter-induced update directions. In this constrained geometry, full-parameter sharpness and effective adapter-update sharpness are no longer equivalent, raising a second question: which perturbation support defines the sharpness term that affect forgetting?
% We answer these two questions through a PAC-Bayesian analysis under constrained update geometry. First, we derive a pathwise hierarchical decomposition that links perturbation-smoothed empirical risk to forgetting by modeling continual learning as a posterior trajectory. Second, we specialize the result to frozen-backbone LoRA and prove that the bound-relevant smoothed empirical risk and KL complexity reduce to the admissible effective adapter-update space. The analysis identifies effective adapter-update sharpness as the relevant flatness object for LoRA-based continual learning. Experiments across perturbation supports, LoRA variants, datasets, and architectures support the resulting predictions.

% We provide a PAC-Bayesian analysis of sharpness under constrained update geometry. Our pathwise hierarchical decomposition connects sharpness-based generalization and forgetting by viewing forgetting as past-task generalization degradation under later posteriors. Specializing this analysis to frozen-backbone LoRA adaptation, we prove a PECL support-reduction theorem showing that the bound-relevant perturbation-smoothed empirical risk and KL complexity reduce to the admissible effective adapter-update geometry. This identifies basis-invariant effective adapter-update sharpness as the relevant sharpness object for LoRA-constrained continual learning. Experiments across perturbation supports, LoRA variants, datasets, and architectures validate the theory-derived predictions.



% Sharpness-aware optimization was originally motivated by the PAC-Bayesian view that flat minima, measured through perturbation-smoothed empirical risk, tend to generalize better in standard single-task learning. 
% Sharpness-aware optimization was originally motivated by PAC-Bayesian generalization theory, where optimizing perturbation-smoothed empirical risk encourages flat solutions in standard single-task learning.
% Although sharpness-based methods have also shown strong performance in continual learning, their connection to forgetting remains theoretically under-explained.
% Existing explanations are mostly geometric: flat old-task basins are assumed to tolerate larger subsequent updates and thus undergo a smaller increase in old-task risk.
% Existing explanations are mostly geometric: flat old-task basins are assumed to tolerate larger subsequent updates and thus incur less forgetting.
% However, this view does not explain why perturbation-smoothed empirical risk should be related to forgetting, which is the degradation of past-task performance rather than a single-task generalization error.
% In LoRA-based parameter-efficient continual learning, constrained update geometry further sharpens this gap: the frozen backbone restricts future updates to adapter-induced directions, so perturbation, update, and interference directions no longer coincide. Consequently, full-parameter-space sharpness is not automatically the relevant theoretical object.
% In LoRA-based parameter-efficient continual learning, this gap is amplified because the frozen backbone restricts future updates to adapter-induced directions, so perturbation, update, and interference directions no longer coincide. Consequently, full-parameter-space sharpness is not automatically the relevant theoretical object.
% Existing exlanations are mostly geometric: a flat old-task basin is expected to tolerate larger subsequent updates and therefore suffer a smaller increase in old-task risk. This intuition is plausible, but it does not explain why perturbation-smoothed empirical risk should be related to forgetting, since forgetting is not a standard single-task generalization error but the degradation of past-task performance. The ambiguity becomes more fundamental in LoRA-based parameter-efficient continual learning, where the pretrained backbone is frozen and future updates are restricted to adapter-induced directions. Perturbation, update, and interference directions therefore no longer coincide, making both full-parameter-space sharpness theoretically questionable.
% We study this problem through a PAC-Bayesian analysis of sharpness under constrained update geometry. We first derive a pathwise hierarchical decomposition that connects sharpness-based generalization and forgetting by interpreting forgetting as past-task generalization degradation under later posteriors. We then specialize the analysis to frozen-backbone LoRA adaptation and prove a PECL support-reduction theorem, showing that the bound-relevant perturbation-smoothed empirical risk and KL complexity reduce to the admissible effective adapter-update geometry. This identifies basis-invariant effective adapter-update sharpness, rather than full-parameter-space or raw factor-space sharpness, as the relevant sharpness object for LoRA-constrained continual learning. Experiments across perturbation scopes, LoRA variants, datasets, and architectures support the resulting theory-derived predictions.


% Flatness has become a central explanation for reduced forgetting in continual learning, yet its theoretical role remains under-specified. The prevailing view is geometric: a flat old-task basin should tolerate larger subsequent updates and thus suffer less forgetting. This intuition is insufficient  because it does not explain why perturbation-smoothed empirical risk should control forgetting in theortical. Forgetting is not a standard single-task generalization error and no esting theorytical expalnation. The ambiguity becomes more fundamental in LoRA-based parameter-efficient continual learning, where the pretrained backbone is frozen and future updates are restricted to adapter-induced directions. Perturbation, update, and interference directions therefore no longer coincide, making both full-parameter-space sharpness and raw LoRA factor-space sharpness theoretically questionable.

% We provide a PAC-Bayesian account of flatness under this constrained update geometry. First, we derive a pathwise hierarchical PAC-Bayesian decomposition that links flatness-based generalization and forgetting through perturbation-smoothed empirical risk, within-task adaptation complexity, and hyperposterior drift complexity. This yields a posterior-trajectory view in which forgetting is past-task generalization degradation under later posteriors. Second, we prove a PECL support-reduction theorem for frozen-backbone LoRA adaptation. By analyzing the pushforward posterior induced by LoRA parameters on effective adapter updates, we show that the bound-relevant smoothed empirical risk and KL complexity reduce to the admissible effective adapter-update geometry. Consequently, the relevant flatness object is basis-invariant effective adapter-update sharpness, not ambient full-space sharpness or raw factor-space sharpness. Experiments across perturbation scopes, LoRA variants, datasets, and architectures support these theory-derived predictions.











% Flat minima are widely associated with improved generalization and reduce forgetting in continual learning. However, it remains unclear why the perturbation-smoothed empirical risk underlying sharpness-aware optimization should explain forgetting, especially in LoRA-based parameter-efficient continual learning, where the pretrained backbone is frozen and future updates are restricted to adapter-induced directions. In this setting, perturbation, update, and interference directions no longer coincide, making full-parameter-space sharpness theoretically ambiguous.
% We study flatness in LoRA-based continual learning from a PAC-Bayesian perspective under constrained update geometry. First, we derive a pathwise hierarchical PAC-Bayesian decomposition that links flatness-based generalization and forgetting through perturbation-smoothed empirical risk, within-task adaptation complexity, and hyperposterior drift complexity. Second, we prove a PECL support-reduction theorem showing that the bound-relevant perturbation-smoothed empirical risk and KL complexity are determined by the admissible effective adapter-update geometry. This identifies basis-invariant effective adapter-update sharpness, rather than full-parameter-space or raw LoRA factor-space sharpness, as the relevant flatness object under LoRA-constrained updates. Experiments across perturbation scopes, LoRA variants, datasets, and architectures support the resulting theory-derived predictions.
%   % Flatness has become a central explanation for generalization, and SAM-style methods optimize this principle through neighborhood-smoothed empirical risk. Although flatter solutions have also been observed to reduce forgetting in continual learning, the theoretical link remains unclear because forgetting is a sequential degradation of past-task population risk, not a standard single-task generalization error. We study this question in LoRA-based parameter-efficient continual learning (PECL), where the frozen backbone further decouples perturbation directions from trainable update and interference directions. We develop a PAC-Bayesian theory of flatness under this constrained update geometry. First, we formulate forgetting as past-task generalization degradation under later posteriors, obtaining a sequential risk decomposition that connects SAM-style smoothed risk with both current-task generalization and forgetting. Second, we specialize the bound to frozen-backbone LoRA adaptation and show that the perturbations, sharpness-smoothed empirical term, and KL complexity reduce to the admissible effective adapter-update subspace. The resulting theory identifies basis-invariant adapter-subspace sharpness as the relevant flatness object in PECL. Experiments across multiple LoRA variants, perturbation scopes, datasets, and architectures support the theory and show that adapter-scoped sharpness-aware optimization improves continual generalization while reducing forgetting.
\end{abstract}



\section{Introduction}
% Flat minima have long been associated with improved generalization and robustness~\cite{FlatMinima1997,wu2017understandinggeneralizationdeeplearning,chaudhariEntropySGDBiasingGradient2019,NEURIPS2021_995f5e03,JMLRAnEmpiricalInvestigatio,yangRevisitingFlatnessAwareOptimization2025}. Sharpness-aware minimization (SAM) and related methods turn this principle into an optimization objective by penalizing losses under local parameter perturbations~\cite{foretSharpnessAwareMinimizationEfficiently2021,pmlr-v151-bisla22a,Zhang_GAM,NEURIPS2024_C-Flat}. In continual learning, studies have also reported that flatter solutions reduce catastrophic forgetting~\cite{,shiOvercomingCatastrophicForgetting2021,NEURIPS2020_518a38cc,liRevisitingCatastrophicForgetting2024,RevisitFlatnessOGD}. However, the theoretical basis linking flatness to reduced forgetting in continual learning remains incomplete, especially from the PAC-Bayesian perspective that originally motivates SAM methods. Existing explanations usually rely on the geometric intuition that a flat old-task basin tolerates larger subsequent updates, and hence suffers a smaller loss increase~\cite{NEURIPS2020_518a38cc,liRevisitingCatastrophicForgetting2024,RevisitFlatnessOGD}. While plausible, this view does not explain why the PAC-Bayesian smoothed-risk objective underlying SAM should also govern forgetting, which is a sequential degradation of past-task performance rather than a standard single-task generalization error.  Thus, the role of flatness in continual learning requires a theory that connects currrent generalization with old-task forgetting.


% Flat minima have long been associated with improved generalization and robustness~\cite{FlatMinima1997,wu2017understandinggeneralizationdeeplearning,chaudhariEntropySGDBiasingGradient2019,NEURIPS2021_995f5e03,JMLRAnEmpiricalInvestigatio,yangRevisitingFlatnessAwareOptimization2025}. Sharpness-aware minimization (SAM) and related methods turn this principle into an optimization objective by penalizing losses under local parameter perturbations~\cite{foretSharpnessAwareMinimizationEfficiently2021,pmlr-v151-bisla22a,Zhang_GAM,NEURIPS2024_C-Flat}. In continual learning, studies have also reported that flat minima can reduce catastrophic forgetting and improve the stability-plasticity trade-off~\cite{shiOvercomingCatastrophicForgetting2021,NEURIPS2020_518a38cc,liRevisitingCatastrophicForgetting2024,RevisitFlatnessOGD}. However, the theoretical basis linking flatness to reduced forgetting in continual learning remains incomplete, especially from the PAC-Bayesian perspective that originally motivates SAM-style methods. Existing explanations usually rely on the geometric intuition that a flat old-task basin tolerates larger subsequent updates, and hence suffers a smaller increase in old-task loss~\cite{NEURIPS2020_518a38cc,liRevisitingCatastrophicForgetting2024,RevisitFlatnessOGD}. While plausible, this view does not explain why the PAC-Bayesian smoothed empirical risk underlying SAM should also appear in the analysis of forgetting, which is a sequential degradation of past-task population risk rather than a standard single-task generalization error. Thus, understanding the role of flatness in continual learning requires a theory that connects neighborhood-smoothed generalization with the degradation of past-task population risk under later updates.
% Flat minima have long been associated with improved generalization and robustness~\cite{FlatMinima1997,wu2017understandinggeneralizationdeeplearning,chaudhariEntropySGDBiasingGradient2019,NEURIPS2021_995f5e03,JMLRAnEmpiricalInvestigatio,yangRevisitingFlatnessAwareOptimization2025}. Sharpness-aware minimization (SAM) and related methods turn this principle into an optimization objective by penalizing losses under local parameter perturbations~\cite{foretSharpnessAwareMinimizationEfficiently2021,pmlr-v151-bisla22a,Zhang_GAM,NEURIPS2024_C-Flat}. In continual learning, studies have also reported that flat minima can reduce catastrophic forgetting and improve the stability-plasticity trade-off~\cite{shiOvercomingCatastrophicForgetting2021,NEURIPS2020_518a38cc,liRevisitingCatastrophicForgetting2024,RevisitFlatnessOGD}. Existing explanations usually rely on the geometric intuition that a flat old-task basin tolerates larger subsequent updates and therefore suffers a smaller increase in old-task loss~\cite{NEURIPS2020_518a38cc,liRevisitingCatastrophicForgetting2024,RevisitFlatnessOGD}. While plausible, this explanation remains incomplete. It does not clarify why the PAC-Bayesian smoothed empirical risk underlying SAM should also appear in the analysis of forgetting, which is a sequential degradation of past-task population risk rather than a standard single-task generalization error. More importantly, it hides a directional assumption: the flat directions being measured must coincide with the directions along which future tasks can move the model and interfere with previous solutions. Thus, understanding the role of flatness in continual learning requires a theory that connects neighborhood-smoothed generalization with past-task risk degradation along future update directions.
% Flat minima() have long been associated with improved generalization and robustness~\cite{FlatMinima1997,wu2017understandinggeneralizationdeeplearning,chaudhariEntropySGDBiasingGradient2019,NEURIPS2021_995f5e03,JMLRAnEmpiricalInvestigatio,yangRevisitingFlatnessAwareOptimization2025}.
% Motivated by PAC-Bayesian analyses, sharpness-aware minimization (SAM) and related methods~\cite{foretSharpnessAwareMinimizationEfficiently2021,pmlr-v151-bisla22a,Zhang_GAM,NEURIPS2024_C-Flat} by optimizing losses under local parameter perturbations (we called it sharpness-based empirical risk) to encouraging flat minima  (i.e., promoting flatness or reducing sharpness; we use sharpness as the unified term in this work). In continual learning, flatness has also been observed to reduce catastrophic forgetting and improve the stability-plasticity trade-off~\cite{shiOvercomingCatastrophicForgetting2021,NEURIPS2020_518a38cc,liRevisitingCatastrophicForgetting2024,RevisitFlatnessOGD}. However, existing explanations usually rely on the geometric intuition that a flat old-task basin tolerates larger subsequent updates and therefore suffers a smaller increase in old-task loss~\cite{NEURIPS2020_518a38cc,liRevisitingCatastrophicForgetting2024,RevisitFlatnessOGD}. While plausible, this intuition does not explain why the PAC-Bayesian sharpness-based empirical risk underlying SAM-style objectives should also be relevant to forgetting. This is nontrivial because forgetting is not a standard single-task generalization error, but a sequential degradation of past-task population risk after later updates. This raises the first question addressed in this work: how can flatness-based generalization and old-task forgetting be described within a unified PAC-Bayesian view of continual learning?

 % Prior work has further suggested that flatter solutions can mitigate catastrophic forgetting in continual learning~\cite{NEURIPS2020_518a38cc,JMLRAnEmpiricalInvestigatio,shiOvercomingCatastrophicForgetting2021,liRevisitingCatastrophicForgetting2024}.



% Flat minima in the loss landscape have long been associated with improved generalization and robustness~\cite{FlatMinima1997,wu2017understandinggeneralizationdeeplearning,chaudhariEntropySGDBiasingGradient2019,NEURIPS2021_995f5e03,JMLRAnEmpiricalInvestigatio,yangRevisitingFlatnessAwareOptimization2025}.
% Motivated by PAC-Bayesian analyses, sharpness-aware minimization (SAM) and related methods~\cite{foretSharpnessAwareMinimizationEfficiently2021,pmlr-v151-bisla22a,Zhang_GAM,NEURIPS2024_C-Flat} turn this principle into an optimization objective by minimizing empirical risk under local parameter perturbations.
% PAC-Bayesian flatness analyses motivate sharpness-aware methods such as SAM, which improve generalization and robustness by minimizing empirical risk under local parameter perturbations~\cite{foretSharpnessAwareMinimizationEfficiently2021,mollenhoffSAMOptimalRelaxation2022}.
% Motivated by PAC-Bayesian analyses, sharpness-aware minimization (SAM) and related methods~\cite{foretSharpnessAwareMinimizationEfficiently2021,pmlr-v151-bisla22a,Zhang_GAM,NEURIPS2024_C-Flat} operationalize flatness by minimizing empirical risk under local parameter perturbations.
% Flat minima have long been associated with improved generalization and robustness~\cite{FlatMinima1997,wu2017understandinggeneralizationdeeplearning,chaudhariEntropySGDBiasingGradient2019,NEURIPS2020_518a38cc,NEURIPS2021_995f5e03,JMLRAnEmpiricalInvestigatio,yangRevisitingFlatnessAwareOptimization2025}, and 
% % sharpness-aware optimization methods such as 
% sharpness-aware minimization (SAM) and related methods~\cite{foretSharpnessAwareMinimizationEfficiently2021,pmlr-v151-bisla22a,Zhang_GAM,NEURIPS2024_C-Flat} operationalize this principle by minimizing empirical risk under local parameter perturbations~\cite{foretSharpnessAwareMinimizationEfficiently2021,pmlr-v151-bisla22a,Zhang_GAM,NEURIPS2024_C-Flat}.
% % These methods have also shown
% Previous work has also shown empirical benefits mitigate catastrophic forgetting in continual learning (CL) ~\cite{shiOvercomingCatastrophicForgetting2021,NEURIPS2020_518a38cc,liRevisitingCatastrophicForgetting2024,RevisitFlatnessOGD}.
% However, the theoretical connection between sharpness-aware objectives and continual learning  performance remains under-explained.
% Flat minima in the loss landscape have long been associated with improved generalization and robustness~\cite{FlatMinima1997,wu2017understandinggeneralizationdeeplearning,chaudhariEntropySGDBiasingGradient2019,NEURIPS2021_995f5e03,JMLRAnEmpiricalInvestigatio,yangRevisitingFlatnessAwareOptimization2025}. 
% Motivated by PAC-Bayesian analyses, sharpness-aware minimization (SAM) and related methods~\cite{foretSharpnessAwareMinimizationEfficiently2021,pmlr-v151-bisla22a,Zhang_GAM,NEURIPS2024_C-Flat} operationalize this principle by minimizing empirical risk under local parameter perturbations. Sharpness-aware learning has also shown empirical benefits in continual learning, including improved average performance and reduced forgetting~\cite{shiOvercomingCatastrophicForgetting2021,NEURIPS2020_518a38cc,liRevisitingCatastrophicForgetting2024,RevisitFlatnessOGD}. However, the theoretical connection between sharpness-aware objectives and continual learning performance remains under-explained.


Flat minima have long been associated with improved generalization and robustness~\cite{FlatMinima1997,wu2017understandinggeneralizationdeeplearning,chaudhariEntropySGDBiasingGradient2019,NEURIPS2021_995f5e03,JMLRAnEmpiricalInvestigatio,yangRevisitingFlatnessAwareOptimization2025}.
Sharpness-aware minimization (SAM) and related methods~\cite{foretSharpnessAwareMinimizationEfficiently2021,pmlr-v151-bisla22a,Zhang_GAM,NEURIPS2024_C-Flat} operationalize this principle by minimizing empirical risk under local parameter perturbations.
Although sharpness-aware learning has been empirically shown to mitigate catastrophic forgetting in continual learning (CL)~\cite{shiOvercomingCatastrophicForgetting2021,NEURIPS2020_518a38cc,liRevisitingCatastrophicForgetting2024,RevisitFlatnessOGD}, why such task-wise objectives improve trajectory-level performance remains theoretically under-explained.

In single-task learning, perturbation-smoothed empirical risk relates flatness to generalization for the posterior learned from one task~\cite{foretSharpnessAwareMinimizationEfficiently2021,mollenhoffSAMOptimalRelaxation2022}.
% In single-task learning, perturbation-smoothed empirical risk relates flatness to generalization for one task-induced posterior~\cite{foretSharpnessAwareMinimizationEfficiently2021,mollenhoffSAMOptimalRelaxation2022}. 
However, CL changes the relevant object from a single posterior to a history-dependent posterior trajectory.
Existing explanations are mainly geometric, arguing that a flat old-task basin can better tolerate later updates and therefore suffer less degradation, as illustrated in Fig.~\ref{fig:flatness_and_scope}(a).
This intuition is limited because it treats flatness as a static property of a past-task solution, while sharpness in CL is local, task-conditioned, and posterior-dependent.
% This intuition is limited because it treats flatness as a static property of an old solution.
% However, sharpness is local, task-conditioned, and posterior-dependent in CL, because it depends on the dataset and the weights around which perturbations are evaluated.
The missing theoretical step is therefore to explain how perturbation-smoothed risks optimized at individual update steps enter a trajectory-level generalization bound.
% Thus, the missing theoretical step is to explain how the perturbation-smoothed risks optimized at individual update steps accumulate into performance along the posterior trajectory.
This leads to our first question Q1:  how can sharpness-based generalization be connected to trajectory-level performance within a unified PAC-Bayesian view?
% Thus, the key missing theory is a unified account of how task-level perturbation-smoothed generalization, optimized during sequential updates, is connected to performance along the posterior trajectory.
% This leads to our first question: how can sharpness-based generalization and trajectory-level continual learning performance be connected within a unified PAC-Bayesian view?
% Existing explanations are mainly geometric, arguing that a flat old-task basin can better tolerate later updates and therefore suffer less degradation, as illustrated in Fig.~\ref{fig:flatness_and_scope}(a). 
% This intuition is incomplete because it treats flatness as a static property of an old solution. 
% Sharpness is local, task-conditioned, and posterior-dependent in CL: it depends on the dataset and the weights around which perturbations are evaluated.
% Since sharpness is optimized when each posterior is produced, it should be understood as a property of the sequential weight-update process rather than a fixed property of one old basin. 
% The central question is therefore how these task-local perturbation-smoothed risks, induced during learning, affect performance along the entire posterior trajectory.


% Existing explanations are mainly geometric, arguing that a flat old-task basin can better tolerate later updates and therefore suffer less degradation, as illustrated in Fig.~\ref{fig:flatness_and_scope}(a). 
% This intuition is incomplete because it treats flatness as a static property of an old solution. 
% Sharpness is local, task-conditioned, and posterior-dependent in CL: it depends on the dataset and the weight around which perturbations are evaluated.
% % , the perturbation distribution. 
% Once the weight update, the relevant object is therefore not the flatness of a single old basin, but the sequence of perturbation-smoothed risks induced along the posterior trajectory. 
% This leads to our first question: how can sharpness-based generalization and trajectory-level continual learning performance be connected within a unified PAC-Bayesian view of posterior trajectories?



% In single-task learning, perturbation-smoothed empirical risk relates flatness to generalization for one task posterior~\cite{foretSharpnessAwareMinimizationEfficiently2021,mollenhoffSAMOptimalRelaxation2022}. However, CL changes the relevant object from a single posterior to a history-dependent posterior trajectory. Existing explanations are mainly geometric, arguing that a flat old-task basin can better tolerate later updates and therefore suffer less degradation, as illustrated in Fig.~\ref{fig:flatness_and_scope}(a). This intuition is incomplete: it does not explain why optimizing current-task perturbation-smoothed empirical risk should improve performance over the entire learning trajectory. Since sharpness is task-conditioned, this missing link requires a theory that explicitly connects task-level sharpness with posterior trajectories induced by sequential updates. This leads to our first question: how can sharpness-based generalization and trajectory-level continual learning performance be connected within a unified PAC-Bayesian view of posterior trajectories?


% These analyses are primarily single-task in nature: sharpness is evaluated at a learned checkpoint with respect to the same dataset~\cite{mollenhoffSAMOptimalRelaxation2022}.
% However, why such task-conditioned sharpness should matter for continual learning remains theoretically under-explained.
% Sharpness-aware learning also has shown empirical benefits in continual learning, including improved average performance and reduced forgetting~\cite{shiOvercomingCatastrophicForgetting2021,NEURIPS2020_518a38cc,liRevisitingCatastrophicForgetting2024,RevisitFlatnessOGD}.
% However, the theoretical connection between sharpness-aware objectives and continual learning performance remains under-explained.
% In standard single-task learning, PAC-Bayesian analyses justify flat minima through perturbation-smoothed empirical risk, where sharpness is evaluated for a posterior under the dataset, loss, and perturbation distribution associated with that task~\cite{foretSharpnessAwareMinimizationEfficiently2021,mollenhoffSAMOptimalRelaxation2022}.
% Continual learning changes this object from a single posterior to a history-dependent posterior trajectory.
% Existing explanations are mainly geometric, arguing that a flatter old-task basin can better tolerate later updates and therefore suffers less degradation, as illustrated in Fig.~\ref{fig:flatness_and_scope}(a).
% Yet this intuition does not explain why optimizing current-task perturbation-smoothed empirical risk should improve trajectory-level continual learning performance.
% Since sharpness is task-conditioned rather than intrinsic to the model, this missing step requires a theory that explicitly relates task-level sharpness to the posterior trajectory induced by continual updates.
% This leads to our first question: how can sharpness-based generalization and trajectory-level continual learning performance be connected within a unified PAC-Bayesian view of posterior trajectories?



% The first gap is that current-task sharpness does not directly explain trajectory-level continual learning performance.
% In standard single-task learning, sharpness is evaluated for one posterior on the same dataset and loss that induce it.
% However, continual learning produces a history-dependent posterior trace, and sharpness becomes a task-conditioned quantity determined by the checkpoint, dataset, loss, and perturbation distribution.
% The relevant question is therefore not only whether each task is learned in a flat region, but how task-conditioned sharpness along the trace affects overall continual learning performance.
% This leaves a missing theoretical step between optimizing current-task perturbation-smoothed empirical risk and explaining sequential performance under continual updates.


% The first gap is that current-task sharpness does not directly explain trajectory-level continual learning performance. 





% In continual learning, the learner does not produce only one posterior, but a history-dependent posterior trace across tasks.
% Accordingly, the central evaluation object is not only whether the current task is learned well or how many earlier tasks are forgotten, but the overall performance induced by the whole learning trace.
% Although flat minima have been empirically associated with improved continual learning performance and reduced forgetting~\cite{shiOvercomingCatastrophicForgetting2021,NEURIPS2020_518a38cc,liRevisitingCatastrophicForgetting2024,RevisitFlatnessOGD}, existing explanations usually remain geometric.
% They argue that a flatter old-task basin can better tolerate later updates and therefore suffers a smaller increase in old-task risk, as illustrated in Fig.~\ref{fig:flatness_and_scope}(a).
% This intuition is plausible, but it does not provide a PAC-Bayesian account of why the current-task perturbation-smoothed empirical risk optimized by SAM-style objectives should control performance along a sequential posterior trace.
% This raises our first question: how can sharpness-based generalization and trajectory-level continual learning performance be connected within a unified PAC-Bayesian view of posterior trajectories?



% Flat minima in the loss landscape have long been associated with improved generalization and robustness~\cite{FlatMinima1997,wu2017understandinggeneralizationdeeplearning,chaudhariEntropySGDBiasingGradient2019,NEURIPS2021_995f5e03,JMLRAnEmpiricalInvestigatio,yangRevisitingFlatnessAwareOptimization2025}. 
% Motivated by PAC-Bayesian analyses, sharpness-aware minimization (SAM) and related methods~\cite{foretSharpnessAwareMinimizationEfficiently2021,pmlr-v151-bisla22a,Zhang_GAM,NEURIPS2024_C-Flat} turn this principle into an optimization objective by minimizing the empirical risk evaluated under local parameter perturbations.
% % , which we refer to as the \emph{perturbation-smoothed empirical risk}. 
% % The excess of this perturbed risk over the unperturbed empirical risk provides a quantitative measure of local sharpness. 
% % Throughout this work, we use \emph{flatness} as a qualitative description of the loss landscape, and \emph{sharpness} as the corresponding quantitative measure of how sensitively the loss changes under local perturbations.
% % Motivated by PAC-Bayesian analyses, sharpness-aware minimization (SAM) and related methods~\cite{foretSharpnessAwareMinimizationEfficiently2021,pmlr-v151-bisla22a,Zhang_GAM,NEURIPS2024_C-Flat} turn this principle into an optimization objective by minimizing perturbation-smoothed empirical risk under local parameter perturbations. 
% % Throughout this work, we use \emph{flatness} as a qualitative description of a loss landscape, and \emph{sharpness} as the corresponding quantitative measure under local perturbations.
% These analyses are primarily single-task in nature: sharpness is evaluated at a learned checkpoint with respect to the dataset and loss that induce the corresponding posterior\cite{mollenhoffSAMOptimalRelaxation2022}.
% However, continual learning changes the object of analysis.
% The learner does not produce one posterior for one task, but a history-dependent posterior trace across tasks.
% Accordingly, the central evaluation object is not only whether the current task is learned well or how many earlier tasks are forgotten, but the overall performance induced by the whole learning trace.


% Flat minima have been empirically associated with improved continual learning performance and reduced forgetting~\cite{shiOvercomingCatastrophicForgetting2021,NEURIPS2020_518a38cc,liRevisitingCatastrophicForgetting2024,RevisitFlatnessOGD}.
% However, existing theoexplanations usually remain geometric.
% They argue that a flatter old-task basin can better tolerate later updates and therefore suffers a smaller increase in old-task risk, as illustrated in Fig.~1(a).
% This intuition is plausible, but it does not provide a PAC-Bayesian account of why the current-task perturbation-smoothed empirical risk should control performance along a sequential posterior trace.
% Indeed, sharpness in continual learning is task-conditioned: it depends on the current checkpoint, dataset, loss, and perturbation distribution, rather than being an intrinsic property of the model alone.
% This leaves a missing theoretical step between optimizing task-conditioned sharpness during current-task learning and explaining trace-level continual learning performance.
% This raises our first question: how can sharpness-based generalization and continual learning performance be connected within a unified PAC-Bayesian view of posterior trajectories?


% In continual learning, flat minima have also been empirically associated with reduced catastrophic forgetting, which is evaluated as the degradation of earlier task performance.~\cite{shiOvercomingCatastrophicForgetting2021,NEURIPS2020_518a38cc,liRevisitingCatastrophicForgetting2024,RevisitFlatnessOGD}.
% However, existing explanations usually rely on the geometric intuition that a flatter old-task basin can better tolerate later updates and therefore suffers a smaller increase in old-task risk, as illustrated in Fig.~1(a). While plausible, this intuition does not yet provide a theoretical connection between sharpness-aware training and forgetting. In CL, sharpness is task-conditioned: it is defined with respect to a specific checkpoint and dataset, rather than being an intrinsic property of the model. Moreover, sharpness-aware objectives such as SAM are optimized only on the current task through its perturbation-smoothed empirical risk. There is therefore still a missing theoretical step between optimizing current-task sharpness and explaining old-task forgetting. This raises the first theoretical question addressed in this work: how can these two quantities be connected within a unified PAC-Bayesian view of continual learning?


% However, existing explanations usually rely on the geometric intuition that a flat old-task basin tolerates larger subsequent updates and therefore suffers a smaller increase in old-task risk, as shown in Fig\ref{fig:flatness_and_scope} (a).
% % However, existing explanations usually rely on the geometric intuition that a flat basin tolerates larger subsequent updates and suffers a smaller increase in the lossland of old-tasks~\cite{NEURIPS2020_518a38cc,liRevisitingCatastrophicForgetting2024,RevisitFlatnessOGD}.
% While plausible, this intuition does not explain why the perturbation-smoothed empirical risk underlying SAM-style objectives should be relevant to forgetting. 
% This is nontrivial because forgetting is not a standard single-task generalization error, but the degradation of previous-task performance after later updates. 
% % From a PAC-Bayesian perspective, forgetting can be formalized as generalization degradation on past tasks under later posteriors. 
% This raises the first theoretical question addressed in this work: how can sharpness-based generalization and forgetting be described within a unified PAC-Bayesian view of continual learning?

% how can perturbation-based generalization and forgetting be described within a unified PAC-Bayesian view of continual learning?


% The central question is therefore whether forgetting can be understood as a form of past-task generalization degradation under later posteriors.

% Flat minima have long been associated with improved generalization and robustness~\cite{FlatMinima1997,wu2017understandinggeneralizationdeeplearning,chaudhariEntropySGDBiasingGradient2019,NEURIPS2021_995f5e03,JMLRAnEmpiricalInvestigatio,yangRevisitingFlatnessAwareOptimization2025}. Sharpness-aware minimization (SAM) and related methods turn this principle into an optimization objective by penalizing losses under local parameter perturbations~\cite{foretSharpnessAwareMinimizationEfficiently2021,pmlr-v151-bisla22a,Zhang_GAM,NEURIPS2024_C-Flat}. In continual learning, studies have also reported that flat minima can reduce catastrophic forgetting and improve the stability-plasticity trade-off~\cite{shiOvercomingCatastrophicForgetting2021,NEURIPS2020_518a38cc,liRevisitingCatastrophicForgetting2024,RevisitFlatnessOGD}. However, the theoretical basis linking flatness to reduced forgetting in continual learning remains incomplete, especially from the PAC-Bayesian perspective that originally motivates SAM-style methods. Existing explanations usually rely on the geometric intuition that a flat old-task basin tolerates larger subsequent updates, and hence suffers a smaller increase in old-task loss~\cite{NEURIPS2020_518a38cc,liRevisitingCatastrophicForgetting2024,RevisitFlatnessOGD}. While plausible, this view does not explain why the PAC-Bayesian smoothed empirical risk underlying SAM should also appear in the analysis of forgetting, which is a sequential degradation of past-task population risk rather than a standard single-task generalization error. Thus, understanding the role of flatness in continual learning requires a theory that connects neighborhood-smoothed generalization with the degradation of past-task population risk under later updates.



% Flat minima have long been associated with improved generalization and robustness~\cite{FlatMinima1997,wu2017understandinggeneralizationdeeplearning,chaudhariEntropySGDBiasingGradient2019,NEURIPS2021_995f5e03,JMLRAnEmpiricalInvestigatio,yangRevisitingFlatnessAwareOptimization2025}. Sharpness-aware minimization (SAM) and related methods turn this principle into an optimization objective by penalizing losses under local parameter perturbations~\cite{foretSharpnessAwareMinimizationEfficiently2021,pmlr-v151-bisla22a,Zhang_GAM,NEURIPS2024_C-Flat}. In continual learning, studies have also reported that flatter solutions can reduce catastrophic forgetting and improve the stability-plasticity trade-off~\cite{shiOvercomingCatastrophicForgetting2021,NEURIPS2020_518a38cc,liRevisitingCatastrophicForgetting2024,RevisitFlatnessOGD}. However, the theoretical basis linking flatness to reduced forgetting in continual learning remains incomplete, especially from the PAC-Bayesian perspective that originally motivates SAM-style methods. Existing explanations usually rely on the geometric intuition that a flat old-task basin tolerates larger subsequent updates, and hence suffers a smaller increase in old-task loss~\cite{NEURIPS2020_518a38cc,liRevisitingCatastrophicForgetting2024,RevisitFlatnessOGD}. While plausible, this view does not explain why the PAC-Bayesian smoothed-risk objective underlying SAM should also relate to forgetting, which is a sequential degradation of past-task performance rather than a standard single-task generalization error. Thus, understanding the role of flatness in continual learning requires a theory that connects neighborhood-smoothed generalization with the degradation of past-task population risk under later updates.




% This gap becomes more fundamental in LoRA-based parameter-efficient continual learning (PECL), where flatness is no longer a property of an unconstrained optimization geometry. In full-parameter training, perturbation directions, optimization directions, and interference directions are usually identified with the same ambient parameter space. PECL breaks this identification. Since the pretrained backbone is frozen, future adaptation is confined to adapter-induced low-rank effective update directions, and any later interference with previous tasks must occur through this admissible update geometry. Consequently, full-space sharpness may measure sensitivity along directions that future tasks can never update, while raw factor-space sharpness may depend on an arbitrary LoRA parameterization rather than on the effective model change, since $\Delta W=BA$. Although Flat-LoRA~\cite{li2024flatlora} addresses the distinction between raw LoRA factors and the effective merged update, it does not answer the continual learning question: under frozen-backbone PECL, should flatness be defined over the full model, over the adapter factors, or over the admissible effective update subspace that carries all future adaptation? This is a theoretical issue, not merely an implementation choice. Until the relevant perturbation support is identified, flatness-based PECL lacks a principled criterion for deciding what should be perturbed, which curvature should be controlled, and how sharpness relates to forgetting under constrained updates.

% This limitation becomes more severe in LoRA-based parameter-efficient continual learning (PECL), where the basin-tolerance intuition above is no longer directly applicable. 
% The common explanation that a flat old-task basin tolerates larger later updates implicitly assumes that the directions used to measure flatness coincide with the directions along which future updates move the model. However, PECL breaks this assumption. Since the pretrained backbone is frozen, future adaptation is confined to adapter-induced low-rank effective update directions, whereas sharpness can still be measured by perturbing the full parameter space, the raw LoRA factors, or the effective update itself. Flat-LoRA~\cite{li2024flatlora} addresses the representational mismatch between raw LoRA factors and the effective merged update, but it does not study how the sharpness governs sequential forgetting under PECL. Therefore, explaining how flatness affects forgetting in PECL requires more than the standard basin-tolerance intuition: it requires a theory that first relates forgetting to sequential generalization, and then identifies the perturbation geometry on which the corresponding sharpness term should be defined. This raises the PECL-specific question: under frozen-backbone adaptation, which perturbation support defines the sharpness term that is relevant to continual generalization and old-task retention? Without this clarification, flatness-based PECL remains partly heuristic: it is unclear whether one should control sensitivity in the full model, in the adapter factors, or in the effective update geometry that actually carries future adaptation.


% This question becomes particularly important in LoRA-based parameter-efficient continual learning (PECL). In PECL, the pretrained backbone is frozen and new tasks are learned through low-rank adapter updates. This makes continual adaptation scalable, but it also changes the geometry in which flatness should be interpreted. Unlike full-parameter training, where perturbations, updates, and interference are usually considered in the same ambient parameter space, LoRA-based PECL admits several possible sharpness objects: one may perturb the full model, the raw LoRA factors, or the effective merged update. These choices are not equivalent. Full-model perturbations can probe directions that are not trainable under frozen-backbone adaptation, while raw factor perturbations may depend on the parameterization of the same effective update. Flat-LoRA~\cite{li2024flatlora} addresses the representational mismatch between raw LoRA factors and the effective merged update, but it does not study how the corresponding sharpness object relates to sequential forgetting under frozen-backbone PECL. Thus, the PAC-Bayesian explanation of flatness in continual learning must be specialized to the PECL update geometry. Without this specialization, flatness-based PECL remains partly heuristic: it is unclear whether one should control sensitivity in the full model, in the adapter factors, or in the effective update geometry that actually carries future adaptation.

% \vspace{-6pt}
\begin{figure}[tp]
    \centering
    \begin{minipage}[t]{0.49\linewidth}
        \centering
        \includegraphics[width=\linewidth,height=3.4cm,keepaspectratio,trim=0 1 0 8,clip]{figures_TRML/losslan_flat.pdf}
        % \vspace{-1pt}
        {\scriptsize (a)}
    \end{minipage}
    \hfill
    \begin{minipage}[t]{0.49\linewidth}
        \centering
        \includegraphics[width=\linewidth,height=3.cm,keepaspectratio,trim=0 6 0 6,clip]{figures_TRML/liuxingshiyitu5.pdf}
        % \vspace{-1pt}
        {\scriptsize (b)}
    \end{minipage}
    
    \vspace{-1pt}
    \caption{\textbf{Static flatness intuition and constrained perturbation geometry in PECL.} 
    (a) A  static explanation for flatness in CL: under the same update magnitude, a flat old-task basin (right) yields a smaller loss increase than a sharp one(left). 
    % However, this view treats sharpness as a fixed basin property, whereas sharpness is local, task-specific, and posterior-dependent along the learning trajectory.
    % A explanation for the benefit of flatness in CL is geometric: under the same update magnitude, a flat loss landscape (right) yields a smaller increase in loss than a sharp one (left). 
    (b) In LoRA-based PECL, future adaptation is confined to task-permissible directions (green), whereas standard sharpness perturbs the full parameter space (blue).}
    \label{fig:flatness_and_scope}
    \vspace{-0.80 cm}
\end{figure}
% \vspace{-6pt}
% This ambiguity becomes sharper in LoRA-based parameter-efficient continual learning (PECL), where the frozen backbone restricts future adaptation to adapter update directions. 
% LoRA-based PECL makes the limitation of the static geometric explanation more explicit.

% The perturbation space remains ambiguous in PECL.

LoRA-based PECL introduces a second ambiguity: the perturbation support that defines sharpness is no longer self-evident.
In full-parameter training, perturbations, updates, and interference directions are often implicitly identified with the same parameter space.
In LoRA-based PECL, this identification breaks down because adaptation is restricted to adapter-induced update directions, while sharpness can be measured in the full parameter space, the raw LoRA factor space, or the effective adapter-update space, as illustrated in Fig.~\ref{fig:flatness_and_scope}(b).
Although Flat-LoRA~\cite{li2024flatlora} identifies the mismatch between raw-factor and merged-weight perturbations, it does not determine which perturbation support is relevant to sequential posterior trajectories under LoRA-based CL.
This motivates our second question Q2: under LoRA-constrained adaptation, what is the correct geometric support of the perturbation-smoothed risk that enters the trajectory-level PAC-Bayesian bound?


% LoRA-based CL generate a new problem: the perturbation space that defines sharpness is no longer self-evident and must be reconsidered.
% In full-parameter training, this issue is often hidden because perturbations, weight updates, and interference directions are all in the same full parameter space.
% In LoRA-based PECL, this identification breaks down because future adaptation occurs only within the adapter update subspace, while sharpness can still be measured in the full parameter space, the raw LoRA factor space, or the merged update space, as illustrated in Fig.~\ref{fig:flatness_and_scope}(b). Flat-LoRA~\cite{li2024flatlora} identifies the mismatch between raw-factor and merged-weight perturbations, but does not study how the corresponding sharpness object interacts with forgetting. This motivates our second question: under LoRA-constrained adaptation, what is the correct geometric support of perturbation-smoothed risk relevant to forgetting?
% In LoRA-based CL, this identification breaks down: the future adaptation directions is restricts, while sharpness can be measured in the full parameter space, the raw LoRA factor space, or the merged adapter-update space.
% These supports induce different smoothed risks.
% Thus, the relevant question is which support determines the bound-relevant smoothed empirical risk and KL complexity along the LoRA-induced posterior trajectory.
% This motivates our second question: under LoRA-constrained adaptation, which perturbation support is theoretically relevant for trajectory-level continual learning performance?



% LoRA-based PECL introduces a second ambiguity: the perturbation support of sharpness is no longer align with the trainable parameter space.
% In full-parameter training, perturbation, update, and interference directions are implicitly identified with the full parameter space. However, LoRA-based PECL breaks this identification. Sharpness can be measured by perturbing the full model, the raw LoRA factors, or the effective adapter update, whereas later adaptation is confined to the adapter update subspace. Thus, global full-parameter sharpness is not automatically the relevant object, since it may probe directions inaccessible to future LoRA updates, as illustrated in Fig.~\ref{fig:flatness_and_scope}(b). Flat-LoRA~\cite{li2024flatlora} identifies the mismatch between raw-factor and merged-weight perturbations, but does not study how the corresponding sharpness object interacts with forgetting. This motivates our second question: under LoRA-constrained adaptation, what is the correct geometric support of perturbation-smoothed risk relevant to forgetting?

PAC-Bayesian analysis provides a natural theoretical lens. However, existing PAC-Bayesian theories of CL do not directly address the setting considered here. Lifelong PAC-Bayes studies transfer across tasks sampled i.i.d. from a fixed hyperdistribution~\cite{pentinaPACBayesianBoundLifelong2014}. Its resulting bound is order invariant and leaves the risk evolution along a sequence of continual updates outside its scope, as shown in Fig.\ref{fig:showPAC_combined}(a). A concurrent sequential PAC-Bayesian analysis derives cumulative risk bounds with data-dependent priors and posteriors~\cite{friedman2026pac}, without addressing flatness-based generalization or the constrained update geometry of PECL.
% A concurrent  PAC-Bayesian analysis derives cumulative risk bounds with data-dependent priors and posteriors~\cite{friedman2026pac}. It consider sequential risk, while leaving unspecified the geometric support on which perturbation-smoothed risk should be evaluated under constrained PECL updates.
% Concurrent sequential PAC-Bayesian analyses derive cumulative risk bounds with data dependent priors and posteriors~\cite{friedman2026pac}. It consider sequential risk, while the geometric support of a perturbation-smoothed risk remains discussed. 
Thus, the missing component is a pathwise analysis that connects perturbation-smoothed risk to overall performance along a posterior trajectory and then specializes this analysis to the constrained update geometry induced by PECL.

% PAC-Bayesian analysis provides a natural theoretical lens for the two questions above. However, existing PAC-Bayesian theories of continual learning do not directly address the setting considered here. Lifelong PAC-Bayes studies transfer across tasks sampled i.i.d. from a fixed hyperdistribution~\cite{pentinaPACBayesianBoundLifelong2014}. Its resulting bound is largely order invariant and not characterizes the risk evolution along a sequence of continual updates. 
% % Consequently, it does not model how a later posterior, produced after learning subsequent tasks, can degrade the population risk of earlier tasks. 
% Concurrent sequential PAC-Bayesian analyses derive cumulative risk bounds with data dependent priors and posteriors~\cite{friedman2026pac}, but they are not designed to identify the geometric sharpness. Thus, the missing component is a pathwise analysis that connects perturbation-smoothed risk to forgetting along a posterior trajectory and then specializes this analysis to the constrained update geometry induced by PECL.

% This leads to two setting-specific technical challenges. First, the prior at task $t$ is no longer fixed in advance, but is produced by the learning history as a process $P_t(S_{1:t-1})$. The bound must therefore be adapted to the filtration generated by the observed task sequence and must accumulate pathwise along the realized trajectory, rather than after averaging over tasks drawn independently from an environment distribution. Second, under frozen-backbone LoRA adaptation, the posterior over trainable adapter parameters induces an effective model update while the backbone remains fixed. Existing continual-learning PAC-Bayesian bounds do not specify how the posterior, the KL complexity, and the perturbation-smoothed empirical term behave under this constrained parameterization. This leaves the bound-relevant sharpness object under PECL theoretically ambiguous.


% This theoretical ambiguity becomes in LoRA-based parameter-efficient continual learning (PECL), where sharpness must be understood under constrained update geometry. 
% % In PECL, the pretrained backbone is frozen and new tasks are learned through low-rank adapter updates. 
% % This theoretical ambiguity becomes sharper in LoRA-based parameter-efficient continual learning.
% Since the pretrained backbone is frozen, subsequent adaptation is restricted to adapter-induced update directions.
% As a result, it also changes the geometry in which the sharpness-forgetting connection must be interpreted. 
% In full-parameter training, the directions used to perturb the model, optimize future tasks, and induce interference are often implicitly identified with the full parameter space. 
% However, LoRA-based PECL breaks this identification: sharpness can be measured by perturbing the full model, the raw LoRA factors, or the effective adapter update, whereas future adaptation is restricted to adapter-induced update directions. 
% Once perturbation, update, and interference directions no longer coincide, global full-parameter-space sharpness is not automatically the relevant object. 
% Full-model perturbations may probe directions that are inaccessible to later LoRA updates, as shown in Fig \ref{fig:flatness_and_scope}(b). 
% % while raw factor perturbations may depend on the parameterization of the same effective adapter update. 
% % Indeed, the same effective adapter update admits equivalent factorizations $BA=(BR)(R^{-1}A)$ for any invertible matrix $R$, so sharpness measured in the raw LoRA factor coordinates can change under a reparameterization that leaves the effective update unchanged. 
% Flat-LoRA~\cite{li2024flatlora} points out the perturbation mismatch between the raw LoRA factors and the merged weight, but it does not study how the resulting sharpness object interacts with sequential forgetting under PECL. 
% The second theoretical question we address is thus: under LoRA-constrained adaptation, what is the correct geometric support of the perturbation-smoothed risk term, and why should this support be relevant to  forgetting?


% This question becomes sharper in LoRA-based parameter-efficient continual learning (PECL), where flatness must be understood under constrained update geometry. In PECL, the pretrained backbone is frozen and new tasks are learned through low-rank adapter updates. This makes continual adaptation scalable, but it also changes the geometry in which the flatness-forgetting connection must be interpreted. In full-parameter training, the directions used to perturb the model, optimize future tasks, and induce interference are often implicitly identified with the full parameter space. However, LoRA-based PECL breaks this identification: sharpness can be measured by perturbing the full model, the raw LoRA factors, or the effective update, whereas future adaptation is restricted to adapter-induced effective update directions. Once perturbation, update, and interference directions no longer coincide, global full-space sharpness is not automatically the relevant object.
% These choices are not equivalent. 
% Full-model perturbations may probe directions that are inaccessible to later LoRA updates, while raw factor perturbations may depend on the parameterization of the same effective update. 
% Flat-LoRA~\cite{li2024flatlora} points out the perturbation mismatch between the raw LoRA factors and the merged weight, but it does not study how the resulting sharpness object interacts with sequential forgetting under PECL.
% More broadly, PECL breaks the usual identification between perturbation, update, and interference directions. 
% This raises a PECL-specific theoretical question: when future adaptation is restricted to LoRA-induced effective update directions, should flatness be defined in the full parameter space, the raw adapter-factor space, or the effective adapter-update geometry?
% More broadly, PECL exposes a hidden assumption behind standard flatness arguments: flatness is relevant to forgetting only insofar as it is measured along directions through which later tasks can update the model and interfere with previous solutions.
 % Therefore, the central PECL-specific hypothesis that adapter-update flatness, rather than ambient full-space flatness, governs continual generalization and forgetting requires a theoretical justification. 
 % The second question we address is thus: under LoRA-constrained adaptation, what is the correct geometric support of the sharpness term, and why should this support be relevant to both future admissible updates and forgetting?
 % The second question we address is thus: under LoRA-constrained adaptation, what is the correct geometric support of the sharpness term, and why should this support be relevant to both future adaptation and old-task retention?


% Flat-LoRA~\cite{li2024flatlora} addresses the representational mismatch between raw LoRA factors and the effective merged update, but it does not study how the resulting sharpness object interacts with sequential forgetting under PECL.
% Thus, a PAC-Bayesian explanation of flatness in continual learning must be specialized to the PECL update geometry. Without this specialization, flatness-based PECL remains theoretically ambiguous.
% More broadly, PECL exposes a hidden assumption behind standard flatness arguments: flatness is relevant to retention only insofar as it is measured along directions through which later tasks can update the model and interfere with previous solutions. Once perturbation, update, and interference directions no longer coincide, global full-space sharpness is not automatically the relevant object.

% More broadly, PECL exposes a hidden assumption behind the standard full-parameter flatness intuition: flatness is relevant to retention only insofar as it is measured along directions through which later tasks can update the model and interfere with previous solutions. Once perturbation, update, and interference directions no longer coincide, global full-space flatness is not automatically the relevant object. A theory of flatness in continual learning must therefore specify the geometry on which sharpness is measured and justify why this geometry governs both future adaptation and old-task retention.

% This question becomes particularly important in LoRA-based parameter-efficient continual learning (PECL). In PECL, the pretrained backbone is frozen and new tasks are learned through low-rank adapter updates. This makes continual adaptation scalable, but it also changes the geometry in which the flatness-forgetting connection must be interpreted. In full-parameter training, the directions used to perturb the model, optimize future tasks, and induce interference are usually treated as the same ambient parameter space. However, LoRA-based PECL breaks this identification. The model can be perturbed in the full parameter space, in the raw LoRA factors, or in the effective merged update, but future adaptation is restricted to adapter-induced effective update directions. These choices are not equivalent: full-model perturbations may probe directions that are not trainable under frozen-backbone adaptation, while raw factor perturbations may depend on the parameterization of the same effective update. Flat-LoRA~\cite{li2024flatlora} addresses the representational mismatch between raw LoRA factors and the effective merged update, but it does not study how the corresponding sharpness object relates to sequential forgetting under PECL. 
% Thus, a PAC-Bayesian explanation of flatness in continual learning must be specialized to the PECL update geometry. Without this specialization, flatness-based PECL remains theoretically ambiguous. More broadly, PECL exposes a hidden assumption behind the standard full-parameter flatness intuition: flatness matters for retention only insofar as it is measured along directions through which later tasks can update and interfere. Once perturbation, update, and interference directions no longer coincide, global full-space flatness is no longer automatically the relevant object. The theory must therefore specify the geometry on which sharpness is measured and explain why that geometry governs future adaptation and old-task retention.


% Thus, a PAC-Bayesian explanation of flatness in continual learning must be specialized to the PECL update geometry. Without this specialization, flatness-based PECL remains theoretically ambiguous. More broadly, PECL exposes a hidden assumption behind the standard full-parameter flatness intuition: flatness matters for retention only insofar as it is measured along directions through which later tasks can update and interfere. Once perturbation, update, and interference directions no longer coincide, global full-space flatness is no longer automatically the relevant object. The theory must therefore specify the geometry on which sharpness is measured and explain why that geometry governs future adaptation and old-task retention.

% Without this specialization, flatness-based PECL remains partly heuristic: it is unclear whether the sharpness term relevant to old-task retention should measure sensitivity in the full model, in the adapter factors, or in the effective update geometry that actually carries future adaptation.


% LoRA-based parameter-efficient continual learning (PECL) makes this hidden directional assumption explicit. In full-parameter continual learning, perturbation directions, optimization directions, and interference directions are often implicitly identified with the same ambient parameter space. Under frozen-backbone PECL, this identification no longer holds. Future adaptation is confined to adapter-induced low-rank effective update directions, whereas sharpness can still be measured by perturbing the full parameter space, the raw LoRA factors, or the effective update itself. Consequently, the standard basin-tolerance intuition is no longer directly applicable: full-space sharpness may measure sensitivity along directions that future tasks cannot update, while raw factor-space sharpness may depend on the parameterization of the same effective update. Flat-LoRA~\cite{li2024flatlora} addresses the representational mismatch between raw LoRA factors and the effective merged update, but it does not study how the corresponding sharpness object governs sequential forgetting under frozen-backbone PECL. Therefore, PECL does not introduce a separate problem from the flatness-forgetting question; rather, it exposes the fact that flatness is support-dependent. The flatness that matters for forgetting is not global loss insensitivity in all parameter directions, but loss insensitivity along the effective update geometry through which future tasks can actually adapt and interfere.


% This limitation becomes more severe in LoRA-based parameter-efficient continual learning (PECL), where the explanations of geometric intuition is no longer directly applicable. The common explanation that a flat old-task basin tolerates larger later updates implicitly assumes that the directions used to measure flatness coincide with the directions along which future updates move the model. However, PECL breaks this assumption. Since the pretrained backbone is frozen, future adaptation is confined to adapter-induced low-rank effective update directions, whereas sharpness can still be measured by perturbing the full parameter space, the raw LoRA factors, or the effective update itself. Flat-LoRA~\cite{li2024flatlora} addresses the distinction between raw LoRA factors and the effective merged update, but it does not answer the continual learning question. Therefore, explaining how flatness affects forgetting in PECL requires more than the standard basin-tolerance intuition: it requires a theory that first relates forgetting to sequential generalization, and then identifies the perturbation geometry on which the corresponding sharpness term should be defined. This raises the PECL-specific question: under frozen-backbone adaptation, which perturbation support defines the sharpness term that is relevant to continual generalization and old-task retention?

% This limitation becomes more severe in LoRA-based parameter-efficient continual learning (PECL), where the geometric intuition above is no longer directly applicable. The common explanation that a flatter old-task basin tolerates larger later updates implicitly assumes that the directions used to measure flatness coincide with the directions along which future updates move the model. PECL breaks this assumption. Since the pretrained backbone is frozen, future adaptation is confined to adapter-induced low-rank update directions, whereas sharpness can still be measured by perturbing the full parameter space, the raw LoRA factors, or the effective update itself. Flat-LoRA~\cite{li2024flatlora} addresses the distinction between raw LoRA factors and the effective merged update, but it does not answer the continual learning question. Therefore, explaining how flatness affects forgetting in PECL requires more than the standard basin-tolerance intuition: it requires a theory that first relates forgetting to sequential generalization, and then identifies the perturbation geometry on which the corresponding sharpness term should be defined.






% Thus, the question is no longer only why flatness affects forgetting, but also which perturbation geometry makes this explanation meaningful.


% This missing link becomes even more consequential in LoRA-based parameter-efficient continual learning (PECL), because the question is no longer only why flatness affects forgetting, but also where flatness should be defined. In PECL, flatness is no longer a property of an unconstrained optimization geometry. In full-parameter training, perturbation directions, optimization directions, and interference directions are usually identified with the same ambient parameter space. PECL breaks this identification. Since the pretrained backbone is frozen, future adaptation is confined to adapter-induced low-rank effective update directions, and any later interference with previous tasks must occur through this admissible update geometry. Consequently, full-space sharpness may measure sensitivity along directions that future tasks can never update, while raw factor-space sharpness may depend on an arbitrary LoRA parameterization rather than on the effective model change, since the same update $\Delta W=BA$ can be represented as $(BR)(R^{-1}A)$ for any invertible matrix $R$. Although 

% Flat-LoRA~\cite{li2024flatlora} addresses the distinction between raw LoRA factors and the effective merged update, it does not answer the continual learning question: under frozen-backbone PECL, should flatness be defined over the full model, over the adapter factors, or over the admissible effective update subspace that carries all future adaptation? This is a theoretical issue, not merely an implementation choice. Until the relevant perturbation support is identified, flatness-based PECL lacks a principled criterion for deciding what should be perturbed, which curvature should be controlled, and how sharpness relates to forgetting under constrained updates.

% PAC-Bayesian analysis is the natural tool for this problem because SAM-style objectives are motivated by neighborhood-smoothed generalization bounds. However, existing PAC-Bayesian continual learning results do not directly yield the theory needed here. The previous lifelong PAC-Bayes bound~\cite{pentinaPACBayesianBoundLifelong2014} treats tasks as i.i.d. draws from a fixed hyperdistribution, producing order-invariant bounds that do not capture trajectory-dependent interference. Recent PAC-Bayesian analyses for continual learning move beyond this assumption by deriving cumulative risk bounds~\cite{friedman2026pac}, but they do not explicitly cast forgetting as past-task population-risk degradation under later posteriors, nor do they specialize the sharpness-smoothed empirical term and KL complexity to the LoRA-admissible effective update subspace required by PECL.
% Recent PAC-Bayesian analyses partially relax this assumption by allowing risk to accumulate over time\cite{friedman2026pac}, but they do not formulate forgetting as past-task generalization degradation under later posteriors, nor do they identify the perturbation support of the sharpness-smoothed empirical term under PECL. 
% The setting considered here requires a different specialization: the prior at task $t$ is history-dependent, the bound must hold pathwise along the realized task trajectory, and both the KL complexity and the SAM-style smoothed empirical term must respect the LoRA-permissible effective update subspace. These requirements define the theoretical problem addressed in this paper.


% PAC-Bayesian analysis is the natural tool for this problem because SAM-style objectives are motivated by neighborhood-smoothed generalization bounds. However, existing PAC-Bayesian results do not directly yield the theory needed here. Lifelong PAC-Bayes learns a hyperposterior over priors for future tasks sampled from a fixed task environment~\cite{pentinaPACBayesianBoundLifelong2014}, which gives a transfer-oriented and order-invariant view of task generalization rather than a trajectory-dependent account of forgetting. Recent PAC-Bayesian analyses for continual learning derive cumulative-risk bounds using data-dependent predictive priors and posteriors~\cite{friedman2026pac}, but they remain at the level of task-indexed posterior transitions. They do not separate the inherited inductive bias from task-specific posterior adaptation through a hierarchical sequential process, nor do they capture the PECL-specific relation between the frozen learned component and the currently trainable LoRA update subspace. As a result, they do not explain forgetting as past-task population-risk degradation under later posteriors, and they do not identify how the SAM-style sharpness-smoothed empirical term and KL complexity should reduce to the LoRA-admissible effective update geometry. This creates three requirements: the prior process must be history-dependent, the bound must hold pathwise along the realized task trajectory, and both the complexity and sharpness terms must respect the update geometry.
% PAC-Bayesian analysis is the natural tool for this problem because SAM-style objectives are motivated by neighborhood-smoothed generalization bounds. However, existing PAC-Bayesian results do not directly provide a theory of how flatness, generalization, and forgetting interact in continual learning. 
% Prior lifelong PAC-Bayes bound~\cite{pentinaPACBayesianBoundLifelong2014}, assumes that tasks are sampled i.i.d. from a fixed environment distribution. As a result, its bound is invariant to task order and does not model trajectory-dependent interference, making it closer to a multi-task formulation than to sequential continual learning (Fig 2). 
% % Lifelong PAC-Bayes learns a hyperposterior over priors for future tasks sampled from a fixed task environment~\cite{pentinaPACBayesianBoundLifelong2014}, giving a transfer-oriented and order-invariant view of task generalization rather than a trajectory-dependent account of how later updates degrade past-task population risk. 
% Recent PAC-Bayesian analyses for continual learning derive cumulative-risk bounds using data-dependent predictive priors and posteriors~\cite{friedman2026pac}, but their objective remains cumulative performance along the task sequence, not the role of neighborhood-smoothed risk in explaining old-task retention under later posteriors. Therefore, existing PAC-Bayesian CL theory does not yet connect the SAM-style flatness term to both current-task generalization and past-task forgetting. PECL further sharpens this gap: the theory must additionally account for the structural decomposition between the frozen learned component and the trainable LoRA update component, and determine whether the perturbations, sharpness-smoothed empirical term, and KL complexity should live in the full parameter space, the raw adapter-factor space, or the LoRA-admissible effective update space. The setting considered here therefore requires a history-dependent prior process, a pathwise bound along the realized task trajectory, and a support-reduction argument showing how the complexity and sharpness terms respect the admissible PECL update geometry.

% PAC-Bayesian analysis provides a natural lens for these questions, but existing continual-learning PAC-Bayes theories are not designed to address them. Prior lifelong PAC-Bayes frameworks typically consider tasks sampled i.i.d. from a fixed hyperdistribution~\cite{pentinaPACBayesianBoundLifelong2014}. It do not directly capture the history-dependent prior process induced by sequential adaptation, nor the pathwise accumulation of generalization error along an observed task sequence.

% .

% First, the prior at each step is no longer fixed or independently chosen, but is produced by the previous learning history. A suitable PAC-Bayesian analysis must therefore handle a time-varying posterior and prior process, and the resulting guarantee must hold along the realized trajectory rather than only after averaging over independently sampled tasks. Second, in LoRA-based PECL, the theoretical object is further constrained by the frozen-backbone parameterization. The posterior is placed on trainable adapter parameters, while the effective model update lies in a lower-dimensional, task-permissible update geometry. Existing PAC-Bayesian continual-learning bounds do not identify how the smoothed empirical term and the KL term should transform under this constrained parameterization, leaving the bound-relevant sharpness object ambiguous.

% PAC-Bayesian analysis provides a natural lens for these questions, but existing continual-learning PAC-Bayes theories do not directly address them. Lifelong PAC-Bayes typically studies transfer across tasks sampled from a fixed hyperdistribution~\cite{pentinaPACBayesianBoundLifelong2014}, whereas continual learning follows a realized, order-dependent trajectory in which later updates may degrade earlier-task performance. This setting requires a pathwise analysis with history-dependent priors, rather than an order-invariant bound averaged over independently sampled tasks. It also requires a PECL-specific support analysis: under frozen-backbone LoRA adaptation, the posterior over trainable adapter parameters induces a posterior over effective adapter updates, while the backbone remains fixed. Hence, the perturbation-smoothed empirical risk and KL complexity should be evaluated on the admissible effective adapter-update geometry, not on full-model perturbations or non-identifiable raw LoRA factors.



% We



% PAC-Bayesian analysis provides a natural theoretical lens for the above questions. 
% However, existing PAC-Bayesian theories of continual learning are not designed to answer the two questions raised here. 
% Lifelong PAC-Bayes studies transfer across tasks sampled from a fixed hyperdistribution~\cite{pentinaPACBayesianBoundLifelong2014}, but its order-invariant formulation does not model how later updates along a realized task trajectory degrade the performance on earlier tasks. 
% Recent PAC-Bayesian bound derive cumulative-risk bounds with data-dependent priors and posteriors~\cite{friedman2026pac}, but they are not designed to identify the geometric support on which a SAM-style perturbation-smoothed empirical risk should be evaluated. More importantly for LoRA-based PECL, these analyses do not address the mismatch among perturbation and trainable update directions under constrained update geometry.
% PAC-Bayesian analysis provides a natural theoretical lens for the above questions because SAM-style flatness objectives optimize sharpness-based empirical risk. However, existing PAC-Bayesian theories of continual learning are not designed to answer the two questions raised here. Lifelong PAC-Bayes studies transfer across tasks sampled from a fixed hyperdistribution~\cite{pentinaPACBayesianBoundLifelong2014}, but its order-invariant formulation does not model how later updates along a realized task trajectory degrade the population risk of earlier tasks. Recent PAC-Bayesian analyses derive cumulative-risk bounds with data-dependent priors and posteriors~\cite{friedman2026pac}, but they are not designed to isolate forgetting as past-task population-risk degradation under later posteriors, nor to identify the perturbation support on.
% These limitations lead to two technical challenges. 
This gap creates two technical challenges.
The first question requires a sequential PAC-Bayesian analysis in which the prior is not fixed, but evolves with the learning history. 
The resulting bound must hold pathwise along the realized task sequence. 
The second question requires a PECL-specific support analysis. 
Under frozen-backbone LoRA adaptation, , the perturbation-smoothed empirical risk and the KL complexity must be related to the admissible effective adapter-update geometry rather than to  full-model perturbations or non-identifiable raw LoRA factor-space perturbations.

% the posterior over trainable adapter parameters induces a posterior over effective adapter updates, while the backbone remains fixed. 
% Therefore

% We develop the analysis in two steps. First, we establish a sequential hierarchical PAC-Bayesian theorem for the posterior trajectory induced by continual learning. The bound decomposes trajectory risk into perturbation-smoothed empirical risk, within-task adaptation complexity, and hyperposterior drift complexity, thereby relating forgetting to past-task generalization degradation under later posteriors. Second, we specialize this result to frozen-backbone LoRA adaptation. Through the pushforward posterior over effective adapter updates, we prove a PECL support-reduction theorem showing that the bound-relevant sharpness and KL complexity are supported on the admissible, basis-invariant effective adapter-update geometry.

% Our contributions are twofold.
% \begin{itemize}
%     \item We provide a pathwise hierarchical PAC-Bayesian decomposition that connects perturbation-smoothed risk with forgetting along posterior trajectories.
%     \item We prove a PECL support-reduction theorem that identifies effective adapter-update geometry as the relevant flatness support under LoRA-constrained updates.
% \end{itemize}


% These limitations lead to two setting-specific technical challenges. The first question requires a sequential PAC-Bayesian analysis in which the prior is not fixed in advance, but evolves with the learning history. The resulting bound must hold pathwise along the realized task sequence, rather than after averaging over tasks sampled independently from an environment distribution. The second question requires a PECL-specific support analysis. Under frozen-backbone LoRA adaptation, the posterior over trainable adapter parameters induces a posterior over effective model updates, while the backbone remains fixed. Therefore, the sharpness and the KL complexity must be related to the admissible effective-update geometry rather than to arbitrary full-model perturbations or non-identifiable raw factor-space perturbations.



% These limitations lead to two setting-specific technical challenges. The first question requires a sequential PAC-Bayesian analysis in which the prior at task $t$ is not fixed, but produced by the learning history as a process $P_t(S_{1:t-1})$. The resulting bound must hold pathwise along the realized task sequence, rather than after averaging over tasks sampled i.i.d. from an environment distribution. The second question requires a PECL-specific support analysis. Under frozen-backbone LoRA adaptation, the posterior over trainable adapter parameters induces a pushforward posterior over effective model updates, while the backbone remains fixed. Therefore, the perturbation-smoothed empirical term and the KL complexity must be related to the admissible effective-update geometry rather than to arbitrary full-model perturbations or non-identifiable raw factor-space perturbations.

% ----------
% PAC-Bayesian analysis provides a natural theoretical lens for the above questions because SAM-style flatness objectives optimize perturbation-smoothed empirical losses. However, existing PAC-Bayesian theories of continual learning are not designed to address the setting considered here. Lifelong PAC-Bayes studies transfer across tasks sampled from a fixed hyperdistribution~\cite{pentinaPACBayesianBoundLifelong2014}, but its order-invariant formulation does not model how later updates along a realized task trajectory degrade the population risk of earlier tasks. Recent PAC-Bayesian analyses derive cumulative-risk bounds with data-dependent priors and posteriors~\cite{friedman2026pac}, but they are not designed to isolate forgetting as past-task population-risk degradation under later posteriors or to identify the perturbation support on which a SAM-style smoothed empirical term should be evaluated. Thus, the missing component is not merely another PAC-Bayesian bound for continual learning, but a pathwise analysis that connects perturbation-smoothed risk to forgetting and then reduces this analysis to the admissible update geometry induced by PECL.





% PAC-Bayesian analysis provides a natural theoretical lens for the above problems, because SAM-style flatness objectives optimize perturbation-smoothed empirical losses. However, existing PAC-Bayesian theories of continual learning do not answer the two questions raised above. Lifelong PAC-Bayes studies transfer across tasks sampled from a fixed hyperdistribution~\cite{pentinaPACBayesianBoundLifelong2014}, but its order-invariant formulation does not model how later updates degrade the forgetting. 
% % Recent analyses derive cumulative-risk bounds with data-dependent priors and posteriors~\cite{friedman2026pac}. These results are not designed to isolate forgetting as the increase of past-task population risk under later posteriors, nor to identify the perturbation support on which a SAM-style smoothed empirical term should be evaluated.
% Recent analyzes derive cumulative-risk bounds with data-dependent priors and posteriors~\cite{friedman2026pac}, but they do not identify how the SAM-style empirical term decomposes past-task forgetting under later posteriors. 
% % Thus, existing PAC-Bayesian CL theory neither resolves the role of flatness in sequential forgetting nor specializes this question to PECL.
% Thus, the missing component is not merely another PAC-Bayesian bound for continual learning, but a pathwise analysis that connects flatness based risk to forgetting and then reduces this analysis to the admissible update geometry induced by PECL.


% This leaves three technical obstacles. First, the prior at task $t$ is no longer fixed but produced by the learning history, giving a history-dependent prior process $P_t(S_{1:t-1})$. Second, the bound must hold pathwise along the realized task trajectory, rather than after averaging over i.i.d. tasks from an environment distribution. Third, under frozen-backbone PECL, the posterior over LoRA parameters induces a pushforward posterior over effective model updates, so the smoothed empirical term and KL complexity must be reduced to the admissible effective-update geometry. We address the first two obstacles using a sequential change-of-measure argument and a martingale-based tail bound, and address the third through a support-reduction argument for the LoRA-induced effective-update posterior.

% In continual learning, the prior used at step $t$ is produced by a history-dependent. This creates three challenges:
% 1. the prior is no longer a fixed object, but a sequential process $P_t$ adapted to the task history
% 2. the generalization analysis must accumulate pathwise over the task sequence
% 3. under PECL, the analysis must further specialize to the task-permissible update subspace $W_{\Delta, t}$.

% To address the history-dependent prior process and the pathwise accumulation over the task sequence, we model continual learning through a timevarying hyperposterior process $P_t$ and analyze it with a sequential change-of-measure argument together with a martingale-based tail bound. This yields Theorem which related flatenss and forgetting. To address the PECL structural constraint, Lemmas  establish the exact PECL reduction of both the KL term and the sharpness-smoothed empirical term to the task-permissible update subspace $W_{\Delta, t}$, which yields Theorem 4.4.


% , where perturbations, sharpness-smoothed empirical risk, and KL complexity may live in the full parameter space or the admissible effective update space.
% PAC-Bayesian analysis is a natural theoretical lens for above problems, since SAM-style flatness objectives are tied to neighborhood-smoothed generalization bounds. 
% PAC-Bayesian analysis provides a natural theoretical lens for the above problems, because SAM-style flatness objectives optimize perturbation-smoothed empirical losses.
% However, existing PAC-Bayesian theories of continual learning do not answer the two questions raised above. Lifelong PAC-Bayes studies transfer across tasks sampled from a fixed hyperdistribution~\cite{pentinaPACBayesianBoundLifelong2014}, but its order-invariant formulation does not model how later updates degrade the population risk of earlier tasks. Recent analyses derive cumulative-risk bounds with data-dependent priors and posteriors~\cite{friedman2026pac}, but they do not identify how the SAM-style smoothed empirical term decomposes past-task forgetting under later posteriors. Nither They do not address the PECL-specific support question: whether the perturbations, sharpness-smoothed empirical risk, and KL complexity should be defined in the full parameter space, the raw LoRA factor space, or the admissible effective update space. 
% % Existing LoRA flatness studies further clarify the mismatch between factor-space and merged-update perturbations, but they remain focused on single-task adaptation and do not explain sequential forgetting under constrained updates. 
% Therefore, the missing theory is not simply another PAC-Bayesian bound for continual learning. What is needed is a pathwise analysis that connects flatness-based smoothing to past-task population-risk degradation and then specializes this connection to the admissible update geometry induced by frozen-backbone LoRA adaptation.

% PAC-Bayesian analysis is a natural theoretical lens for this problem, since SAM-style objectives are motivated by neighborhood-smoothed generalization bounds. However, existing PAC-Bayesian analyses do not resolve the two issues central to this work: they do not connect the SAM-style flatness term to both current-task generalization and past-task forgetting, and they do not identify the perturbation support relevant to LoRA-constrained PECL.
% Prior lifelong PAC-Bayes bounds~\cite{pentinaPACBayesianBoundLifelong2014}, assume that tasks are sampled i.i.d. from a fixed environment distribution. Consequently, their bounds are invariant to task order and do not model trajectory-dependent interference, which makes them closer to multi-task transfer than to sequential continual learning. 
% % They therefore cannot characterize forgetting as the degradation of past-task population risk caused by later updates.
% A concurrent PAC-Bayesian analysis for sequential continual learning~\cite{friedman2026pac} derives a cumulative-loss bound using data-dependent predictive priors and posteriors. Our contribution differs in two respects. First, we use a sequential hierarchical PAC-Bayesian formulation to expose how past-task forgetting can be decomposed through neighborhood-smoothed empirical risk and complexity terms under later posteriors. Second, we specialize the analysis to PECL, where the frozen backbone and trainable LoRA updates induce a restricted effective update geometry. Under this geometry, both the KL complexity and the sharpness-smoothed empirical term are determined by the LoRA-admissible effective update subspace rather than by arbitrary full-space perturbations.

% Thus, the missing theory is not merely a PAC-Bayesian bound for sequential prediction. What is needed is a pathwise analysis that connects flatness-based smoothing to past-task degradation and then identifies the support on which this smoothing is meaningful under PECL. This is the gap addressed in this work.

% PAC-Bayesian analysis is the natural tool for this problem because SAM-style objectives are motivated by neighborhood-smoothed generalization bounds. However, existing PAC-Bayesian results do not directly yield the theory needed here. Lifelong PAC-Bayes learns a hyperposterior over priors for future tasks sampled from a fixed task environment~\cite{pentinaPACBayesianBoundLifelong2014}, which gives a transfer-oriented and order-invariant view of task generalization rather than a trajectory-dependent account of forgetting. Recent PAC-Bayesian analyses for CL derive cumulative-risk bounds using data-dependent predictive priors and posteriors~\cite{friedman2026pac}, but they remain at the level of task-indexed posterior transitions. They do not separate the inherited inductive bias from task-specific posterior adaptation through a hierarchical sequential process, nor do they capture the PECL-specific relation between the frozen learned component and the currently trainable LoRA update subspace. As a result, they do not explain forgetting as past-task population-risk degradation under later posteriors, and they do not identify how the SAM-style sharpness-smoothed empirical term and KL complexity should reduce to the LoRA-admissible effective update geometry.

% PAC-Bayesian analysis is the natural tool for this problem because SAM-style objectives are motivated by neighborhood-smoothed generalization bounds. However, existing PAC-Bayesian continual learning results do not directly yield the theory needed here. Lifelong PAC-Bayes usually treats tasks as i.i.d. draws from a fixed environment or hyperdistribution, producing order-invariant bounds that do not capture trajectory-dependent interference~\cite{pentinaPACBayesianBoundLifelong2014,JMLR:v24:22-1254}. Rencent paper relax part of this assumption by allowing risk to accumulate over time, but they do not identify the perturbation support of the sharpness term  and the architrechor desing of  PECL. The setting considered here requires a different specialization: the prior at task $t$ is history-dependent, the bound must hold pathwise along the task trajectory, and both the KL complexity and the SAM-style smoothed empirical term must respect the LoRA-permissible update subspace. These requirements define the theoretical problem addressed in this paper.

% Existing PAC-Bayesian theory provides important tools, but it does not resolve the flatness-forgetting question under PECL. Classical PAC-Bayes relates empirical risk, posterior perturbations, and generalization complexity~\cite{mcallesterPACBayesianModelAveraging1999,mcallesterPACBayesianStochasticModel2003,germainPACBayesianTheoryMeets2016,lotfiPACBayesCompressionBounds2022}. Extensions to online, lifelong, and meta-learning settings allow data-dependent or task-dependent priors~\cite{pentinaPACBayesianBoundLifelong2014,NEURIPS2022_onlineBayes,NEURIPS2024_xinrui,JMLR:v24:22-1254}, but their objectives differ from ours. Lifelong PAC-Bayes typically models tasks as draws from a fixed environment or hyperdistribution, leading to order-invariant guarantees that do not capture trajectory-dependent interference. Online PAC-Bayes allows priors to evolve over time, but is usually formulated over example-level sequences rather than task-level posterior inheritance. More recent sequential CL analyses bound cumulative risk, but they do not identify forgetting as past-task generalization degradation induced by later posteriors, nor do they connect this degradation to a SAM-style smoothed-risk term. At the same time, SAM theory explains flatness-based generalization mainly in single-task or full-parameter regimes, while LoRA flatness theory resolves the distinction between raw factor perturbations and effective update perturbations without studying continual forgetting. Consequently, no existing theory specifies the PECL-relevant sharpness object that simultaneously connects flatness, sequential generalization, and old-task retention.


% We resolve this issue by treating flatness in continual learning as both a sequential generalization problem and a constrained-geometry problem. First, we show that forgetting can be formulated as past-task generalization degradation under later task posteriors. This gives a sequential PAC-Bayesian risk decomposition in which neighborhood-smoothed empirical risk, the quantity underlying SAM-style flatness objectives, controls not only current-task generalization but also old-task retention. Second, we specialize the analysis to frozen-backbone LoRA-based PECL and prove a support-reduction result: the perturbations, sharpness-smoothed empirical term, and KL complexity that enter the bound are supported on the admissible effective adapter-update subspace $\mathcal W_{\Delta,t}$. Consequently, the theoretically relevant notion of flatness in PECL is not ambient full-space sharpness, nor raw LoRA factor-space sharpness, but basis-invariant sharpness of the effective adapter-update geometry. This provides a principled explanation for why adapter-scoped sharpness-aware optimization can improve both continual generalization and forgetting, and it yields testable predictions that we validate across LoRA variants, perturbation scopes, datasets, and architectures.

% We resolve these issues by treating flatness in continual learning as both a sequential generalization problem and a constrained-geometry problem. First, we formulate forgetting as past-task generalization degradation under later task posteriors. This gives a sequential PAC-Bayesian risk decomposition in which neighborhood-smoothed empirical risk, the quantity underlying SAM-style flatness objectives, appears in the bound for both current-task generalization and old-task retention. Second, we specialize the analysis to frozen-backbone LoRA-based PECL and prove a support-reduction result: the perturbations, sharpness-smoothed empirical term, and KL complexity that enter the bound are supported on the admissible effective adapter-update subspace $\mathcal W_{\Delta,t}$. Consequently, the theoretically relevant notion of flatness in PECL is not ambient full-space sharpness, nor raw LoRA factor-space sharpness, but basis-invariant sharpness of the effective adapter-update geometry. This helps explain the empirical effectiveness of adapter-scoped sharpness-aware optimization for improving continual generalization and reducing forgetting, and it yields testable predictions that we validate across LoRA variants, perturbation scopes, datasets, and architectures.



% \begin{figure}[tbp]
%     \centering
%     \begin{minipage}[t]{0.49\linewidth}
%         \centering
% \includegraphics[width=\linewidth,height=3.4cm,keepaspectratio]{figures_TRML/losslan_flat.pdf}
% % {figures_TRML/lossland.pdf}
%         \vspace{2pt}
%         {\footnotesize (a)}
%     \end{minipage}
%     \hfill
%     \begin{minipage}[t]{0.49\linewidth}
%         \centering
% \includegraphics[width=\linewidth,height=3.4cm,keepaspectratio]{figures_TRML/liuxingshiyitu5.pdf}
%         \vspace{2pt}
%         {\footnotesize (b)}
%     \end{minipage}
%     \vspace{-2pt}
%     \caption{\textbf{Flatness and admissible perturbation geometry in PECL.} 
%     (a) Under the same amount of weight update, the loss change of the flat loss landscape (right) is smaller than a sharp loss landscape (left) .
%     (b) PECL updates are confined to task-permissible directions (green), while standard sharpness perturbs the full parameter space (blue).}
%     \label{fig:flatness_and_scope}
%     \vspace{-8pt}
% \end{figure}

We develop the analysis in two steps, each corresponding to one of the questions above. First, to connect sharpness-based generalization with forgetting, we develop a sequential hierarchical PAC-Bayesian analysis of the posterior trajectory in CL. 
% Using a sequential change-of-measure argument and a martingale-based tail bound, 
Our first main theorem gives a pathwise decomposition of trajectory risk into perturbation-smoothed empirical risk, within-task adaptation complexity, and hyperposterior drift complexity. This links sharpness and CL performance through a shared perturbation-smoothed risk view.
% First, to connect flatness-based generalization with forgetting, we use a sequential change-of-measure argument and a martingale-based tail bound to handle history-dependent priors and pathwise accumulation along the realized task trajectory. 
% This yields our first main theorem: a sequential PAC-Bayesian risk decomposition showing that perturbation-smoothed empirical risk bounds generalization on past tasks under later posteriors. 
% First, to connect flatness-based generalization with forgetting, we use a sequential change-of-measure argument and a martingale-based tail bound to handle history-dependent priors and pathwise accumulation along the realized task trajectory. 
% This gives a sequential PAC-Bayesian risk decomposition in which perturbation-smoothed empirical risk bounds generalization on past tasks under later posteriors. 
% As a result, forgetting can be decomposed into old-task perturbation-smoothed empirical-risk degradation and a hierarchical KL divergence term reflecting posterior complexity. 
% This result explains why the same perturbation-smoothed empirical risk that motivates SAM-style generalization objectives is also relevant to forgetting in continual learning. 
Second, we specialize the sequential PAC-Bayesian decomposition to frozen-backbone LoRA adaptation. 
% The key structural constraint is that the posterior is parameterized in the trainable LoRA factors, while the induced model perturbations act through effective adapter updates and the pretrained backbone remains fixed. 
We analyze the pushforward posterior induced by the LoRA parameterization on effective adapter updates. 
This yields our second main theorem, showing that both the perturbation-smoothed empirical risk and the KL complexity in the bound reduce to quantities supported on the admissible adapter-update geometry. 
This result answers the second question of the paper: under LoRA-constrained updates, the bound-relevant perturbation support is the basis-invariant effective adapter-update geometry. 
% Consequently, perturbations outside this geometry are theoretically misaligned with PECL, while adapter-scoped perturbations target the sharpness object that is relevant to sequential generalization and forgetting. 
% This contribution is analytical rather than algorithmic: it clarifies which perturbation geometry is meaningful under subspace-constrained PECL and why flatness-aware methods can remain effective even when trainable parameters occupy only a small adapter-induced subspace.
% Second, to identify the PECL-specific support of this perturbation-smoothed risk term, we specialize the decomposition PECL. 
% We analyze the pushforward posterior induced by LoRA parameters on effective adapter updates and prove a support-reduction result: the perturbation-smoothed empirical risk, and the KL complexity are determined by the admissible effective adapter-update geometry. 
% Consequently, for explaining sequential generalization and forgetting under LoRA-constrained updates, the relevant flatness object is not full-parameter-space sharpness or raw LoRA factor-space sharpness, but basis-invariant sharpness of the effective adapter-update geometry.




% We develop the analysis in two steps, each corresponding to one of the questions above. First, to connect flatness-based generalization with forgetting, we use a sequential change-of-measure argument and a martingale-based tail bound to handle history-dependent priors and pathwise accumulation along the realized task trajectory. This gives a sequential PAC-Bayesian risk decomposition in which sharpness empirical risk bounds past-task population risk under later posteriors. This result explains why the same flatness empirical risk that motivates SAM-style generalization objectives is also relevant to old-task forgetting in continual learning. Second, to identify the PECL-specific support of this sharpness term, we specialize the decomposition to frozen-backbone LoRA adaptation. We analyze the pushforward posterior induced by LoRA parameters on effective model updates and prove a support-reduction result: the PAC-Bayesian perturbations, the sharpness, and the KL complexity in the effective model space are determined by the admissible adapter-update geometry. Consequently, for explaining sequential generalization and past-task degradation under LoRA-constrained updates, the relevant flatness object is not full-space sharpness or raw factor-space sharpness, but basis-invariant sharpness of the effective adapter-update geometry.

% We resolve these issues by treating flatness in continual learning as both a sequential generalization problem and a constrained-geometry problem. First, we formulate forgetting as past-task generalization degradation under later task posteriors.
% This gives a sequential PAC-Bayesian risk decomposition in which neighborhood-smoothed empirical risk bounds past-task population risk under later posteriors, thereby expressing forgetting as a combination of old-task empirical degradation and PAC-Bayesian complexity terms.
% % This gives a sequential PAC-Bayesian risk decomposition in which neighborhood-smoothed empirical risk, the quantity underlying SAM-style flatness objectives, appears in the bound for both current-task generalization and old-task retention. 
% Second, we specialize the analysis to frozen-backbone LoRA-based PECL and prove a support-reduction result for the pushforward posterior induced by LoRA parameters: the perturbations, sharpness-smoothed empirical term, and KL complexity in the effective model space are determined by the admissible adapter-update geometry.
% % prove a support-reduction result: the perturbations, sharpness-smoothed empirical term, and KL complexity that enter the bound are supported on the admissible effective adapter-update subspace $\mathcal W_{\Delta,t}$. 
% Consequently, for explaining sequential generalization and past-task degradation under LoRA-constrained updates, the theoretically relevant notion of flatness is not full-space sharpness or raw factor-space sharpness, but basis-invariant sharpness of the effective adapter-update geometry.
% Consequently, the theoretically relevant notion of flatness in PECL is not full-space sharpness, nor raw LoRA factor-space sharpness, but basis-invariant sharpness of the effective adapter-update geometry. This helps explain the empirical effectiveness of adapter-scoped sharpness-aware optimization for improving continual generalization and reducing forgetting, and it yields testable predictions that we validate across LoRA variants, perturbation scopes, datasets, and architectures.

% Our contributions are threefold.
% \begin{itemize}
% \item We formulate forgetting as generalization degradation on past tasks under later posteriors and derive a sequential hierarchical PAC-Bayesian decomposition of the posterior trajectory. The resulting bound separates perturbation-smoothed empirical risk, within-task adaptation complexity, and hyperposterior drift complexity, making cross-task interference explicit.
%     % \item We formulate forgetting as past-task generalization degradation under later posteriors and derive a sequential PAC-Bayesian decomposition that links SAM-style neighborhood-smoothed empirical risk to both current-task generalization and old-task retention.
%     \item We identify a PECL-specific geometric mismatch among perturbation directions, trainable update directions, and interference directions, showing why full-space or raw factor-space sharpness is not automatically the relevant object under frozen-backbone LoRA adaptation. We prove a support-reduction result showing that the PAC-Bayesian perturbations, sharpness-smoothed empirical term, and KL complexity reduce to the admissible effective adapter-update subspace, and validate the resulting predictions across LoRA variants, perturbation scopes, datasets, and architectures.
% \end{itemize}

% Second, we specialize this result to frozen-backbone LoRA adaptation. Through the pushforward posterior over effective adapter updates, we prove a PECL support-reduction theorem showing that the bound-relevant sharpness and KL complexity are supported on the admissible, basis-invariant effective adapter-update geometry.

Our contributions are twofold.
\begin{itemize}
    \item We provide a pathwise hierarchical PAC-Bayesian decomposition that connects perturbation-smoothed risk with forgetting along posterior trajectories.
    \item We prove a PECL support-reduction theorem that identifies effective adapter-update geometry as the relevant flatness support under LoRA-constrained updates.
    
\end{itemize}



% Our contributions are twofold.
% \begin{itemize}
%     \item We derive a pathwise hierarchical PAC-Bayesian decomposition for continual learning. 
%     The bound separates perturbation-smoothed empirical risk, within-task adaptation complexity, and hyperposterior drift complexity, thereby linking flatness-based generalization and forgetting through a shared posterior-trajectory view.
% \vspace{-2pt}
%     \item We prove a PECL support-reduction theorem for frozen-backbone LoRA adaptation. 
%     By analyzing the pushforward posterior induced by LoRA parameters on effective adapter updates, we show that the bound-relevant perturbation-smoothed empirical risk and KL complexity are determined by the admissible effective adapter-update geometry. 
%     This identifies basis-invariant effective adapter-update sharpness as the relevant flatness object under LoRA-constrained updates.
% % \vspace{-2pt}
% %     \item We empirically validate the theory-derived predictions through perturbation-support comparisons, adapter-subspace sharpness diagnostics, forgetting analysis, and cross-architecture experiments across LoRA-based continual learning settings.
% \end{itemize}

\section{Related Work}

\minisection{LoRA-Based CL.}
% PECL has emerged as a practical solution for incrementally adapting large pre-trained models
% % , with LoRA becoming a representative approach due to its efficiency and compatibility with frozen backbones
% ~\cite{pmlr-v97-houlsby19a,Wang_2022_CVPR,dingParameterefficientFinetuningLargescale2023,xu2023parameterefficientfinetuningmethodspretrained}. 
By introducing trainable low-rank matrices, LoRA enables continual adaptation with minimal overhead, and has demonstrated strong performance~\cite{hu2022lora}.
However, LoRA-based methods remain vulnerable to forgetting and task interference
% , particularly in long and diverse task sequence
s~\cite{Gao_2023_ICCV,10.1145/3730436.3730456,yang2024parametercollisionhinderingcontinual}. 
 % Gated variants further select or combine adapters with task-dependent weights~\citep{Gao_2023_ICCV,10.1145/3730436.3730456,zhao-etal-2024-sapt,cheng2025exploiting,wu2025sdlorascalabledecoupledlowrank}.
Recent variants such as SDLoRA~\cite{wu2025sdlorascalabledecoupledlowrank} and OLoRA~\cite{NEURIPS2020_70d85f35,wang2023orthogonal} aim to alleviate these issues via adapter modularity and orthogonal constraints. 
Although such designs yield empirical improvements, they lack a principled understanding of how constrained update geometry affects performance in CL .
% In particular, the interaction between adapter subspace design and solution stability remains insufficiently explored.

\minisection{Flat Minima.}
Prior work has explored flat minima in CL through sharpness-aware optimization~\cite{pmlr-v151-bisla22a,Zhang_GAM,NEURIPS2024_C-Flat,yangRevisitingFlatnessAwareOptimization2025,chen2026sharpnessflatnessdecompositionframework}, mode connectivity~\cite{izmailov2019averagingweightsleadswider,mirzadeh2020linearmodeconnectivitymultitask}, and improved representation learning~\cite{hu2022doesselfsupervisedpretrainingperform,ramasesh2022effect}.
These methods proposed for full model training do not account for continual training dynamics in the low-rank subspace~\cite{yangRevisitingFlatnessAwareOptimization2025}. Whether flatness in the constrained update geometry can ensure generalization in CL remains unclear, motivating further investigation.


\minisection{PAC-Bayesian Theory.}
The PAC-Bayesian framework provides a principled foundation for analyzing generalization by balancing empirical risk and posterior complexity~\cite{mcallesterPACBayesianModelAveraging1999,mcallesterPACBayesianStochasticModel2003,germainPACBayesianTheoryMeets2016,lotfiPACBayesCompressionBounds2022}.
Online PAC-Bayes~\citep{NEURIPS2022_onlineBayes,NEURIPS2024_xinrui} allows time-varying, data-dependent priors, but it is formulated at the level of individual examples rather than tasks and does not explicitly model task-level inheritance in CL.
Hierarchical PAC-Bayes bounds for meta learning~\citep{pentinaPACBayesianBoundLifelong2014,JMLR:v24:22-1254} introduce hyperposteriors over task priors, but they assume a single stationary hyperposterior and do not model the sequential, history-dependent evolution of the inductive bias that underlies forgetting in CL. A concurrent PAC-Bayesian for CL analysis derives cumulative risk bounds with data-dependent priors and posteriors~\cite{friedman2026pac}, without addressing flatness-based generalization or the constrained update geometry of PECL.
 % Concurrent PAC-Bayes  bound ~\cite{friedman2026pac} derives a cumulative-loss bound, but not However, it does not distinguish the hyperposterior process from task-specific priors, nor does it specialize the analysis to parameter-efficient continual learning.
% In continual learning, the inductive bias must adapt to the observed task history, and this non-stationarity is particularly pronounced in LoRA-based PECL, where the update scope and the initialization of incremnental adapters could change with training.







\section{Preliminaries}


% \paragraph{Risks and perturbation-induced local posteriors.}
% Let $\mathcal X$ denote the input space, $\mathcal Y$ the label space, and $\mathcal Z=\mathcal X\times\mathcal Y$ the data space.
% A model parameterized by $W\in\mathbb R^d$ induces a predictor $f_W$.
% We consider a bounded loss $\ell(W,z)\in[0,1]$.
% For a data distribution $\mathcal D$ and a sample $S=\{z_j\}_{j=1}^{m}$, the population and empirical risks are
% $$
% L_{\mathcal D}(W)
% =
% \mathbb E_{z\sim\mathcal D}\ell(W,z),
% \qquad
% \widehat L_S(W)
% =
% \frac{1}{m}\sum_{j=1}^{m}\ell(W,z_j).
% $$



% PAC-Bayesian analysis evaluates risk under a posterior distribution over parameters.






% In flatness-based analysis, such a posterior can be locally induced by perturbing a deterministic solution $\widehat W$.
% Given a perturbation law $\epsilon\sim\Pi$, the corresponding perturbation-induced posterior is the distribution of
% $$
% W=\widehat W+\epsilon .
% $$
% Equivalently, if $Q_{\widehat W}^{\Pi}$ denotes this local posterior, then
% $$
% Q_{\widehat W}^{\Pi}
% =
% \mathcal L(\widehat W+\epsilon),
% \qquad
% \epsilon\sim\Pi ,
% $$
% where $\mathcal L(\cdot)$ denotes the induced law.
% The posterior-averaged empirical risk is therefore the perturbation-smoothed empirical risk
% $$
% \widehat L_S(Q_{\widehat W}^{\Pi})
% =
% \mathbb E_{W\sim Q_{\widehat W}^{\Pi}}
% \widehat L_S(W)
% =
% \mathbb E_{\epsilon\sim\Pi}
% \widehat L_S(\widehat W+\epsilon).
% $$







% PAC-Bayesian analysis evaluates risks under a posterior distribution over parameters.
% A perturbation law $\Pi(\cdot\mid W)$ provides a concrete way to construct such a local posterior by sampling
% $$
% W\sim Q,
% \qquad
% \epsilon\sim\Pi(\cdot\mid W),
% \qquad
% \widetilde W=W+\epsilon .
% $$
% The corresponding perturbation-smoothed empirical risk is
% $$
% \widehat L_S^\Pi(Q)
% =
% \mathbb E_{W\sim Q}
% \mathbb E_{\epsilon\sim\Pi(\cdot\mid W)}
% \widehat L_S(W+\epsilon).
% $$
% When $Q$ is concentrated at a deterministic solution $\widehat W$, this reduces to the empirical risk averaged under a local perturbation-induced parameter distribution:
% $$
% \widehat L_S^\Pi(\widehat W)
% =
% \mathbb E_{\epsilon\sim\Pi(\cdot\mid \widehat W)}
% \widehat L_S(\widehat W+\epsilon).
% $$
% The corresponding distributional sharpness is the excess empirical risk under this local posterior:
% $$
% \mathrm{ES}_S^\Pi(\widehat W)
% =
% \mathbb E_{\epsilon\sim\Pi(\cdot\mid \widehat W)}
% \left[
% \widehat L_S(\widehat W+\epsilon)
% -
% \widehat L_S(\widehat W)
% \right].
% $$

% SAM uses an adversarial version of this local posterior view, replacing the expected perturbation loss by the worst-case loss within a radius $\rho$:
% $$
% \widehat L_S^{\mathrm{SAM}}(\widehat W,\rho)
% =
% \max_{\|\epsilon\|\le \rho}
% \widehat L_S(\widehat W+\epsilon)
% =
% \widehat L_S(\widehat W)
% +
% \mathrm{AS}^{(0)}_S(\widehat W,\rho),
% $$
% where
% $$
% \mathrm{AS}^{(0)}_S(\widehat W,\rho)
% =
% \max_{\|\epsilon\|\le\rho}
% \left[
% \widehat L_S(\widehat W+\epsilon)
% -
% \widehat L_S(\widehat W)
% \right].
% $$






% Let $\mathcal X$ be the input space, $\mathcal Y$ the label space, and $\mathcal Z=\mathcal X\times\mathcal Y$ the data space. A model parameterized by $W\in\mathbb R^d$ induces a predictor $f_W$, and the loss $\ell(W,z)$ is bounded in $[0,1]$. For a distribution $\mathcal D$ and a sample $S=\{z_j\}_{j=1}^{m}$, define
% $$
% L_{\mathcal D}(W)
% =
% \mathbb E_{z\sim\mathcal D}\ell(W,z),
% \qquad
% \widehat L_S(W)
% =
% \frac{1}{m}\sum_{j=1}^{m}\ell(W,z_j).
% $$

% PAC-Bayesian analysis evaluates risks under a posterior distribution over parameters.
% A perturbation law $\Pi(\cdot\mid W)$ provides a concrete way to construct such a local posterior by sampling
% $$
% W\sim Q,
% \qquad
% \epsilon\sim\Pi(\cdot\mid W),
% \qquad
% \widetilde W=W+\epsilon .
% $$
% The corresponding perturbation-smoothed empirical risk is
% $$
% \widehat L_S^\Pi(Q)
% =
% \mathbb E_{W\sim Q}
% \mathbb E_{\epsilon\sim\Pi(\cdot\mid W)}
% \widehat L_S(W+\epsilon).
% $$
% When $Q$ is concentrated at a deterministic solution $\widehat W$, this reduces to the empirical risk averaged under a local perturbation-induced parameter distribution:
% $$
% \widehat L_S^\Pi(\widehat W)
% =
% \mathbb E_{\epsilon\sim\Pi(\cdot\mid \widehat W)}
% \widehat L_S(\widehat W+\epsilon).
% $$
% The corresponding distributional sharpness is the excess empirical risk under this local posterior:
% $$
% \mathrm{ES}_S^\Pi(\widehat W)
% =
% \mathbb E_{\epsilon\sim\Pi(\cdot\mid \widehat W)}
% \left[
% \widehat L_S(\widehat W+\epsilon)
% -
% \widehat L_S(\widehat W)
% \right].
% $$

% SAM uses an adversarial version of this local posterior view, replacing the expected perturbation loss by the worst-case loss within a radius $\rho$:
% $$
% \widehat L_S^{\mathrm{SAM}}(\widehat W,\rho)
% =
% \max_{\|\epsilon\|\le \rho}
% \widehat L_S(\widehat W+\epsilon)
% =
% \widehat L_S(\widehat W)
% +
% \mathrm{AS}^{(0)}_S(\widehat W,\rho),
% $$
% where
% $$
% \mathrm{AS}^{(0)}_S(\widehat W,\rho)
% =
% \max_{\|\epsilon\|\le\rho}
% \left[
% \widehat L_S(\widehat W+\epsilon)
% -
% \widehat L_S(\widehat W)
% \right].
% $$




% PAC-Bayesian analysis considers a prior $P$ and a posterior $Q$ over parameters, with posterior-averaged risks
% $$
% L_{\mathcal D}(Q)
% =
% \mathbb E_{W\sim Q}L_{\mathcal D}(W),
% \qquad
% \widehat L_S(Q)
% =
% \mathbb E_{W\sim Q}\widehat L_S(W).
% $$
% A perturbation law $\Pi(\cdot\mid W)$ induces a perturbation-smoothed posterior by sampling
% $$
% W\sim Q,
% \qquad
% \epsilon\sim\Pi(\cdot\mid W),
% \qquad
% \widetilde W=W+\epsilon .
% $$
% The corresponding perturbation-smoothed empirical risk is
% $$
% \widehat L_S^\Pi(Q)
% =
% \mathbb E_{W\sim Q}
% \mathbb E_{\epsilon\sim\Pi(\cdot\mid W)}
% \widehat L_S(W+\epsilon).
% $$
% When $Q$ is concentrated at a deterministic solution $\widehat W$, this becomes
% $$
% \widehat L_S^\Pi(\widehat W)
% =
% \mathbb E_{\epsilon\sim\Pi(\cdot\mid \widehat W)}
% \widehat L_S(\widehat W+\epsilon).
% $$








\minisection{PAC-Bayesian bound and Sharpness}
PAC-Bayesian analysis considers a prior $P$ and a posterior $Q$ over model parameters.
Let $\mathcal X$ denote the input space, $\mathcal Y$ the label space, and $\mathcal Z=\mathcal X\times\mathcal Y$ the data space.
A model parameterized by $W\in\mathbb R^d$ induces a predictor $f_W$.
We consider a bounded loss $\ell(W,z)\in[0,1]$.
For a data distribution $\mathcal D$ and a sample $S=\{z_j\}_{j=1}^m$, the population and empirical risks are defined as
$
L_{\mathcal D}(W)
=
\mathbb E_{z\sim\mathcal D}\ell(W,z),
\ 
\hat L_S(W)
=
\frac{1}{m}\sum_{j=1}^m \ell(W,z_j).
$
The corresponding posterior-averaged risks are
$
L_{\mathcal D}(Q)
=
\mathbb E_{W\sim Q}L_{\mathcal D}(W)
$
and
$
\widehat L_S(Q)
=
\mathbb E_{W\sim Q}\widehat L_S(W).
$
A PAC-Bayesian bound controls $L_{\mathcal D}(Q)$ by the empirical term $\widehat L_S(Q)$ and a complexity term depending on $\mathrm{KL}(Q\Vert P)$.


Sharpness enters this framework when the posterior is induced by perturbing a deterministic solution $\widehat W$.
Let $\epsilon\sim\Pi$ and let $Q^\Pi$ denote the distribution of $W=\widehat W+\epsilon$.
Then the posterior-averaged empirical risk becomes \emph{the perturbation-smoothed empirical risk}:
$
\widehat L_S(Q^\Pi)
=
\mathbb E_{\epsilon\sim\Pi}
\widehat L_S(\widehat W+\epsilon).
$
This term admits the decomposition
$$
    \widehat L_S(Q^\Pi)
=
\widehat L_S(\widehat W)
+
\mathrm{ES}^{\Pi}_S(\widehat W), \quad
\mathrm{ES}^{\Pi}_S(\widehat W)
=
\mathbb E_{\epsilon\sim\Pi}
[
\widehat L_S(\widehat W+\epsilon)
-
\widehat L_S(\widehat W)
].
$$
% $$
$\mathrm{ES}^{\Pi}_S(\widehat W)$ is the distributional excess sharpness under the perturbation law $\Pi$.
Thus, perturbation-smoothed empirical risk consists of the clean empirical fit and a task flatness term.
When $\Pi=\mathcal N(0,\Sigma)$, this recovers \emph{the expected sharpness}~\citep{liRevisitingRandomWeight2024,jiang2019fantasticgeneralizationmeasures}:
$
\mathrm{ES}^{(0)}_S(\widehat W,\Sigma)
=
\mathbb E_{\epsilon\sim\mathcal N(0,\Sigma)}
[
\widehat L_S(\widehat W+\epsilon)
-
\widehat L_S(\widehat W)
].
$
SAM uses \emph{the adversarial sharpness} by replacing the expectation with the worst-case perturbation within a radius $\rho$~\citep{foretSharpnessAwareMinimizationEfficiently2021}:
$
\mathrm{AS}^{(0)}_S(\widehat W,\rho)
=
\max_{\|\epsilon\|\le\rho}
[
\widehat L_S(\widehat W+\epsilon)
-
\widehat L_S(\widehat W)
],
$
% For any perturbation distribution $\Pi$ supported in the radius-$\rho$ neighborhood, the random excess sharpness is upper-bounded by this adversarial quantity:
% $$
% \mathrm{ES}^{\Pi}_S(\widehat W)
% \le
% \mathrm{AS}^{(0)}_S(\widehat W,\rho).
% $$
% Therefore, 
which can be interpreted as controlling an adversarial upper of the perturbation-smoothed empirical risk that appears in PAC-Bayesian bounds\cite{mollenhoffSAMOptimalRelaxation2022}.
% GAM~\citep{Zhang_GAM} further considers \emph{the first-order adversarial sharpness} by maximizing the gradient norm within the perturbation neighborhood:
% $
% \mathrm{AS}^{(1)}_S(\widehat W,\rho)
% =
% \max_{\|\epsilon\|\le\rho}
% \left\|
% \nabla_W \widehat L_S(\widehat W+\epsilon)
% \right\|.
% $

% Compared with zeroth-order sharpness, which measures the local increase in loss value, first-order sharpness measures the local sensitivity of the gradient field around $\widehat W$.
% Both quantities are task-conditioned because they depend on the dataset $S$, the loss $\widehat L_S$, the perturbation distribution or radius, and the parameter point around which perturbations are evaluated.



% PAC-Bayesian analysis considers a prior $P$ and a posterior $Q$ over model parameters.
% The corresponding posterior-averaged risks are
% $
% L_{\mathcal D}(Q)
% =
% \mathbb E_{W\sim Q}L_{\mathcal D}(W),
% \ 
% \hat L_S(Q)
% =
% \mathbb E_{W\sim Q}\hat L_S(W).
% $
% A PAC-Bayesian bound controls $L_{\mathcal D}(Q)$ by $\hat L_S(Q)$ and a complexity term depending on $\mathrm{KL}(Q\Vert P)$.
% If the posterior is induced by perturbing a deterministic solution $W$, namely $W+\epsilon$ with $\epsilon\sim\Pi$, then the posterior-averaged empirical risk becomes \emph{the perturbation-smoothed empirical risk}
% $
% \hat L_S(Q)= \hat L_S^{\Pi}(Q)
% =
% \mathbb E_{\epsilon\sim\Pi}
% \hat L_S(W+\epsilon).
% $
% % The corresponding sharpness is defined as the excess empirical risk under the same perturbation law:
% % $$
% % \mathrm{Sh}_S^{\Pi}(W)
% % =
% % \hat L_S^{\Pi}(W)-\hat L_S(W).
% % $$
% The corresponding distributional sharpness is the excess empirical risk under this perturbation-induced posterior.
% % $$
% % \mathrm{ES}_S^{\Pi}(\widehat W)
% % =
% % \mathbb E_{\epsilon\sim\Pi}
% % \left[
% % \widehat L_S(\widehat W+\epsilon)
% % -
% % \widehat L_S(\widehat W)
% % \right].
% % $$
% When $\Pi=\mathcal N(0,\Sigma)$, this recovers the expected sharpness~\citep{liRevisitingRandomWeight2024,jiang2019fantasticgeneralizationmeasures}:
% $
% \mathrm{ES}^{(0)}_S(\widehat W,\Sigma)
% =
% \mathbb E_{\epsilon\sim\mathcal N(0,\Sigma)}
% \left[
% \widehat L_S(\widehat W+\epsilon)
% -
% \widehat L_S(\widehat W)
% \right].
% $
% SAM uses an adversarial counterpart, replacing the expectation by the worst-case perturbation within a radius $\rho$ ~\cite{bahri-etal-2022-sharpness}:
% $
% \mathrm{AS}^{(0)}_S(\widehat W,\rho)
% =
% \max_{\|\epsilon\|\le\rho}
% \left[
% \widehat L_S(\widehat W+\epsilon)
% -
% \widehat L_S(\widehat W)
% \right].
% $ GAM~\citep{Zhang_GAM} define the first order adversarial in the maximum gradient norm within a radius $\rho$ on the training set:
% $
% \mathrm{AS}^{(1)}_S(W,\rho) = \max_{ \| \epsilon \|\leq \rho } \| \nabla_{W} \hat{L}_S(W+\epsilon)\|.
% $



% Equivalently, SAM minimizes the worst-case perturbation-smoothed empirical risk
% $$
% \max_{\|\epsilon\|\le\rho}
% \widehat L_S(\widehat W+\epsilon)
% =
% \widehat L_S(\widehat W)
% +
% \mathrm{AS}^{(0)}_S(\widehat W,\rho).
% $$



% For adversarial perturbations within radius $\rho$, SAM minimizes the worst-case perturbation-smoothed empirical risk
% $$
% \max_{\|\epsilon\|\le \rho}\hat L_S(W+\epsilon)
% =
% \hat L_S(W)
% +
% \mathrm{ASh}^{(0)}_S(W,\rho),\quad
% \mathrm{ASh}^{(0)}_S(W,\rho)
% =
% \max_{\|\epsilon\|\le\rho}
% \left[
% \hat L_S(W+\epsilon)-\hat L_S(W)
% \right].
% $$
% % Sharpness is not an intrinsic property of a model alone, but a dataset- and loss-dependent quantity evaluated at a given parameter point. 
% % In continual learning, we treat sharpness as task-conditioned, since each task induces its own empirical risk and may assign different local sensitivities to the same model.
% % Throughout this work, we use  \emph{sharpness} to denote a quantitative measure of local sensitivity under a specified perturbation set or perturbation distribution.




% \subsection{PAC-Bayesian Risk and Perturbation-Smoothed Sharpness}

% Let $\mathcal X$ be the input space, $\mathcal Y$ the label space, and $\mathcal Z=\mathcal X\times\mathcal Y$ the data space.
% A model with parameter $W\in\mathbb R^d$ induces a predictor $f_W$.
% We consider a bounded loss $\ell(W,z)\in[0,1]$.
% For a distribution $\mathcal D$ and a sample $S=\{z_j\}_{j=1}^m$, define
% $$
% L_{\mathcal D}(W)=\mathbb E_{z\sim\mathcal D}\ell(W,z),
% \qquad
% \hat L_S(W)=\frac{1}{m}\sum_{j=1}^m \ell(W,z_j).
% $$

% For a prior $P$ and a posterior $Q$ over parameters, the posterior-averaged risks are
% $$
% L_{\mathcal D}(Q)=\mathbb E_{W\sim Q}L_{\mathcal D}(W),
% \qquad
% \hat L_S(Q)=\mathbb E_{W\sim Q}\hat L_S(W).
% $$ A standard PAC-Bayesian bound controls $L_{\mathcal D}(Q)$ by $\hat L_S(Q)$ plus a complexity term depending on $\mathrm{KL}(Q\Vert P)$.
% Thus, the empirical term in the bound is not the loss of a single deterministic model, but the loss averaged under a posterior distribution.


% When the posterior is induced by perturbing a deterministic solution $W$, namely $W+\epsilon$ with $\epsilon\sim\Pi$, the posterior-averaged empirical risk becomes the perturbation-smoothed empirical risk
% $$
% \hat L_S^{\Pi}(W)
% =
% \mathbb E_{\epsilon\sim\Pi}\hat L_S(W+\epsilon).
% $$ The corresponding sharpness is the excess sensitivity under the same perturbation law:
% $$
% \mathrm{Sh}_S^{\Pi}(W)
% =
% \hat L_S^{\Pi}(W)-\hat L_S(W).
% $$ 
% For adversarial perturbations within radius $\rho$, SAM minimizes
% $$
% \max_{\|\epsilon\|\le \rho}\hat L_S(W+\epsilon)
% =
% \hat L_S(W)+\mathrm{AS}^{(0)}_S(W,\rho),
% $$
% where
% $$
% \mathrm{AS}^{(0)}_S(W,\rho)
% =
% \max_{\|\epsilon\|\le\rho}
% \left[
% \hat L_S(W+\epsilon)-\hat L_S(W)
% \right].
% $$
% Throughout this work, we use \emph{flat} to denote a qualitative property of the loss landscape, and \emph{sharpness} to denote a quantitative measure of local sensitivity under a specified perturbation set or perturbation distribution.

% % Throughout this work, we use \emph{flatness} to denote a qualitative property of the loss landscape, and \emph{sharpness} to denote a quantitative measure of local sensitivity under a specified perturbation set or perturbation distribution.


\minisection{Continual Learing and LoRA-Constrained Updates}
In CL, the learner observes a sequence of tasks indexed by $t=1,\ldots,T$.
Each task provides a training sample $S_t=\{z_{tj}\}_{j=1}^{m_t}$ drawn i.i.d. from a task distribution $\mathcal D_t$.
After learning tasks $1,\ldots,t$, the learner outputs a model, or more generally a posterior $Q_t$, whose performance is evaluated along the learning trajectory.
This trajectory based evaluation reflects both the acquisition of new tasks and the retention of previously learned tasks, and is the primary object of our analysis.



% In continual learning, the learner observes a sequence of tasks indexed by $t=1,\ldots,T$.
% Each task provides a training sample $S_t=\{z_{tj}\}_{j=1}^{m_t}$ drawn i.i.d. from a task distribution $\mathcal D_t$.
% After learning tasks $1,\ldots,t$, the learner outputs a model, or more generally a posterior $Q_t$, which is evaluated not only on the current task but also on previously encountered tasks.
% This trajectory-level evaluation gives rise to two complementary metrics. The first is forgetting, commonly quantified by backward transfer, which measures the change in earlier-task performance after later updates. The second is overall continual learning performance, commonly summarized by average-accuracy metrics at the final stage or along the learning trajectory. 
% In our analysis, we adopt trajectory-level average performance as the primary object because it captures both the acquisition of new tasks and the retention of previously learned tasks.

% Compared with forgetting alone, average performance provides a broader evaluation criterion because it reflects both the acquisition of new tasks and the retention of previously learned tasks. In our theoretical analysis, we use this trajectory-level view as the starting point.


% The central problem is catastrophic forgetting, namely the degradation of performance on earlier tasks. 
% after subsequent tasks have been learned.
% In continual learning evaluation, forgetting is commonly quantified by backward transfer, which measures the change in earlier-task performance after later tasks.
% Overall performance is commonly evaluated by average-accuracy metrics, which summarize accuracy at the final stage or across the learning trajectory.
% Overall continual learning performance is usually summarized by average accuracy based metrics, such as the last average accuracy over all tasks and trajectory level averages over learning stages.
% These metrics motivate the risk based formulation introduced later, where forgetting is interpreted as a degradation of past task generalization along the posterior trajectory.


% For any previously observed task $i\le t$, we denote the posterior-averaged population risk of $Q_t$ on $\mathcal D_i$ by
% $$
% R_{t,i}
% =
% L_{\mathcal D_i}(Q_t)
% =
% \mathbb E_{W\sim Q_t}L_{\mathcal D_i}(W).
% $$


% In continual learning, the learner observes a sequence of tasks $t=1,\ldots,T$, where each task provides a sample $S_t=\{z_{tj}\}_{j=1}^{m_t}$ drawn i.i.d. from a task distribution $\mathcal D_t$.
% Let $Q_t$ denote the posterior after learning tasks up to $t$.



% For an earlier task $i\le t$, its population risk under the later posterior $Q_t$ is
% $$
% L_{\mathcal D_i}(Q_t)
% =
% \mathbb E_{W\sim Q_t}L_{\mathcal D_i}(W).
% $$
% From a PAC-Bayesian perspective, continual learning can be understood as maintaining generalization along a sequence of posteriors $\{Q_t\}_{t=1}^{T}$ over previously encountered task distributions.
% Under this view, forgetting corresponds to the degradation of generalization on earlier tasks after learning later ones.
% For $i<t$, we write this risk-form forgetting as
% $$
% \mathcal F_{i,t}
% =
% L_{\mathcal D_i}(Q_t)
% -
% L_{\mathcal D_i}(Q_i),
% $$
% where $Q_i$ is the posterior immediately after learning task $i$.
% This definition is consistent with the standard backward transfer(BWT) criterion in continual learning, up to the sign convention and the use of risk rather than accuracy.
% However, our theory does not focus on a single pairwise forgetting gap.
% Instead, it studies the learner's overall behavior along the posterior trajectory.
% The central object is the trajectory-averaged hierarchical test risk
% $$
% \mathcal L_{\mathrm{traj}}(\mathcal Q_{1:T})
% =
% \frac{1}{T}
% \sum_{t=1}^{T}
% \frac{1}{t}
% \sum_{i=1}^{t}
% \mathbb E_{P_t\sim\mathcal Q_t}
% \mathbb E_{W\sim Q_t(P_t)}
% L_{\mathcal D_i}(W).
% $$
% This quantity is the risk-based counterpart of average anytime accuracy, which averages the accuracy over all observed tasks at each learning stage and then averages over stages.
% Thus, minimizing $\mathcal L_{\mathrm{traj}}$ aligns with maximizing the standard trajectory-level continual learning performance.
% Theorem~\ref{thm:pathwise_pac} bounds this trajectory-averaged hierarchical test risk by the corresponding trajectory-averaged empirical risk plus a complexity penalty.
% We define forgetting as the degradation of generalization on previously learned tasks after learning new ones.
% Equivalently, for $i<t$, the forgetting of task $i$ after task $t$ is measured by
% $$
% \mathcal F_{i,t}
% =
% L_{\mathcal D_i}(Q_t)
% -
% L_{\mathcal D_i}(Q_i),
% $$
% where $Q_i$ is the posterior immediately after learning task $i$.
% Thus, forgetting is treated as past-task generalization degradation along the posterior trajectory.
% \subsection{LoRA-Constrained Updates}
LoRA-based PECL adapts a
frozen backbone by learning task-wise low-rank updates.
% In LoRA based CL, task adaptation is performed through trainable low rank updates while the pretrained backbone remains fixed.
% At task $t$, the model can be written as
% $
% W_t = W_t^{\mathrm{froz}} + \Delta W_t,
% $
% where $W_t^{\mathrm{froz}}$ denotes the frozen part during task $t$ and 
$\Delta W_t =A_tB_t$ is generated by low rank factors, and the effective update is restricted to a task dependent geometry.
Different PECL designs can be understood as different rules for organizing and evolving this update geometry across tasks.


% \section{Why Flatness Matters for Continual Learning}
% \label{sec:flatness_cl}

% Sharpness-aware optimization is well understood through PAC-Bayesian generalization theory as a way to control perturbation-smoothed empirical risk. However, continual learning raises a different question. Forgetting is not the generalization error of the current posterior on the current task, but the degradation of an earlier task under a later posterior. Let $Q_t$ denote the posterior produced after learning task $t$. For an earlier task $i<t$, forgetting can be viewed as
% $$
% F_{i\to t}=L_{\mathcal D_i}(Q_t)-L_{\mathcal D_i}(Q_i).
% $$
% Thus, the central question is not merely whether flatness improves single-task generalization, but why the perturbation-smoothed empirical risk optimized during sequential learning should be relevant to the old-task risks evaluated under later posteriors.

% This perspective exposes a gap in existing explanations of flatness in continual learning. A common intuition is that flat old-task basins tolerate larger subsequent updates and therefore suffer less forgetting. While plausible, this geometric argument does not identify the PAC-Bayesian quantity that connects sharpness-aware training to forgetting. In particular, it does not explain how a smoothed empirical risk term, originally introduced for single-task generalization, can control risk changes along a posterior trajectory. We therefore require a sequential analysis in which the risks of previous tasks are evaluated under the evolving posteriors produced by later tasks.
% \section{Why Flatness Matters for Continual Learning}
% \label{sec:flatness_cl}

{%
\setlength{\intextsep}{8pt}        % 浮动体与正文上下间距
\setlength{\textfloatsep}{6pt}     % 页顶/页底浮动体与正文间距
\setlength{\abovecaptionskip}{2pt} % 图与caption之间间距
\setlength{\belowcaptionskip}{0pt} % caption与正文之间间距
\captionsetup{skip=4pt}

\begin{figure}[t]
    \centering
    \begin{minipage}[t]{0.49\linewidth}
        \centering
        \includegraphics[width=\linewidth,height=4.2cm,keepaspectratio]{figures_TRML/ouranalysis.pdf}
        \vspace{2pt}
        {\footnotesize (a)}
    \end{minipage}
    \hfill
    \begin{minipage}[t]{0.49\linewidth}
        \centering
        \includegraphics[width=\linewidth,height=4.2cm,keepaspectratio]{figures_TRML/KLshiyi1.pdf}
        
        \vspace{2pt}
        {\footnotesize (b)}
    \end{minipage}

    \vspace{-2pt}
    \caption{\textbf{Time-varying hyperposterior and KL decomposition in PECL.}
    (a) In previous bound, task priors are drawn from a single time-invariant hyperdistribution, so there is no sequential drift. In continual learning, the hyperposterior is updated over time, inducing a history-dependent cumulative drift over tasks.
    (b) The posterior and prior defined on the adapter update subspace, $(Q_t^\Delta, P_t^\Delta)$, induce full-model distributions $(Q_t, P_t)$ with $\mathrm{KL}(Q_t \Vert P_t)=\mathrm{KL}(Q_t^\Delta \Vert P_t^\Delta)$. Therefore, the complexity term depends only on the task update subspace.}
    \label{fig:showPAC_combined}
    \vspace{-0.4 cm}
\end{figure}
} % end local spacing group
 % \vspace{-0.6 cm}
% % The previous section introduced flatness through perturbation-smoothed empirical risk in single task. 
% We now ask why this quantity should matter in continual learning. 
% % The issue is not merely whether a flatt posterior generalizes better on the current task. 
% In continual learning, the learner produces a trajectory of posteriors, and forgetting occurs when later posteriors degrade the population risk of earlier tasks. 
% This section shows that perturbation-smoothed empirical risk enters continual learning through a pathwise PAC-Bayesian decomposition, where it is coupled with task-level adaptation complexity and hyperposterior drift.

% \subsection{From Task-Conditioned Sharpness to Sequential Generalization}
% \label{sec:flatness_to_forgetting}


% Continual learning changes flatness analysis from a task-local posterior problem to a pathwise posterior-trajectory problem. A flatness term models the local behavior of a task loss around the posterior induced at a given learning step. In a sequential setting, this local object is embedded in a history-dependent trajectory $Q_1,\ldots,Q_T$, where each posterior $Q_t$ is produced from the current task sample $S_t$ and the learning history before task $t$. Hence, the sharpness term at task $t$ is a history-conditioned local model of the current task loss around $Q_t$.

% This trajectory view has three consequences. First, perturbation-smoothed risk is task-conditioned, since each term is tied to the current distribution $\mathcal D_t$, loss, perturbation distribution, and posterior $Q_t$. Second, the prior-generating mechanism is history dependent: the prior before task $t$ is produced after observing $S_{1:t-1}$, rather than sampled from a fixed task environment. Third, trajectory-level performance depends jointly on local smoothed risks and posterior movement. Sharpness-aware optimization controls the current-task smoothed risk when $Q_t$ is produced, while within-task adaptation and hyperposterior drift enter the bound as complexity terms that measure the cost of the posterior trajectory. 
% Therefore, The first step is  to show how task-conditioned perturbation-smoothed empirical risks enter a pathwise PAC-Bayesian bound over the posterior trajectory. 

% The next subsection provides this decomposition through a time-varying hyperposterior process.


% Sharpness is task-conditioned quantity, its role in continual learning cannot be inferred from a single posterior alone. Continual learning produces a history-dependent posterior trajectory $Q_1,\ldots,Q_T$, where each posterior $Q_t$ is induced by the current task sample $S_t$ and the learning history before task $t$. When task $t$ is learned, a sharpness-aware objective optimizes perturbation-smoothed empirical risk on $S_t$. Its PAC-Bayesian counterpart is therefore the population risk of $Q_t$ on the current task distribution $\mathcal D_t$.

% This quantity should not be confused with forgetting. Forgetting is evaluated by applying later posteriors to earlier tasks, whereas sharpness-aware learning acts locally when each posterior $Q_t$ is produced. The first step is therefore to show how current-task perturbation-smoothed risks control the sequential generalization components along the posterior trajectory. The next subsection provides this pathwise PAC-Bayesian decomposition.


% Thus, before explaining trajectory-level degradation geometrically, one must first explain why the current-task smoothed empirical risks optimized along the sequence are theoretically meaningful for continual learning. The next subsection provides this pathwise PAC-Bayesian decomposition.


% Sharpness is a task-conditioned quantity rather than an intrinsic property of a model alone. It depends on the task loss, the data distribution, the perturbation distribution, and the parameter point at which it is evaluated. Thus, the same posterior may exhibit different local sensitivities under different task distributions. In standard single-task learning, PAC-Bayesian flatness analyses relate perturbation-smoothed empirical risk to the population risk of the same task-induced posterior.
% Continual learning changes this object of analysis. The learner does not produce a single posterior for one task, but a history-dependent posterior trajectory $Q_1,\ldots,Q_T$ induced by sequential updates. When task $t$ is learned, a sharpness-aware objective is optimized on the current sample $S_t$, and its immediate PAC-Bayesian counterpart is the population risk of $Q_t$ on the current task distribution $\mathcal D_t$. This quantity should not be confused with backward transfer or forgetting, which are evaluated by applying later posteriors to earlier tasks. Rather, it is the task-conditioned sequential generalization component through which perturbation-smoothed empirical risks first enter the analysis of a continual learning trajectory.

% This distinction is important. Geometric explanations usually argue that flatter old-task basins are more tolerant to later updates, but this does not explain why minimizing current-task perturbation-smoothed empirical risk should be theoretically meaningful for performance along the entire posterior trajectory. A sequential theory must first account for how task-level smoothed risks accumulate along the realized learning path, and only then can posterior drift and task-conditioned sensitivity be used to interpret degradation on earlier tasks. The next subsection provides this pathwise PAC-Bayesian decomposition.





% Previous sharpness analyses 
% % provide a principled way to explain why sharpness-aware learning can improve generalization in standard supervised learning. 
% % For a single task, the learner outputs a posterior $Q$ over hypotheses, and the population risk $L_{\mathcal D}(Q)$ is controlled by a perturbation-smoothed empirical risk together with a KL complexity term. 
% % This view 
% formalizes flatness as local stability of the loss around the learned posterior: if nearby perturbations do not substantially increase the empirical loss, then the posterior is expected to generalize better.
% However, continual learning changes the object of analysis. 
% The learner no longer produces one posterior for one task, but a history-dependent posterior trajectory $Q_1,\ldots,Q_T$. 
% After learning task $t$, the resulting posterior $Q_t$ is evaluated not only on the current task, but also on earlier tasks. 
% For task $i$ evaluated after learning task $t$, we define
% \begin{equation}
% R_{i,t}
% =
% L_{\mathcal D_i}(Q_t)
% =
% \mathbb E_{W\sim Q_t}
% L_{\mathcal D_i}(W),
% \qquad i\leq t.
% \end{equation}
% The diagonal term $R_{t,t}$ measures current-task generalization, while the off-diagonal term $R_{i,t}$ with $i<t$ measures the performance of an earlier task under a later posterior. 
% Accordingly, risk-form forgetting from task $i$ to task $t$ is
% \begin{equation}
% \mathcal F_{i,t}
% =
% R_{i,t}-R_{i,i}
% =
% L_{\mathcal D_i}(Q_t)-L_{\mathcal D_i}(Q_i),
% \qquad i<t.
% \end{equation}
% The average seen-task risk after learning task $t$ is
% \begin{equation}
% \mathcal R_t^{\mathrm{seen}}
% =
% \frac{1}{t}
% \sum_{i=1}^{t}
% R_{i,t}
% =
% \frac{1}{t}
% \sum_{i=1}^{t}
% L_{\mathcal D_i}(Q_t).
% \end{equation}
%  
% This trajectory view also clarifies how flatness can be connected to forgetting. 
% For any earlier task $i<t$, forgetting can be written as the accumulated effect of later posterior transitions:
% \begin{equation}
% \mathcal F_{i,t}
% =
% R_{i,t}-R_{i,i}
% =
% \sum_{k=i+1}^{t}
% \left(
% R_{i,k}-R_{i,k-1}
% \right).
% \end{equation}
% Therefore, the relevant question is not only whether $Q_i$ generalizes well on task $i$, but also how much each subsequent update from $Q_{k-1}$ to $Q_k$ increases the old-task risk. 
% Flatness becomes relevant precisely through this transition. 
% If the old-task loss remains locally stable around the posterior trajectory, then a later update is less likely to amplify posterior movement into a large increase in $R_{i,k}$. 
% In this sense, flatness in continual learning should be understood as a local sensitivity term that modulates how sequential posterior drift is converted into forgetting.

% This distinction explains why existing PAC-Bayesian views do not directly answer the flatness-in-continual-learning question. 
% PAC-Bayesian flatness analyses justify sharpness-aware learning through perturbation-smoothed empirical risk, but they primarily concern the population risk of a posterior on the task that induces it. 
% Hierarchical PAC-Bayesian bounds for lifelong learning separate transferable inductive bias from task-specific adaptation, but their exchangeable-task formulation assumes a fixed task-generating environment and therefore does not model the realized posterior trajectory along a specific task order. 
% Recent cumulative PAC-Bayesian bounds for continual learning are sequential and allow data-dependent priors, but their objective is to certify cumulative prediction risk along the task sequence. 
% A naive insertion of perturbation-smoothed losses into such a cumulative bound would smooth the current-task or cumulative prediction terms, but it would not by itself expose how the transition from $Q_i$ to $Q_t$ changes the old-task risk $R_{i,t}$ for $i<t$.


% \subsection{From Flatness-Based Generalization to Continual Forgetting}
% \label{sec:flatness_to_forgetting}

% PAC-Bayesian flatness analyses in standard supervised learning relate the population risk of a posterior $Q$ to its perturbation-smoothed empirical risk and a KL complexity term. 
% This explains why sharpness-aware learning can improve single-task generalization. 
% However, continual learning changes the object of analysis. 
% The learner does not produce one posterior for one task, but a history-dependent posterior trajectory $Q_1,\ldots,Q_T$.
% For an earlier task $i\leq t$, we define the risk of task $i$ under the posterior after learning task $t$ as
% \begin{equation}
% R_{i,t}
% =
% L_{\mathcal D_i}(Q_t)
% =
% \mathbb E_{W\sim Q_t}L_{\mathcal D_i}(W).
% \end{equation}
% The diagonal term $R_{t,t}$ measures current-task generalization, while the off-diagonal term $R_{i,t}$ with $i<t$ measures the performance of an earlier task under a later posterior. 
% Accordingly, risk-form forgetting after learning task $t$ is
% \begin{equation}
% \mathcal F_{i,t}
% =
% R_{i,t}-R_{i,i}
% =
% L_{\mathcal D_i}(Q_t)-L_{\mathcal D_i}(Q_i),
% \qquad i<t.
% \end{equation}
% The average seen-task risk after task $t$ is
% \begin{equation}
% \mathcal R_t^{\mathrm{seen}}
% =
% \frac{1}{t}\sum_{i=1}^t R_{i,t}
% =
% \frac{1}{t}\sum_{i=1}^t L_{\mathcal D_i}(Q_t).
% \end{equation}
% Thus, both forgetting and overall continual-learning performance depend on off-diagonal risks along the posterior trajectory, rather than only on the current-task risks $\{R_{t,t}\}_{t=1}^T$.

% % This trajectory view also changes the role of sharpness. 
% % In single-task analysis, sharpness is evaluated on the same task that produces the posterior. 
% % In continual learning, forgetting depends on whether a later posterior remains locally stable for earlier task losses. 
% % Let $\widehat W_t$ denote the deterministic parameter obtained after task $t$, and let $\rho_t$ be a local perturbation distribution around $\widehat W_t$. 
% % For task $i$ evaluated at time $t$, we define the perturbation-smoothed empirical risk as
% % \begin{equation}
% % \widehat R_{i,t}^{\rho}
% % =
% % \mathbb E_{\epsilon\sim\rho_t}
% % \widehat L_{S_i}(\widehat W_t+\epsilon),
% % \qquad i\leq t,
% % \end{equation}
% % and the corresponding task-conditioned sharpness as
% % \begin{equation}
% % \mathcal S_{i,t}^{\rho}
% % =
% % \widehat R_{i,t}^{\rho}
% % -
% % \widehat L_{S_i}(\widehat W_t).
% % \end{equation}
% % When $i=t$, this recovers the usual current-task sharpness. 
% % When $i<t$, it measures the local sensitivity of an earlier task loss around a later posterior, which is the sharpness quantity directly relevant to forgetting.

% This distinction clarifies why existing PAC-Bayesian views do not directly answer the flatness-in-continual-learning question. 
% PAC-Bayesian flatness analyses justify sharpness-aware learning through perturbation-smoothed empirical risk, but they primarily control single-task generalization. 
% Hierarchical PAC-Bayesian bounds for lifelong learning separate transferable inductive bias from task-specific adaptation, but usually assume a stationary task environment and are therefore exchangeable over the observed tasks. 
% Recent cumulative PAC-Bayesian bounds for continual learning are sequential, but they certify cumulative prediction risk rather than old-task degradation under later posteriors. 
% Even replacing the empirical loss in such a cumulative bound with a perturbation-smoothed loss would control smoothed current-task risk, not the forgetting quantity $R_{i,t}$ for $i<t$.

% Therefore, explaining flatness in continual learning requires a bound on old-task risk under later posteriors. 
% Such a bound must be sequential, because the posterior evolves along the task order; hierarchical, because the history-dependent inductive bias should be separated from task-specific adaptation; and sharpness-preserving, because the perturbation-smoothed empirical risk is the formal PAC-Bayesian object corresponding to flatness. 
% We derive this decomposition next.


% % \subsection{From Flatness-Based Generalization to Forgetting}
% % \label{sec:flatness_to_forgetting}

% % In standard supervised learning, PAC-Bayesian flatness analyses relate the population risk of a posterior $Q$ to its perturbation-smoothed empirical risk and a KL complexity term. 
% % This gives a single-task explanation of why sharpness-aware learning can improve generalization. 
% % Continual learning changes the object of analysis. 
% % The learner observes a sequence of tasks $1,\ldots,T$, where task $t$ provides a dataset $S_t=\{z_{t,j}\}_{j=1}^{m_t}$ drawn from a task distribution $\mathcal D_t$, and produces a posterior trajectory $Q_1,\ldots,Q_T$.

% % The central quantity in continual learning is not only the current-task risk $L_{\mathcal D_t}(Q_t)$, but also the risk of earlier tasks under later posteriors. 
% % For an earlier task $i<t$, we define the risk-form forgetting after learning task $t$ as
% % \begin{equation}
% % \mathcal F_{i,t}
% % =
% % L_{\mathcal D_i}(Q_t)
% % -
% % L_{\mathcal D_i}(Q_i),\quad
% % \text{where } 
% % L_{\mathcal D_i}(Q_t)
% % =
% % \mathbb E_{W\sim Q_t}
% % L_{\mathcal D_i}(W).
% % \end{equation}
% % and the average perofrmance we care is about 
% % \begin{equation}
% % \mathcal A_{i,t}
% % =
% % L_{\mathcal D_i}(Q_t)
% % +
% % L_{\mathcal D_i}(Q_i),\quad
% % \end{equation}
% % but the problem is the later is depended on the pervious result, is not exchangable.

% % This definition measures whether replacing the posterior obtained after task $i$ with the later posterior obtained after task $t$ increases the population risk on the same earlier distribution $\mathcal D_i$. 
% % It is the risk-based analogue of backward-transfer forgetting metrics (BWT), up to the sign convention and the use of risk rather than accuracy.

% % This perspective clarifies the theoretical gap. 
% % Flatness-based generalization explains why the posterior $Q_t$ may generalize well around task $t$, but forgetting asks whether the later posterior $Q_t$ remains reliable on $\mathcal D_i$ for $i<t$. 
% % Therefore, explaining flatness in continual learning requires a bound that does not only control current-task generalization, but also describes how risks evolve along the posterior trajectory.



% % \subsection{Why Existing PAC-Bayesian Views Are Insufficient}
% % \label{sec:limits_existing_pac_cl}

% % Existing PAC-Bayesian results provide important ingredients, but they do not directly answer the flatness-in-continual-learning question studied here. PAC-Bayesian flatness analyses justify sharpness-aware objectives by replacing empirical risk with perturbation-smoothed empirical risk \cite{bahri-etal-2022-sharpness,mollenhoffSAMOptimalRelaxation2022}. This explains why flatter solutions can improve single-task generalization, but it does not model the sequential posterior trajectory that gives rise to forgetting. In continual learning, the central concern is not only the risk of $Q_t$ on the current task, but also how later posteriors affect earlier task risks after subsequent updates.

% % PAC-Bayesian analyses for lifelong or continual learning address parts of this sequential problem, but remain insufficient for our purpose. Hierarchical lifelong-learning bounds introduce a hyperposterior over task-level priors, which is useful for separating transferable inductive bias from task-specific adaptation \cite{pentinaPACBayesianBoundLifelong2014}. However, they typically assume tasks are sampled from a fixed task environment, yielding an order-invariant view that does not capture history-dependent drift. Recent cumulative-loss bounds for continual learning move closer to the sequential setting \cite{friedman2026pac}, but they primarily certify cumulative generalization loss and do not isolate perturbation-smoothed empirical risk as the mechanism through which flatness may affect forgetting.

% % Therefore, explaining flatness in continual learning requires a framework that is simultaneously sequential, hierarchical, and sharpness-preserving. It must track how posteriors evolve along a task sequence, separate within-task adaptation from the drift of transferable inductive bias, and retain the perturbation-smoothed empirical risk term optimized by sharpness-aware learning. This motivates our pathwise hierarchical PAC-Bayesian decomposition, which links flatness, posterior drift, and forgetting within a single continual-learning bound.


% % \subsection{Why Existing PAC-Bayesian Views Are Insufficient}
% % \label{sec:limits_existing_pac_cl}

% % Existing PAC-Bayesian results provide important ingredients, but they do not directly answer the flatness-in-continual-learning question. 
% % SAM and Bayesian based flatness analyses justify sharpness-aware objectives by replacing empirical risk with perturbation-smoothed empirical risk \cite{bahri-etal-2022-sharpness, mollenhoffSAMOptimalRelaxation2022}. 
% % These results explain why flatter solutions can generalize better in a single-task setting, but they do not model the posterior drift responsible for forgetting.

% % % Hierarchical PAC-Bayesian analyses for lifelong or meta-learning \cite{pentinaPACBayesianBoundLifelong2014}introduce a hyperposterior over task-level priors, which is useful for separating transferable inductive bias from task-specific adaptation. 
% % % However, the formulation assumes that tasks are sampled exchangeably from a fixed hyperdistribution. 
% % % Consequently, the resulting bound is order-invariant and does not describe how the inductive bias changes along a realized task sequence. 
% % % In our framework, this setting corresponds to the special case where the hyperposterior drift is zero.

% % % Recent PAC-Bayesian analyses for sequential continual learning\cite{friedman2026pac} move closer to the continual-learning setting by considering data-dependent priors and cumulative loss along a task sequence. 
% % % These results are important for certifying sequential learning, but they are not designed to explain flatness. 
% % % They bound cumulative generalization loss, rather than isolating perturbation-smoothed empirical risk as the mechanism through which sharpness-aware learning affects forgetting. 
% % % Moreover, they do not explicitly separate the hyper-level evolution of inductive bias from task-level posterior adaptation.

% % % These limitations suggest that explaining flatness in continual learning requires a framework with three properties. 
% % % First, it must be sequential, so that posterior trajectories and history-dependent priors are explicitly modeled. 
% % % Second, it must be hierarchical, so that task-level adaptation and the evolution of transferable inductive bias are separated. 
% % % Third, it must retain the perturbation-smoothed empirical risk term, since this is the formal quantity optimized by sharpness-aware learning.



% % % Section~\ref{sec:setup} recalled that PAC-Bayesian flatness analysis naturally leads to perturbation-smoothed empirical risk, and that continual learning evaluates a trajectory of posteriors rather than a single posterior. 
% % % We now ask whether the same smoothed-risk object can explain forgetting. The key change is that continual learning does not only require controlling the current-task risk $L_{\mathcal D_t}(Q_t)$, but also the old-task risks $L_{\mathcal D_i}(Q_t)$ induced by later posteriors for $i<t$.

% % % Let $Q_1,\dots,Q_T$ denote the posterior trajectory produced by continual learning. For each task $i$ and stage $t\ge i$, define
% $$
% R_{i,t}=L_{\mathcal D_i}(Q_t).
% $$
% The diagonal terms $R_{t,t}$ correspond to current-task generalization, whereas the off-diagonal terms $R_{i,t}$ measure how later posteriors behave on earlier task distributions. Forgetting can therefore be written as
% $$
% F_{i\to t}
% =
% R_{i,t}-R_{i,i}
% =
% L_{\mathcal D_i}(Q_t)-L_{\mathcal D_i}(Q_i),
% \quad i<t.
% $$
% Thus, explaining forgetting requires controlling past-task population risks under later posteriors, rather than only controlling the risk of each posterior on its own training distribution.
% \subsection{A Sequential Hierarchical Decomposition}
% \label{sec:seq_hier_decomp}

% % Continual learning requires a pathwise flatness analysis because sharpness is evaluated locally at each task posterior, whereas performance is accumulated along a history-dependent posterior trajectory.
% At task $t$, the perturbation-smoothed risk is tied to the current sample $S_t$, task distribution $\mathcal D_t$, perturbation law, and posterior produced at that step.
% Meanwhile, the prior before task $t$ is generated from the previous learning history rather than sampled from a fixed task environment.
% Thus, a PAC-Bayesian analysis for continual learning must account for both the task-conditioned smoothed risk and the cost of moving the posterior process along the task sequence.

% We therefore need a decomposition that keeps track of both the local geometry of task losses and the sequential evolution of the learner. 


% \subsection{From Task-Conditioned Sharpness to Posterior Trajectories}

% Sharpness-aware learning optimizes a task-conditioned quantity, whereas continual learning evaluates a posterior trajectory. 
% In standard single-task PAC-Bayesian analyses, perturbation-smoothed empirical risk is associated with the population risk of one posterior on the task that induces it. 
% Continual learning changes this object: after task $t$, the learner has produced a history-dependent posterior sequence rather than an isolated posterior. 
% Therefore, the relevant theoretical question is not whether sharpness improves one task in isolation, but whether the task-conditioned smoothed risks optimized along the sequence enter a bound on trajectory-level performance.




% \section{A Pathwise PAC-Bayesian Answer to Q1}
% To formalize this view, let $\mathcal F_{t-1}=\sigma(S_1,\ldots,S_{t-1})$ denote the information available before observing task $t$. 


\section{A Pathwise PAC-Bayesian Answer to Q1}
\label{sec:pathwise_pac_bayes_q1}

We answer Q1 by showing that the perturbation-smoothed empirical risks optimized at individual tasks enter a pathwise generalization bound over the posterior trajectory.
This is the missing step in existing explanations of flatness in continual learning.
Single-task PAC-Bayesian flatness analyses connect perturbation-smoothed empirical risk to the population risk of one task-induced posterior, while lifelong PAC-Bayesian analyses typically consider tasks sampled from a fixed environment distribution.
Neither view directly captures the realized sequence of posterior updates produced by continual learning.
We need a sequential PAC-Bayesian formulation in which the prior-generating process evolves with the learning history and the bound holds along the realized task trajectory.

Let $\mathcal F_{t-1}=\sigma(S_1,\ldots,S_{t-1})$ denote the information available before observing task $t$.
Before task $t$, the hyperposterior $\mathcal P_{t-1}$ is $\mathcal F_{t-1}$-measurable and generates a task prior $P_t$.
After observing $S_t$, the learner produces a perturbation-induced task posterior $Q_t^\Pi(\cdot\mid P_t,S_t)$, where $\Pi_t$ denotes the perturbation law used to define the corresponding smoothed posterior.
The hyperposterior is then updated from $\mathcal P_{t-1}$ to $\mathcal P_t$.
This construction separates within-task adaptation from the history-dependent drift of the prior-generating mechanism.

For each task $t$, define
$
\widehat L_{S_t}
\bigl(
Q_t^\Pi(\cdot\mid P_t,S_t)
\bigr)
=
\mathbb E_{W_t\sim Q_t^\Pi(\cdot\mid P_t,S_t)}
\widehat L_{S_t}(W_t),
$
and
$
L_{\mathcal D_t}
\bigl(
Q_t^\Pi(\cdot\mid P_t,S_t)
\bigr)
=
\mathbb E_{W_t\sim Q_t^\Pi(\cdot\mid P_t,S_t)}
L_{\mathcal D_t}(W_t).
$
The corresponding task-conditioned empirical and population risks are
$
\widehat{\mathcal L}_t^\Pi
=
\mathbb E_{P_t\sim\mathcal P_{t-1}}
\widehat L_{S_t}
\bigl(
Q_t^\Pi(\cdot\mid P_t,S_t)
\bigr),
$
and
$
\mathcal L_t^\Pi
=
\mathbb E_{P_t\sim\mathcal P_{t-1}}
L_{\mathcal D_t}
\bigl(
Q_t^\Pi(\cdot\mid P_t,S_t)
\bigr).
$
To move from a single-task quantity to the continual-learning trace, we aggregate these task-conditioned risks along the posterior trajectory:
$
\widehat{\mathcal L}_{1:T}^\Pi
=
\frac{1}{T}
\sum_{t=1}^{T}
\widehat{\mathcal L}_t^\Pi,
\quad
\mathcal L_{1:T}^\Pi
=
\frac{1}{T}
\sum_{t=1}^{T}
\mathcal L_t^\Pi .
$
These trajectory-level quantities summarize the task-conditioned generalization components induced by sequential learning.
They are the PAC-Bayesian objects through which perturbation-smoothed empirical risk at each update step enters a bound over the posterior trace.






% We answer Q1 by showing that the perturbation-smoothed empirical risks optimized at individual tasks enter a pathwise generalization bound over the posterior trajectory. This is the missing step in existing explanations of flatness in continual learning. Single-task PAC-Bayesian flatness analyses connect perturbation-smoothed empirical risk to the population risk of one task-induced posterior, while lifelong PAC-Bayesian analyses assume tasks sampled from a fixed environment distribution. Neither view directly captures the realized sequence of posterior updates produced by continual learning. We therefore need a sequential PAC-Bayesian formulation in which the prior-generating process evolves with the learning history and the bound holds along the realized task trajectory.

% % We model this process through a time-varying hyperposterior.
% Let $\mathcal F_{t-1}=\sigma(S_1,\ldots,S_{t-1})$ denote the information available before observing task $t$. Before task $t$, the hyperposterior $\mathcal P_{t-1}$ is $\mathcal F_{t-1}$-measurable and generates a task prior $P_t$. After observing $S_t$, the learner produces a perturbation-induced task posterior $Q_t^\Pi(\cdot\mid P_t,S_t)$, where $\Pi_t$ denotes the perturbation law used to define the corresponding smoothed posterior. The hyperposterior is then updated from $\mathcal P_{t-1}$ to $\mathcal P_t$. This construction separates within-task adaptation from the history-dependent drift of the prior-generating mechanism.
% For each task $t$, we define the perturbation-smoothed empirical risk and the corresponding population risk as
% $
% \widehat{\mathcal{L}}_t^\Pi
% =
% \mathbb E_{P_t\sim\mathcal P_{t-1}}  \widehat L_S(Q^\Pi)
% ,
% $
% and
% $
% \mathcal{L}_t^\Pi
% =
% \mathbb E_{P_t\sim\mathcal P_{t-1}}L_{\mathcal D_t}(Q^\Pi).
% $ To move from a single-task quantity to the continual-learning trace, we aggregate these task-conditioned risks along the posterior trajectory:
% $
% \widehat{\mathcal L}_{1:T}^\Pi
% =
% \frac{1}{T}
% \sum_{t=1}^{T}
% \widehat{\mathcal L}_t^\Pi,
% \quad
% \mathcal L_{1:T}^\Pi
% =
% \frac{1}{T}
% \sum_{t=1}^{T}
% \mathcal L_t^\Pi .
% $
% These trajectory-level quantities summarize the task-conditioned generalization components induced by sequential learning.
% They are the PAC-Bayesian objects that connect perturbation-smoothed empirical risk at each update step to overall performance along the posterior trace.


% To move from a single-task quantity to the continual-learning trace, we aggregate these task-conditioned risks along the posterior trajectory: $
% \widehat{\mathcal L}_{1:T}^\Pi
% =
% \frac{1}{T}
% \sum_{t=1}^{T}
% \widehat{\mathcal L}_t^\Pi,
% \quad
% \mathcal L_{1:T}^\Pi
% =
% \frac{1}{T}
% \sum_{t=1}^{T}
% \mathcal L_t^\Pi .
% $
To make the connection to flatness explicit, suppose that for each sampled prior $P_t$, the posterior $Q_t^\Pi(\cdot\mid P_t,S_t)$ is induced by perturbing a learned solution $\widehat W_t$ with $\epsilon_t\sim\Pi_t$, namely $W_t=\widehat W_t+\epsilon_t$.
Then the perturbation-smoothed empirical risk admits the exact decomposition
\begin{equation}
\widehat L_{S_t}
\bigl(
Q_t^\Pi(\cdot\mid P_t,S_t)
\bigr)
=
\mathbb E_{\epsilon_t\sim\Pi_t}
\widehat L_{S_t}(\widehat W_t+\epsilon_t)
=
\widehat L_{S_t}(\widehat W_t)
+
\mathrm{ES}_{S_t}^{\Pi_t}(\widehat W_t),
\end{equation}
where
$
\mathrm{ES}_{S_t}^{\Pi_t}(\widehat W_t)
=
\mathbb E_{\epsilon_t\sim\Pi_t}
[
\widehat L_{S_t}(\widehat W_t+\epsilon_t)
-
\widehat L_{S_t}(\widehat W_t)
]
$
is the task-conditioned excess sharpness under $\Pi_t$.
Thus, the empirical term in the trajectory-level bound contains both current-task empirical fit and task-conditioned sharpness accumulated along the posterior trajectory.
For adversarial sharpness-aware objectives such as SAM, the corresponding worst-case sharpness can be viewed as an upper envelope of this perturbation-smoothed term.




% To make the connection to flatness explicit, suppose that, conditional on $P_t$ and $S_t$, the posterior $Q_t^\Pi$ is induced by perturbing a learned solution $\widehat W_t=\widehat W_t(P_t,S_t): 
% W_t=\widehat W_t+\epsilon_t,
% \quad
% \epsilon_t\sim\Pi_t.
% $
% Then the posterior-averaged empirical risk admits the exact decomposition
% \begin{equation}
% \begin{aligned}
% \hat{L}(Q_t^\Pi)=
% \mathbb E_{W_t\sim Q_t^\Pi(\cdot\mid P_t,S_t)}
% \widehat L_{S_t}(W_t)
% =
% \mathbb E_{\epsilon_t\sim\Pi_t}
% \widehat L_{S_t}(\widehat W_t+\epsilon_t)
% =
% \widehat L_{S_t}(\widehat W_t)
% +
% \mathrm{ES}_{S_t}^{\Pi_t}(\widehat W_t),
% \end{aligned}
% \end{equation}
% where
% $
% \mathrm{ES}_{S_t}^{\Pi_t}(\widehat W_t)
% =
% \mathbb E_{\epsilon_t\sim\Pi_t}
% \left[
% \widehat L_{S_t}(\widehat W_t+\epsilon_t)
% -
% \widehat L_{S_t}(\widehat W_t)
% \right]
% $
% is the task-conditioned excess sharpness under the perturbation law $\Pi_t$. 

% This decomposition is an exact identity, not a local Taylor approximation.





% We answer the first question by showing that task-conditioned perturbation-smoothed empirical risks enter a pathwise generalization bound over the posterior trajectory.
% This is the missing step in existing explanations of flatness in continual learning. Single-task PAC-Bayesian flatness analyses\cite{bahri-etal-2022-sharpness, mollenhoffSAMOptimalRelaxation2022} connect perturbation-smoothed empirical risk to the population risk of one task-induced posterior, while lifelong PAC-Bayesian analyses\cite{pentinaPACBayesianBoundLifelong2014} assume tasks sampled from a fixed environment distribution.
% Neither view directly captures the realized sequence of posterior updates produced by continual learning.
% Therefore, we need a sequential PAC-Bayesian formulation in which the prior-generating process evolves with the learning history and the bound holds along the realized task trajectory. 



% We model this process through a time-varying hyperposterior.
% Let $\mathcal F_{t-1}=\sigma(S_1,\ldots,S_{t-1})$ denote the information available before observing task $t$.
% Before task $t$, the hyperposterior $\mathcal P_{t-1}$ is $\mathcal F_{t-1}$-measurable and generates a task prior $P_t$.
% After observing $S_t$, the learner produces a task posterior $Q_t^\Pi(\cdot \mid P_t,S_t)$, where $\Pi_t$ denotes the perturbation distribution used to define the corresponding smoothed posterior.
% The hyperposterior is then updated from $\mathcal P_{t-1}$ to $\mathcal P_t$.
% This construction separates within-task adaptation from the history-dependent drift of the prior-generating mechanism.
% For each task $t$, we define the perturbation-smoothed empirical risk as
% $
% \widehat L_t^\Pi
% =
% \mathbb E_{P_t\sim \mathcal P_{t-1}}
% \mathbb E_{W_t\sim Q_t^\Pi(\cdot \mid P_t,S_t)}
% \widehat L_{S_t}(W_t),
% $
% and the corresponding task-conditioned population risk as
% $
% L_t^\Pi
% =
% \mathbb E_{P_t\sim \mathcal P_{t-1}}
% \mathbb E_{W_t\sim Q_t^\Pi(\cdot \mid P_t,S_t)}
% L_{\mathcal D_t}(W_t).
% $


\begin{theorem}[Sequential hierarchical PAC-Bayesian decomposition]
\label{thm:sequential_pac_bayes}
Assume that the loss is bounded in $[0,1]$.
For any $\delta\in(0,1]$, with probability at least $1-\delta$ over the task samples $S_{1:T}$, the following bound holds:
$$
\mathcal L_{1:T}^\Pi
\le
\widehat{\mathcal L}_{1:T}^\Pi
+
\sqrt{
\frac{
\sum_{t=1}^{T}
\mathbb E_{P_t\sim \mathcal P_{t-1}}
\mathrm{KL}\!\left(
Q_t^\Pi(\cdot \mid P_t,S_t)
\,\middle\|\,
P_t
\right)
+
\sum_{t=1}^{T}
\mathrm{KL}\!\left(
\mathcal P_t
\,\middle\|\,
\mathcal P_{t-1}
\right)
+
\log \frac{1}{\delta}
}{
2\sum_{t=1}^{T}(m_t-1)
}
}.
$$
\end{theorem}


Theorem~\ref{thm:sequential_pac_bayes} identifies how task-conditioned sharpness enters continual learning theory.
The first term on the right-hand side is the average perturbation-smoothed empirical risk along the posterior trajectory.
By the decomposition above, this term consists of clean empirical fit and task-conditioned excess sharpness at each update step.
This is the component most directly optimized by sharpness-aware learning.
The two complexity terms determine whether such locally flat solutions are compatible with sequential adaptation.
The task-level KL term measures the cost of adapting from the task prior $P_t$ to the perturbation-induced posterior $Q_t^\Pi$.
The hyperposterior KL term measures how much the prior-generating mechanism drifts across the realized task sequence.
Thus, flatness is not sufficient by itself.
Trajectory-level generalization depends jointly on low perturbation-smoothed empirical risk, moderate within-task adaptation complexity, and controlled hyperposterior drift.

This theorem answers Q1 at the level of trajectory-level generalization.
It does not claim that a flat solution at one checkpoint automatically prevents forgetting.
Rather, it shows that the perturbation-smoothed risks optimized during current-task learning are precisely the empirical terms that enter a sequential PAC-Bayesian bound over the learning trace.
Consequently, sharpness-aware updates can affect continual learning performance through the posterior states they produce and through the amount of adaptation and hyperposterior drift required along the sequence.

% Theorem~\ref{thm:sequential_pac_bayes} identifies how flatness enters continual learning theory.
% The first term on the right-hand side is the average perturbation-smoothed empirical risk optimized by sharpness-aware learning at individual tasks.
% The first KL term measures the cost of adapting from the task prior $P_t$ to the posterior $Q_t^\Pi$.
% The second KL term measures the drift of the hyperposterior process across the realized task sequence.
% Thus, task-conditioned sharpness does not appear as an isolated single-task quantity, but as one component of a pathwise bound over the posterior trajectory.

% This theorem answers Q1 at the level of trajectory-level generalization.
% It does not claim that a flat solution at one checkpoint automatically prevents forgetting.
% Rather, it shows that the perturbation-smoothed risks optimized during current-task learning are precisely the empirical terms that enter a sequential PAC-Bayesian bound over the learning trace.
% Consequently, sharpness-aware updates can affect continual learning performance through the posterior states they produce and through the amount of adaptation and hyperposterior drift required along the sequence.

% This result also clarifies the limitation of existing PAC-Bayesian formulations.
% If $\mathcal P_t=\mathcal P_{t-1}$ for all $t$, the hyperposterior drift term vanishes and the formulation reduces to an exchangeable-task view similar to lifelong or meta-learning analyses.
% Continual learning corresponds to the nonzero-drift regime, where the prior-generating process evolves with the realized history.
% Compared with sequential cumulative-risk bounds, Theorem~\ref{thm:sequential_pac_bayes} explicitly isolates the perturbation-smoothed empirical risk term optimized by sharpness-aware learning and places it in the same pathwise bound as within-task adaptation and hyperposterior drift.

% The proof follows a sequential change-of-measure argument combined with a martingale concentration bound and is provided in Appendix~\ref{app:proof_seq_pac_bayes}.




% Let $\mathcal F_{t-1}=\sigma(S_1,\ldots,S_{t-1})$ denote the information available before observing task $t$.
% We model continual learning through a time-varying hyperposterior process $\{\mathcal P_t\}_{t=0}^{T}$, where $\mathcal P_{t-1}$ is $\mathcal F_{t-1}$-measurable and generates the task prior $P_t$.
% After observing $S_t$, the learner produces a task posterior $Q_t(\cdot\mid P_t,S_t)$, and the hyperposterior is updated from $\mathcal P_{t-1}$ to $\mathcal P_t$.
% This hierarchy separates within-task adaptation from the history-dependent evolution of the prior-generating mechanism.
% Let $\mathcal F_{t-1}=\sigma(S_1,\ldots,S_{t-1})$ denote the information available before observing task $t$.
% We model continual learning through a time-varying hyperposterior process $\{\mathcal P_t\}_{t=0}^{T}$. Before observing task $t$, the hyperposterior $\mathcal P_{t-1}$ is $\mathcal F_{t-1}$-measurable and generates a task prior $P_t$. After observing $S_t$, the learner produces a task posterior $Q_t(\cdot \mid P_t,S_t)$, and the hyperposterior is updated from $\mathcal P_{t-1}$ to $\mathcal P_t$. This construction separates two sources of variation along the learning trajectory: the task posterior $Q_t$ captures within-task adaptation from the prior available before task $t$, while the hyperposterior process captures the history-dependent evolution of the prior-generating mechanism.



% This decomposition is designed to address the two sequential challenges identified above. Since the prior before task $t$ is produced from the previous learning history, it must be treated as an adapted process rather than as a fixed distribution sampled from an exchangeable task environment. Moreover, because performance is evaluated along the realized sequence of updates, the guarantee must accumulate pathwise over $S_{1:T}$. In this setting, the perturbation-smoothed empirical risk is the term through which task-conditioned sharpness enters the trajectory-level analysis, while the KL terms account for within-task adaptation and hyperposterior drift.
% For a perturbation distribution $\Pi_t$, let $Q_t^\Pi(\cdot\mid P_t,S_t)$ denote the posterior induced at task $t$. The perturbation-smoothed empirical risk is defined as
% $
% \widehat L_t^{\Pi}
% =
% \mathbb E_{P_t\sim \mathcal P_{t-1}}
% \mathbb E_{W_t\sim Q_t^\Pi(\cdot\mid P_t,S_t)}
% \widehat L_{S_t}(W_t).
% $
% The corresponding task-conditioned population risk is
% $
% L_t
% =
% \mathbb E_{P_t\sim \mathcal P_{t-1}}
% \mathbb E_{W_t\sim Q_t^\Pi(\cdot\mid P_t,S_t)}
% L_{\mathcal D_t}(W_t).
% $


% We model continual learning through a time-varying hyperposterior process $\{\mathcal P_t\}_{t=0}^{T}$, where $\mathcal P_{t-1}$ is $\mathcal F_{t-1}$-measurable and generates the task prior used before observing $S_t$. 
% A task prior $P_t$ is sampled from $\mathcal P_{t-1}$, and after observing $S_t$, the learner produces a task posterior $Q_t(\cdot\mid P_t,S_t)$. 
% The hyperposterior is then updated from $\mathcal P_{t-1}$ to $\mathcal P_t$.
% % Given a task prior $P_t\sim\mathcal P_{t-1}$, the learner produces a task posterior $Q_t(\cdot\mid P_t,S_t)$. 
% % The resulting bound separates the trajectory-averaged perturbation-smoothed empirical risk, the within-task adaptation complexity $\mathrm{KL}(Q_t\Vert P_t)$, and the hyperposterior drift complexity $\mathrm{KL}(\mathcal P_t\Vert\mathcal P_{t-1})$. 
% % When the hyperposterior does not drift, the formulation reduces to the exchangeable-task hierarchical setting; when the perturbation-smoothed term is removed, it reduces to a sequential PAC-Bayesian control of prediction risk. 
% The hierarchical view is useful  because it separates two distinct sources of change. 
% The task-level posterior $Q_t$ captures adaptation from the prior available before task $t$, whereas the hyperposterior process captures how this prior-generating mechanism evolves with the history $S_{1:t-1}$. 
% % This separation allows us to distinguish within-task adaptation complexity from hyperposterior drift complexity. 
% The perturbation-smoothed empirical risk then plays a specific role: it measures the local sensitivity of task losses along the posterior trajectory, and therefore determines how much sequential drift is amplified into old-task risk increase. 
% Thus, flatness is not added as an auxiliary term to a cumulative bound. 
% It is the PAC-Bayesian object that connects local loss geometry with the pathwise forgetting quantity $R_{i,t}$.
% Our focus is the intermediate regime specific to continual learning, where history-dependent posterior drift and local flatness jointly determine how current learning affects old-task risk.

% We now formulate a sequential hierarchical PAC-Bayesian view for continual learning. 
% Let $\mathcal F_{t-1}=\sigma(S_1,\ldots,S_{t-1})$ denote the information available before observing task $t$. 
% The task prior used at step $t$ is not treated as a fixed external object. 
% Instead, it is sampled from a history-dependent hyperposterior process. 
% Specifically, $\mathcal P_{t-1}$ denotes a hyperposterior over task-level priors before task $t$, and it is $\mathcal F_{t-1}$ measurable. 
% A task prior $P_t$ is sampled from $\mathcal P_{t-1}$, and after observing $S_t$, the learner produces a task posterior $Q_t(\cdot\mid P_t,S_t)$. 
% The hyperposterior is then updated from $\mathcal P_{t-1}$ to $\mathcal P_t$.

% For a perturbation distribution $\Pi_t$, the perturbation-smoothed empirical risk of posterior $Q_t$ on task $t$ is defined as
% \begin{equation}
% \widehat{\mathcal L}^{\Pi}_t
% =
% \mathbb E_{P_t\sim\mathcal P_{t-1}}
% \mathbb E_{W_t\sim Q_t(\cdot\mid P_t,S_t)}
% \widehat L_{S_t}(W_t).
% \end{equation}
% The corresponding hierarchical population risk is
% \begin{equation}
% \mathcal L_t
% =
% \mathbb E_{P_t\sim\mathcal P_{t-1}}
% \mathbb E_{W_t\sim Q_t(\cdot\mid P_t,S_t)}
% L_{\mathcal D_t}(W_t).
% \end{equation}

% We obtain the following pathwise decomposition.

% \begin{theorem}[Sequential hierarchical PAC-Bayesian decomposition]
% \label{thm:pathwise_pac}
% Assume the loss is bounded in $[0,1]$. 
% For any $\delta\in(0,1]$, with probability at least $1-\delta$ over the task samples $S_{1:T}$, the following bound holds:
% \begin{equation}
% \frac{1}{T}
% \sum_{t=1}^T
% \mathcal L_t
% \le
% \frac{1}{T}
% \sum_{t=1}^T
% \widehat{\mathcal L}^{\Pi}_t
% +
% \sqrt{
% \frac{
% \sum_{t=1}^T
% \mathbb E_{P_t\sim\mathcal P_{t-1}}
% \mathrm{KL}\bigl(Q_t(\cdot\mid P_t,S_t)\Vert P_t\bigr)
% +
% \sum_{t=1}^T
% \mathrm{KL}\bigl(\mathcal P_t\Vert\mathcal P_{t-1}\bigr)
% +
% \log\frac{1}{\delta}
% }
% {
% 2\sum_{t=1}^T(m_t-1)
% }
% }.
% \end{equation}
% \end{theorem}

% This decomposition clarifies how flatness enters continual learning theory. Sharpness-aware learning acts on the perturbation-smoothed empirical term at each task, and Theorem~4.1 shows that these task-conditioned terms control the sequential generalization components along the posterior trajectory, together with two complexity costs. The first KL term measures the cost of adapting from the task prior $P_t$ to the posterior $Q_t$, while the second KL term measures the drift of the hyperposterior process across tasks.

% The result also clarifies the relation to existing PAC-Bayesian analyses. If $\mathcal P_t=\mathcal P_{t-1}$ for all $t$, the hyperposterior drift term vanishes, and the formulation reduces to an exchangeable-task hierarchical view similar to lifelong or meta-learning analyses. Continual learning corresponds to the nonzero-drift regime, where the prior-generating process evolves with the realized history. Compared with sequential cumulative-risk bounds, Theorem~4.1 explicitly isolates the perturbation-smoothed empirical term optimized by sharpness-aware learning and places it in the same pathwise bound as within-task adaptation and hyperposterior drift.

% The proof follows a sequential change-of-measure argument combined with a martingale concentration bound, and is provided in Appendix~??.


% Theorem~\ref{thm:pathwise_pac} decomposes continual-learning generalization into three terms. 
% The first term is the perturbation-smoothed empirical risk, which is the quantity directly controlled by sharpness-aware learning. 
% The second term,
% \begin{equation}
% \sum_{t=1}^T
% \mathbb E_{P_t\sim\mathcal P_{t-1}}
% \mathrm{KL}\bigl(Q_t(\cdot\mid P_t,S_t)\Vert P_t\bigr),
% \end{equation}
% measures within-task adaptation complexity. 
% It quantifies how far the task posterior must move away from the task prior in order to fit $S_t$. 
% The third term,
% \begin{equation}
% \sum_{t=1}^T
% \mathrm{KL}\bigl(\mathcal P_t\Vert\mathcal P_{t-1}\bigr),
% \end{equation}
% measures hyperposterior drift. 
% It quantifies how much the mechanism that generates task priors changes as new tasks are incorporated.

% This decomposition explains why flatness is relevant to continual learning. 
% Sharpness-aware learning can reduce the perturbation-smoothed empirical term, but forgetting also depends on how the posterior trajectory moves across tasks. 
% A flat solution is therefore not sufficient by itself. 
% It is useful when the smoothed risk remains low along the trajectory and when the induced posterior movement does not require excessive task-level adaptation or hyperposterior drift.

% The result also clarifies the relation to existing PAC-Bayesian analyses. 
% If $\mathcal P_t=\mathcal P_{t-1}$ for all $t$, the hyperposterior drift term vanishes, and the formulation reduces to the exchangeable-task hierarchical view used in lifelong or meta-learning analyses. 
% In continual learning, however, this zero-drift assumption is generally too restrictive, because the prior used at task $t$ is produced by the previous learning history. 
% Conversely, sequential cumulative-risk bounds account for task order, but they do not isolate the perturbation-smoothed empirical risk term that sharpness-aware learning optimizes. 
% Theorem~\ref{thm:pathwise_pac} combines these two aspects by making flatness and posterior drift appear in the same pathwise bound.
% A direct implication is that forgetting can be interpreted through two coupled effects. 
% For an earlier task $i<t$, the risk-form forgetting
% \begin{equation}
% \mathfrak F_{i,t}
% =
% L_{\mathcal D_i}(Q_t)
% -
% L_{\mathcal D_i}(Q_i)
% \end{equation}
% is controlled by how much the smoothed empirical risk of task $i$ increases under the later posterior, together with the complexity accumulated by moving from the earlier posterior to the later posterior. 
% Thus, flatness matters for continual learning because it controls the local stability of empirical risk under perturbations, while the KL terms quantify the cost of maintaining this stability along a sequential posterior path.

% The proof of Theorem~\ref{thm:pathwise_pac} follows a sequential change-of-measure argument together with a martingale concentration bound, and is provided in Appendix~\ref{app:proof_pathwise_pac}.

% \begin{corollary}[Seen-task risk, forgetting, and perturbation-smoothed risk]
% \label{cor:seen_task_forgetting}
% Define the trajectory-averaged seen-task risk as
% \begin{equation}
% \mathcal R_{\mathrm{seen}}
% =
% \frac{1}{T}
% \sum_{t=1}^{T}
% \frac{1}{t}
% \sum_{i=1}^{t}
% L_{\mathcal D_i}(Q_t)
% =
% \frac{1}{T}
% \sum_{t=1}^T
% \mathcal R_t^{\mathrm{seen}}.
% \label{eq:seen_task_risk}
% \end{equation}
% This quantity includes both diagonal terms $R_{i,i}$ (current-task generalization) and off-diagonal terms $R_{i,t}$ with $i<t$ (old-task population risk under later posteriors).
% Risk-form forgetting appears as the off-diagonal increment:
% \begin{equation}
% \mathcal R_{\mathrm{seen}}
% =
% \frac{1}{T}\sum_{t=1}^T R_{t,t}
% +
% \frac{1}{T}\sum_{t=1}^T \frac{1}{t}\sum_{i=1}^{t-1}
% \underbrace{(R_{i,t}-R_{i,i})}_{\mathcal F_{i,t}\ (\text{forgetting})}
% +
% \frac{1}{T}\sum_{t=1}^T \frac{1}{t}\sum_{i=1}^{t-1}
% R_{i,i}.
% \label{eq:seen_task_forgetting_decomp}
% \end{equation}
% In replay-free continual learning, the posterior $Q_t$ is learned from $(S_t, S_{1:t-1}^{\mathrm{buffer}})$ and is approximately independent of $S_i$ for $i < t$.
% Applying the standard PAC-Bayesian bound to each pair $(Q_t, S_i)$ with $Q_t$ as the evaluated posterior and $Q_i$ as the reference prior, with probability at least $1-\delta$ over the draws of $S_{1:T}$,
% \begin{equation}
% R_{i,t}
% \le
% \widehat L_{S_i}^{\Pi_t}(Q_t)
% +
% \sqrt{
% \frac{
% \mathrm{KL}(Q_t\Vert Q_i) + \log\frac{1}{\delta}
% }{
% 2(m_i-1)
% }
% },
% \qquad i \le t,
% \label{eq:off_diagonal_bound}
% \end{equation}
% where $\widehat L_{S_i}^{\Pi_t}(Q_t) = \mathbb E_{W\sim Q_t}\mathbb E_{\epsilon\sim\Pi_t}\widehat L_{S_i}(W+\epsilon)$ is the perturbation-smoothed empirical risk of old dataset $S_i$ evaluated at the later posterior $Q_t$.
% Eq.~\ref{eq:off_diagonal_bound} shows that forgetting $\mathcal F_{i,t}$ is bounded by the increase in perturbation-smoothed old-task risk under the new posterior, plus a KL complexity term $\mathrm{KL}(Q_t\Vert Q_i)$ that reflects the total posterior movement from task $i$ to task $t$.
% \end{corollary}

% Corollary~\ref{cor:seen_task_forgetting} closes the loop between the PAC-Bayesian framework and forgetting.
% The same perturbation-smoothed empirical risk that Theorem~\ref{thm:pathwise_pac} identifies as the bound-relevant quantity for current-task generalization also controls old-task risk under later posteriors.
% Specifically, $\widehat L_{S_i}^{\Pi_t}(Q_t)$ measures how stably old-task $i$'s loss behaves under the later posterior $Q_t$ and its associated perturbations.
% When the old-task loss remains insensitive to perturbations along the directions explored by $Q_t$, this term stays close to $\widehat L_{S_i}(Q_t)$, which in turn bounds $R_{i,t}$.
% This provides the theoretical grounding for why sharpness-aware training, which reduces the perturbation-smoothed empirical risk on the current task, may also reduce forgetting: if the same loss geometry is locally stable under all posterior perturbations along the task trajectory, the off-diagonal risks $R_{i,t}$ are better controlled.

% \subsection{The Missing Link in Continual Learning}
% \label{sec:missing_link_cl}

% Continual learning differs from the standard PAC-Bayesian setting because the prior at task $t$ is not a fixed external object. It is produced by the previous history $S_{1:t-1}$. This history dependence is inherent to many CL mechanisms. 

% % Curvature-based methods update the effective prior using Fisher or Hessian summaries of earlier tasks~\citep{EWC2017,SI2017}, while scope-allocation methods change which parameters remain plastic over time~\citep{progressive2016,HAT2018,Supermasks2020,pmlr-v97-houlsby19a}. 
% Hence the prior sequence $(P_t)_{t=1}^T$ is a nonstationary process, and its drift is part of the forgetting mechanism.

% Existing PAC-Bayesian analyses provide the necessary ingredients, but they do not directly answer the flatness-in-continual-learning question above.
% Bayesian analyses explain perturbation-smoothed empirical risk in single-task learning, but not posterior drift \cite{bahri-etal-2022-sharpness,mollenhoffSAMOptimalRelaxation2022}. Hierarchical PAC-Bayesian bound\cite{pentinaPACBayesianBoundLifelong2014} introduce a hyperposterior over task-level priors, which is useful for separating transferable inductive bias from task-specific adaptation. However, the standard formulation assumes that tasks are exchangeable draws from a fixed hyperdistribution. The resulting bound is order-invariant and does not describe how the inductive bias changes along a realized task sequence. Recent PAC-Bayesian analyses for sequential continual learning\cite{friedman2026pac} move closer to this setting by considering data-dependent priors and cumulative losses. However, they are not designed to isolate perturbation-smoothed empirical risk as the mechanism through which sharpness-aware objectives affect forgetting. Moreover, without an explicit hyperposterior level, they do not separate task-specific posterior adaptation from the evolution of transferable inductive bias.

% We therefore require a decomposition with three properties. It must be sequential, because forgetting is defined along a posterior trajectory. It must be hierarchical, because continual learning updates both task-level posteriors and the transferable inductive bias that generates future priors. Finally, it must retain the perturbation-smoothed





% % \subsection{Limitations of Existing PAC-Bayesian Views}
% % \label{sec:limitations_existing_pac_cl}

% % Existing PAC-Bayesian results provide important ingredients, but they do not directly answer the flatness-in-continual-learning question considered here. PAC-Bayesian flatness analyses justify sharpness-aware objectives by replacing empirical risk with perturbation-smoothed empirical risk. These results explain why flatter solutions can generalize better in a single-task setting, but they do not model the posterior drift responsible for forgetting.

% % Hierarchical PAC-Bayesian analyses for lifelong or meta-learning introduce a hyperposterior over task-level priors, which is conceptually useful for separating transferable inductive bias from task-specific adaptation. However, the standard formulation assumes that tasks are sampled exchangeably from a fixed hyperdistribution. Consequently, the resulting bound is order-invariant and does not describe how the inductive bias changes along a realized task sequence. In our framework, this setting corresponds to the special case in which the hyperposterior drift is zero. Continual learning requires the more general case where the prior used at task $t$ is produced by the history $S_{1:t-1}$ and may therefore evolve over time.

% % Recent PAC-Bayesian analyses for sequential continual learning move closer to this setting by considering data-dependent priors and cumulative loss along a task sequence. These results are important for certifying learning plasticity, but they are not designed to explain flatness. They bound cumulative generalization loss, rather than isolating perturbation-smoothed empirical risk as the mechanism through which sharpness-aware objectives affect forgetting. Moreover, they do not explicitly separate the hyper-level evolution of inductive bias from task-level posterior adaptation.

% % These limitations indicate that explaining flatness in continual learning requires a framework with three properties. First, it must be sequential, so that posterior trajectories and history-dependent priors are explicitly modeled. Second, it must be hierarchical, so that task-level adaptation and the evolution of transferable inductive bias are separated. Third, it must retain the perturbation-smoothed empirical risk term, since this is the formal object optimized by sharpness-aware learning. We develop such a decomposition next.








% % \section{Why Flatness Matters for Continual Learning}




% \section{From Flatness-Based Generalization to Forgetting}


% \section{Sequential Hierarchical PAC-Bayesian Bound for Continual Learning }
% \subsection{Question I: Can Perturbation-Smoothed Risk Explain Forgetting?}

% The PAC-Bayesian view relates the population risk of a posterior to its empirical risk and a KL complexity term. 
% In standard supervised learning, this gives a single-posterior statement: a posterior $Q$ learned from a sample $S$ is evaluated on the same underlying distribution $\mathcal D$.
% Continual learning changes the object of analysis. 
% The learner produces a posterior trajectory $Q_1,\ldots,Q_T$, and a later posterior $Q_t$ must remain reliable not only on the current task distribution $\mathcal D_t$, but also on previous task distributions $\mathcal D_i$ for $i<t$.

% This requires a pathwise view of the task sequence.
% Within each task, the sample $S_t$ is drawn i.i.d. from $\mathcal D_t$, but the sequence of task distributions $\mathcal D_1,\ldots,\mathcal D_T$ is treated as a realized continual-learning path rather than as an exchangeable collection of tasks from a fixed environment.
% Accordingly, the relevant prior at step $t$ may depend on the learning history before task $t$, while the generalization question concerns how the posterior trajectory behaves on the expanding set of encountered task distributions.

% Under this view, forgetting is the degradation of past-task generalization under later posteriors.
% For an earlier task $i<t$, its risk-form forgetting after learning task $t$ is
% $$
% \mathfrak F_{i,t}
% =
% L_{\mathcal D_i}(Q_t)
% -
% L_{\mathcal D_i}(Q_i).
% $$
% This quantity compares the population risk on the same earlier distribution $\mathcal D_i$ before and after subsequent updates.
% It is therefore the risk-based analogue of backward-transfer style forgetting metrics, up to the sign convention and the use of risk rather than accuracy.

% This perspective also clarifies trajectory-level continual learning metrics.
% Average-accuracy metrics evaluate each intermediate model on the tasks observed so far and then aggregate performance across stages.
% Their risk-based counterpart asks whether the posterior trajectory maintains low population risk over the expanding set of task distributions.
% Thus, explaining continual learning performance requires more than controlling the current-task risk $L_{\mathcal D_t}(Q_t)$; it requires controlling old-task risks $L_{\mathcal D_i}(Q_t)$ under later posteriors.

% Perturbation-smoothed empirical risk is central to this question because SAM-style objectives optimize empirical loss under local posterior perturbations.
% However, existing PAC-Bayesian analyses do not directly show how such smoothed empirical risk controls past-task population risk under later posteriors.
% Classical lifelong PAC-Bayes mainly studies transfer under shared or exchangeable task environments, while recent cumulative-risk analyses do not isolate how old-task smoothed empirical risk and posterior drift jointly contribute to forgetting along a realized task sequence.

% This leads to our first theoretical question:

% \textbf{Q1.}
% Can one derive a pathwise PAC-Bayesian bound for continual learning in which past-task population risks under later posteriors are controlled by perturbation-smoothed empirical risks and posterior-trajectory complexity?

% Section~\ref{sec:seq-pac} answers this question through a sequential hierarchical PAC-Bayesian decomposition with history-dependent task priors.
% The resulting bound separates within-task adaptation complexity from cross-task posterior drift complexity, thereby connecting sharpness-based generalization and forgetting through a shared posterior-trajectory view.

% \section{Which Flatness Matters in LoRA-Based PECL}
% \label{sec:pecl_support}

% % Theorem~\ref{thm:pathwise_pac} shows that perturbation-smoothed empirical risk and KL complexity are the two bound-relevant quantities through which flatness enters the analysis of a continual learning trajectory. We now identify the geometric support of these quantities in frozen-backbone PECL. The key point is that LoRA does not directly update the ambient parameter space. Instead, trainable LoRA coordinates induce an effective adapter update, while the pretrained backbone and other frozen components remain fixed. Therefore, the relevant flatness object should be defined on the effective update geometry reachable by the current adapter.

% \subsection{Effective Update Posterior under Frozen-Backbone LoRA}
% \label{sec:effective_update_posterior}

% At task $t$, let $W_t^{\mathrm{froz}}$ denote the parameters that are fixed before learning $S_t$. This includes the pretrained backbone and, depending on the PECL method, previously frozen adapter components. The trainable LoRA coordinates are denoted by $\theta_t\in\Theta_t$, and the corresponding effective adapter update is given by a task-dependent map
% $$
% \phi_t:\Theta_t\rightarrow \mathbb R^d,
% \qquad
% \Delta W_t=\phi_t(\theta_t).
% $$
% Thus the model parameter used at task $t$ can be written as
% $$
% W_t
% =
% W_t^{\mathrm{froz}}+\Delta W_t.
% $$
% The admissible effective update geometry is
% $$
% \mathcal W_{\Delta,t}^{\mathrm{train}}
% =
% \{\phi_t(\theta_t):\theta_t\in\Theta_t\}.
% $$
% All effective update distributions are viewed as probability measures on the ambient parameter space $\mathbb R^d$, with support restricted to $\mathcal W_{\Delta,t}^{\mathrm{train}}$.

% A posterior over raw LoRA coordinates $Q_t^\theta$ induces a pushforward posterior over effective updates,
% $$
% Q_t^\Delta
% =
% (\phi_t)_\# Q_t^\theta,
% $$
% meaning that for any measurable set $A\subseteq\mathbb R^d$,
% $$
% Q_t^\Delta(A)
% =
% Q_t^\theta\bigl(\{\theta_t:\phi_t(\theta_t)\in A\}\bigr).
% $$
% Similarly, a prior $P_t^\theta$ over LoRA coordinates induces an effective-update prior
% $$
% P_t^\Delta
% =
% (\phi_t)_\# P_t^\theta.
% $$

% This pushforward view is important because LoRA factors are not identifiable. Different factor coordinates may induce the same effective update $\Delta W_t$. Hence, quantities defined directly on raw factors may depend on the chosen parameterization, whereas quantities defined on $\Delta W_t$ are invariant to equivalent LoRA factor representations.

% The perturbation law can also be expressed directly in effective-update space. If the learner obtains a deterministic effective update $\widehat{\Delta W}_t$, then sampling a local perturbation
% $$
% \epsilon_t\sim \Pi_{\Delta,t}
% $$
% induces a local posterior over effective updates through
% $$
% \Delta W_t
% =
% \widehat{\Delta W}_t+\epsilon_t.
% $$
% Thus, in the PAC-Bayesian view, the perturbation distribution is a way of specifying a local posterior around the learned effective update.

% For a task prior $P_t^\Delta$ sampled from the hyperposterior $\mathcal P_{\Delta,t-1}$, the perturbation-smoothed empirical risk in effective-update space is
% $$
% \widehat{\mathcal L}_{\Delta,t}^{\Pi}
% =
% \mathbb E_{P_t^\Delta\sim \mathcal P_{\Delta,t-1}}
% \mathbb E_{\Delta W_t\sim Q_t^\Delta(\cdot\mid P_t^\Delta,S_t)}
% \widehat L_{S_t}
% \left(
% W_t^{\mathrm{froz}}+\Delta W_t
% \right).
% $$
% The corresponding task-conditioned population risk is
% $$
% \mathcal L_{\Delta,t}^{\Pi}
% =
% \mathbb E_{P_t^\Delta\sim \mathcal P_{\Delta,t-1}}
% \mathbb E_{\Delta W_t\sim Q_t^\Delta(\cdot\mid P_t^\Delta,S_t)}
% L_{\mathcal D_t}
% \left(
% W_t^{\mathrm{froz}}+\Delta W_t
% \right).
% $$

% For a deterministic effective update $\widehat{\Delta W}_t$, the expected adapter-update sharpness is
% $$
% \mathrm{ES}_{\Delta,t}^{\Pi}
% (\widehat{\Delta W}_t)
% =
% \mathbb E_{\epsilon_t\sim\Pi_{\Delta,t}}
% \left[
% \widehat L_{S_t}
% \left(
% W_t^{\mathrm{froz}}
% +
% \widehat{\Delta W}_t
% +
% \epsilon_t
% \right)
% -
% \widehat L_{S_t}
% \left(
% W_t^{\mathrm{froz}}
% +
% \widehat{\Delta W}_t
% \right)
% \right].
% $$
% The corresponding adversarial zeroth-order adapter-update sharpness is
% $$
% \mathrm{AS}_{\Delta,t}^{(0)}
% (\widehat{\Delta W}_t,\rho)
% =
% \max_{\epsilon_t\in\mathcal W_{\Delta,t}^{\mathrm{train}},\ \|\epsilon_t\|\le \rho}
% \left[
% \widehat L_{S_t}
% \left(
% W_t^{\mathrm{froz}}
% +
% \widehat{\Delta W}_t
% +
% \epsilon_t
% \right)
% -
% \widehat L_{S_t}
% \left(
% W_t^{\mathrm{froz}}
% +
% \widehat{\Delta W}_t
% \right)
% \right].
% $$
% These definitions make the perturbation support explicit: the perturbation is applied to the effective adapter update, not to arbitrary full-model directions or non-identifiable raw LoRA factors.

% \begin{proposition}[Effective-update reduction]
% \label{prop:effective_update_reduction}
% Assume that the task loss depends on the LoRA coordinates only through the effective parameter
% $$
% W_t^{\mathrm{froz}}+\phi_t(\theta_t).
% $$
% Let $Q_t^\Delta=(\phi_t)_\#Q_t^\theta$ and $P_t^\Delta=(\phi_t)_\#P_t^\theta$. Then, for any task sample $S_t$ and distribution $\mathcal D_t$,
% $$
% \mathbb E_{\theta_t\sim Q_t^\theta}
% \widehat L_{S_t}
% \left(
% W_t^{\mathrm{froz}}+\phi_t(\theta_t)
% \right)
% =
% \mathbb E_{\Delta W_t\sim Q_t^\Delta}
% \widehat L_{S_t}
% \left(
% W_t^{\mathrm{froz}}+\Delta W_t
% \right),
% $$
% and
% $$
% \mathbb E_{\theta_t\sim Q_t^\theta}
% L_{\mathcal D_t}
% \left(
% W_t^{\mathrm{froz}}+\phi_t(\theta_t)
% \right)
% =
% \mathbb E_{\Delta W_t\sim Q_t^\Delta}
% L_{\mathcal D_t}
% \left(
% W_t^{\mathrm{froz}}+\Delta W_t
% \right).
% $$
% Moreover, the pushforward KL satisfies
% $$
% \mathrm{KL}
% \left(
% Q_t^\Delta\Vert P_t^\Delta
% \right)
% \le
% \mathrm{KL}
% \left(
% Q_t^\theta\Vert P_t^\theta
% \right),
% $$
% whenever both sides are well-defined.
% \end{proposition}

% Proposition~\ref{prop:effective_update_reduction} shows that the empirical and population risk terms are exactly determined by the effective-update posterior. It also shows that a raw-factor KL may contain parameterization-dependent information that disappears after mapping to effective updates. Therefore, the PAC-Bayesian analysis should be formulated directly on the effective-update posterior when the goal is to identify the flatness object relevant to LoRA-constrained adaptation.

% \subsection{Support-Reduced PAC-Bayesian Bound for PECL}
% \label{sec:pecl_support_bound}

% Applying Theorem~\ref{thm:pathwise_pac} to the effective-update hypothesis space yields the following PECL-specific bound.

% \begin{theorem}[Support-reduced PAC-Bayesian bound for PECL]
% \label{thm:pecl_support_reduction}
% Assume the loss is bounded in $[0,1]$. For any $\delta\in(0,1]$, with probability at least $1-\delta$ over the task samples $S_{1:T}$, the following bound holds:
% $$
% \frac{1}{T}
% \sum_{t=1}^{T}
% \mathcal L_{\Delta,t}^{\Pi}
% \le
% \frac{1}{T}
% \sum_{t=1}^{T}
% \widehat{\mathcal L}_{\Delta,t}^{\Pi}
% +
% \sqrt{
% \frac{
% \sum_{t=1}^{T}
% \mathbb E_{P_t^\Delta\sim\mathcal P_{\Delta,t-1}}
% \mathrm{KL}
% \left(
% Q_t^\Delta(\cdot\mid P_t^\Delta,S_t)
% \Vert
% P_t^\Delta
% \right)
% +
% \sum_{t=1}^{T}
% \mathrm{KL}
% \left(
% \mathcal P_{\Delta,t}
% \Vert
% \mathcal P_{\Delta,t-1}
% \right)
% +
% \log\frac{1}{\delta}
% }
% {
% 2\sum_{t=1}^{T}(m_t-1)
% }
% }.
% $$
% \end{theorem}

% Theorem~\ref{thm:pecl_support_reduction} identifies the support on which flatness enters the PECL bound. Both the perturbation-smoothed empirical risk and the task-level KL complexity are evaluated on the admissible effective adapter-update geometry $\mathcal W_{\Delta,t}^{\mathrm{train}}$. The frozen backbone still determines the loss through the offset $W_t^{\mathrm{froz}}$, but the stochastic quantity controlled by the posterior is the adapter-induced effective update $\Delta W_t$.

% This result gives a precise answer to the support question raised in the introduction. In LoRA-based PECL, the bound-relevant flatness object is not ambient full-parameter sharpness, because full-space perturbations may probe directions that future LoRA updates cannot realize. It is also not raw-factor sharpness, because raw LoRA factors are non-identifiable and basis dependent. The relevant object is basis-invariant effective adapter-update sharpness, measured over the update geometry reachable by the current adapter.

% The theorem also clarifies how the PECL result relates to the sequential decomposition in Theorem~\ref{thm:pathwise_pac}. Theorem~\ref{thm:pathwise_pac} explains why perturbation-smoothed risks and posterior drift jointly control the trajectory-level sequential generalization components. Theorem~\ref{thm:pecl_support_reduction} specifies where these quantities live under frozen-backbone LoRA adaptation. Thus, Chapter~\ref{sec:flatness_cl} establishes the pathwise role of flatness in continual learning, while the present chapter identifies the correct flatness support in PECL.

% The proof of Theorem~\ref{thm:pecl_support_reduction} follows by applying the sequential hierarchical PAC-Bayesian decomposition to the effective-update hypothesis space and using Proposition~\ref{prop:effective_update_reduction}. Details are provided in Appendix~??.

% \section{A PECL Support-Reduction Answer to Q2}
% \label{sec:pecl_support_reduction}

% \section{A PECL Support-Reduction Answer to Q2}
% \label{sec:pecl_support_reduction}


\section{A PECL Support-Reduction Answer to Q2}

Section 4 shows that perturbation-smoothed empirical risks enter a pathwise PAC-Bayesian bound over the posterior trajectory. We now identify the perturbation support on which these quantities should be evaluated in LoRA-based PECL. This is the point at which PECL differs from the standard geometric explanation of flatness. Classical accounts usually treat flatness as a property of the ambient parameter basin. In contrast, our analysis identifies a support-specific notion of flatness: in frozen-backbone LoRA PECL, the bound-relevant sharpness is the task-local curvature of the effective adapter update $\Delta W_t$ along the admissible perturbation support.

At task $t$, write the model as
$
W_t = W_t^{\mathrm{froz}} + \Delta W_t,
$
where $W_t^{\mathrm{froz}}$ is frozen and $\Delta W_t= A_tB_t$ is the effective adapter update induced by the trainable LoRA factors. Let $\mathcal W_{\Delta,t}$ denote the admissible effective-update geometry and let $\mathcal T_{\Delta,t}\subseteq \mathcal W_{\Delta,t}$ denote the local perturbation support around the learned update $\widehat{\Delta W}_t$. A PECL-admissible perturbation law is a distribution $\Pi_{\Delta,t}$ satisfying
$
\mathrm{supp}(\Pi_{\Delta,t})\subseteq \mathcal T_{\Delta,t}.
$
It induces the effective-update posterior
$$
Q_t^\Delta = \widehat{\Delta W}_t+\Pi_{\Delta,t},
\quad
\Delta W_t=\widehat{\Delta W}_t+\epsilon_t,
\quad
\epsilon_t\sim \Pi_{\Delta,t}.
$$
The corresponding PECL perturbation-smoothed empirical risk is
$
\widehat L_{S_t}^{\Delta}
(
\widehat{\Delta W}_t,\Pi_{\Delta,t}
)
=
\mathbb E_{\epsilon_t\sim\Pi_{\Delta,t}}
[
\widehat L_{S_t}
(
W_t^{\mathrm{froz}}+\widehat{\Delta W}_t+\epsilon_t
)
],
$
and the associated adapter-update sharpness is
$$
TS_{\Delta,t}
(
\widehat{\Delta W}_t,\Pi_{\Delta,t}
)
=
\mathbb E_{\epsilon_t\sim\Pi_{\Delta,t}}
[
\widehat L_{S_t}
(
W_t^{\mathrm{froz}}+\widehat{\Delta W}_t+\epsilon_t
)
-
\widehat L_{S_t}
(
W_t^{\mathrm{froz}}+\widehat{\Delta W}_t
)
].
$$
For small perturbations, when the first-order term is negligible, this term admits the local quadratic approximation
$$
TS_{\Delta,t}
(
\widehat{\Delta W}_t,\Pi_{\Delta,t}
)
\approx
\frac{1}{2}
\mathbb E_{\epsilon_t\sim\Pi_{\Delta,t}}
\left[
\mathrm{vec}(\epsilon_t)^\top
H_{\Delta,t}
\mathrm{vec}(\epsilon_t)
\right],
$$
where
$
H_{\Delta,t}
=
\nabla_{\mathrm{vec}(\Delta W)}^2
\widehat L_{S_t}
(
W_t^{\mathrm{froz}}+\Delta W
)
\big|_{\Delta W=\widehat{\Delta W}_t}.
$
Thus, the relevant curvature is not ambient full-parameter curvature, but curvature restricted to the locally admissible effective adapter-update directions.

This distinction also affects the complexity term. Let $T_t(\Delta)=W_t^{\mathrm{froz}}+\Delta$ be the translation from effective updates to full model weights. For an effective-update posterior $Q_t^\Delta$ and prior $P_t^\Delta$, define the induced model-level distributions as
$
Q_t=(T_t)_\# Q_t^\Delta,
\quad
P_t=(T_t)_\# P_t^\Delta.
$

\textbf{Lemma 5.1.}
Assume $Q_t^\Delta\ll P_t^\Delta$. Then $Q_t\ll P_t$ and
$$
\mathrm{KL}(Q_t\|P_t)
=
\mathrm{KL}(Q_t^\Delta\|P_t^\Delta).
$$

Lemma 5.1 shows that, once the frozen backbone is fixed, the task-level KL complexity is exactly the divergence between distributions over effective adapter updates, as shown in Fig\ref{fig:showPAC_combined}(b).
% Raw factor-space sharpness over $(A_t,B_t)$ is not canonical because the same effective update may admit multiple factor representations, for example $A_tB_t=(A_tR)(R^{-1}B_t)$. Therefore, factor-space perturbations and norms can be representation-dependent, whereas $\Delta W_t$ defines the basis-invariant effective update.
Combining Lemma 5.1 with Theorem 4.1 gives the PECL support-reduction result.

% \textbf{Theorem 5.2.}
\begin{theorem}[PECL support-reduction bound]
\label{thm:pecl-bound}
Under the setup above, assume that the loss is bounded in $[0,1]$.
For any $\delta\in(0,1]$, with probability at least $1-\delta$ over the task samples $S_{1:T}$,
$$
\begin{aligned}
\mathcal L_{1:T}^{\Delta}
\le
&
\frac{1}{T}
\sum_{t=1}^{T}
\mathbb E_{P_t^\Delta\sim\mathcal P_{t-1}}
\left[
\widehat L_{S_t}
\left(
W_t^{\mathrm{froz}}
+
\widehat{\Delta W}_t
\right)
+
\mathrm{TS}_t^\Delta
\left(
\widehat{\Delta W}_t,
\Pi_{\Delta,t}
\right)
\right]
\\
&+
\sqrt{
\frac{
\sum_{t=1}^{T}
\mathbb E_{P_t^\Delta\sim \mathcal P_{t-1}}
\mathrm{KL}
\left(
Q_t^\Delta
\,\Vert\,
P_t^\Delta
\right)
+
\sum_{t=1}^{T}
\mathrm{KL}
\left(
\mathcal P_t
\,\Vert\,
\mathcal P_{t-1}
\right)
+
\log \frac{1}{\delta}
}{
2\sum_{t=1}^{T}(m_t-1)
}
},
\end{aligned}
$$
where
$
\mathcal L_{1:T}^{\Delta}
=
\frac{1}{T}
\sum_{t=1}^{T}
\mathbb E_{P_t^\Delta\sim\mathcal P_{t-1}}
\mathbb E_{\Delta W_t\sim Q_t^\Delta}
L_{\mathcal D_t}
\left(
W_t^{\mathrm{froz}}
+
\Delta W_t
\right).
$
\end{theorem}



Under frozen-backbone LoRA adaptation, the perturbation-smoothed empirical term and the task-level KL complexity in the pathwise PAC-Bayesian bound reduce to quantities supported on the effective adapter-update geometry. In particular, the empirical term decomposes into the clean empirical loss and $TS_{\Delta,t}$, while the task-level complexity becomes $\mathrm{KL}(Q_t^\Delta\|P_t^\Delta)$.

Theorem 5.2 refines the classical flat-basin explanation rather than rejecting it. Both views rely on local curvature, but they differ in the geometric support on which curvature is evaluated. The classical view implicitly uses ambient parameter-space flatness. In frozen-backbone LoRA PECL, however, the posterior trajectory can only perturb, optimize, and reuse directions induced by the effective adapter update. Hence, the bound-relevant flatness object is neither full-space sharpness nor raw factor-space sharpness, but task-local effective adapter-update sharpness.



% \section{A PECL Support-Reduction Answer to Q2}

% Section 4 shows that perturbation-smoothed empirical risks enter a pathwise PAC-Bayesian bound over the posterior trajectory. We now identify the perturbation support on which these quantities should be evaluated in LoRA-based PECL. This is the point at which PECL differs from the standard geometric explanation of flatness. Classical accounts usually treat flatness as a property of the ambient parameter basin. In contrast, our analysis identifies a support-specific notion of flatness: in frozen-backbone LoRA PECL, the bound-relevant sharpness is the task-local curvature of the effective adapter update $\Delta W_t$ along the admissible perturbation support.

% % % In LoRA-based PECL, the support of the perturbation law is no longer self-evident.、
% % Section~\ref{sec:pathwise_pac_bayes_q1} answers Q1 by showing that perturbation-smoothed empirical risks enter a trajectory-level PAC-Bayesian bound.
% % We now address Q2 to identify the perturbation support on which the smoothed empirical risk and the task-level KL complexity should be evaluated in LoRA-based PECL.


% % identify the perturbation support on which both the smoothed empirical risk and the task-level KL complexity in Theorem~\ref{thm:sequential_pac_bayes} should be evaluated.



% At task $t$, the model is written as
% $
% W_t
% =
% W_t^{\mathrm{froz}}
% +
% \Delta W_t,
% $
% where $W_t^{\mathrm{froz}}$ is the frozen backbone and $\Delta W_t=A_tB_t$ is the effective adapter update induced by the trainable LoRA coordinates. 
% Let $\mathcal W_{\Delta,t}$ denote the admissible effective update geometry, and let $\mathcal T_{\Delta,t}\subseteq \mathcal W_{\Delta,t}$ denote the local perturbation support around the learned update $\widehat{\Delta W}_t$. Define the translation map from effective updates to model weights as
% $
% T_t(\Delta)
% =
% W_t^{\mathrm{froz}}
% +
% \Delta.
% $
% For an posterior $Q_t^\Delta$ and prior $P_t^\Delta$, the corresponding model-level distributions are
% $
% Q_t=(T_t)_\#Q_t^\Delta,
% \quad
% P_t=(T_t)_\#P_t^\Delta.
% $

% A PECL-admissible perturbation law is a distribution $\Pi_{\Delta,t}$ with
% $\mathrm{supp}(\Pi_{\Delta,t})\subseteq \mathcal T_{\Delta,t}$. It induces the
% effective-update posterior $Q_t^\Delta=\widehat{\Delta W}_t+\Pi_{\Delta,t}$,
% namely $\Delta W_t=\widehat{\Delta W}_t+\varepsilon_t$ with
% $\varepsilon_t\sim\Pi_{\Delta,t}$.
% The PECL perturbation-smoothed empirical risk is
% $
% \widehat L_{S_t}^{\Delta}(\widehat{\Delta W}_t,\Pi_{\Delta,t})
% =
% \mathbb E_{\varepsilon_t\sim\Pi_{\Delta,t}}
% [
% \widehat L_{S_t}(W_t^{\mathrm{froz}}+\widehat{\Delta W}_t+\varepsilon_t)
% ],
% $
% and the associated adapter-update sharpness is
% $
% \mathrm{TS}_{\Delta,t}(\widehat{\Delta W}_t,\Pi_{\Delta,t})
% =
% \mathbb E_{\varepsilon_t\sim\Pi_{\Delta,t}}
% [
% \widehat L_{S_t}(W_t^{\mathrm{froz}}+\widehat{\Delta W}_t+\varepsilon_t)
% -
% \widehat L_{S_t}(W_t^{\mathrm{froz}}+\widehat{\Delta W}_t)
% ].
% $
% For small perturbations, if the first-order term is negligible, this quantity admits a local quadratic approximation. Let
% $
% H_{\Delta,t}
% :=
% \nabla_{\mathrm{vec}(\Delta W)}^2
% \widehat L_{S_t}(W_t^{\mathrm{froz}}+\Delta W)
% \big|_{\Delta W=\widehat{\Delta W}_t}.
% $
% Then
% $
% \mathrm{TS}_{\Delta,t}(\widehat{\Delta W}_t,\Pi_{\Delta,t})
% \approx
% \frac{1}{2}
% \mathbb E_{\varepsilon_t\sim\Pi_{\Delta,t}}
% [
% \mathrm{vec}(\varepsilon_t)^\top
% H_{\Delta,t}
% \mathrm{vec}(\varepsilon_t)
% ].
% $
% Thus, the relevant curvature is not measured in the full parameter space, but along the locally admissible effective adapter-update directions.




% % A PECL-admissible perturbation law is a distribution $\Pi_{\Delta,t}$ with
% % $
% % \mathrm{supp}(\Pi_{\Delta,t})
% % \subseteq
% % \mathcal T_{\Delta,t}.
% % $
% % It induces the effective-update posterior
% % $$
% % Q_t^\Delta
% % =
% % \widehat{\Delta W}_t+\Pi_{\Delta,t},
% % \qquad
% % \Delta W_t
% % =
% % \widehat{\Delta W}_t+\varepsilon_t,
% % \quad
% % \varepsilon_t\sim\Pi_{\Delta,t}.
% % $$
% % % If Gaussian perturbations are used, this corresponds to $\Pi_{\Delta,t}=\mathcal N(0,\Sigma_{\Delta,t})$ with $\mathrm{Range}(\Sigma_{\Delta,t})\subseteq\mathcal T_{\Delta,t}$.
% % % We first isolate the empirical term selected by this support. 
% % Accordingly, the PECL perturbation-smoothed empirical risk is
% % $
% % \widehat L_{S_t}^{\Delta}
% % (
% % \widehat{\Delta W}_t,
% % \Pi_{\Delta,t}
% % )
% % =
% % \mathbb E_{\varepsilon_t\sim\Pi_{\Delta,t}}
% % [
% % \widehat L_{S_t}
% % (
% % W_t^{\mathrm{froz}}
% % +
% % \widehat{\Delta W}_t
% % +
% % \varepsilon_t
% % )
% % ],
% % $
% % and the associated adapter-update sharpness is
% % $
% % \mathrm{TS}_t^\Delta
% % (
% % \widehat{\Delta W}_t,
% % \Pi_{\Delta,t}
% % t)
% % =
% % \mathbb E_{\varepsilon_t\sim\Pi_{\Delta,t}}
% % [
% % \widehat L_{S_t}
% % (
% % W_t^{\mathrm{froz}}
% % +
% % \widehat{\Delta W}_t
% % +
% % \varepsilon_t
% % )
% % -
% % \widehat L_{S_t}
% % (
% % W_t^{\mathrm{froz}}
% % +
% % \widehat{\Delta W}_t
% % )
% % ].
% % $
% % For small perturbations, if the first-order term is negligible, this quantity admits the local approximation
% % $$
% % \mathrm{TS}_{\Delta,t}
% % (\widehat{\Delta W}_t,\Pi_{\Delta,t})
% % \approx
% % \frac{1}{2}
% % \mathbb E_{\epsilon\sim\Pi_{\Delta,t}}
% % [
% % \mathrm{vec}(\epsilon)^\top
% % H_{\Delta,t}
% % \mathrm{vec}(\epsilon)
% % ],
% % H_{\Delta,t}
% % =
% % \nabla_{\mathrm{vec}(\Delta W)}^2
% % \widehat L_{S_t}
% % (
% % W_t^{\mathrm{froz}}+\Delta W
% % )
% % \bigg|_{\Delta W=\widehat{\Delta W}_t}.
% % $$
% % Thus, the relevant curvature is not measured in the full parameter space, but along the locally admissible effective adapter-update directions.

% % \begin{lemma}[Empirical reduction to adapter-update sharpness]
% % \label{lem:subspace-loss-pecl}
% % For each task $t$, let $Q_t^\Delta=\widehat{\Delta W}_t+\Pi_{\Delta,t}$ with $\mathrm{supp}(\Pi_{\Delta,t})\subseteq\mathcal T_{\Delta,t}$, and let $Q_t=(T_t)_\#Q_t^\Delta$. Then
% % $$
% % \mathbb E_{W_t\sim Q_t}
% % \left[
% % \widehat L_{S_t}(W_t)
% % \right]
% % =
% % \widehat L_{S_t}
% % \left(
% % W_t^{\mathrm{froz}}
% % +
% % \widehat{\Delta W}_t
% % \right)
% % +
% % \mathrm{TS}_t^\Delta
% % \left(
% % \widehat{\Delta W}_t,
% % \Pi_{\Delta,t}
% % \right).
% % $$
% % \end{lemma}

% % Lemma~\ref{lem:subspace-loss-pecl} shows that the empirical term in the PAC-Bayesian bound becomes clean empirical loss plus adapter-update sharpness once the perturbation law is restricted to the local PECL-admissible support. We refer to Appendix~\ref{proof:lemma2} for the proof.

% % We next state the corresponding reduction for the KL term. Define the translation map from effective updates to model weights as
% % $
% % T_t(\Delta W_t)
% % =
% % W_t^{\mathrm{froz}}
% % +
% % \Delta W_t.
% % $
% % For an effective-update posterior $Q_t^\Delta$ and an effective-update prior $P_t^\Delta$, the corresponding model-level distributions are
% % $
% % Q_t=(T_t)_\#Q_t^\Delta,
% % \quad
% % P_t=(T_t)_\#P_t^\Delta.
% % $
% Not only the support perturabtion is local on $\Delta W$, the task KL divergence is also rely on $\Delta W$, as shown in Fig\ref{fig:showPAC_combined}(b).And we have following conclusion.
% \begin{lemma}[KL reduction under frozen-backbone translation]
% \label{lem:kl-decomp-pecl}
% Assume $Q_t^\Delta \ll P_t^\Delta$. Then $Q_t \ll P_t$ and
% $$
% \mathrm{KL}
% (
% Q_t
% \,\Vert\,
% P_t
% )
% =
% \mathrm{KL}
% (
% Q_t^\Delta
% \,\Vert\,
% P_t^\Delta
% ).
% $$
% \end{lemma}

% Lemma~\ref{lem:kl-decomp-pecl} shows that the task-level KL complexity is exactly the divergence between the effective-update posterior and the effective-update prior. We refer to Appendix~\ref{proof:Lemma1} for the proof.






% % -------------------


% % at the level of support identification.
% % In LoRA-based PECL, the frozen backbone restricts adaptation to adapter-induced effective update directions, while sharpness can still be measured in several inequivalent spaces: the ambient parameter space, the raw LoRA factor space, or the effective adapter-update space.
% % These spaces are not interchangeable.
% % Full-parameter perturbations include directions unavailable to the PECL learner, whereas raw-factor perturbations depend on a parameterization of the same effective update.
% % Therefore, the task is 

% % We show that this support is the basis-invariant effective adapter-update geometry.

% % We now answer Q2 by identifying the perturbation support that is relevant to the trajectory-level PAC-Bayesian bound under LoRA-constrained adaptation.
% % Section~\ref{sec:pathwise_pac_bayes_q1} shows that the bound-relevant empirical term is the perturbation-smoothed empirical risk induced by a task-specific perturbation law $\Pi_t$.
% % In LoRA-based PECL, the remaining question is which perturbation law is admissible.
% % The key issue is that LoRA-based PECL separates three geometric objects that are usually conflated in full-parameter training: the ambient parameter space, the raw LoRA factor space, and the effective adapter-update geometry.
% % Only the last one corresponds to directions through which the frozen-backbone model can actually be modified by future tasks.
% % Thus, the relevant question is not whether flatness is useful in general, but which support defines the perturbation-smoothed empirical risk and KL complexity that enter Theorem~\ref{thm:sequential_pac_bayes}.

% % At task $t$, the PECL model is written as
% % $
% % W_t
% % =
% % W_t^{\mathrm{froz}}
% % +
% % \Delta W_t,
% % \quad
% % \Delta W_t=A_tB_t,
% % $
% % where $\Delta W_t$ is the effective adapter update induced by the trainable LoRA coordinates.
% % Let $\mathcal W_{\Delta,t}$ denote the admissible effective adapter-update geometry, and let $\mathcal T_{\Delta,t}$ denote the local perturbation support around the learned update $\widehat{\Delta W}_t$.
% % A PECL-admissible perturbation law is a distribution $\Pi_{\Delta,t}$ satisfying
% % $
% % \mathrm{supp}(\Pi_{\Delta,t})
% % \subseteq
% % \mathcal T_{\Delta,t}.
% % $
% % It induces the effective-update posterior
% % $
% % Q_t^\Delta
% % =
% % \widehat{\Delta W}_t+\Pi_{\Delta,t},
% % \quad
% % \Delta W_t
% % =
% % \widehat{\Delta W}_t+\varepsilon_t,
% % \quad
% % \varepsilon_t\sim\Pi_{\Delta,t}.
% % $ The following lemma~\ref{lem:subspace-loss-pecl} isolates an additive loss elevation caused by perturbations supported on $\mathcal W_{\Delta,t}$.
% % These notations will be used in the empirical-risk and KL reduction lemmas below.
% % When Gaussian perturbations are used, we write
% % $
% % \Pi_{\Delta,t}
% % =
% % \mathcal N(0,\Sigma_{\Delta,t}),
% % \qquad
% % \mathrm{Range}(\Sigma_{\Delta,t})
% % \subseteq
% % \mathcal T_{\Delta,t},
% % $
% % which corresponds to a degenerate Gaussian on the local admissible update support.



% % At task $t$, the model has the form
% % $
% % W_t
% % =
% % W_t^{\mathrm{froz}}
% % +
% % \Delta W_t,
% % $
% % where $W_t^{\mathrm{froz}}$ denotes the frozen backbone and $\Delta W_t=A_tB_t$ denotes the adapter update induced by the trainable LoRA coordinates.
% % Let $\mathcal W_{\Delta,t}$ denote the admissible  adapter-update geometry induced by the LoRA parameterization.
% % For perturbation-based sharpness, we use $\mathcal T_{\Delta,t}$ to denote the local admissible perturbation support around the learned update $\widehat{\Delta W}_t$.
% % % This support can be understood as the local tangent support, or equivalently the linear support spanned by the adapter-update coordinates that are perturbable in the implementation.
% % % When $\mathcal W_{\Delta,t}$ is already treated as a local linear support, we identify $\mathcal T_{\Delta,t}$ with $\mathcal W_{\Delta,t}$ for simplicity.

% % A PECL-admissible perturbation law is a distribution $\Pi_{\Delta,t}$ over perturbations $\varepsilon_t$ such that
% % $
% % \mathrm{supp}(\Pi_{\Delta,t})
% % \subseteq
% % \mathcal T_{\Delta,t}.
% % $
% % The corresponding update posterior is
% % $
% % Q_t^\Delta
% % =
% % \widehat{\Delta W}_t
% % +
% % \Pi_{\Delta,t},
% % $
% % meaning that
% % $
% % \Delta W_t
% % =
% % \widehat{\Delta W}_t
% % +
% % \varepsilon_t, \quad
% % \varepsilon_t\sim \Pi_{\Delta,t}.
% % $
% % % The induced model-level posterior is the pushforward of $Q_t^\Delta$ under the translation map
% % % $$
% % % T_t(\Delta W_t)
% % % =
% % % W_t^{\mathrm{froz}}
% % % +
% % % \Delta W_t.
% % % $$
% % % Thus,
% % % $$
% % % Q_t
% % % =
% % % (T_t)_\# Q_t^\Delta .
% % % $$
% % % A Gaussian perturbation is one possible instantiation:
% % % $$
% % % \Pi_{\Delta,t}
% % % =
% % % \mathcal N(0,\Sigma_{\Delta,t}),
% % % \qquad
% % % \mathrm{Range}(\Sigma_{\Delta,t})
% % % \subseteq
% % % \mathcal T_{\Delta,t}.
% % % $$
% % % In this case, $Q_t^\Delta=\mathcal N(\widehat{\Delta W}_t,\Sigma_{\Delta,t})$ is a degenerate Gaussian on the local admissible adapter-update support.
% % % The analysis below depends on the support restriction rather than on Gaussianity itself.
% % We define the PECL perturbation-smoothed empirical risk at task $t$ as
% % $
% % \widehat L_{S_t}^{\Delta}
% % \left(
% % \widehat{\Delta W}_t,
% % \Pi_{\Delta,t}
% % \right)
% % =
% % \mathbb E_{\varepsilon_t\sim\Pi_{\Delta,t}}
% % \left[
% % \widehat L_{S_t}
% % \left(
% % W_t^{\mathrm{froz}}
% % +
% % \widehat{\Delta W}_t
% % +
% % \varepsilon_t
% % \right)
% % \right].
% % $


% % \begin{lemma}[Task adapter subspace sharpness]
% % \label{lem:subspace-loss-pecl}
% % For each task $t$,
% % {%
% % \setlength{\abovedisplayskip}{1pt}
% % \setlength{\belowdisplayskip}{2pt}
% % \setlength{\abovedisplayshortskip}{1pt}
% % \setlength{\belowdisplayshortskip}{2pt}
% % \begin{equation*}
% % \label{eq:subspace-loss-pecl}
% % \begin{aligned}
% % \mathbb E_{W_t\sim Q_t}\bigl[\hat L_{S_t}(W_t)\bigr] 
% % = \mathbb{E}_{\varepsilon\sim\mathcal N(0,\Sigma_t)}
% %     \big[\hat{L}_{S_t}(W^{\mathrm{froz}}_t+\widehat{\Delta W}_t+\varepsilon)\big]
% %  =\hat L_{S_t}(W^{\mathrm{froz}}_t + \widehat{\Delta W}_t)
% % +
% % \mathrm{TS}_t^{\Delta}(\widehat{\Delta W}_t,\Sigma_t).
% % \end{aligned}
% % \end{equation*}
% % }where the task sharpness is $
% % \mathrm{TS}_t^{\Delta}(\widehat{\Delta W}_t,\Sigma_t)
% % = \mathbb{E}_{\varepsilon\sim\mathcal N(0,\Sigma_t)}
% %    \big[\hat{L}_{S_t}(W^{\mathrm{froz}}_t+\widehat{\Delta W}_t+\varepsilon)
% %    - \hat{L}_{S_t}(W^{\mathrm{froz}}_t+\widehat{\Delta W}_t)\big] \text{\,and\, } 
% %      \varepsilon\sim\Pi_{\Delta,t}$.
% %      % \begin{equation*}
% %      % \begin{aligned}
% %      % &\mathrm{RS}_t^{\Delta}(\hat{\Delta W}_t,\Sigma_t):= \\
% %      % &\qquad \mathbb{E}_{\varepsilon\sim\mathcal N(0,\Sigma_t)}
% %      %    \big[\hat{L}_{S_t}(W_\mathrm {frozen }+\hat{\Delta W}_t+\varepsilon) \\
% %      % & \qquad- \hat{L}_{S_t}(W_\mathrm {frozen }+\hat{\Delta W}_t)\big] \\
% %      % &\qquad \varepsilon\sim\mathcal N(0,\Sigma_t) \in \mathcal W_\Delta
% %      % \end{aligned}
% %      % \end{equation*}
% %      \end{lemma}





% % This quantity admits the exact decomposition
% % $$
% % \widehat L_{S_t}^{\Delta}
% % \left(
% % \widehat{\Delta W}_t,
% % \Pi_{\Delta,t}
% % \right)
% % =
% % \widehat L_{S_t}
% % \left(
% % W_t^{\mathrm{froz}}
% % +
% % \widehat{\Delta W}_t
% % \right)
% % +
% % \mathrm{Sh}_t^\Delta
% % \left(
% % \widehat{\Delta W}_t,
% % \Pi_{\Delta,t}
% % \right),
% % $$
% % where the adapter-update sharpness is
% % $$
% % \mathrm{Sh}_t^\Delta
% % \left(
% % \widehat{\Delta W}_t,
% % \Pi_{\Delta,t}
% % \right)
% % =
% % \mathbb E_{\varepsilon_t\sim\Pi_{\Delta,t}}
% % \left[
% % \widehat L_{S_t}
% % \left(
% % W_t^{\mathrm{froz}}
% % +
% % \widehat{\Delta W}_t
% % +
% % \varepsilon_t
% % \right)
% % -
% % \widehat L_{S_t}
% % \left(
% % W_t^{\mathrm{froz}}
% % +
% % \widehat{\Delta W}_t
% % \right)
% % \right].
% % $$
% % This is the sharpness object selected by PECL geometry.
% % It measures loss elevation under perturbations of the effective adapter update, not under perturbations of frozen full-model weights.

% % The following lemma makes this empirical reduction explicit.

% % \begin{lemma}[Empirical risk reduction to adapter-update sharpness]
% % \label{lem:subspace-loss-pecl}
% % For each task $t$, let $Q_t^\Delta=\widehat{\Delta W}_t+\Pi_{\Delta,t}$ be an effective-update posterior with $\mathrm{supp}(\Pi_{\Delta,t})\subseteq\mathcal T_{\Delta,t}$, and let $Q_t=(T_t)_\#Q_t^\Delta$ be its pushforward to model weights.
% % Then
% % $$
% % \begin{aligned}
% % \mathbb E_{W_t\sim Q_t}
% % \left[
% % \widehat L_{S_t}(W_t)
% % \right]
% % &=
% % \mathbb E_{\varepsilon_t\sim\Pi_{\Delta,t}}
% % \left[
% % \widehat L_{S_t}
% % \left(
% % W_t^{\mathrm{froz}}
% % +
% % \widehat{\Delta W}_t
% % +
% % \varepsilon_t
% % \right)
% % \right]
% % \\
% % &=
% % \widehat L_{S_t}
% % \left(
% % W_t^{\mathrm{froz}}
% % +
% % \widehat{\Delta W}_t
% % \right)
% % +
% % \mathrm{Sh}_t^\Delta
% % \left(
% % \widehat{\Delta W}_t,
% % \Pi_{\Delta,t}
% % \right).
% % \end{aligned}
% % $$
% % % \end{lemma}

% % Lemma~\ref{lem:subspace-loss-pecl} shows that the empirical term in the PAC-Bayesian bound reduces to clean empirical loss plus adapter-update sharpness when the perturbation law is supported on the admissible PECL update geometry. We refer to Appendix~\ref{proof:lemma2} for the proof.
% % This decomposition is an exact identity for the perturbation-induced posterior.
% % It is not a local Taylor approximation.

% % A local second-order approximation further clarifies the geometric meaning of this term.
% % Let $H_t$ be the Hessian of $\widehat L_{S_t}$ at $\widehat W_t$.
% % Let $\Pi_{\Delta,t}^{\mathrm{proj}}$ denote the projector onto the local admissible adapter-update support $\mathcal T_{\Delta,t}$.
% % The restricted Hessian is
% % $
% % H_{\Delta,t}
% % =
% % \Pi_{\Delta,t}^{\mathrm{proj}}
% % H_t
% % \Pi_{\Delta,t}^{\mathrm{proj}}.
% % $
% % If $\varepsilon_t$ has zero mean and covariance $\Sigma_{\Delta,t}$ supported on $\mathcal T_{\Delta,t}$, then
% % $$
% % \mathrm{Sh}_t^\Delta
% % \left(
% % \widehat{\Delta W}_t,
% % \Pi_{\Delta,t}
% % \right)
% % \approx
% % \frac{1}{2}
% % \mathrm{tr}
% % \left(
% % H_{\Delta,t}
% % \Sigma_{\Delta,t}
% % \right).
% % $$
% % Therefore, PECL sharpness is governed by curvature and perturbation variance inside the update-carrying geometry, rather than by curvature over the full ambient parameter space.


% % Because $T_t$ is a translation by the frozen backbone, it preserves the divergence between the effective-update posterior and prior.

% % \begin{lemma}[KL support reduction under frozen-backbone adaptation]
% % \label{lem:kl-decomp-pecl}
% % Let $Q_t^\Delta$ and $P_t^\Delta$ be posterior and prior distributions on the effective adapter-update support, and let $Q_t=(T_t)_\#Q_t^\Delta$ and $P_t=(T_t)_\#P_t^\Delta$ be their pushforward distributions on $W_t^{\mathrm{froz}}+\mathcal W_{\Delta,t}$.
% % If $Q_t^\Delta \ll P_t^\Delta$, then $Q_t \ll P_t$ and
% % $$
% % \mathrm{KL}
% % \left(
% % Q_t
% % \,\middle\|\,
% % P_t
% % \right)
% % =
% % \mathrm{KL}
% % \left(
% % Q_t^\Delta
% % \,\middle\|\,
% % P_t^\Delta
% % \right).
% % $$
% % \end{lemma}

% % We next show that the KL complexity term in the same bound reduces to the same support.
% % Let
% % $
% % T_t(\Delta W_t)
% % =
% % W_t^{\mathrm{froz}}
% % +
% % \Delta W_t
% % $
% % be the translation map from effective updates to model weights.
% % The induced model-level posterior is
% % $
% % Q_t=(T_t)_\# Q_t^\Delta.
% % $
% % Similarly, for a task prior $P_t^\Delta$ over effective adapter updates, we write
% % $
% % P_t=(T_t)_\# P_t^\Delta.
% % $
% % We propose the following lemma to formalize the resulting translation-invariance of the KL divergence:
% % % , and its proof is included in Appendix~\ref{proof:Lemma1}.
% % \begin{lemma}
% % \label{lem:kl-decomp-pecl}
% % Assume $Q_t^\Delta \ll P_t^\Delta$. Then $Q_t \ll P_t$ and
% % \[
% % \mathrm{KL}\bigl(Q_t \,\Vert\, P_t\bigr)
% % =
% % \mathrm{KL}\bigl(Q_t^\Delta \,\Vert\, P_t^\Delta\bigr).
% % \]
% % \end{lemma}
% % % Lemma~\ref{lem:kl-decomp-pecl} shows that, under a frozen backbone, the  KL divergence in the full parameter space is exactly the KL divergence on the task-permissible update subspace. 
% % Lemma~\ref{lem:kl-decomp-pecl} shows that the PAC-Bayesian complexity term does not depend on the ambient full-parameter space once the backbone is frozen.
% % The task-level KL cost is exactly the cost of moving from the adapter-update prior to the adapter-update posterior inside the effective update support. We refer to Appendix~\ref{proof:Lemma1} for the proof.

% Combining the previous perturabtion analyse and Lemma~\ref{lem:kl-decomp-pecl} with Theorem~\ref{thm:sequential_pac_bayes} yields the following PECL-specific support-reduction bound.
% % Combining Lemma~\ref{lem:subspace-loss-pecl} and Lemma~\ref{lem:kl-decomp-pecl} with Theorem~\ref{thm:sequential_pac_bayes} yields the following PECL support-reduction bound.
% \begin{theorem}[PECL support-reduction bound]
% \label{thm:pecl-bound}
% Under the setup above, assume that the loss is bounded in $[0,1]$.
% For any $\delta\in(0,1]$, with probability at least $1-\delta$ over the task samples $S_{1:T}$,
% $$
% \begin{aligned}
% \mathcal L_{1:T}^{\Delta}
% \le
% &
% \frac{1}{T}
% \sum_{t=1}^{T}
% \mathbb E_{P_t^\Delta\sim\mathcal P_{t-1}}
% \left[
% \widehat L_{S_t}
% \left(
% W_t^{\mathrm{froz}}
% +
% \widehat{\Delta W}_t
% \right)
% +
% \mathrm{TS}_t^\Delta
% \left(
% \widehat{\Delta W}_t,
% \Pi_{\Delta,t}
% \right)
% \right]
% \\
% &+
% \sqrt{
% \frac{
% \sum_{t=1}^{T}
% \mathbb E_{P_t^\Delta\sim \mathcal P_{t-1}}
% \mathrm{KL}
% \left(
% Q_t^\Delta
% \,\Vert\,
% P_t^\Delta
% \right)
% +
% \sum_{t=1}^{T}
% \mathrm{KL}
% \left(
% \mathcal P_t
% \,\Vert\,
% \mathcal P_{t-1}
% \right)
% +
% \log \frac{1}{\delta}
% }{
% 2\sum_{t=1}^{T}(m_t-1)
% }
% },
% \end{aligned}
% $$
% where
% $
% \mathcal L_{1:T}^{\Delta}
% =
% \frac{1}{T}
% \sum_{t=1}^{T}
% \mathbb E_{P_t^\Delta\sim\mathcal P_{t-1}}
% \mathbb E_{\Delta W_t\sim Q_t^\Delta}
% L_{\mathcal D_t}
% \left(
% W_t^{\mathrm{froz}}
% +
% \Delta W_t
% \right).
% $
% \end{theorem}

% % Theorem~\ref{thm:pecl-bound} answers Q2 by showing that the perturbation support selected by the trajectory-level PAC-Bayesian bound is the effective adapter-update support.
% % Both bound-relevant quantities reduce to this support.
% % The empirical term becomes clean empirical loss plus adapter-update sharpness, while the task-level KL term becomes the divergence between the effective-update posterior and the effective-update prior.

% % This result identifies the flatness object relevant to LoRA-based PECL.
% % It is not ambient full-parameter sharpness, because full-space perturbations include frozen directions that cannot be used by future LoRA updates.
% % It is also not raw factor-space sharpness over $A_t$ and $B_t$, because different factorizations can induce the same effective update while yielding different perturbation scales and norms.
% % The bound instead selects the basis-invariant effective update geometry of $\Delta W_t$ as the natural support for perturbation-smoothed empirical risk.

% % Consequently, sharpness-aware methods in PECL should be evaluated by whether they reduce loss sensitivity along the admissible adapter-update directions that carry future sequential adaptation and potential interference.
% % The proof of Theorem~\ref{thm:pecl-bound} follows by applying Theorem~\ref{thm:sequential_pac_bayes} to the pushforward posterior induced by frozen-backbone translation, together with Lemma~\ref{lem:subspace-loss-pecl} and Lemma~\ref{lem:kl-decomp-pecl}.
% Theorem~5.2 provides a support-identification result. Under frozen-backbone LoRA adaptation, both bound-relevant quantities reduce to the effective adapter-update geometry: the empirical term becomes clean empirical loss plus $\mathrm{TS}_t^\Delta$, while the task-level KL term becomes $\mathrm{KL}(Q_t^\Delta\|P_t^\Delta)$. Hence, the relevant flatness object in PECL is neither ambient full-parameter sharpness nor raw factor-space sharpness over $(A_t,B_t)$, but task-local sharpness of the effective update $\Delta W_t$ along the support through which the posterior trajectory can move.
% We provide the proof in Appendix~\ref{proof:theorem}.

% % \begin{theorem}[PECL support-reduction bound]
% % \label{thm:pecl-bound}
% % Consider a LoRA-based PECL model with frozen backbone $W_t^{\mathrm{froz}}$ at task $t$.
% % Let $Q_t^\Delta=\widehat{\Delta W}_t+\Pi_{\Delta,t}$ be an effective-update posterior whose perturbation law satisfies $\mathrm{supp}(\Pi_{\Delta,t})\subseteq\mathcal T_{\Delta,t}$.
% % Let $P_t^\Delta$ be a task prior sampled from the history-dependent hyperposterior $\mathcal P_{t-1}$ over adapter-update priors.
% % Let $Q_t=(T_t)_\#Q_t^\Delta$ and $P_t=(T_t)_\#P_t^\Delta$ be the corresponding model-level pushforwards.
% % Assume that the loss is bounded in $[0,1]$.
% % For any $\delta\in(0,1]$, with probability at least $1-\delta$ over the task samples $S_{1:T}$,
% % $$
% % \begin{aligned}
% % \mathcal L_{1:T}^{\Delta}
% % \le
% % &
% % \frac{1}{T}
% % \sum_{t=1}^{T}
% % \mathbb E_{P_t^\Delta\sim\mathcal P_{t-1}}
% % \left[
% % \widehat L_{S_t}
% % \left(
% % W_t^{\mathrm{froz}}
% % +
% % \widehat{\Delta W}_t
% % \right)
% % +
% % \mathrm{Sh}_t^\Delta
% % \left(
% % \widehat{\Delta W}_t,
% % \Pi_{\Delta,t}
% % \right)
% % \right]
% % \\
% % &+
% % \sqrt{
% % \frac{
% % \sum_{t=1}^{T}
% % \mathbb E_{P_t^\Delta\sim \mathcal P_{t-1}}
% % \mathrm{KL}
% % \left(
% % Q_t^\Delta
% % \,\middle\|\,
% % P_t^\Delta
% % \right)
% % +
% % \sum_{t=1}^{T}
% % \mathrm{KL}
% % \left(
% % \mathcal P_t
% % \,\middle\|\,
% % \mathcal P_{t-1}
% % \right)
% % +
% % \log \frac{1}{\delta}
% % }{
% % 2\sum_{t=1}^{T}(m_t-1)
% % }
% % }.
% % \end{aligned}
% % $$
% % Here,
% % $$
% % \mathcal L_{1:T}^{\Delta}
% % =
% % \frac{1}{T}
% % \sum_{t=1}^{T}
% % \mathbb E_{P_t^\Delta\sim\mathcal P_{t-1}}
% % \mathbb E_{\Delta W_t\sim Q_t^\Delta}
% % L_{\mathcal D_t}
% % \left(
% % W_t^{\mathrm{froz}}
% % +
% % \Delta W_t
% % \right).
% % $$
% % Equivalently, by defining
% % $$
% % \widehat{\mathcal L}_{1:T}^{\Delta}
% % =
% % \frac{1}{T}
% % \sum_{t=1}^{T}
% % \mathbb E_{P_t^\Delta\sim\mathcal P_{t-1}}
% % \widehat L_{S_t}^{\Delta}
% % \left(
% % \widehat{\Delta W}_t,
% % \Pi_{\Delta,t}
% % \right),
% % $$
% % the bound can be written compactly as
% % $$
% % \mathcal L_{1:T}^{\Delta}
% % \le
% % \widehat{\mathcal L}_{1:T}^{\Delta}
% % +
% % \sqrt{
% % \frac{
% % \sum_{t=1}^{T}
% % \mathbb E_{P_t^\Delta\sim \mathcal P_{t-1}}
% % \mathrm{KL}
% % \left(
% % Q_t^\Delta
% % \,\middle\|\,
% % P_t^\Delta
% % \right)
% % +
% % \sum_{t=1}^{T}
% % \mathrm{KL}
% % \left(
% % \mathcal P_t
% % \,\middle\|\,
% % \mathcal P_{t-1}
% % \right)
% % +
% % \log \frac{1}{\delta}
% % }{
% % 2\sum_{t=1}^{T}(m_t-1)
% % }
% % }.
% % $$
% % \end{theorem}

% % Theorem~\ref{thm:pecl-bound} answers Q2 by specializing the generic perturbation law in Theorem~\ref{thm:sequential_pac_bayes} to the PECL-admissible perturbation law $\Pi_{\Delta,t}$.
% % Under frozen-backbone LoRA adaptation, both bound-relevant quantities reduce to the effective adapter-update support.
% % The empirical term becomes the clean empirical loss plus adapter-update sharpness, while the task-level KL term becomes the divergence between the effective-update posterior and the effective-update prior.

% % This result identifies the flatness object relevant to LoRA-based PECL.
% % It is not ambient full-parameter sharpness, because full-space perturbations include frozen directions that cannot be used by future LoRA updates.
% % It is also not raw factor-space sharpness over $A_t$ and $B_t$, because different factorizations can induce the same effective update while yielding different perturbation scales and norms.
% % The bound instead selects the basis-invariant effective update geometry of $\Delta W_t$ as the natural support for perturbation-smoothed empirical risk.

% % Consequently, sharpness-aware methods in PECL should not be evaluated only by whether they reduce loss sensitivity in the full parameter space.
% % They should be evaluated by whether they reduce loss sensitivity along the admissible adapter-update directions that carry future sequential adaptation and potential interference.
% % The proof of Theorem~\ref{thm:pecl-bound} follows by applying Theorem~\ref{thm:sequential_pac_bayes} to the pushforward posterior induced by $T_t$, together with Lemma~\ref{lem:subspace-loss-pecl} and Lemma~\ref{lem:kl-decomp-pecl}, and is provided in Appendix~\ref{app:proof_pecl_bound}.









% % ---------------------

% % We now answer Q2 by identifying the perturbation support that is relevant to the trajectory-level PAC-Bayesian bound under LoRA-constrained adaptation.
% % Section~\ref{sec:pathwise_pac_bayes_q1} shows that the bound-relevant empirical term is the perturbation-smoothed empirical risk induced by a task-specific perturbation law $\Pi_t$.
% % In LoRA-based PECL, the remaining question is which perturbation law is admissible.
% % The key issue is that LoRA-based PECL separates three geometric objects that are usually conflated in full-parameter training: the ambient parameter space, the raw LoRA factor space, and the effective adapter-update geometry.
% % Only the last one corresponds to directions through which the frozen-backbone model can actually be modified by future tasks.
% % Thus, the relevant question is not whether flatness is useful in general, but which support defines the perturbation-smoothed empirical risk and KL complexity that enter Theorem~\ref{thm:sequential_pac_bayes}.

% % We now answer Q2 by identifying the perturbation support that is relevant to the trajectory-level PAC-Bayesian bound under LoRA-constrained adaptation. 
% % The key issue is that LoRA-based PECL separates three geometric objects that are usually conflated in full-parameter training: the ambient parameter space, the raw LoRA factor space, and the effective adapter-update space. 
% % Only the last one corresponds to directions along which the model can actually be modified under a frozen backbone. 
% % Thus, the question is which perturbation support defines the perturbation-smoothed empirical risk and KL complexity that enter the bound in Theorem~\ref{thm:sequential_pac_bayes}.

% % At task $t$, let $W_t^{\mathrm{froz}}$ denote the frozen backbone and let $\mathcal W_{\Delta,t}$ denote the admissible effective adapter-update support induced by the LoRA parameterization. 
% % A task posterior over full model weights is generated by first sampling an effective update $\Delta W_t$ and then translating it by the frozen backbone:
% % $
% % W_t = W_t^{\mathrm{froz}} + \Delta W_t,
% % \quad
% % \Delta W_t \sim Q_t^\Delta .
% % $
% % We consider a Gaussian posterior on the effective adapter-update support,
% % $
% % Q_t^\Delta = \mathcal N(\widehat{\Delta W}_t,\Sigma_t),
% % $
% % where $\widehat{\Delta W}_t$ is the learned effective adapter update and $\Sigma_t$ defines perturbations supported on $\mathcal W_{\Delta,t}$. 

% % \begin{lemma}[KL support reduction under frozen-backbone adaptation]
% % \label{lem:kl-decomp-pecl}
% % Let $Q_t^\Delta$ and $P_t^\Delta$ be posterior and prior distributions supported on $\mathcal W_{\Delta,t}$, and let $Q_t$ and $P_t$ be their pushforward distributions on $W_t^{\mathrm{froz}}+\mathcal W_{\Delta,t}$. 
% % If $Q_t^\Delta \ll P_t^\Delta$, then $Q_t \ll P_t$ and
% % $$
% % \mathrm{KL}\bigl(Q_t \,\Vert\, P_t\bigr)
% % =
% % \mathrm{KL}\bigl(Q_t^\Delta \,\Vert\, P_t^\Delta\bigr).
% % $$
% % \end{lemma}

% % Lemma~\ref{lem:kl-decomp-pecl} shows that the PAC-Bayesian complexity term does not depend on the ambient full-parameter space once the backbone is frozen. 
% % The KL cost is exactly the cost of moving from the adapter prior to the adapter posterior inside the effective update support $\mathcal W_{\Delta,t}$.

% % Next, we show that the empirical term in the same bound also reduces to the same support.  Equivalently,
% % $
% % \Delta W_t = \widehat{\Delta W}_t + \varepsilon_t,
% % \qquad
% % \varepsilon_t \sim \mathcal N(0,\Sigma_t),
% % \qquad
% % \varepsilon_t \in \mathcal W_{\Delta,t}.
% % $
% % The induced full-model posterior $Q_t$ is therefore the pushforward of $Q_t^\Delta$ under the translation map $\Delta W_t \mapsto W_t^{\mathrm{froz}}+\Delta W_t$.
% % For each task $t$, define the adapter-update sharpness as the loss elevation induced by perturbations supported on $\mathcal W_{\Delta,t}$:
% % $$
% % \mathrm{TS}_t^{\Delta}(\widehat{\Delta W}_t,\Sigma_t)
% % =
% % \mathbb{E}_{\varepsilon_t\sim\mathcal N(0,\Sigma_t)}
% % \left[
% % \widehat L_{S_t}
% % \left(
% % W_t^{\mathrm{froz}}+\widehat{\Delta W}_t+\varepsilon_t
% % \right)
% % -
% % \widehat L_{S_t}
% % \left(
% % W_t^{\mathrm{froz}}+\widehat{\Delta W}_t
% % \right)
% % \right],
% % \qquad
% % \varepsilon_t \in \mathcal W_{\Delta,t}.
% % $$

% % \begin{lemma}[Empirical risk decomposition on the adapter-update support]
% % \label{lem:subspace-loss-pecl}
% % For each task $t$, the perturbation-smoothed empirical risk induced by $Q_t^\Delta$ satisfies
% % $$
% % \begin{aligned}
% % \mathbb E_{W_t\sim Q_t}
% % \left[
% % \widehat L_{S_t}(W_t)
% % \right]
% % &=
% % \mathbb{E}_{\varepsilon_t\sim\mathcal N(0,\Sigma_t)}
% % \left[
% % \widehat L_{S_t}
% % \left(
% % W_t^{\mathrm{froz}}+\widehat{\Delta W}_t+\varepsilon_t
% % \right)
% % \right] \\
% % &=
% % \widehat L_{S_t}
% % \left(
% % W_t^{\mathrm{froz}}+\widehat{\Delta W}_t
% % \right)
% % +
% % \mathrm{TS}_t^{\Delta}(\widehat{\Delta W}_t,\Sigma_t).
% % \end{aligned}
% % $$
% % \end{lemma}

% % Lemma~\ref{lem:subspace-loss-pecl} identifies the sharpness term that appears in PECL. 
% % It is not full-parameter sharpness and it is not raw-factor sharpness over $A_t$ and $B_t$. 
% % It is the loss elevation caused by perturbing the effective adapter update $\Delta W_t$ inside the support through which future task updates are actually realized.

% % Combining Lemma~\ref{lem:kl-decomp-pecl} and Lemma~\ref{lem:subspace-loss-pecl} with Theorem~\ref{thm:sequential_pac_bayes} yields the following PECL-specific support-reduction bound.

% % \begin{theorem}[PECL support-reduction bound]
% % \label{thm:pecl-bound}
% % Consider a LoRA-based PECL model with frozen backbone $W_t^{\mathrm{froz}}$ at task $t$. 
% % Let $Q_t^\Delta=\mathcal N(\widehat{\Delta W}_t,\Sigma_t)$ be a Gaussian posterior supported on the effective adapter-update support $\mathcal W_{\Delta,t}$, and let $P_t^\Delta$ be a task prior sampled from the history-dependent hyperposterior $\mathcal P_{t-1}$. 
% % Assume that the loss is bounded in $[0,1]$. 
% % For any $\delta\in(0,1]$, with probability at least $1-\delta$ over $S_{1:T}$,
% % % $$
% % % \begin{aligned}
% % % \frac{1}{T}\sum_{t=1}^T
% % % \mathbb E_{P_t^\Delta\sim\mathcal P_{t-1}}
% % % &\mathbb E_{W_t\sim Q_t}
% % % L_{\mathcal D_t}(W_t)
% % % \le
% % % \frac{1}{T}\sum_{t=1}^T
% % % \mathbb E_{P_t^\Delta\sim\mathcal P_{t-1}}
% % % \left[
% % % \widehat L_{S_t}
% % % \left(
% % % W_t^{\mathrm{froz}}+\widehat{\Delta W}_t
% % % \right)
% % % +
% % % \mathrm{TS}_t^{\Delta}(\widehat{\Delta W}_t,\Sigma_t)
% % % \right] \\
% % % &\quad+
% % % \sqrt{
% % % \frac{
% % % \sum_{t=1}^{T}
% % % \mathbb E_{P_t^\Delta\sim \mathcal P_{t-1}}
% % % \mathrm{KL}
% % % \left(
% % % Q_t^\Delta
% % % \,\middle\|\,
% % % P_t^\Delta
% % % \right)
% % % +
% % % \sum_{t=1}^{T}
% % % \mathrm{KL}
% % % \left(
% % % \mathcal P_t
% % % \,\middle\|\,
% % % \mathcal P_{t-1}
% % % \right)
% % % +
% % % \log \frac{1}{\delta}
% % % }{
% % % 2\sum_{t=1}^{T}(m_t-1)
% % % }
% % % }.
% % % \end{aligned}
% % % $$
% % $$
% % \begin{aligned}
% % % \frac{1}{T}\sum_{t=1}^T
% % % \mathbb E_{P_t^\Delta\sim\mathcal P_{t-1}}
% % % \mathbb E_{W_t\sim Q_t}
% % % L_{\mathcal D_t}(W_t)
% % \mathcal L_{1:T}^\Pi
% % \le
% % % % \frac{1}{T}\sum_{t=1}^T
% % % % \mathbb E_{P_t^\Delta\sim\mathcal P_{t-1}}
% % % % \left[
% % % % \widehat L_{S_t}
% % % % \left(
% % % % W_t^{\mathrm{froz}}+\widehat{\Delta W}_t
% % % % \right)
% % % +
% % % \mathrm{TS}_t^{\Delta}(\widehat{\Delta W}_t,\Sigma_t)
% % % \right] \\
% % \widehat{\mathcal L}_{1:T}^\Pi
% % +
% % \sqrt{
% % \frac{
% % \sum_{t=1}^{T}
% % \mathbb E_{P_t^\Delta\sim \mathcal P_{t-1}}
% % \mathrm{KL}
% % \left(
% % Q_t^\Delta
% % \,\middle\|\,
% % P_t^\Delta
% % \right)
% % +
% % \sum_{t=1}^{T}
% % \mathrm{KL}
% % \left(
% % \mathcal P_t
% % \,\middle\|\,
% % \mathcal P_{t-1}
% % \right)
% % +
% % \log \frac{1}{\delta}
% % }{
% % 2\sum_{t=1}^{T}(m_t-1)
% % }
% % }.
% % \end{aligned}
% % $$
% % \end{theorem}

% % Theorem~\ref{thm:pecl-bound} answers Q2. 
% % Under frozen-backbone LoRA adaptation, the bound-relevant perturbation support is the effective adapter-update geometry $\mathcal W_{\Delta,t}$. 
% % Both components of the PAC-Bayesian bound reduce to this support: the empirical term becomes clean training loss plus adapter-update sharpness, and the KL complexity becomes the divergence between adapter-update posterior and adapter-update prior. 
% % Therefore, the relevant flatness object in LoRA-based PECL is not ambient full-parameter flatness, because full-space perturbations include directions that the learner cannot use for future adaptation. 
% % It is also not raw LoRA factor-space flatness, because different factorizations of $A_tB_t$ can represent the same effective update while inducing different norms and perturbation scales in factor space. 
% % The support identified by the bound is instead the basis-invariant effective update geometry of $\Delta W_t$.

% % This result refines the flatness explanation for continual learning under constrained update geometry. 
% % In full-parameter training, perturbation directions, trainable update directions, and interference directions are often treated as the same object. 
% % In LoRA-based PECL, this equivalence fails. 
% % Theorem~\ref{thm:pecl-bound} shows that the perturbation-smoothed empirical risk relevant to trajectory-level generalization is the one measured along the admissible adapter-update support. 
% % Consequently, sharpness-aware methods should be evaluated not merely by whether they reduce loss sensitivity in the full parameter space, but by whether they reduce sensitivity along the effective LoRA update directions that govern sequential adaptation.




% % \section{Which Flatness Matters in LoRA Based PECL}
% % \label{sec:which_flatness_pecl}

% % % Section~\ref{sec:flatness_cl} shows that flatness enters continual learning through perturbation-smoothed empirical risk along a sequential posterior trajectory. We now address the PECL-specific question left open by that result: when the pretrained backbone is frozen and only LoRA-induced updates are trainable, which perturbation support defines the flatness term that enters the bound?

% % % This question is nontrivial because LoRA based PECL separates three geometries that are usually conflated in full-parameter training. The loss is evaluated on the full model, optimization moves only through trainable adapter-induced directions, and raw LoRA factors provide a non-unique parameterization of the same effective weight update. Therefore, full-parameter sharpness, raw-factor sharpness, and effective adapter-update sharpness are not equivalent. The relevant flatness object must be derived from the posterior support through which the learner can actually move.

% % % The main result of this section is a support-reduction statement. When the sequential PAC-Bayesian decomposition of Section~\ref{sec:flatness_cl} is specialized to frozen-backbone LoRA adaptation, both the task-level KL complexity and the perturbation-smoothed empirical risk reduce to the trainable effective adapter-update support. Hence, the flatness entering the LoRA PECL bound is not full-space flatness and not raw-factor flatness. It is effective adapter-update flatness on the current task's trainable support.
% % % Theorem~\ref{thm:pathwise_pac} shows that perturbation-smoothed empirical risk and KL complexity are the two bound-relevant quantities through which flatness enters the analysis of a continual learning trajectory. We now identify the geometric support of these quantities in frozen-backbone PECL. The key point is that LoRA does not directly update the ambient parameter space. Instead, trainable LoRA coordinates induce an effective adapter update, while the pretrained backbone and other frozen components remain fixed. Therefore, the relevant flatness object should be defined on the effective update geometry reachable by the current adapter.

% % % \subsection{Effective Update Posterior under Frozen-Backbone LoRA}
% % % \subsection{Trainable Effective Update Support in PECL}
% % % \label{sec:pecl_geometry}




% % % At task $t$, the model can be written as
% % % $
% % % W_t
% % % =
% % % W_t^{\mathrm{froz}}
% % % +
% % % \Delta W_t ,
% % % $
% % % where $W_t^{\mathrm{froz}}$ collects all parameters that are fixed during task-$t$ adaptation, while $\Delta W_t$ denotes the effective update induced by the LoRA coordinates that are trainable at task $t$.
% % % % For a LoRA layer, let
% % % % $
% % % % \Delta W_t=A_tB_t
% % % % $
% % % % where $(A_t,B_t)$ are raw LoRA factors. 
% % % % The theoretical object of interest is not the raw factor pair itself, but the effective update that changes the model function. 
% % % We define the trainable effective update support as
% % % $
% % % \mathcal W_{\Delta,t}^{\mathrm{train}}
% % % =
% % % \left\{
% % % \Delta W:
% % % \Delta W=A_tB_t,
% % % \ (A_t,B_t)\in\mathcal A_t^{\mathrm{train}}
% % % \right\},
% % % % \label{eq:trainable_effective_update_support}
% % % $
% % % where $\mathcal A_t^{\mathrm{train}}$ is the feasible set of LoRA factors actively trained on task $t$. 
% % % For local perturbation analysis, $\mathcal W_{\Delta,t}^{\mathrm{train}}$ can be understood as the tangent update support induced by the Jacobian of $\phi_t$ around the learned factors.

% % % This definition separates three non-equivalent spaces. The first is the ambient full parameter space of $W$. Perturbations in this space include frozen coordinates that the PECL learner cannot update. The second is the raw factor space of $(A,B)$. Perturbations in this space depend on the chosen LoRA parameterization. Under a bilinear parameterization, for example,
% % % \begin{equation}
% % % (A,B)\mapsto (cA,B/c)
% % % \label{eq:lora_rescaling}
% % % \end{equation}
% % % can preserve the same effective update while changing distances and perturbation geometry in raw factor coordinates. The third is the effective adapter-update support $\mathcal W_{\Delta,t}^{\mathrm{train}}$, which directly describes the model-level directions that the current task can realize.

% % % A further distinction is needed for LoRA variants that accumulate adapter components across tasks. The accumulated model adapter support after task $t$ is
% % % \begin{equation}
% % % \mathcal W_{\Delta,1:t}^{\mathrm{model}}
% % % =
% % % \mathcal W_{\Delta,1}^{\mathrm{train}}
% % % +
% % % \mathcal W_{\Delta,2}^{\mathrm{train}}
% % % +
% % % \cdots
% % % +
% % % \mathcal W_{\Delta,t}^{\mathrm{train}} .
% % % \label{eq:accumulated_model_support}
% % % \end{equation}
% % % This support describes all adapter-induced directions present in the model, but it is not necessarily the posterior support of the learner at task $t$. In SeqLoRA, the same adapter branch is reused, so the trainable and accumulated supports largely coincide. In IncLoRA-style methods, a new branch is allocated for each task, while previous branches are frozen. In that case,
% % % \begin{equation}
% % % W_t^{\mathrm{froz}}
% % % =
% % % W_0
% % % +
% % % \sum_{k<t}\widehat{\Delta W}_k,
% % % \qquad
% % % \Delta W_t\in\mathcal W_{\Delta,t}^{\mathrm{train}} .
% % % \label{eq:inclora_frozen_decomposition}
% % % \end{equation}
% % % Thus, previous adapters contribute to the frozen model state, whereas the current posterior only moves through the new trainable support.

% % % This distinction is central to the bound. The perturbation distribution associated with task-$t$ adaptation should be supported on $\mathcal W_{\Delta,t}^{\mathrm{train}}$, not on the full parameter space and not necessarily on $\mathcal W_{\Delta,1:t}^{\mathrm{model}}$. Full-model probes may still diagnose accumulated sensitivity, but they are not the support-reduced sharpness term induced by the PECL posterior trajectory.

% % We formalize the LoRA adapter posterior on a Gaussian distribution in the permissible space $\mathcal W_{\Delta,t}$ as
% % $
% % Q_t^\Delta = \mathcal N(\widehat{\Delta W}_t,\Sigma_t)$ where
% % $\widehat{\Delta W}_t$ refers to the learned weight update 
% % and $\Sigma_t$ is a covariance around $\widehat{\Delta W}_t$ on $\mathcal W_{\Delta,t}$ .
% % Equivalently, sampling $W_t\sim Q_t$ on $W^{\mathrm{froz}}_t+\mathcal W_{\Delta,t}$ amounts to 
% % % sampling $ \Delta W_t \sim Q_t^{\Delta}$ \text{ on } $W_{\Delta,t}$ and 
% % setting $W_t=W^{\mathrm{froz}}_t+\Delta W_t$ and sampling $ \Delta W_t \sim Q_t^{\Delta}$ where $\Delta W_t = \widehat{\Delta W}_t + \varepsilon$
% % and
% % $
% % \varepsilon \sim \mathcal N(0,\Sigma_t).$

% % % \text{ on }  \mathcal W_{\Delta,t}.$

% % % and setting $W=W_{\mathrm{frozen}}+\Delta W$.


% % We propose the following lemma to formalize the resulting translation-invariance of the KL divergence:
% % % , and its proof is included in Appendix~\ref{proof:Lemma1}.
% % \begin{lemma}
% % \label{lem:kl-decomp-pecl}
% % Assume $Q_t^\Delta \ll P_t^\Delta$. Then $Q_t \ll P_t$ and
% % \[
% % \mathrm{KL}\bigl(Q_t \,\Vert\, P_t\bigr)
% % =
% % \mathrm{KL}\bigl(Q_t^\Delta \,\Vert\, P_t^\Delta\bigr).
% % \]
% % \end{lemma}
% % Lemma~\ref{lem:kl-decomp-pecl} shows that, under a frozen backbone, the  KL divergence in the full parameter space is exactly the KL divergence on the task-permissible update subspace. We refer to Appendix~\ref{proof:Lemma1} for the proof.


% % Now, we define the adapter subspace sharpness as the expected loss elevation under perturbations supported on the adapter update subspace $\mathcal W_{\Delta,t}$ for task $t$. The following lemma~\ref{lem:subspace-loss-pecl} isolates an additive loss elevation caused by perturbations supported on $\mathcal W_{\Delta,t}$.

% % \begin{lemma}[Task adapter subspace sharpness]
% % \label{lem:subspace-loss-pecl}
% % For each task $t$,
% % {%
% % \setlength{\abovedisplayskip}{1pt}
% % \setlength{\belowdisplayskip}{2pt}
% % \setlength{\abovedisplayshortskip}{1pt}
% % \setlength{\belowdisplayshortskip}{2pt}
% % \begin{equation*}
% % \label{eq:subspace-loss-pecl}
% % \begin{aligned}
% % \mathbb E_{W_t\sim Q_t}\bigl[\hat L_{S_t}(W_t)\bigr] 
% % = \mathbb{E}_{\varepsilon\sim\mathcal N(0,\Sigma_t)}
% %     \big[\hat{L}_{S_t}(W^{\mathrm{froz}}_t+\widehat{\Delta W}_t+\varepsilon)\big]  =\hat L_{S_t}(W^{\mathrm{froz}}_t + \widehat{\Delta W}_t)
% % +
% % \mathrm{TS}_t^{\Delta}(\widehat{\Delta W}_t,\Sigma_t).
% % \end{aligned}
% % \end{equation*}
% % }where the task adapter subspace sharpness is $
% % \mathrm{TS}_t^{\Delta}(\widehat{\Delta W}_t,\Sigma_t)
% % = \mathbb{E}_{\varepsilon\sim\mathcal N(0,\Sigma_t)}
% %    \big[\hat{L}_{S_t}(W^{\mathrm{froz}}_t+\widehat{\Delta W}_t+\varepsilon)
% %    - \hat{L}_{S_t}(W^{\mathrm{froz}}_t+\widehat{\Delta W}_t)\big] \text{\,and\, } 
% %      \varepsilon\sim\mathcal N(0,\Sigma_t) \in  \mathcal W_{\Delta,t}$.
% %      % \begin{equation*}
% %      % \begin{aligned}
% %      % &\mathrm{RS}_t^{\Delta}(\hat{\Delta W}_t,\Sigma_t):= \\
% %      % &\qquad \mathbb{E}_{\varepsilon\sim\mathcal N(0,\Sigma_t)}
% %      %    \big[\hat{L}_{S_t}(W_\mathrm {frozen }+\hat{\Delta W}_t+\varepsilon) \\
% %      % & \qquad- \hat{L}_{S_t}(W_\mathrm {frozen }+\hat{\Delta W}_t)\big] \\
% %      % &\qquad \varepsilon\sim\mathcal N(0,\Sigma_t) \in \mathcal W_\Delta
% %      % \end{aligned}
% %      % \end{equation*}
% %      \end{lemma}

% % We refer to Appendix~\ref{proof:lemma2} for the proof.
% % % \subsection{KL Reduction to Effective Adapter Updates}
% % % \label{sec:pecl_support_reduction}

% % % We first formalize the posterior induced by LoRA adaptation. Let $q_t$ be a posterior over the trainable raw LoRA factors at task $t$. The corresponding pushforward posterior over effective updates is
% % % \begin{equation}
% % % \bar Q_t
% % % =
% % % (\phi_t)_{\#}q_t .
% % % \label{eq:pushforward_posterior}
% % % \end{equation}
% % % Similarly, a prior $p_t$ over raw LoRA factors induces an effective-update prior
% % % \begin{equation}
% % % \bar P_t
% % % =
% % % (\phi_t)_{\#}p_t .
% % % \label{eq:pushforward_prior}
% % % \end{equation}

% % % Let $T_t$ be the translation map from effective updates to full model parameters:
% % % \begin{equation}
% % % T_t(\Delta)
% % % =
% % % W_t^{\mathrm{froz}}+\Delta .
% % % \label{eq:translation_map}
% % % \end{equation}
% % % The model-level posterior and prior are therefore
% % % \begin{equation}
% % % Q_t
% % % =
% % % (T_t)_{\#}\bar Q_t,
% % % \qquad
% % % P_t
% % % =
% % % (T_t)_{\#}\bar P_t .
% % % \label{eq:model_level_pushforward}
% % % \end{equation}
% % % Both distributions are supported on the affine set
% % % \begin{equation}
% % % W_t^{\mathrm{froz}}+\mathcal W_{\Delta,t}^{\mathrm{train}} .
% % % \label{eq:affine_trainable_support}
% % % \end{equation}

% % % \begin{lemma}[Translation invariance of the PECL KL term]
% % % \label{lem:translation_kl_pecl}
% % % Assume $\bar Q_t\ll \bar P_t$. Let $Q_t=(T_t)_{\#}\bar Q_t$ and $P_t=(T_t)_{\#}\bar P_t$, where $T_t(\Delta)=W_t^{\mathrm{froz}}+\Delta$. Then
% % % \begin{equation}
% % % \mathrm{KL}(Q_t\Vert P_t)
% % % =
% % % \mathrm{KL}(\bar Q_t\Vert \bar P_t).
% % % \label{eq:translation_kl}
% % % \end{equation}
% % % \end{lemma}

% % % Lemma~\ref{lem:translation_kl_pecl} shows that the frozen component contributes no task-level KL complexity. Since the same frozen state is shared by the task prior and posterior, the divergence is entirely determined by distributions over effective adapter updates.

% % % This equality should not be confused with an equality in raw LoRA factor space. Because LoRA factorization is generally non-identifiable, raw factor coordinates may contain variation that is invisible at the effective model-update level.

% % % \begin{lemma}[Raw factor complexity upper bounds effective update complexity]
% % % \label{lem:factor_update_kl}
% % % Let $\bar Q_t=(\phi_t)_{\#}q_t$ and $\bar P_t=(\phi_t)_{\#}p_t$. Then
% % % \begin{equation}
% % % \mathrm{KL}(\bar Q_t\Vert \bar P_t)
% % % \le
% % % \mathrm{KL}(q_t\Vert p_t).
% % % \label{eq:data_processing_lora}
% % % \end{equation}
% % % \end{lemma}

% % % Lemma~\ref{lem:factor_update_kl} follows from the data processing inequality. It implies that raw-factor KL is not the invariant model-level complexity. The appropriate PAC-Bayesian object is the pushforward posterior over effective updates.



% % % \subsection{Task-Conditioned Sharpness on the Adapter Update Support}
% % % \label{sec:pecl_smoothed_risk}

% % % We now show that the empirical term carrying flatness admits the same support reduction. Let $\bar Q_t$ be an effective-update posterior supported on $\mathcal W_{\Delta,t}^{\mathrm{train}}$. The PECL empirical term for task $t$ is
% % % \begin{equation}
% % % \widehat L_{S_t}^{\Delta}(\bar Q_t)
% % % =
% % % \mathbb E_{\Delta W_t\sim \bar Q_t}
% % % \left[
% % % \widehat L_{S_t}
% % % \left(
% % % W_t^{\mathrm{froz}}+\Delta W_t
% % % \right)
% % % \right].
% % % \label{eq:pecl_smoothed_empirical_risk_general}
% % % \end{equation}

% % % When the posterior is instantiated as a perturbation distribution around the learned effective update $\widehat{\Delta W}_t$, we write
% % % \begin{equation}
% % % \Delta W_t
% % % =
% % % \widehat{\Delta W}_t+\epsilon_t,
% % % \qquad
% % % \epsilon_t\sim\nu_{\Delta,t},
% % % \qquad
% % % \operatorname{supp}(\nu_{\Delta,t})
% % % \subseteq
% % % \mathcal W_{\Delta,t}^{\mathrm{train}} .
% % % \label{eq:update_perturbation_distribution}
% % % \end{equation}
% % % Then
% % % \begin{equation}
% % % \widehat L_{S_t}^{\Delta}
% % % (\widehat{\Delta W}_t,\nu_{\Delta,t})
% % % =
% % % \mathbb E_{\epsilon_t\sim\nu_{\Delta,t}}
% % % \left[
% % % \widehat L_{S_t}
% % % \left(
% % % W_t^{\mathrm{froz}}
% % % +
% % % \widehat{\Delta W}_t
% % % +
% % % \epsilon_t
% % % \right)
% % % \right].
% % % \label{eq:pecl_smoothed_empirical_risk}
% % % \end{equation}

% % % This smoothed risk decomposes into empirical loss at the learned solution and an effective adapter-update sharpness term:
% % % \begin{equation}
% % % \begin{aligned}
% % % \widehat L_{S_t}^{\Delta}
% % % (\widehat{\Delta W}_t,\nu_{\Delta,t})
% % % &=
% % % \widehat L_{S_t}
% % % \left(
% % % W_t^{\mathrm{froz}}+\widehat{\Delta W}_t
% % % \right)
% % % +
% % % \mathrm{Sh}_{t\mid t}^{\Delta}
% % % (\widehat W_t,\nu_{\Delta,t}),
% % % \\
% % % \mathrm{Sh}_{t\mid t}^{\Delta}
% % % (\widehat W_t,\nu_{\Delta,t})
% % % &=
% % % \mathbb E_{\epsilon_t\sim\nu_{\Delta,t}}
% % % \left[
% % % \widehat L_{S_t}(\widehat W_t+\epsilon_t)
% % % -
% % % \widehat L_{S_t}(\widehat W_t)
% % % \right],
% % % \end{aligned}
% % % \label{eq:pecl_risk_sharpness_decomposition}
% % % \end{equation}
% % % where $\widehat W_t=W_t^{\mathrm{froz}}+\widehat{\Delta W}_t$.

% % % The notation $\mathrm{Sh}_{t\mid t}^{\Delta}$ emphasizes that sharpness is task-conditioned. More generally, for any evaluation task $i\le t$, we define
% % % \begin{equation}
% % % \mathrm{Sh}_{i\mid t}^{\Delta}
% % % (\widehat W_t,\nu_{\Delta,t})
% % % =
% % % \mathbb E_{\epsilon_t\sim\nu_{\Delta,t}}
% % % \left[
% % % \widehat L_{S_i}(\widehat W_t+\epsilon_t)
% % % -
% % % \widehat L_{S_i}(\widehat W_t)
% % % \right],
% % % \qquad
% % % \operatorname{supp}(\nu_{\Delta,t})
% % % \subseteq
% % % \mathcal W_{\Delta,t}^{\mathrm{train}} .
% % % \label{eq:task_conditioned_adapter_sharpness}
% % % \end{equation}
% % % Thus, the same checkpoint may be flat for one task and sharp for another, because each task defines a different empirical risk. For current-task generalization, the relevant term is $\mathrm{Sh}_{t\mid t}^{\Delta}$. For forgetting, the relevant diagnostic is the old-task sharpness under the later posterior support, namely $\mathrm{Sh}_{i\mid t}^{\Delta}$ for $i<t$.

% % % A local second-order approximation clarifies the geometry. Let $H_i(\widehat W_t)$ be the Hessian of $\widehat L_{S_i}$ at $\widehat W_t$, and let $\Pi_{\Delta,t}^{\mathrm{train}}$ be the orthogonal projector onto the local trainable adapter-update support. If $\epsilon_t$ has zero mean and covariance $\Sigma_{\Delta,t}$ supported on $\mathcal W_{\Delta,t}^{\mathrm{train}}$, then
% % % \begin{equation}
% % % \mathrm{Sh}_{i\mid t}^{\Delta}
% % % (\widehat W_t,\nu_{\Delta,t})
% % % \approx
% % % \frac{1}{2}
% % % \operatorname{tr}
% % % \left(
% % % \Pi_{\Delta,t}^{\mathrm{train}}
% % % H_i(\widehat W_t)
% % % \Pi_{\Delta,t}^{\mathrm{train}}
% % % \Sigma_{\Delta,t}
% % % \right).
% % % \label{eq:local_task_conditioned_adapter_sharpness}
% % % \end{equation}
% % % Hence, PECL sharpness is governed by task-specific curvature projected onto the trainable adapter-update geometry, rather than by curvature in the full ambient parameter space.



% % % \subsection{A Support-Reduced PAC-Bayesian Bound for LoRA PECL}
% % % \label{sec:pecl_bound}

% % % We now combine the KL support reduction and the smoothed-risk support reduction with the sequential PAC-Bayesian decomposition from Section~\ref{sec:flatness_cl}. Let $\bar\Pi_t$ denote the hyperposterior over effective-update priors after task $t$. For each task, a task prior $\bar P_t$ is drawn from $\bar\Pi_{t-1}$, and the task posterior $\bar Q_t$ is supported on $\mathcal W_{\Delta,t}^{\mathrm{train}}$.

% % % \begin{theorem}[Support reduction for LoRA based PECL]
% % % \label{thm:pecl_support_reduction}
% % % Assume the conditions of Theorem~\ref{thm:pathwise_pac}. Consider a frozen-backbone LoRA based PECL model in which, for each task $t$, the effective-update posterior $\bar Q_t$ and prior $\bar P_t$ are supported on $\mathcal W_{\Delta,t}^{\mathrm{train}}$, and their model-level pushforwards are
% % % \begin{equation}
% % % Q_t=(T_t)_{\#}\bar Q_t,
% % % \qquad
% % % P_t=(T_t)_{\#}\bar P_t .
% % % \label{eq:theorem_pushforward}
% % % \end{equation}
% % % Then, for any $\delta\in(0,1]$, with probability at least $1-\delta$ over $S_{1:T}$,
% % % \begin{equation}
% % % \begin{aligned}
% % % &
% % % \frac{1}{T}
% % % \sum_{t=1}^{T}
% % % \mathbb E_{\bar P_t\sim\bar\Pi_{t-1}}
% % % \mathbb E_{\Delta W_t\sim\bar Q_t}
% % % \left[
% % % L_{\mathcal D_t}
% % % \left(
% % % W_t^{\mathrm{froz}}+\Delta W_t
% % % \right)
% % % \right]
% % % \\
% % % &\le
% % % \frac{1}{T}
% % % \sum_{t=1}^{T}
% % % \mathbb E_{\bar P_t\sim\bar\Pi_{t-1}}
% % % \left[
% % % \widehat L_{S_t}^{\Delta}(\bar Q_t)
% % % \right]
% % % \\
% % % &\quad+
% % % \sqrt{
% % % \frac{
% % % \displaystyle
% % % \sum_{t=1}^{T}
% % % \mathbb E_{\bar P_t\sim\bar\Pi_{t-1}}
% % % \mathrm{KL}
% % % \left(
% % % \bar Q_t
% % % \Vert
% % % \bar P_t
% % % \right)
% % % +
% % % \sum_{t=1}^{T}
% % % \mathrm{KL}
% % % \left(
% % % \bar\Pi_t
% % % \Vert
% % % \bar\Pi_{t-1}
% % % \right)
% % % +
% % % \log\frac{1}{\delta}
% % % }{
% % % 2\displaystyle\sum_{t=1}^{T}(m_t-1)
% % % }
% % % }.
% % % \end{aligned}
% % % \label{eq:pecl_support_reduced_bound}
% % % \end{equation}
% % % \end{theorem}

% % % \paragraph{Proof sketch.}
% % % The sequential PAC-Bayesian result in Section~\ref{sec:flatness_cl} applies to the model-level posterior trajectory $\{Q_t\}_{t=1}^{T}$. Under frozen-backbone LoRA adaptation, Lemma~\ref{lem:translation_kl_pecl} gives
% % % \begin{equation}
% % % \mathrm{KL}(Q_t\Vert P_t)
% % % =
% % % \mathrm{KL}(\bar Q_t\Vert\bar P_t).
% % % \label{eq:proof_kl_substitution}
% % % \end{equation}
% % % Moreover, since $Q_t=(T_t)_{\#}\bar Q_t$ with $T_t(\Delta)=W_t^{\mathrm{froz}}+\Delta$, the empirical term becomes
% % % \begin{equation}
% % % \mathbb E_{W_t\sim Q_t}
% % % \left[
% % % \widehat L_{S_t}(W_t)
% % % \right]
% % % =
% % % \mathbb E_{\Delta W_t\sim\bar Q_t}
% % % \left[
% % % \widehat L_{S_t}
% % % \left(
% % % W_t^{\mathrm{froz}}+\Delta W_t
% % % \right)
% % % \right]
% % % =
% % % \widehat L_{S_t}^{\Delta}(\bar Q_t).
% % % \label{eq:proof_empirical_substitution}
% % % \end{equation}
% % % Substituting Eq.~\ref{eq:proof_kl_substitution} and Eq.~\ref{eq:proof_empirical_substitution} into the pathwise bound yields Eq.~\ref{eq:pecl_support_reduced_bound}.

% % % \subsection{PAC-Bayesian Bound for PECL}

% % In PECL, Lemma~\ref{lem:kl-decomp-pecl} and Lemma~\ref{lem:subspace-loss-pecl} reformulates the complexity and empirical terms entirely in the task adpter subspace $\mathcal W_{\Delta,t}$ through the task KL divergence $\mathrm{KL}\bigl(Q_t^\Delta \,\Vert\, P_t^\Delta\bigr)$ and the adapter subspace sharpness $\mathrm{TS}_t^{\Delta}(\widehat{\Delta W}_t,\Sigma_t)$ respectively. 
% % % See Fig.~\ref{fig: showPAC3} and Fig.~\ref{fig: KL_decompose} for an overview of the subspace reduction and the hierarchical construction.
% % Substituting these reductions into the bound from Theorem~\ref{thm:pathwise_pac}, we obtain the PAC-Bayesian bound for PECL as follows:
% % \begin{theorem}
% % \label{thm:pecl-bound}
% % Consider a PECL model with frozen backbone $W^{\mathrm{froz}}_t$ and Gaussian posteriors 
% % $Q_t^\Delta = \mathcal N(\widehat{\Delta W}_t,\Sigma_t)$ supported on $\mathcal W_{\Delta,t}$. 
% % Let $P_t^\Delta$ be task priors sampled from a hyper posterior $\mathcal P_{t-1}$ over adapter 
% % distributions. For any $\delta\in(0,1]$, with probability at least $1-\delta$ over $S_{1:T}$,
% % {%
% % \setlength{\abovedisplayskip}{2pt}
% % \setlength{\belowdisplayskip}{3pt}
% % \setlength{\abovedisplayshortskip}{3pt}
% % \setlength{\belowdisplayshortskip}{3pt}
% % \begin{equation}
% % \label{eq:pecl-main-bound}
% % % \resizebox{\linewidth}{!}{$
% % \begin{aligned}
% % \frac{1}{T}\sum_{t=1}^T
% % \mathbb E_{P_t^\Delta\sim\mathcal P_{t-1}}
% % \mathbb E_{W_t\sim Q_t}
% % L_{\mathcal D_t}(W_t) &
% % \le
% % \frac{1}{T}\sum_{t=1}^T
% % \mathbb E_{P_t^\Delta\sim\mathcal P_{t-1}}
% % \bigl[
% % \hat L_{S_t}(W^{\mathrm{froz}}_t + \widehat{\Delta W}_t)
% % +
% % \mathrm{TS}_t^{\Delta}(\widehat{\Delta W}_t,\Sigma_t)
% % \bigr] \\
% % &+
% % \sqrt{
% % \frac{
% % \displaystyle
% % \sum_{t=1}^T
% % \mathbb E_{P_t^\Delta\sim\mathcal P_{t-1}}
% % \mathrm{KL}\bigl(Q_t^\Delta \,\Vert\, P_t^\Delta\bigr)
% % +
% % \sum_{t=1}^T
% % \mathrm{KL}\bigl(\mathcal P_t\Vert\mathcal P_{t-1}\bigr)
% % +
% % \log\frac{1}{\delta}}
% % {2\displaystyle\sum_{t=1}^T (m_t - 1)}
% % }.
% % \end{aligned}
% % % $}
% % \end{equation}
% % }
% % \end{theorem}
% % Theorem~\ref{thm:pecl-bound} specializes the Theorem~\ref{thm:pathwise_pac} to PECL by exploiting the fact that all task dependent updates are confined to the permissible task subspace $\mathcal W_{\Delta,t}$. 
% % % This restriction is not merely technical.
% % Under a frozen backbone $\mathcal W_{\Delta,t}$ is exactly the set of directions along which the model can be modified by future tasks.
% % Therefore, in PECL the relevant notion of flatness for generalization is the effective geometry induced on $\mathcal W_{\Delta,t}$, rather than a global property of the full parameter space. We provide the proof in Appendix~\ref{proof:theorem}.

% % % Theorem~\ref{thm:pecl_support_reduction} identifies the flatness object that enters the LoRA PECL bound. The empirical term is the perturbation-smoothed empirical risk on $\mathcal W_{\Delta,t}^{\mathrm{train}}$, and the task-level complexity is the KL divergence between effective-update distributions on the same support. Therefore, the bound-relevant flatness in LoRA PECL is effective adapter-update flatness.


% % % \subsection{Implications for Forgetting and Evaluation}
% % % \label{sec:pecl_implications}

% % % The support-reduction result has three implications for analyzing flatness in LoRA based continual learning.

% % % First, full-parameter sharpness is not the primary bound-relevant object under a frozen backbone. Full-space perturbations may probe coordinates that are frozen during PECL adaptation. Such perturbations can be useful as a diagnostic of global model sensitivity, but they do not match the posterior support through which the learner moves at task $t$.

% % % Second, raw-factor sharpness is not invariant to the representation of the same effective update. Since different factor pairs can induce the same $\Delta W$, perturbing $(A,B)$ may change the measured sharpness without changing the model function. Lemma~\ref{lem:factor_update_kl} expresses the same issue at the complexity level: raw-factor KL can upper bound effective-update KL, but it is not necessarily the invariant model-level complexity.

% % % Third, the relevant flatness criterion is task-conditioned effective adapter-update sharpness. For current-task generalization, this corresponds to $\mathrm{Sh}_{t\mid t}^{\Delta}$. For forgetting, the more informative quantity is old-task sensitivity under later trainable supports, namely $\mathrm{Sh}_{i\mid t}^{\Delta}$ for $i<t$. This quantity measures how much the old-task loss changes along the directions that the current or future task can actually update.

% % % This distinction also clarifies how different LoRA organizations should be evaluated. SeqLoRA reuses a shared adapter branch, so the trainable support remains largely shared across tasks. In this case, old-task sharpness along the shared support is directly related to potential forgetting. IncLoRA allocates new branches across tasks, so previous branches become part of $W_t^{\mathrm{froz}}$ and the current posterior support is the newly trainable branch. In this case, the relevant question is whether the new branch's support intersects high-curvature directions of previous tasks. Orthogonal or projection-based variants aim to reduce such intersections by constraining $\mathcal W_{\Delta,t}^{\mathrm{train}}$ away from old sensitive directions.

% % % The theorem also yields falsifiable experimental controls. A raw-factor rescaling control should preserve the model function and effective update while changing raw-factor sharpness, which would show that raw-factor sharpness is not invariant. A frozen-coordinate perturbation control should test whether perturbing directions outside the trainable support explains forgetting. A random matched-subspace control should test whether dimensionality alone is sufficient, or whether alignment with $\mathcal W_{\Delta,t}^{\mathrm{train}}$ is necessary. These controls are more diagnostic than accuracy improvements alone because they directly test which flatness object is theoretically meaningful.

% % % The conclusion is therefore precise. Full-space sharpness and raw-factor sharpness may still be useful diagnostics, but they are not the flatness terms naturally induced by the LoRA PECL bound. The bound identifies $\mathcal W_{\Delta,t}^{\mathrm{train}}$ as the support on which flatness should be evaluated, and it identifies $\mathrm{Sh}_{i\mid t}^{\Delta}$ as the task-conditioned quantity most directly aligned with forgetting.



% % % % \section{PAC-Bayesian Bound for PECL}

% % % % \section{Which Flatness Matters in LoRA Based Continual Learning}
% % % % \label{sec:which_flatness_pecl}

% % % % Section~\ref{sec:flatness_cl} shows that flatness enters continual learning through perturbation smoothed empirical risk along a sequential posterior trajectory. We now ask the PECL specific question left open by that result: on which support should this smoothed risk be evaluated when the pretrained backbone is frozen and only LoRA induced updates are trainable?

% % % % This question is nontrivial because LoRA based PECL breaks the implicit alignment between perturbation directions, update directions, and interference directions. In full parameter training, these directions are often identified with the same ambient parameter space. In PECL, however, future adaptation is restricted to adapter induced update directions. A full model perturbation may probe frozen coordinates that no later task can update, while a raw factor perturbation may depend on the non unique parameterization of the same effective adapter update. Therefore, the relevant flatness object cannot be chosen by default. It must be derived from the support of the constrained posterior induced by PECL.

% % % % The main result of this section is a support reduction statement. When the sequential PAC Bayesian decomposition of Section~\ref{sec:flatness_cl} is specialized to frozen backbone LoRA adaptation, both the task level KL complexity and the perturbation smoothed empirical risk reduce to the effective adapter update support. This shows that the flatness relevant to LoRA based continual learning is adapter update flatness, rather than full space sharpness or raw factor space sharpness.

% % % % \subsection{Constrained Update Geometry in PECL}
% % % % \label{sec:pecl_geometry}

% % % % Let the model used after learning task $t$ be written as
% % % % \begin{equation}
% % % % W_t = W_t^{\mathrm{froz}} + \Delta W_t ,
% % % % \label{eq:pecl_affine_model}
% % % % \end{equation}
% % % % where $W_t^{\mathrm{froz}}$ denotes the parameters that are fixed during task $t$, including the frozen pretrained backbone and possibly previously frozen adapter components, and $\Delta W_t$ denotes the effective update induced by the trainable LoRA coordinates at task $t$.

% % % % For a LoRA layer, we write the effective update as
% % % % \begin{equation}
% % % % \Delta W_t = \phi_t(A_t,B_t),
% % % % \label{eq:lora_effective_update}
% % % % \end{equation}
% % % % where $(A_t,B_t)$ are the raw LoRA factors and $\phi_t$ maps them to the model level update injected into the corresponding weight matrix. 
% % % % % We deliberately use $\phi_t$ rather than fixing a specific convention such as $A_tB_t$ or $B_tA_t$, since different implementations may adopt different matrix layouts. 
% % % % The theoretical object of interest is not the raw factor pair itself, but the effective update that changes the model function.
% % % % We denote the admissible effective update support by
% % % % \begin{equation}
% % % % \mathcal W_{\Delta,t}
% % % % =
% % % % \left\{
% % % % \Delta W:\Delta W=\phi_t(A,B),\ (A,B)\in\mathcal A_t
% % % % \right\},
% % % % \label{eq:effective_update_support}
% % % % \end{equation}
% % % % where $\mathcal A_t$ is the feasible set of trainable LoRA factors at task $t$. 
% % % % % In local sharpness analysis, $\mathcal W_{\Delta,t}$ may be replaced by the tangent update support around the learned factors $(\widehat A_t,\widehat B_t)$:
% % % % % \begin{equation}
% % % % % \mathcal T_{\Delta,t}
% % % % % =
% % % % % \operatorname{Range}
% % % % % \left(
% % % % % J_{\phi_t}(\widehat A_t,\widehat B_t)
% % % % % \right),
% % % % % \label{eq:lora_tangent_support}
% % % % % \end{equation}
% % % % % where $J_{\phi_t}$ is the Jacobian of the effective update map. 
% % % % For notational simplicity, we use $\mathcal W_{\Delta,t}$ to refer to the admissible effective update support, with the understanding that local perturbation based sharpness is evaluated on its local tangent support when needed.

% % % % This distinction separates three different spaces. The first is the ambient full parameter space of $W$. The second is the raw LoRA factor space of $(A,B)$. The third is the effective adapter update support $\mathcal W_{\Delta,t}$. These spaces are not equivalent. Perturbing the full parameter space includes directions in frozen coordinates that are inaccessible to future LoRA updates. Perturbing raw factors may depend on the chosen factorization, since different factor pairs can induce the same effective update. For example, under a bilinear LoRA parameterization, the rescaling $(A,B)\mapsto (cA,B/c)$ can preserve the same $\Delta W$ while changing the geometry of the raw factor space.

% % % % \paragraph{Trainable support versus accumulated model support.}
% % % % A further distinction is essential for LoRA variants that expand adapter capacity over tasks.
% % % % We define the \emph{trainable effective update support at task $t$} as
% % % % \begin{equation}
% % % % \mathcal W_{\Delta,t}^{\mathrm{train}}
% % % % =
% % % % \bigl\{
% % % % \Delta W:\Delta W=\phi_t(A,B),\ (A,B)\in\mathcal A_t^{\mathrm{train}}
% % % % \bigr\},
% % % % \label{eq:trainable_support}
% % % % \end{equation}
% % % % where $\mathcal A_t^{\mathrm{train}}$ is the set of LoRA factors that are \emph{actively trained} at task $t$.
% % % % This is distinct from the \emph{accumulated model adapter support}
% % % % \begin{equation}
% % % % \mathcal W_{\Delta,1:t}^{\mathrm{model}}
% % % % =
% % % % \mathcal W_{\Delta,1}^{\mathrm{train}}
% % % % +
% % % % \mathcal W_{\Delta,2}^{\mathrm{train}}
% % % % +
% % % % \cdots
% % % % +
% % % % \mathcal W_{\Delta,t}^{\mathrm{train}},
% % % % \label{eq:accumulated_support}
% % % % \end{equation}
% % % % which spans all adapter-induced directions that have ever been trained and are present in the model at time $t$.

% % % % In methods such as SeqLoRA that reuse a single adapter branch, $\mathcal W_{\Delta,t}^{\mathrm{train}} \approx \mathcal W_{\Delta,k}^{\mathrm{train}}$ for all $k$, so the trainable and accumulated supports essentially coincide.
% % % % In IncLoRA-style methods that allocate a new adapter branch per task, $\mathcal W_{\Delta,t}^{\mathrm{train}}$ is the support of the newly created branch, while previously created adapters are frozen and absorbed into $W_t^{\mathrm{froz}}$, so
% % % % \begin{equation}
% % % % W_t^{\mathrm{froz}}
% % % % =
% % % % W_0 + \sum_{k<t}\Delta W_k,
% % % % \qquad
% % % % \Delta W_t \in \mathcal W_{\Delta,t}^{\mathrm{train}}.
% % % % \label{eq:inclora_frozen_decomp}
% % % % \end{equation}
% % % % This makes $W_t^{\mathrm{froz}}$ and $\mathcal W_{\Delta,t}^{\mathrm{train}}$ cleanly separated, but the two accumulated supports may still overlap in model-space directions.

% % % % For the PAC-Bayesian bound, the perturbation distribution $\nu_{\Delta,t}$ should be supported on $\mathcal W_{\Delta,t}^{\mathrm{train}}$, not on $\mathcal W_{\Delta,1:t}^{\mathrm{model}}$, because only the trainable support is the posterior support of the learner at task $t$.
% % % % Full-model diagnostic perturbations may probe $\mathcal W_{\Delta,1:t}^{\mathrm{model}}$ and are useful for measuring accumulated sensitivity, but they are not the bound-relevant object.
% % % % This distinction matters especially for forgetting analysis, where the relevant quantity is how sensitive old-task $i$'s loss is to perturbations along \emph{future trainable directions} $\mathcal W_{\Delta,t}^{\mathrm{train}}$ for $t>i$, as formalized in Section~\ref{sec:pecl_implications}.

% % % % Thus, the central question is not whether flatness matters, which was addressed in Section~\ref{sec:flatness_cl}, but which flatness matters under the constrained update geometry of PECL. The answer must be determined by the posterior support through which the learner can actually move at each task step.

% % % % \subsection{Posterior Support Reduction to Effective Adapter Updates}
% % % % \label{sec:pecl_support_reduction}

% % % % We first formalize the posterior induced by LoRA adaptation. Let $q_t$ be a posterior distribution over raw LoRA factors. The corresponding pushforward posterior over effective updates is
% % % % \begin{equation}
% % % % \bar Q_t = (\phi_t)_{\#} q_t .
% % % % \label{eq:pushforward_posterior}
% % % % \end{equation}
% % % % Similarly, if $p_t$ is a prior over raw LoRA factors, the induced prior over effective updates is
% % % % \begin{equation}
% % % % \bar P_t = (\phi_t)_{\#} p_t .
% % % % \label{eq:pushforward_prior}
% % % % \end{equation}

% % % % Let $T_t$ be the translation map from effective updates to model parameters:
% % % % \begin{equation}
% % % % T_t(\Delta)=W_t^{\mathrm{froz}}+\Delta .
% % % % \label{eq:translation_map}
% % % % \end{equation}
% % % % The model level posterior and prior are then
% % % % \begin{equation}
% % % % Q_t = (T_t)_{\#}\bar Q_t,
% % % % \qquad
% % % % P_t = (T_t)_{\#}\bar P_t .
% % % % \label{eq:model_level_pushforward}
% % % % \end{equation}
% % % % Therefore, both $Q_t$ and $P_t$ are supported on the same affine set
% % % % \begin{equation}
% % % % W_t^{\mathrm{froz}}+\mathcal W_{\Delta,t}.
% % % % \label{eq:affine_support}
% % % % \end{equation}

% % % % The following lemma states that the frozen part does not contribute to the KL complexity.

% % % % \begin{lemma}[Translation invariance of the PECL KL term]
% % % % \label{lem:translation_kl_pecl}
% % % % Assume $\bar Q_t\ll \bar P_t$. Let $Q_t=(T_t)_{\#}\bar Q_t$ and $P_t=(T_t)_{\#}\bar P_t$, where $T_t(\Delta)=W_t^{\mathrm{froz}}+\Delta$. Then
% % % % \begin{equation}
% % % % \mathrm{KL}(Q_t\Vert P_t)
% % % % =
% % % % \mathrm{KL}(\bar Q_t\Vert \bar P_t).
% % % % \label{eq:translation_kl}
% % % % \end{equation}
% % % % \end{lemma}

% % % % % \paragraph{Proof sketch.}
% % % % % The map $T_t$ is a measurable bijection between $\mathcal W_{\Delta,t}$ and the affine support $W_t^{\mathrm{froz}}+\mathcal W_{\Delta,t}$. Since $T_t$ is only a translation by the frozen parameter vector, it preserves the Radon Nikodym derivative between the posterior and prior. Therefore, the KL divergence is unchanged under this pushforward transformation.

% % % % Lemma~\ref{lem:translation_kl_pecl} shows that the task level PAC Bayesian complexity is entirely determined by the effective adapter update posterior. The frozen backbone is shared by prior and posterior, so it adds no divergence.

% % % % However, this equality should not be confused with a KL equality in the raw LoRA factor space. Since the LoRA factorization is generally not identifiable, the factor posterior may contain complexity that is invisible at the model level. The next lemma makes this precise.

% % % % \begin{lemma}[Raw factor complexity upper bounds effective update complexity]
% % % % \label{lem:factor_update_kl}
% % % % Let $\bar Q_t=(\phi_t)_{\#}q_t$ and $\bar P_t=(\phi_t)_{\#}p_t$. Then
% % % % \begin{equation}
% % % % \mathrm{KL}(\bar Q_t\Vert \bar P_t)
% % % % \le
% % % % \mathrm{KL}(q_t\Vert p_t).
% % % % \label{eq:data_processing_lora}
% % % % \end{equation}
% % % % \end{lemma}

% % % % % \paragraph{Proof sketch.}
% % % % % The result follows directly from the data processing inequality for KL divergence under the measurable map $\phi_t$. Mapping raw factors to effective updates can only discard information, and therefore cannot increase the divergence.

% % % % Lemma~\ref{lem:factor_update_kl} is important for interpreting LoRA flatness. Raw factor space complexity is not the invariant model level complexity. It can overestimate the effective update complexity because multiple factor pairs may correspond to the same $\Delta W$. Therefore, the pushforward posterior $\bar Q_t$ over effective updates is the appropriate object for PAC Bayesian analysis.

% % % % \subsection{Perturbation Smoothed Empirical Risk on the Adapter Update Support}
% % % % \label{sec:pecl_smoothed_risk}

% % % % We now show that the same support reduction applies to the empirical term that carries flatness. Let $\bar Q_t$ be a posterior over effective updates supported on $\mathcal W_{\Delta,t}$. The PECL perturbation smoothed empirical risk is
% % % % \begin{equation}
% % % % \widehat L_{S_t}^{\Delta}(\bar Q_t)
% % % % =
% % % % \mathbb E_{\Delta W_t\sim \bar Q_t}
% % % % \left[
% % % % \widehat L_{S_t}
% % % % \left(
% % % % W_t^{\mathrm{froz}}+\Delta W_t
% % % % \right)
% % % % \right].
% % % % \label{eq:pecl_smoothed_empirical_risk_general}
% % % % \end{equation}

% % % % When the posterior is instantiated as a perturbation distribution around the learned effective update $\widehat{\Delta W}_t$, we may write
% % % % \begin{equation}
% % % % \Delta W_t = \widehat{\Delta W}_t+\epsilon_t,
% % % % \qquad
% % % % \epsilon_t\sim \nu_{\Delta,t},
% % % % \qquad
% % % % \operatorname{supp}(\nu_{\Delta,t})\subseteq \mathcal W_{\Delta,t}.
% % % % \label{eq:update_perturbation_distribution}
% % % % \end{equation}
% % % % Then Eq.~\ref{eq:pecl_smoothed_empirical_risk_general} becomes
% % % % \begin{equation}
% % % % \widehat L_{S_t}^{\Delta}
% % % % (\widehat{\Delta W}_t,\nu_{\Delta,t})
% % % % =
% % % % \mathbb E_{\epsilon_t\sim \nu_{\Delta,t}}
% % % % \left[
% % % % \widehat L_{S_t}
% % % % \left(
% % % % W_t^{\mathrm{froz}}+\widehat{\Delta W}_t+\epsilon_t
% % % % \right)
% % % % \right].
% % % % \label{eq:pecl_smoothed_empirical_risk}
% % % % \end{equation}

% % % % This smoothed risk decomposes into the empirical loss at the learned solution and an adapter update sharpness term:
% % % % \begin{equation}
% % % % \widehat L_{S_t}^{\Delta}
% % % % (\widehat{\Delta W}_t,\nu_{\Delta,t})
% % % % =
% % % % \widehat L_{S_t}
% % % % \left(
% % % % W_t^{\mathrm{froz}}+\widehat{\Delta W}_t
% % % % \right)
% % % % +
% % % % \mathrm{Sh}_{t}^{\Delta}
% % % % (\widehat{\Delta W}_t,\nu_{\Delta,t}),
% % % % \label{eq:pecl_risk_sharpness_decomposition}
% % % % \end{equation}
% % % % where
% % % % \begin{equation}
% % % % \begin{aligned}
% % % % \mathrm{Sh}_{t}^{\Delta}
% % % % (\widehat{\Delta W}_t,\nu_{\Delta,t})
% % % % =
% % % % \mathbb E_{\epsilon_t\sim \nu_{\Delta,t}}
% % % % \bigg[
% % % % &
% % % % \widehat L_{S_t}
% % % % \left(
% % % % W_t^{\mathrm{froz}}+\widehat{\Delta W}_t+\epsilon_t
% % % % \right)
% % % % \\
% % % % &
% % % % -
% % % % \widehat L_{S_t}
% % % % \left(
% % % % W_t^{\mathrm{froz}}+\widehat{\Delta W}_t
% % % % \right)
% % % % \bigg].
% % % % \end{aligned}
% % % % \label{eq:adapter_update_sharpness}
% % % % \end{equation}

% % % % Eq.~\ref{eq:adapter_update_sharpness} defines the flatness object that enters the PECL bound. It evaluates loss elevation only along admissible effective update directions. Perturbations outside $\mathcal W_{\Delta,t}$ correspond to changes that the PECL learner cannot realize through future LoRA updates, and therefore do not form the natural support of the posterior trajectory.

% % % % A local second order approximation further clarifies this point. Let $H_t$ be the Hessian of $\widehat L_{S_t}$ at $\widehat W_t=W_t^{\mathrm{froz}}+\widehat{\Delta W}_t$, and let $\Pi_{\Delta,t}$ be the projector onto the local adapter update support. The restricted curvature is
% % % % \begin{equation}
% % % % H_{\Delta,t}
% % % % =
% % % % \Pi_{\Delta,t} H_t \Pi_{\Delta,t}.
% % % % \label{eq:restricted_hessian}
% % % % \end{equation}
% % % % If $\epsilon_t$ has zero mean and covariance $\Sigma_{\Delta,t}$ supported on $\mathcal W_{\Delta,t}$, then
% % % % \begin{equation}
% % % % \mathrm{Sh}_{t}^{\Delta}
% % % % (\widehat{\Delta W}_t,\nu_{\Delta,t})
% % % % \approx
% % % % \frac{1}{2}
% % % % \operatorname{tr}
% % % % \left(
% % % % H_{\Delta,t}\Sigma_{\Delta,t}
% % % % \right).
% % % % \label{eq:local_adapter_sharpness}
% % % % \end{equation}
% % % % Thus, PECL sharpness is governed by curvature and perturbation variance inside the update carrying geometry, rather than by curvature over the full ambient parameter space.

% % % % \subsection{A Support Reduction Bound for LoRA Based PECL}
% % % % \label{sec:pecl_bound}

% % % % We now combine the KL support reduction and the smoothed risk support reduction with the sequential PAC Bayesian decomposition of Section~\ref{sec:flatness_cl}. Let $\bar \Pi_t$ denote the hyper posterior over effective update priors after task $t$. For each task, a task prior $\bar P_t$ is drawn from $\bar \Pi_{t-1}$, and the task posterior $\bar Q_t$ is supported on $\mathcal W_{\Delta,t}$.

% % % % \begin{theorem}[Support reduction for LoRA based PECL]
% % % % \label{thm:pecl_support_reduction}
% % % % Assume the conditions of the sequential PAC Bayesian bound in Section~\ref{sec:flatness_cl}. Consider a frozen backbone LoRA based PECL model in which, for each task $t$, the effective update posterior $\bar Q_t$ and prior $\bar P_t$ are supported on $\mathcal W_{\Delta,t}$, and their model level pushforwards are defined by
% % % % \begin{equation}
% % % % Q_t=(T_t)_{\#}\bar Q_t,
% % % % \qquad
% % % % P_t=(T_t)_{\#}\bar P_t.
% % % % \label{eq:theorem_pushforward}
% % % % \end{equation}
% % % % Then, for any $\delta\in(0,1]$, with probability at least $1-\delta$ over the task samples $S_{1:T}$,
% % % % \begin{equation}
% % % % \begin{aligned}
% % % % &
% % % % \frac{1}{T}
% % % % \sum_{t=1}^{T}
% % % % \mathbb E_{\bar P_t\sim \bar \Pi_{t-1}}
% % % % \mathbb E_{\Delta W_t\sim \bar Q_t}
% % % % \left[
% % % % L_{\mathcal D_t}
% % % % \left(
% % % % W_t^{\mathrm{froz}}+\Delta W_t
% % % % \right)
% % % % \right]
% % % % \\
% % % % &\le
% % % % \frac{1}{T}
% % % % \sum_{t=1}^{T}
% % % % \mathbb E_{\bar P_t\sim \bar \Pi_{t-1}}
% % % % \left[
% % % % \widehat L_{S_t}^{\Delta}(\bar Q_t)
% % % % \right]
% % % % \\
% % % % &\quad+
% % % % \sqrt{
% % % % \frac{
% % % % \displaystyle
% % % % \sum_{t=1}^{T}
% % % % \mathbb E_{\bar P_t\sim \bar \Pi_{t-1}}
% % % % \mathrm{KL}
% % % % \left(
% % % % \bar Q_t
% % % % \Vert
% % % % \bar P_t
% % % % \right)
% % % % +
% % % % \sum_{t=1}^{T}
% % % % \mathrm{KL}
% % % % \left(
% % % % \bar \Pi_t
% % % % \Vert
% % % % \bar \Pi_{t-1}
% % % % \right)
% % % % +
% % % % \log\frac{1}{\delta}
% % % % }{
% % % % 2\displaystyle\sum_{t=1}^{T}(m_t-1)
% % % % }
% % % % }.
% % % % \end{aligned}
% % % % \label{eq:pecl_support_reduced_bound}
% % % % \end{equation}
% % % % \end{theorem}

% % % % \paragraph{Proof sketch.}
% % % % The sequential PAC Bayesian result in Section~\ref{sec:flatness_cl} applies to the model level posterior trajectory $\{Q_t\}_{t=1}^{T}$. Under frozen backbone LoRA adaptation, Lemma~\ref{lem:translation_kl_pecl} gives
% % % % \begin{equation}
% % % % \mathrm{KL}(Q_t\Vert P_t)
% % % % =
% % % % \mathrm{KL}(\bar Q_t\Vert \bar P_t).
% % % % \label{eq:proof_kl_substitution}
% % % % \end{equation}
% % % % Moreover, because $Q_t$ is the pushforward of $\bar Q_t$ through $T_t(\Delta)=W_t^{\mathrm{froz}}+\Delta$, its empirical term becomes
% % % % \begin{equation}
% % % % \mathbb E_{W_t\sim Q_t}
% % % % \left[
% % % % \widehat L_{S_t}(W_t)
% % % % \right]
% % % % =
% % % % \mathbb E_{\Delta W_t\sim \bar Q_t}
% % % % \left[
% % % % \widehat L_{S_t}
% % % % \left(
% % % % W_t^{\mathrm{froz}}+\Delta W_t
% % % % \right)
% % % % \right]
% % % % =
% % % % \widehat L_{S_t}^{\Delta}(\bar Q_t).
% % % % \label{eq:proof_empirical_substitution}
% % % % \end{equation}
% % % % Substituting Eq.~\ref{eq:proof_kl_substitution} and Eq.~\ref{eq:proof_empirical_substitution} into the trajectory level bound yields Eq.~\ref{eq:pecl_support_reduced_bound}.

% % % % Theorem~\ref{thm:pecl_support_reduction} specializes the continual learning bound to the constrained geometry of LoRA based PECL. It shows that the bound does not require controlling sharpness in arbitrary full parameter directions. The empirical term is the perturbation smoothed empirical risk on the effective adapter update support, and the task level complexity is the KL divergence between effective update distributions. Therefore, the flatness relevant to the PECL bound is adapter update flatness.

% % % % \subsection{Implications for Flatness, Forgetting, and Evaluation}
% % % % \label{sec:pecl_implications}

% % % % The support reduction result has three direct implications.

% % % % First, full parameter sharpness is not the primary theoretical object in frozen backbone PECL. Full space perturbations may still be useful as a diagnostic of global sensitivity after training. However, they include frozen coordinates that the learner cannot update through future LoRA adaptation. Such directions are outside the support of the constrained posterior trajectory and therefore are not the sharpness term naturally induced by Theorem~\ref{thm:pecl_support_reduction}.

% % % % Second, raw LoRA factor sharpness is not invariant to the parameterization of the same effective update. Since different factor pairs can induce the same $\Delta W$, perturbing $(A,B)$ can change the measured sharpness without changing the model level update. This makes raw factor sharpness an implementation dependent quantity. Lemma~\ref{lem:factor_update_kl} formalizes the same issue at the level of complexity: raw factor KL upper bounds effective update KL, but it is not necessarily the invariant model level complexity.

% % % % Third, the appropriate flatness criterion for LoRA based PECL is effective adapter update flatness. This criterion evaluates perturbations along the same geometry through which the learner adapts to new tasks and through which later tasks can interfere with earlier solutions. In other words, the relevant support is not the ambient parameter space, but the admissible update support $\mathcal W_{\Delta,t}$.

% % % % This conclusion also clarifies the relation to prior LoRA flatness analyses. Prior work correctly observes that perturbing raw LoRA factors is not equivalent to perturbing the effective merged update. Our result further shows why this distinction is essential in continual learning. The effective update support is not merely a parameterization issue. It is the support on which posterior movement, task adaptation, and future interference are realized.

% % % % Fourth, for forgetting analysis the relevant sharpness object is the \emph{cross-task interference sharpness}, not the old-task's own sharpness.
% % % % Corollary~\ref{cor:seen_task_forgetting} shows that forgetting of task $i$ after task $t$ is bounded through $\widehat L_{S_i}^{\Pi_t}(Q_t)$, where $\Pi_t$ is supported on $\mathcal W_{\Delta,t}^{\mathrm{train}}$.
% % % % This suggests measuring
% % % % \begin{equation}
% % % % \mathrm{Sh}_{i\leftarrow t}^{\Delta}
% % % % =
% % % % \max_{\epsilon\in \mathcal W_{\Delta,t}^{\mathrm{train}},\ \|\epsilon\|\le \rho}
% % % % \left[
% % % % \widehat L_{S_i}(\widehat W_{t-1}+\epsilon)
% % % % -
% % % % \widehat L_{S_i}(\widehat W_{t-1})
% % % % \right],
% % % % \label{eq:cross_task_sharpness}
% % % % \end{equation}
% % % % which quantifies how sensitively old task $i$'s loss changes under perturbations along the \emph{future task $t$'s trainable update directions}.
% % % % A local second-order approximation gives
% % % % \begin{equation}
% % % % \mathrm{Sh}_{i\leftarrow t}^{\Delta}
% % % % \approx
% % % % \frac{1}{2}
% % % % \operatorname{tr}
% % % % \!\left(
% % % % \Pi_{\Delta,t}^{\mathrm{train}}\,
% % % % H_i\,
% % % % \Pi_{\Delta,t}^{\mathrm{train}}
% % % % \right)
% % % % \rho^2,
% % % % \label{eq:cross_task_sharpness_hessian}
% % % % \end{equation}
% % % % where $H_i$ is the Hessian of $\widehat L_{S_i}$ at $\widehat W_{t-1}$ and $\Pi_{\Delta,t}^{\mathrm{train}}$ is the orthogonal projector onto the local tangent of $\mathcal W_{\Delta,t}^{\mathrm{train}}$.
% % % % Thus, forgetting is governed by how much old-task curvature leaks into the directions that the new task is actually able to update.
% % % % When $\Pi_{\Delta,t}^{\mathrm{train}} H_i \Pi_{\Delta,t}^{\mathrm{train}}$ has small trace, the new task's updates are approximately orthogonal to the old-task's high-curvature directions, and forgetting is controlled.
% % % % This motivates a more precise empirical question: beyond measuring current-task sharpness $\mathrm{Sh}_t^{\Delta}$, one should measure $\mathrm{Sh}_{i\leftarrow t}^{\Delta}$ and examine its correlation with actual forgetting $\mathcal F_{i,t}$.

% % % % Different LoRA organizations instantiate different evolution rules for $\mathcal W_{\Delta,t}^{\mathrm{train}}$.
% % % % SeqLoRA keeps an approximately fixed update support, so $\mathcal W_{\Delta,t}^{\mathrm{train}} \approx \mathcal W_{\Delta,k}^{\mathrm{train}}$ and the cross-task sharpness $\mathrm{Sh}_{i\leftarrow t}^{\Delta} \approx \mathrm{Sh}_{i}^{\Delta}$ reduces to the old-task's own sharpness along the shared support.
% % % % IncLoRA expands the admissible support by allocating new adapter branches, reducing direct overwrite; however, cross-task interference can still occur when the new branch's support intersects old-task high-curvature directions, making $\mathrm{Sh}_{i\leftarrow t}^{\Delta}$ the informative diagnostic rather than old-task sharpness.
% % % % Orthogonal or projection-based variants such as OLoRA and InfLoRA explicitly aim to reduce $\Pi_{\Delta,t}^{\mathrm{train}} H_i \Pi_{\Delta,t}^{\mathrm{train}}$ by constraining new update directions to avoid previously sensitive regions.
% % % % Theorem~\ref{thm:pecl_support_reduction} provides a common criterion: their effect should be assessed through the support-reduced sharpness on $\mathcal W_{\Delta,t}^{\mathrm{train}}$ and through the cross-task interference $\mathrm{Sh}_{i\leftarrow t}^{\Delta}$.

% % % % The conclusion of this section is therefore not that full space sharpness is meaningless, but that it is not the support-reduced sharpness term naturally induced by LoRA-based PECL.
% % % % Full-space probes may diagnose global sensitivity, and raw-factor probes may reveal implementation-level behavior.
% % % % However, the PAC-Bayesian trajectory bound identifies the effective adapter update support $\mathcal W_{\Delta,t}^{\mathrm{train}}$ as the space in which sharpness should be controlled during continual adaptation, and the cross-task interference sharpness $\mathrm{Sh}_{i\leftarrow t}^{\Delta}$ as the quantity most directly aligned with forgetting.
% % % % This gives a principled empirical criterion: a flatness measure should be judged by whether it evaluates perturbations along the trainable geometry through which the learner adapts and future tasks can interfere.


% % % % % In frozen-backbone parameter-efficient continual learning, the model at task $t$ is written as
% % % % % $$
% % % % % W_t
% % % % % =
% % % % % W_t^{\mathrm{froz}}
% % % % % +
% % % % % \Delta W_t,
% % % % % $$
% % % % % where $W_t^{\mathrm{froz}}$ is fixed during task-$t$ adaptation and $\Delta W_t$ is the trainable low-rank update.
% % % % % The update is restricted to a task-dependent admissible geometry
% % % % % $$
% % % % % \Delta W_t\in \mathcal W_{\Delta,t}.
% % % % % $$
% % % % % For LoRA, with adapter coordinates $\theta_\Delta=(A,B)$ and effective update $\Delta W(\theta_\Delta)=AB$, the locally admissible update geometry around the learned solution $\hat\theta_{\Delta,t}$ is
% % % % % $$
% % % % % \mathcal W_{\Delta,t}
% % % % % =
% % % % % \operatorname{Im}
% % % % % \left(
% % % % % J_{\mathrm{LoRA}}(\hat\theta_{\Delta,t})
% % % % % \right),
% % % % % $$
% % % % % where $J_{\mathrm{LoRA}}(\hat\theta_{\Delta,t})$ is the Jacobian of the map $\theta_\Delta\mapsto\Delta W(\theta_\Delta)$.
% % % % % Different LoRA-based continual learning methods, such as SeqLoRA, IncLoRA, and constrained variants, can be viewed as different rules for evolving the sequence of admissible geometries $\{\mathcal W_{\Delta,t}\}_{t=1}^T$.




% % % % % \subsection{Continual Learning with LoRA-Constrained Updates}

% % % % % In continual learning, the learner observes a sequence of tasks $t=1,\ldots,T$.
% % % % % Each task provides a sample $S_t=\{z_{tj}\}_{j=1}^{m_t}$ drawn i.i.d. from a task distribution $\mathcal D_t$.
% % % % % % We do not require the task distributions to be sampled exchangeably from a fixed environment, since the realized task order is central to forgetting.
% % % % % Let $Q_t$ denote the posterior over model parameters after learning task $t$.
% % % % % The risk of an earlier task $i\le t$ under the later posterior $Q_t$ is
% % % % % $$
% % % % % L_{\mathcal D_i}(Q_t)
% % % % % =
% % % % % \mathbb E_{W\sim Q_t}L_{\mathcal D_i}(W).
% % % % % $$

% % % % % For $i<t$, the forgetting of task $i$ after learning task $t$ can be measured by
% % % % % $$
% % % % % \mathfrak F_{i,t}
% % % % % =
% % % % % L_{\mathcal D_i}(Q_t)-L_{\mathcal D_i}(Q_i),
% % % % % $$
% % % % % where $Q_i$ is the posterior immediately after task $i$.

% % % % % In frozen-backbone PECL, task-dependent adaptation is confined to a low-rank update added to a frozen backbone:
% % % % % $$
% % % % % W_t
% % % % % =
% % % % % W_t^{\mathrm{froz}}+\Delta W_t.
% % % % % $$
% % % % % The trainable increment is not arbitrary, but is restricted to a task-dependent admissible update geometry
% % % % % $$
% % % % % \Delta W_t\in\mathcal W_{\Delta,t}.
% % % % % $$

% % % % % Let $\theta_{\Delta}=(A,B)$ denote LoRA coordinates and let $\Delta W(\theta_{\Delta})=AB$ be the effective weight update.
% % % % % Around the learned solution $\hat\theta_{\Delta,t}$, the locally admissible update directions are given by
% % % % % $$
% % % % % \mathcal W_{\Delta,t}
% % % % % =
% % % % % \operatorname{Im}\left(J_{\mathrm{LoRA}}(\hat\theta_{\Delta,t})\right),
% % % % % $$
% % % % % where $J_{\mathrm{LoRA}}(\hat\theta_{\Delta,t})$ is the Jacobian of the map $\theta_{\Delta}\mapsto\Delta W(\theta_{\Delta})$.

% % % % % Different LoRA-based CL methods induce different evolution rules for $\mathcal W_{\Delta,t}$.
% % % % % SeqLoRA reuses a single adapter and therefore maximizes subspace sharing across tasks.
% % % % % IncLoRA expands the adapter set over time and enlarges the available update geometry.
% % % % % Constraint-based variants such as OLoRA, InfLoRA, and SDLoRA further regularize or decompose this geometry to reduce cross-task interference.
% % % % % Our analysis applies to these organizations through the sequence of admissible geometries $\{\mathcal W_{\Delta,t}\}_{t=1}^T$.




% % % % % \subsection{Question I: Can Perturbation-Smoothed Risk Explain Forgetting?}

% % % % % The PAC-Bayesian view explains why perturbation-smoothed empirical risk is relevant to generalization.
% % % % % However, continual learning requires a different risk comparison.
% % % % % The learner is not evaluated only on the current posterior $Q_t$ and current distribution $\mathcal D_t$, but also on earlier distributions $\mathcal D_i$ under later posteriors $Q_t$ for $i<t$.
% % % % % Forgetting is therefore a pathwise phenomenon along the posterior trajectory, rather than a single-task generalization gap.

% % % % % Existing lifelong PAC-Bayesian analyses typically study transfer across tasks sampled from a fixed hyperdistribution, while recent cumulative-risk bounds control risks along a task sequence.
% % % % % However, these formulations do not directly isolate how perturbation-smoothed empirical risk controls old-task risk under later posteriors, nor do they decompose forgetting into empirical degradation and posterior-trajectory complexity.

% % % % % This leads to our first theoretical question:

% % % % % \textbf{Q1.}
% % % % % Can perturbation-smoothed empirical risk provide a PAC-Bayesian control of old-task risk under later posteriors, and can forgetting be decomposed into smoothed empirical-risk degradation and a complexity term along the posterior trajectory?

% % % % % Section~\ref{sec:seq-pac} answers this question by deriving a sequential PAC-Bayesian decomposition with history-dependent priors and a hierarchical drift penalty.

% % % % % \subsection{Question I: Can Perturbation-Smoothed Risk Explain Forgetting?}

% % % % % The PAC-Bayesian view links posterior-averaged generalization to posterior-averaged empirical risk and a KL complexity term.
% % % % % In continual learning, however, the relevant comparison is not only the current-task risk $L_{\mathcal D_t}(Q_t)$.
% % % % % The learner must also maintain performance on previous task distributions $\mathcal D_i$ under later posteriors $Q_t$ for $i<t$.
% % % % % Thus, forgetting is a pathwise phenomenon along the posterior trajectory $Q_i\rightarrow Q_t$.

% % % % % For a previous task $i$, forgetting after learning task $t$ is measured by
% % % % % $$
% % % % % \mathfrak F_{i,t}
% % % % % =
% % % % % L_{\mathcal D_i}(Q_t)
% % % % % -
% % % % % L_{\mathcal D_i}(Q_i).
% % % % % $$
% % % % % This quantity asks whether the later posterior $Q_t$ still generalizes to the earlier distribution $\mathcal D_i$.
% % % % % Therefore, explaining forgetting requires more than a single-task PAC-Bayesian bound.
% % % % % It requires a sequential bound that compares old-task risks under changing posteriors and separates empirical degradation from the complexity of moving along the posterior trajectory.

% % % % % Existing lifelong PAC-Bayesian analyses mainly study transfer under a shared or exchangeable task environment.
% % % % % Recent cumulative-risk bounds control risks along a task sequence, but they do not explicitly identify how perturbation-smoothed empirical risk contributes to old-task risk under later posteriors, nor do they isolate the complexity induced by cross-task posterior drift.

% % % % % This leads to our first theoretical question:

% % % % % \textbf{Q1.}
% % % % % Can perturbation-smoothed empirical risk provide PAC-Bayesian control of past-task generalization under later posteriors, and can forgetting be decomposed into smoothed empirical-risk degradation and posterior-trajectory complexity?

% % % % % Section~\ref{sec:seq-pac} answers this question by deriving a sequential hierarchical PAC-Bayesian decomposition with history-dependent task priors.
% % % % % The resulting bound separates within-task adaptation complexity from hyperposterior drift complexity, making cross-task interference explicit across task transitions.


% % % % % % \subsection{Question II: Which Perturbation Support Is Relevant in PECL?}

% % % % % % The second issue is specific to PECL.
% % % % % % In full-parameter training, the perturbation space used to define sharpness largely coincides with the space in which optimization can update the model.
% % % % % % In frozen-backbone PECL, this equivalence breaks down.
% % % % % % The loss is defined on the full model, but future adaptation is constrained to the admissible update geometry $\mathcal W_{\Delta,t}$.

% % % % % % This creates three non-equivalent perturbation scopes:
% % % % % % full-parameter perturbations in the ambient space,
% % % % % % factor-coordinate perturbations in the LoRA parameters $(A,B)$,
% % % % % % and effective-update perturbations supported on $\mathcal W_{\Delta,t}$.
% % % % % % The first includes directions that are frozen during future adaptation.
% % % % % % The second depends on the chosen LoRA parameterization.
% % % % % % The third is defined directly in weight-increment space and captures the directions that PECL can locally realize.

% % % % % % For a perturbation law $\Pi_{\Delta,t}$ supported on $\mathcal W_{\Delta,t}$, the PECL-specific perturbation-smoothed empirical risk is
% % % % % % $$
% % % % % % \hat L_{S_t}^{\Delta,\Pi}(W_t)
% % % % % % =
% % % % % % \mathbb E_{\epsilon\sim\Pi_{\Delta,t}}
% % % % % % \hat L_{S_t}(W_t+\epsilon),
% % % % % % \qquad
% % % % % % \operatorname{supp}(\Pi_{\Delta,t})\subseteq \mathcal W_{\Delta,t}.
% % % % % % $$

% % % % % % This leads to our second theoretical question:

% % % % % % \textbf{Q2.}
% % % % % % Under frozen-backbone LoRA adaptation, which perturbation support determines the bound-relevant smoothed empirical risk and KL complexity: the ambient parameter space, the raw LoRA factor space, or the effective adapter-update geometry?


% % % % % \subsection{Question I: Can Perturbation-Smoothed Risk Explain Forgetting?}

% % % % % The PAC-Bayesian view links posterior-averaged population risk to posterior-averaged empirical risk and a KL complexity term.
% % % % % In standard supervised learning, this gives a single-posterior generalization statement: a posterior $Q$ learned from a sample $S$ is evaluated on the same underlying distribution $\mathcal D$.
% % % % % Continual learning changes this object of analysis.
% % % % % The learner does not output a single posterior, but a posterior trajectory $Q_1,\ldots,Q_T$, and each later posterior $Q_t$ must remain reliable on all previously encountered task distributions $\mathcal D_i$ for $i\le t$.

% % % % % Under this view, forgetting is the degradation of past-task generalization along the posterior trajectory.
% % % % % For an earlier task $i<t$, its risk-form forgetting after learning task $t$ can be written as
% % % % % $$
% % % % % \mathfrak F_{i,t}
% % % % % =
% % % % % L_{\mathcal D_i}(Q_t)
% % % % % -
% % % % % L_{\mathcal D_i}(Q_i).
% % % % % $$
% % % % % This quantity asks whether replacing the task-specific posterior $Q_i$ by the later posterior $Q_t$ increases the population risk on the same earlier distribution $\mathcal D_i$.
% % % % % It is therefore the risk-based analogue of standard backward-transfer style forgetting metrics, up to the sign convention and the use of risk rather than accuracy.

% % % % % This perspective also clarifies the role of trajectory-level metrics such as average anytime accuracy.
% % % % % Such metrics evaluate each intermediate model on all tasks observed so far and then average across learning stages.
% % % % % Their risk-based counterpart measures how well the posterior trajectory preserves generalization over the expanding set of encountered task distributions.
% % % % % Thus, explaining continual learning performance requires more than bounding current-task risk $L_{\mathcal D_t}(Q_t)$.
% % % % % It requires controlling old-task risks $L_{\mathcal D_i}(Q_t)$ under later posteriors.

% % % % % Perturbation-smoothed empirical risk is central to this question because SAM-style objectives optimize posterior-averaged empirical loss under local perturbations.
% % % % % However, existing PAC-Bayesian analyses do not directly explain how such smoothed empirical risk controls past-task population risk under later posteriors.
% % % % % Classical lifelong PAC-Bayes studies transfer under a shared or exchangeable task environment, but they do not model pathwise degradation along a realized task sequence.
% % % % % Recent cumulative-risk bounds control risks over task sequences, but they do not explicitly decompose forgetting into old-task smoothed empirical-risk degradation and the complexity induced by posterior drift across tasks.

% % % % % This leads to our first theoretical question:

% % % % % \textbf{Q1.}
% % % % % Can perturbation-smoothed empirical risk provide PAC-Bayesian control of past-task generalization under later posteriors, and can forgetting be decomposed into smoothed empirical-risk degradation and posterior-trajectory complexity?

% % % % % Section~\ref{sec:seq-pac} answers this question by deriving a sequential hierarchical PAC-Bayesian decomposition with history-dependent task priors.
% % % % % The resulting bound separates within-task adaptation complexity from hyperposterior drift complexity, making cross-task interference explicit across task transitions.



% % % % % \subsection{Question II: Which Perturbation Support Is Relevant in PECL?}

% % % % % The second question is specific to frozen-backbone PECL.
% % % % % In full-parameter training, the perturbation space used to define sharpness largely coincides with the space in which optimization can update the model.
% % % % % In PECL, this equivalence no longer holds.
% % % % % The loss is evaluated on the full model, but future adaptation can only move the model through the admissible update geometry $\mathcal W_{\Delta,t}$.

% % % % % This creates three non-equivalent perturbation scopes.
% % % % % Full-parameter perturbations act in the ambient parameter space and include directions that are frozen during adaptation.
% % % % % LoRA factor-coordinate perturbations act on $(A,B)$, but depend on a non-unique parameterization of the same effective update.
% % % % % Effective-update perturbations act directly in weight-increment space and are supported on the locally realizable geometry $\mathcal W_{\Delta,t}$.

% % % % % For a perturbation law $\Pi_{\Delta,t}$ supported on $\mathcal W_{\Delta,t}$, the PECL-specific perturbation-smoothed empirical risk is
% % % % % $$
% % % % % \hat L_{S_t}^{\Delta,\Pi}(W_t)
% % % % % =
% % % % % \mathbb E_{\epsilon\sim\Pi_{\Delta,t}}
% % % % % \hat L_{S_t}(W_t+\epsilon),
% % % % % \qquad
% % % % % \operatorname{supp}(\Pi_{\Delta,t})\subseteq \mathcal W_{\Delta,t}.
% % % % % $$

% % % % % This leads to our second theoretical question:

% % % % % \textbf{Q2.}
% % % % % Under frozen-backbone LoRA adaptation, do the bound-relevant perturbation-smoothed empirical risk and KL complexity reduce to the effective adapter-update geometry $\mathcal W_{\Delta,t}$, rather than the ambient parameter space or the raw LoRA factor space?

% % % % % Section~\ref{sec:pecl-spec} answers this question by analyzing the pushforward posterior induced by LoRA parameters on effective weight updates.
% % % % % The result shows that, under the frozen-backbone constraint, both the smoothed empirical risk term and the KL complexity are supported on the admissible adapter-update geometry.
% % % % % Consequently, the theoretically relevant sharpness object in PECL is the basis-invariant sharpness of effective adapter updates.




% % % % % \section{PECL Support Reduction under LoRA-Constrained Updates}












% % % \section{Theory-Derived Predictions and Empirical Protocol}
% % % \label{sec:predictions}

% % % Sections~\ref{sec:flatness_cl} and~\ref{sec:which_flatness_pecl} yield four concrete predictions about how sharpness should behave in LoRA-based PECL.
% % % These predictions translate the theoretical analysis into testable hypotheses that organize the empirical investigation in Section~\ref{sec:experiments}.

% % % \paragraph{Prediction 1 (Effective-update sharpness aligns with forgetting).}
% % % The cross-task interference sharpness $\mathrm{Sh}_{i\leftarrow t}^{\Delta}$ (Eq.~\ref{eq:cross_task_sharpness}), which perturbs along the \emph{trainable support} $\mathcal W_{\Delta,t}^{\mathrm{train}}$ and measures old-task sensitivity in that direction, should be the sharpness object most predictive of forgetting $\mathcal F_{i,t}$.
% % % Current-task adapter-update sharpness $\mathrm{Sh}_t^{\Delta}$ should be more predictive of current-task generalization.
% % % Full-parameter sharpness $\mathrm{Sh}^{\mathrm{full}}$ and raw-factor sharpness $\mathrm{Sh}^{A,B}$ should be less predictive of forgetting than support-reduced objects, because they measure sensitivity along directions that include frozen coordinates or depend on non-invariant parameterization choices.

% % % \paragraph{Prediction 2 (Raw-factor sharpness is not parameterization-invariant).}
% % % The same effective adapter update $\Delta W$ can be realized by a family of factor pairs $(cA,B/c)$ for any scalar $c\ne0$.
% % % Since this rescaling leaves the model function and loss landscape unchanged, both the effective-update sharpness $\mathrm{Sh}_t^{\Delta}$ and model-level quantities should be stable under rescaling.
% % % However, raw-factor sharpness $\mathrm{Sh}^{A,B}$, which evaluates perturbations in $(A,B)$-space, will change because the local geometry of the factor space depends on $c$.
% % % This provides a clean negative control: invariance under rescaling is a necessary condition for a sharpness measure to reflect an intrinsic model-level property.

% % % \paragraph{Prediction 3 (Full-space perturbations are misaligned with PECL forgetting).}
% % % Full-parameter perturbations probe directions in $W_t^{\mathrm{froz}}$ that no future LoRA update can realize.
% % % Although full-space sharpness may measure global loss sensitivity, Theorem~\ref{thm:pecl_support_reduction} shows that it is not the PAC-Bayesian object induced by the PECL posterior.
% % % Consequently, frozen-coordinate perturbations may produce large loss increases, but they should be systematically less predictive of actual LoRA-induced forgetting than perturbations restricted to $\mathcal W_{\Delta,t}^{\mathrm{train}}$.
% % % Furthermore, a random subspace of the same dimension as $\mathcal W_{\Delta,t}^{\mathrm{train}}$ should be less predictive of forgetting than the true trainable support, showing that it is not merely \emph{any} low-dimensional reduction but specifically the task-permissible geometry that matters.

% % % \paragraph{Prediction 4 (LoRA organization changes the support-forgetting relation).}
% % % Different LoRA designs induce different trainable support processes $\{\mathcal W_{\Delta,t}^{\mathrm{train}}\}_{t=1}^T$.
% % % \begin{itemize}
% % % \item \textbf{SeqLoRA:} The support is approximately fixed across tasks.
% % % Cross-task sharpness $\mathrm{Sh}_{i\leftarrow t}^{\Delta}$ is close to old-task sharpness $\mathrm{Sh}_i^{\Delta}$, so current-task adapter-update flatness should directly predict forgetting.
% % % \item \textbf{IncLoRA:} Each task introduces a new adapter branch with support $\mathcal W_{\Delta,t}^{\mathrm{train}}$ disjoint from previous branches in the factor space.
% % % The relevant diagnostic is $\mathrm{Sh}_{i\leftarrow t}^{\Delta}$ (how much the new branch's support intersects old-task curvature), not old-task sharpness.
% % % \item \textbf{OLoRA / InfLoRA:} These methods impose orthogonality or projection constraints to reduce $\Pi_{\Delta,t}^{\mathrm{train}}H_i\Pi_{\Delta,t}^{\mathrm{train}}$.
% % % If successful, $\mathrm{Sh}_{i\leftarrow t}^{\Delta}$ should be lower than in SeqLoRA or IncLoRA even when current-task sharpness is similar, leading to reduced forgetting.
% % % \end{itemize}
% % % Analyzing each method through the lens of support evolution and cross-task sharpness rather than raw performance benchmarking provides a mechanistic explanation for when and why each design reduces forgetting.



\section{Experiments}
\label{sec:experiments}

% Our experiments are organized as empirical probes of the support-reduction claim in Theorem~\ref{thm:pecl_support_reduction}, not as a method-centric benchmark.
% Each RQ targets a specific aspect of whether the adapter-update geometry $\mathcal W_{\Delta,t}^{\mathrm{train}}$ governs forgetting:
% RQ1 tests the causal role of perturbation support during training;
% RQ2 tests whether adapter-update sharpness is the most explanatory post-hoc diagnostic;
% RQ3 tests how perturbation types (SAM, RWP, GAM) behave once the support is fixed to $\mathcal W_{\Delta,t}^{\mathrm{train}}$;
% RQ4 tests robustness of the support principle across LoRA organizations, datasets, and model families.

\subsection{Experimental Setup}
\label{sec:exp_setup}

\noindent\textbf{Datasets}\quad
We evaluate PECL methods on vision benchmarks derived from ImageNet. ImageNet-R~\cite{ImagenetR_Hendrycks_2021_ICCV} serves as our main class-incremental benchmark following prior work~\cite{liangInfLoRAInterferenceFreeLowRank2024a}.
To assess robustness under corruptions and perturbations, we use the Tiny ImageNet variants of ImageNet-C and ImageNet-P (200 classes)~\cite{hendrycks2018benchmarking}.

\noindent\textbf{Models}\quad 
All experiments use a ViT-B/16 backbone~\cite{dosovitskiy2020image} pre-trained on ImageNet-21K.
We attach LoRA adapters to each transformer block and consider three organizations: SeqLoRA, IncLoRA, and OLoRA~\cite{wang2023orthogonal}.

% To study perturbation \emph{scope}, we adopt {full-parameter} setting~\cite{li2024flatlora} by injecting isotropic Gaussian perturbation $\mathrm{RS}_S^{(0)}({W}, \sigma^2 
% I)$ into all model parameters.

% To study perturbation \emph{type}, we compare SAM~\cite{foretSharpnessAwareMinimizationEfficiently2021}, RWP~\cite{pmlr-v151-bisla22a}, and GAM~\cite{Zhang_GAM}, which optimize different sharpness objectives: SAM minimizes $\mathrm{AS}^{(0)}_S(W,\rho)$, RWP targets $\mathrm{RS}^{(0)}_S(W,\Sigma)$, and GAM targets $\mathrm{AS}^{(1)}_S(W,\rho)$. For a controlled comparison in PECL, all methods apply their perturbations only to the trainable LoRA coordinates. Unless otherwise stated, we fix the perturbation radius to $\rho=0.05$, following the standard empirical setting used in SAM, and use a dedicated radius sweep only when diagnosing the scope effect.



\noindent\textbf{Metrics}\quad
The average class-incremental accuracy at step $t$ is
$
\operatorname{Acc}_t = \frac{1}{t} \sum_{j=1}^{t} a_{tj}
$,  where $a_{tj}$ denotes the accuracy on task $j$ after learning task $t$. 
We report accuracy $\operatorname{Acc}_t^{\mathrm{cls}}$ using the linear classifier and
% , and also evaluate representation quality using Nearest Class Mean (NCM) classifier, denoted by $\operatorname{Acc}_t^{\mathrm{nme}}$.
compute the corresponding Nearest-Class-Mean Classifier (NCM) results
% Nearest Mean of Exemplars (NME) metrics 
$(\operatorname{Acc}_t^{\mathrm{ncm}})$ using class prototypes in the feature space.
For the controlled trajectory experiments, let $R_{t,i}$ denote the linear-classifier accuracy on task $i$ after learning task $t$.
We report final average accuracy $\mathrm{FAA}=\frac{1}{T}\sum_{i=1}^{T}R_{T,i}$, average anytime accuracy $\mathrm{AAA}=\frac{2}{T(T+1)}\sum_{t=1}^{T}\sum_{i=1}^{t}R_{t,i}$, and final backward transfer $\mathrm{BWT}=\frac{1}{T-1}\sum_{i=1}^{T-1}(R_{T,i}-R_{i,i})$.
Higher $\mathrm{FAA}$, $\mathrm{AAA}$, and $\mathrm{BWT}$ are better.

\noindent\textbf{Controlled perturbation variants.}\quad
Our main empirical probes use SeqLoRA on ImageNet-R with rank $r=16$, $T=20$ tasks of 10 classes each, and ViT-B/16.
All variants update the same trainable variables, namely the LoRA factors and the classifier head.
They differ only in the temporary perturbation used to form the perturbed loss.
We use the notation \texttt{direction-support}: \texttt{sam} denotes the adversarial SAM ascent direction, \texttt{random} denotes a normalized Gaussian direction, and the supports are raw LoRA factors (\texttt{factor}), effective adapter-update tangent directions (\texttt{delta}), merged effective weights (\texttt{full}), all raw parameters (\texttt{all}), or frozen backbone coordinates (\texttt{frozen}).
Thus, \texttt{sam_delta} and \texttt{random_delta} use the same effective adapter-update support but different perturbation directions.
This separation is important: \texttt{sam_random}, if used, would mean a SAM direction projected onto a random support and is not the same object as \texttt{random_delta}.

% \noindent\textbf{Experimental objective.}\quad
% The central question is whether, in LoRA-based PECL, the flatness term relevant to continual generalization should be evaluated over the effective adapter-update support $\mathcal W_{\Delta,t}^{\mathrm{train}}$, rather than over the full parameter space or raw LoRA factor coordinates.
% The manipulated variable across RQ1--RQ3 is either the temporary perturbation support used during training (RQ1) or the support over which post-hoc sharpness is measured (RQ2, RQ3).
% In all variants, the actual gradient update is applied only to LoRA parameters $(A_t, B_t)$; the frozen pretrained backbone is never updated.
% This design isolates the effect of perturbation support from confounds such as parameter count, learning rate, or LoRA rank.

% \noindent\textbf{Datasets.}\quad
% % \emph{Primary (RQ1--RQ3).}\;
% ImageNet-R~\cite{ImagenetR_Hendrycks_2021_ICCV} (200 classes, $T=20$ tasks of 10 classes each with random class shuffle) serves as the main class-incremental benchmark, following prior LoRA-based CL work~\cite{liangInfLoRAInterferenceFreeLowRank2024a}.
% The 20-task structure yields 19 per-task BWT measurements per run, sufficient for reliable Spearman correlation in RQ2.
% To assess robustness under distribution shift, we additionally evaluate on the Tiny ImageNet variants of ImageNet-C and ImageNet-P~\cite{hendrycks2018benchmarking}.

% % \emph{Generalization (RQ4).}\;
% Five fine-grained benchmarks---CUB-200-2011, FGVC-Aircraft, Stanford Cars, Oxford 102 Flowers, and Oxford-IIIT Pet~\cite{Gao_2023_ICCV,wu2025sdlorascalabledecoupledlowrank}---are each split into $T=10$ equal tasks (Oxford-IIIT Pet: $T=9$).
% These provide cross-domain coverage well beyond the ImageNet rendition space.

% \noindent\textbf{Models and LoRA organizations.}\quad
% All vision experiments use ViT-B/16~\cite{dosovitskiy2020image} pretrained on ImageNet-21K, with LoRA adapters of rank $r=16$ attached to the query and value projection matrices of all attention blocks.
% We study three LoRA organizations that induce different admissible update geometries:
% \textbf{SeqLoRA} (shared adapter trained sequentially; $\mathcal W_{\Delta,t}^{\mathrm{train}}$ is the tangent space of the current adapter factors, geometrically fixed across tasks),
% \textbf{IncLoRA} (one new adapter branch per task; previous branches frozen and accumulated, so $\mathcal W_{\Delta,t}^{\mathrm{train}}$ is disjoint from previous supports in factor space), and
% \textbf{OLoRA}~\cite{wang2023orthogonal} (orthogonality-constrained updates that directly reduce cross-task curvature overlap $c_t$).
% RQ1--RQ3 use SeqLoRA, which provides the simplest and most tractable support structure for the causal probes.
% RQ4 extends to all three organizations.

% \noindent\textbf{Perturbation supports (RQ1).}\quad
% We hold the LoRA update rule fixed and vary only which subspace receives the perturbation during training.
% All five variants share the same architecture, learning rate, rank, and training epochs.
% Table~\ref{tab:training_variants} defines the five variants.

% \begin{table}[ht]
% \caption{Training variants for the RQ1 causal experiment (SeqLoRA, ImageNet-R, rank $r=16$, ViT-B/16, $\rho=0.05$).
% $(A_t, B_t)$ receive the gradient update in every variant.
% \emph{Perturbation support} specifies which subspace the temporary sharpness-seeking step targets before the gradient step.}
% \label{tab:training_variants}
% \centering
% \small
% \setlength{\tabcolsep}{3.5pt}
% \begin{tabular}{@{}lp{5.0cm}ll@{}}
% \toprule
% Variant & Perturbation support & Trainable update & Purpose \\
% \midrule
% SGD
%   & None
%   & $(A_t,\,B_t)$
%   & Baseline \\[2pt]
% Adapter-param SAM
%   & $(A_t,B_t)$ factor space \newline (norm in $\mathbb{R}^{r\times d_{\mathrm{in}}}\!\times\!\mathbb{R}^{d_{\mathrm{out}}\times r}$)
%   & $(A_t,\,B_t)$
%   & Existing optimizer; approximate $\mathcal W_{\Delta,t}^{\mathrm{train}}$ \\[2pt]
% Adapter-update SAM
%   & $\mathcal W_{\Delta,t}^{\mathrm{train}}$,\; $\|\epsilon\|_F \le \rho$ in $\mathbb{R}^{d_{\mathrm{out}}\times d_{\mathrm{in}}}$
%   & $(A_t,\,B_t)$
%   & Main causal test (P1, P3) \\[2pt]
% Frozen-coord perturb.
%   & Frozen backbone directions, \newline $\epsilon \perp \mathcal W_{\Delta,t}^{\mathrm{train}}$
%   & $(A_t,\,B_t)$
%   & Negative control (P3) \\[2pt]
% Random matched SAM
%   & Random subspace of $\mathbb{R}^{d_{\mathrm{out}}\times d_{\mathrm{in}}}$, \newline $\dim = \dim\mathcal W_{\Delta,t}^{\mathrm{train}}$
%   & $(A_t,\,B_t)$
%   & Dimension-matched control (P3) \\
% \bottomrule
% \end{tabular}
% \end{table}

% The key distinction between adapter-param SAM and adapter-update SAM is the norm constraint: adapter-param SAM searches $\max_{\|(\delta A,\delta B)\|_F \le \rho}$ in factor coordinates (so the effective $\Delta W$ perturbation depends on the scaling of $A_t, B_t$), while adapter-update SAM searches $\max_{\|\epsilon\|_F \le \rho,\;\epsilon \in \mathcal W_{\Delta,t}^{\mathrm{train}}}$ directly in weight space, removing this parameterization dependence.
% The frozen-coordinate and random-subspace perturbations are temporary: they are applied only during the perturbed-loss forward pass and do not make backbone parameters trainable.

% \noindent\textbf{Perturbation types (RQ3).}\quad
% After RQ1--RQ2 identify the correct perturbation support, RQ3 asks how different sharpness objectives behave \emph{within} $\mathcal W_{\Delta,t}^{\mathrm{train}}$.
% We compare three types, all restricted to the adapter-update support:
% $\mathrm{RS}^{(0),\Delta}$ (RWP~\cite{pmlr-v151-bisla22a}: isotropic Gaussian smoothing within $\mathcal W_{\Delta,t}^{\mathrm{train}}$),
% $\mathrm{AS}^{(0),\Delta}$ (SAM~\cite{foretSharpnessAwareMinimizationEfficiently2021}: adversarial ascent within $\mathcal W_{\Delta,t}^{\mathrm{train}}$), and
% $\mathrm{AS}^{(1),\Delta}$ (GAM~\cite{Zhang_GAM}: gradient-norm-aware smoothing within $\mathcal W_{\Delta,t}^{\mathrm{train}}$).
% For all SAM-style variants we fix $\rho = 0.05$ following the standard setting~\cite{foretSharpnessAwareMinimizationEfficiently2021}.

% \noindent\textbf{Cross-architecture checks (RQ4).}\quad
% To test whether the support principle extends beyond ViT-based vision models, we additionally evaluate T5-small and T5-large on a standard NLP CL benchmark (AG News, Amazon Reviews, DBpedia, Yahoo Answers, three task orders), using Adam as the common base optimizer.
% We further include Llama-3.2-1B and Llama-3.2-3B as decoder-only LLM checks under the same protocol.
% These serve as targeted structural checks, not a full second ablation suite.

% \noindent\textbf{Sharpness and curvature metrics (RQ2).}\quad
% Post-training, at each task-$t$ checkpoint, we evaluate (batch size 16, dataset fraction 0.1):
% \begin{itemize}[leftmargin=1.5em,itemsep=2pt]
%     \item $\mathrm{Sh}^{\mathrm{full}}_t$: $\mathrm{AS}^{(0)}$ in the full model parameter space (global sensitivity).
%     \item $\mathrm{Sh}^{A,B}_t$: $\mathrm{AS}^{(0)}$ in raw LoRA factor $(A_t,B_t)$-coordinates (parameterization-dependent; P2).
%     \item $\mathrm{Sh}^{\Delta}_t$: $\mathrm{AS}^{(0),\Delta}$ with $\epsilon = \rho\,g_{\Delta W}/\|g_{\Delta W}\|_F$, where $g_{\Delta W} = B_t g_A + g_B A_t$; bound-relevant, parameterization-invariant.
%     \item $\mathrm{Sh}^{\mathrm{rand}}_t$: $\mathrm{AS}^{(0)}$ along a random unit vector projected onto the $\dim(\mathcal W_{\Delta,t}^{\mathrm{train}})$-dimensional subspace (negative control).
% \end{itemize}
% For auxiliary curvature diagnostics we additionally estimate $\lambda_{\max}$ and $\mathrm{tr}(\cdot)$ of the Hessian and empirical Fisher information matrix via Lanczos~\cite{ghorbaniInvestigationNeuralNet2019}.

% \begin{table}[ht]
% \caption{Overview of the four empirical probes.
% Experiments are designed to test whether the adapter-update support $\mathcal W_{\Delta,t}^{\mathrm{train}}$ is the relevant locus of flatness in LoRA-based PECL.}
% \label{tab:exp_overview}
% \centering
% \small
% \setlength{\tabcolsep}{4pt}
% \begin{tabular}{@{}lp{3.5cm}p{3.5cm}p{3.5cm}@{}}
% \toprule
% RQ & Manipulated variable & Measured outcome & Purpose \\
% \midrule
% RQ1 & Perturbation support during training (5 variants) & LAA, AAA, BWT, Forgetting & Causal: does support alignment reduce forgetting? \\
% \addlinespace[2pt]
% RQ2 & Sharpness measurement support (4 objects) & Spearman $\rho$ with $\mathrm{BWT}_t$; rescaling invariance & Diagnostic: which sharpness object explains forgetting? \\
% \addlinespace[2pt]
% RQ3 & Perturbation type within $\mathcal W_{\Delta,t}^{\mathrm{train}}$ (SAM/RWP/GAM) & Restricted sharpness, curvature, forgetting & Within-support: do different smoothing objectives agree? \\
% \addlinespace[2pt]
% RQ4 & LoRA organization, dataset, model family & Directional consistency of support effect & Robustness: is the conclusion structural, not implementation-specific? \\
% \bottomrule
% \end{tabular}
% \end{table}

\subsection{RQ1: Does Sharpness-Aware Optimization Affect the Posterior Trajectory?}
\label{sec:rq1_trajectory}

Theorem~\ref{thm:sequential_pac_bayes} suggests that sharpness-aware optimization should be interpreted as acting on a posterior trajectory rather than only on an isolated checkpoint.
We therefore test the following empirical implication: after the same posterior state, changing the optimizer used for later tasks should change the subsequent trajectory and the retention of earlier tasks.
To isolate this effect, we use a forked-trajectory design on ImageNet-R.
We first train a prefix up to task 9 and then copy the identical checkpoint into multiple suffix branches.
Within each prefix group, all suffix variants start from the same task-9 checkpoint; hence differences after task 9 are attributable to the suffix optimizer rather than the prefix solution.
We use \texttt{random_factor} as a control for generic perturbation noise.

\begin{table}[t]
\centering
\small
\setlength{\tabcolsep}{4.2pt}
\caption{Forked trajectory experiment on ImageNet-R with SeqLoRA ($r=16$, $T=20$).
The columns ``t0--t10'' use the first ten-task prefix only, while ``t0--t20'' uses the full twenty-task trajectory.
Within each prefix group, the t0--t10 numbers are identical because the suffix branches start from the same task-9 checkpoint.
}
\label{tab:forked_prefix_full}
\begin{tabular}{lrrrrrr}
\toprule
Method & \multicolumn{3}{c}{t0--t10} & \multicolumn{3}{c}{t0--t20} \\
\cmidrule(lr){2-4} \cmidrule(lr){5-7}
 & AAA $\uparrow$ & FAA $\uparrow$ & BWT $\uparrow$ & AAA $\uparrow$ & FAA $\uparrow$ & BWT $\uparrow$ \\
\midrule
SGD $\rightarrow$ SGD & 72.75 & 69.91 & -7.09 & 67.60 & 61.72 & -15.32 \\
SGD $\rightarrow$ SAM-factor & 72.75 & 69.91 & -7.09 & 69.51 & 66.38 & -9.97 \\
SGD $\rightarrow$ Random-factor & 72.75 & 69.91 & -7.09 & 67.59 & 61.72 & -15.32 \\
\midrule
SAM-factor $\rightarrow$ SGD & 73.47 & 71.90 & -4.55 & 69.84 & 65.29 & -10.38 \\
SAM-factor $\rightarrow$ SAM-factor & 73.47 & 71.90 & -4.55 & 70.42 & 67.50 & -7.94 \\
\bottomrule
\end{tabular}
\end{table}

Table~\ref{tab:forked_prefix_full} shows two effects.
First, from the same SGD prefix, suffix SAM-factor improves full-trajectory FAA from $61.72$ to $66.38$, improves AAA from $67.60$ to $69.51$, and increases BWT from $-15.32$ to $-9.97$.
In contrast, suffix Random-factor is essentially identical to suffix SGD.
This rules out the explanation that the gain comes merely from injecting perturbation noise.
Second, with a SAM-factor prefix, continuing with SAM-factor further improves the full trajectory over continuing with SGD, increasing FAA from $65.29$ to $67.50$ and BWT from $-10.38$ to $-7.94$.
These results support the trajectory-level interpretation of the theory: sharpness-aware optimization changes the posterior states produced by later updates, and these altered states retain earlier tasks better.
The conclusion is not that SAM universally dominates every current-task trade-off, but that its effect is not confined to the task on which it is applied.

\subsection{RQ2: Which Perturbation Geometry Is Relevant in PECL?}
\label{sec:rq2_geometry}

The support-reduction result in Theorem~\ref{thm:pecl-bound} identifies the effective adapter-update geometry as the bound-relevant support in frozen-backbone LoRA PECL.
Empirically, this raises two separate questions: whether the perturbation direction must be adversarial, and whether the perturbation support is aligned with the trainable adapter-update geometry.
We therefore run a $2\times5$ decomposition on ImageNet-R using the same SeqLoRA setting as above.
The direction axis compares an adversarial SAM direction with a normalized Gaussian random direction.
The support axis compares raw factor coordinates, effective adapter-update tangent directions, merged effective weights, all raw parameters, and frozen backbone coordinates.

\begin{table}[t]
\centering
\small
\setlength{\tabcolsep}{4.0pt}
\caption{Direction--support decomposition on ImageNet-R with SeqLoRA ($r=16$, $T=20$).
For each perturbation support, we compare an adversarial SAM direction against a normalized Gaussian random direction.
$\Delta$ reports SAM minus random; higher FAA and BWT are better.
}
\label{tab:support_direction}
\begin{tabular}{lrrrrrr}
\toprule
Support & SAM FAA & Random FAA & $\Delta$ FAA & SAM BWT & Random BWT & $\Delta$ BWT \\
\midrule
factor & 67.94 & 58.79 & +9.14 & -6.81 & -18.47 & +11.67 \\
delta  & 65.41 & 58.26 & +7.15 & -11.56 & -19.08 & +7.52 \\
full   & 64.67 & 58.15 & +6.52 & -12.01 & -19.24 & +7.23 \\
all    & 63.19 & 58.32 & +4.87 & -12.28 & -19.01 & +6.72 \\
frozen & 62.91 & 58.24 & +4.67 & -12.37 & -19.10 & +6.73 \\
\midrule
SGD baseline & 58.29 & -- & -- & -18.99 & -- & -- \\
\bottomrule
\end{tabular}
\end{table}

Table~\ref{tab:support_direction} shows that the perturbation direction is decisive.
For every support, the adversarial SAM direction substantially outperforms the corresponding random direction.
The FAA gain ranges from $+4.67$ for frozen-coordinate perturbations to $+9.14$ for raw factor perturbations, and the BWT gain ranges from $+6.72$ to $+11.67$.
Moreover, all random variants remain close to the SGD baseline, showing that generic Gaussian perturbation is not sufficient to improve the trajectory.

The support comparison gives a more nuanced picture.
Under the default LoRA parameterization, \texttt{sam_factor} gives the strongest performance, while \texttt{sam_delta} remains clearly better than SGD and its random counterpart.
This should not be read as evidence that raw factor-space sharpness is the theoretically correct object.
Raw factor perturbations are parameterization-dependent: the same effective update $\Delta W$ can be represented by rescaled factors without changing the model function.
Thus, Table~\ref{tab:support_direction} supports the necessity of adversarial sharpness-aware perturbations, while the basis-invariant role of the effective adapter-update support should be further checked by a rescaling control and post-hoc $\mathrm{Sh}^{\Delta}$ diagnostics.

\subsection{RQ3: How Do Sharpness Objectives Behave within the Adapter-Update Support?}

\subsection{RQ4: Is the Support Principle Robust Across LoRA Organizations and Model Families?}

% \subsection{RQ1: verify }


% \subsection{RQ2: Does Perturbation Support Causally Affect Forgetting?}
% \label{sec:rq1}

% We train the five variants from Table~\ref{tab:training_variants} on ImageNet-R with SeqLoRA and report FAA, AAA, and $\overline{\mathrm{BWT}} = \frac{1}{T}\sum_t \mathrm{BWT}_t$ in Table~\ref{tab:rq1_results}.

% \begin{table}[ht]
% \caption{RQ1 results: effect of perturbation support on SeqLoRA continual learning.
% ImageNet-R, $T=20$, rank\,16, ViT-B/16, averaged over 3 seeds.
% $\uparrow$ denotes higher is better.}
% \label{tab:rq1_results}
% \centering
% \small
% \setlength{\tabcolsep}{5pt}
% \begin{tabular}{@{}lccc@{}}
% \toprule
% Variant & FAA (\%)$\uparrow$ & AAA (\%)$\uparrow$ & $\overline{\mathrm{BWT}}$ (\%)$\uparrow$ \\
% \midrule
% SGD                      & --- & --- & --- \\
% Adapter-param SAM        & --- & --- & --- \\
% Adapter-update SAM       & --- & --- & --- \\
% Frozen-coord perturb.    & --- & --- & --- \\
% Random matched SAM       & --- & --- & --- \\
% \bottomrule
% \end{tabular}
% \end{table}

% \paragraph{Expected pattern.}
% Predictions P1 and P3 jointly imply the following ordering in forgetting reduction:
% \begin{equation*}
% \text{Adapter-update SAM} \;\ge\; \text{Adapter-param SAM} \;\gg\; \text{Random matched SAM} \;\approx\; \text{Frozen-coord} \;\approx\; \text{SGD}.
% \end{equation*}
% Adapter-update SAM is expected to show the most consistent reduction in $\overline{\mathrm{BWT}}$ because its perturbation is exactly aligned with $\mathcal W_{\Delta,t}^{\mathrm{train}}$, the subspace the PAC-Bayesian bound identifies as the locus of forgetting risk.
% Adapter-param SAM is expected to perform similarly but slightly less stably, because it implicitly searches within $\mathcal W_{\Delta,t}^{\mathrm{train}}$ but with a factor-space norm that is sensitive to the scaling of $A_t$ and $B_t$.
% Frozen-coordinate perturbation is expected to yield no consistent improvement over SGD: although it seeks a flat loss landscape along the frozen backbone, those directions are outside the reachable adapter geometry and therefore do not affect the bound-relevant sharpness.
% Random matched SAM is expected to show an intermediate, unreliable benefit: matching the dimension of $\mathcal W_{\Delta,t}^{\mathrm{train}}$ is not sufficient; it is the specific geometry that matters.

% This contrast---adapter-update geometry reduces forgetting; frozen-coordinate and random do not---provides causal evidence for P3: it is not perturbation \emph{magnitude} but perturbation \emph{support} that determines forgetting.


% \subsection{RQ2: Does Adapter-Update Sharpness Best Explain Forgetting?}
% \label{sec:rq2}

% RQ2 takes the same five sets of trained checkpoints from RQ1 and asks a diagnostic question: among the four sharpness objects defined in Section~\ref{sec:exp_setup}, which is most predictive of per-task forgetting?
% This separates the causal question (RQ1: does the training perturbation support affect forgetting?) from the explanatory question (RQ2: does the measurement support explain it?).

% \paragraph{Spearman correlation analysis.}
% At each task-$t$ checkpoint, we measure all four sharpness objects and compute the Spearman rank correlation $\rho$ with $\mathrm{BWT}_t$, pooled over all $T-1$ transitions and 3 seeds.
% Table~\ref{tab:rq2_spearman} reports these correlations.

% \begin{table}[ht]
% \caption{RQ2: Spearman $\rho$ between sharpness objects and per-task forgetting $\mathrm{BWT}_t$.
% SeqLoRA on ImageNet-R, $T=20$, 3 seeds ($3\times19=57$ data points per cell).
% Bold: highest $|\rho|$ in each row.
% Negative $\rho$: higher sharpness predicts more forgetting.}
% \label{tab:rq2_spearman}
% \centering
% \small
% \setlength{\tabcolsep}{5pt}
% \begin{tabular}{@{}lcccc@{}}
% \toprule
% Training variant & $\mathrm{Sh}^{\mathrm{full}}$ & $\mathrm{Sh}^{A,B}$ & $\mathrm{Sh}^{\Delta}$ & $\mathrm{Sh}^{\mathrm{rand}}$ \\
% \midrule
% SGD                   & --- & --- & --- & --- \\
% Adapter-param SAM     & --- & --- & --- & --- \\
% Adapter-update SAM    & --- & --- & --- & --- \\
% Frozen-coord perturb. & --- & --- & --- & --- \\
% Random matched SAM    & --- & --- & --- & --- \\
% \bottomrule
% \end{tabular}
% \end{table}

% Prediction P1 expects $\mathrm{Sh}^{\Delta}_t$ to achieve the highest $|\rho|$ consistently across all five training variants, because it measures sensitivity in precisely the subspace through which the LoRA update causes forgetting.
% $\mathrm{Sh}^{\mathrm{full}}$ may show moderate correlation---it subsumes $\mathcal W_{\Delta,t}^{\mathrm{train}}$ but dilutes the signal with frozen-coordinate sensitivity.
% $\mathrm{Sh}^{\mathrm{rand}}$ is expected to be weak and inconsistent: the random subspace does not respect the task-permissible geometry.
% $\mathrm{Sh}^{A,B}$ may track $\mathrm{Sh}^{\Delta}$ for adapter-param SAM (where the two are closely linked) but should be less stable across variants because its value depends on the scaling of $A_t$ and $B_t$.

% % \paragraph{Parameterization invariance control (P2).}
% % To confirm that $\mathrm{Sh}^{A,B}$ is not an intrinsic model property, we take a SeqLoRA/SGD checkpoint after task $t$ and apply the rescaling $(A_t, B_t)\mapsto(c A_t,\, B_t/c)$ for $c\in\{0.1,\,0.5,\,1,\,2,\,10\}$.
% % Since $\Delta W = B_t A_t$ is preserved by this rescaling, the model function, empirical loss, and $\mathrm{Sh}^{\Delta}_t$ are invariant by construction.
% % Prediction P2 expects $\mathrm{Sh}^{A,B}$ to vary monotonically with $c$ while $\mathrm{Sh}^{\Delta}$ remains flat.
% % Any sharpness object that changes under such a rescaling cannot serve as the PAC-Bayesian flatness term of Theorem~\ref{thm:pecl_support_reduction}.


% \subsection{RQ3: How Does Perturbation Type Behave Within the Correct Support?}
% \label{sec:rq3}

% RQ1 and RQ2 identify adapter-update SAM as causally effective and $\mathrm{Sh}^{\Delta}$ as the most explanatory diagnostic.
% RQ3 asks a subordinate question: given that the perturbation support is fixed to $\mathcal W_{\Delta,t}^{\mathrm{train}}$, do different sharpness objectives---zeroth-order adversarial (SAM), random smoothing (RWP), and first-order gradient-norm-aware (GAM)---all yield consistent forgetting reduction, or does the \emph{type} of smoothing within the correct support also matter?

% \paragraph{Setup.}
% We train three additional SeqLoRA variants on ImageNet-R, all using the adapter-update support but differing in the perturbation type:
% $\mathrm{AS}^{(0),\Delta}$ (SAM restricted to $\mathcal W_{\Delta,t}^{\mathrm{train}}$, already run in RQ1),
% $\mathrm{RS}^{(0),\Delta}$ (RWP: isotropic Gaussian smoothing within $\mathcal W_{\Delta,t}^{\mathrm{train}}$, $\sigma = 0.05$), and
% $\mathrm{AS}^{(1),\Delta}$ (GAM: gradient-norm-aware smoothing within $\mathcal W_{\Delta,t}^{\mathrm{train}}$).
% We compare them against SGD using the same metrics as RQ1.
% Additionally, for each variant we measure $\mathrm{Sh}^{\Delta}_t$ post-training to confirm that all three actually reduce adapter-update sharpness relative to SGD.

% \begin{table}[ht]
% \caption{RQ3: Effect of perturbation type within the adapter-update support $\mathcal W_{\Delta,t}^{\mathrm{train}}$.
% SeqLoRA on ImageNet-R, rank\,16, ViT-B/16, 3 seeds.
% All variants use the adapter-update support; only the smoothing objective differs.}
% \label{tab:rq3_perttype}
% \centering
% \small
% \setlength{\tabcolsep}{5pt}
% \begin{tabular}{@{}lcccc@{}}
% \toprule
% Variant & FAA (\%)$\uparrow$ & AAA (\%)$\uparrow$ & $\overline{\mathrm{BWT}}$ (\%)$\uparrow$ & $\mathrm{Sh}^{\Delta}$ \\
% \midrule
% SGD (baseline)          & --- & --- & --- & --- \\
% $\mathrm{AS}^{(0),\Delta}$ (SAM) & --- & --- & --- & --- \\
% $\mathrm{RS}^{(0),\Delta}$ (RWP) & --- & --- & --- & --- \\
% $\mathrm{AS}^{(1),\Delta}$ (GAM) & --- & --- & --- & --- \\
% \bottomrule
% \end{tabular}
% \end{table}

% \paragraph{Expected pattern.}
% All three perturbation types are expected to reduce forgetting relative to SGD, because they all target the theoretically relevant support $\mathcal W_{\Delta,t}^{\mathrm{train}}$.
% The relative ranking (SAM vs.\ RWP vs.\ GAM) is a secondary question; any consistent ordering would indicate that, within the correct support, zeroth-order adversarial smoothing is empirically strongest (as suggested by the broader SAM literature).
% The key prediction is that $\mathrm{Sh}^{\Delta}_t$ is reduced for all three---confirming that the sharpness reduction is the common mechanism---and that the reduction in $\mathrm{Sh}^{\Delta}_t$ correlates with the reduction in forgetting across the three types.
% This provides within-support evidence that it is genuinely sharpness in $\mathcal W_{\Delta,t}^{\mathrm{train}}$ that drives the forgetting reduction, not any other implicit effect of the training procedure.


% \subsection{RQ4: Is the Support Principle Robust Across LoRA Organizations and Model Families?}
% \label{sec:rq4}

% RQ1--RQ3 establish the support principle for SeqLoRA on ImageNet-R.
% RQ4 tests whether the conclusion is structural---following from the frozen-backbone PECL setup---or an artifact of a specific LoRA organization or vision model.

% \paragraph{Across LoRA organizations.}
% For each of SeqLoRA, IncLoRA, and OLoRA, we compare SGD against adapter-update SAM (where the support $\mathcal W_{\Delta,t}^{\mathrm{train}}$ is defined by each method's own current adapter).
% We apply the same perturbation radius $\rho$ across all organizations and report the BWT improvement $\Delta\mathrm{BWT} = \mathrm{BWT}_{\mathrm{Adapter\text{-}update\,SAM}} - \mathrm{BWT}_{\mathrm{SGD}}$ in Table~\ref{tab:rq4_organizations}.

% \begin{table}[ht]
% \caption{RQ4: BWT improvement from adapter-update SAM over SGD, across LoRA organizations, datasets, and model families.
% $\Delta\mathrm{BWT} > 0$ indicates less forgetting.
% Rank\,16, ViT-B/16 (vision) or Adam base optimizer (NLP), averaged over 3 seeds.}
% \label{tab:rq4_organizations}
% \centering
% \small
% \setlength{\tabcolsep}{4pt}
% \begin{tabular}{@{}llccc@{}}
% \toprule
% Model & LoRA org. & $\Delta\mathrm{BWT}$ on ImageNet-R & $\Delta\mathrm{BWT}$ on Fine-grained (avg.) & $\Delta\mathrm{BWT}$ on NLP bench. \\
% \midrule
% \multirow{3}{*}{ViT-B/16}
%   & SeqLoRA & --- & --- & --- \\
%   & IncLoRA & --- & --- & --- \\
%   & OLoRA   & --- & --- & --- \\
% \addlinespace[2pt]
% \multirow{2}{*}{T5}
%   & small & --- & --- & --- \\
%   & large & --- & --- & --- \\
% \addlinespace[2pt]
% \multirow{2}{*}{Llama-3.2}
%   & 1B & --- & --- & --- \\
%   & 3B & --- & --- & --- \\
% \bottomrule
% \end{tabular}
% \end{table}

% \paragraph{Expected pattern.}
% All three LoRA organizations should show $\Delta\mathrm{BWT} > 0$, validating Prediction P4: the support principle is not tied to a specific LoRA structure.
% The magnitude of improvement may differ systematically.
% IncLoRA's support $\mathcal W_{\Delta,t}^{\mathrm{train}}$ is already disjoint from previous branches in factor space, so the dominant forgetting mechanism shifts from within-support sharpness to cross-task curvature alignment $c_t$; consequently, adapter-update SAM may provide smaller absolute BWT gains for IncLoRA.
% OLoRA imposes orthogonality to reduce $c_t$ structurally, so its SGD baseline already exhibits lower forgetting, leaving less room for improvement.
% These differences in magnitude---not in sign---are informative: they confirm that each organization's specific support geometry governs the forgetting dynamics, consistent with the general bound in Theorem~\ref{thm:pecl_support_reduction}.

% For T5 and Llama, the prediction is directional: adapter-update SAM should consistently improve over SGD or Adam baseline, even when the base optimizer changes.
% This would indicate that the support principle is structural (following from frozen-backbone PECL) rather than implementation-specific.


% % \subsection{Summary}
% % \label{sec:exp_summary}

% % \begin{table}[ht]
% % \caption{Summary: how each RQ maps to a theory-derived prediction, the experimental design, and the expected outcome.}
% % \label{tab:prediction_summary}
% % \centering
% % \small
% % \setlength{\tabcolsep}{3pt}
% % \begin{tabular}{@{}lp{3.2cm}p{3.0cm}p{3.5cm}@{}}
% % \toprule
% % RQ & Prediction tested & Experiment & Expected outcome \\
% % \midrule
% % RQ1 & P1, P3: perturbation support causally affects forgetting & 5 support variants $\times$ SeqLoRA, ImageNet-R & Adapter-update SAM reduces BWT; frozen-coord and random do not \\
% % \addlinespace[3pt]
% % RQ2 & P1: $\mathrm{Sh}^{\Delta}$ best explains forgetting;\newline P2: $\mathrm{Sh}^{A,B}$ is not invariant & Spearman $\rho$ on 5 checkpoint sets; rescaling sweep $c\in\{0.1,\ldots,10\}$ & $\mathrm{Sh}^{\Delta}$ highest $|\rho|$; $\mathrm{Sh}^{A,B}$ varies with $c$, $\mathrm{Sh}^{\Delta}$ flat \\
% % \addlinespace[3pt]
% % RQ3 & P1 within correct support: type matters less than support & SAM / RWP / GAM, all in $\mathcal W_{\Delta,t}^{\mathrm{train}}$ & All three reduce BWT and $\mathrm{Sh}^{\Delta}$; $\mathrm{Sh}^{\Delta}$ reduction correlates with BWT reduction \\
% % \addlinespace[3pt]
% % RQ4 & P4: support principle is structural, not implementation-specific & SeqLoRA/IncLoRA/OLoRA $\times$ 2 vision datasets $\times$ T5 $\times$ Llama & Consistent $\Delta\mathrm{BWT}>0$ across organizations and model families \\
% % \bottomrule
% % \end{tabular}
% % \end{table}

% % Together, the four RQs form a layered argument.
% % RQ1 establishes that perturbation support \emph{causally} governs forgetting.
% % RQ2 confirms that $\mathrm{Sh}^{\Delta}$---the bound-relevant sharpness object---is also the most \emph{explanatory} diagnostic, and that $\mathrm{Sh}^{A,B}$ is not intrinsic.
% % RQ3 shows that within the correct support, different smoothing objectives all work, suggesting the support alignment is the primary driver.
% % RQ4 checks that the conclusion is not an artifact of a particular LoRA organization or model family.
% % The convergence of causal, diagnostic, and generalization evidence supports the paper's central claim: in LoRA-based PECL, it is not the \emph{magnitude} of flatness but \emph{where} it is measured---specifically, within the adapter-update geometry $\mathcal W_{\Delta,t}^{\mathrm{train}}$---that governs forgetting.




% % Our contributions are threefold.
% % \begin{itemize}
% %     \item We derive a sequential hierarchical PAC-Bayesian decomposition for continual learning. 
% %     The resulting pathwise bound decomposes the posterior trajectory into perturbation-smoothed empirical risk, within-task adaptation complexity, and hyperposterior drift complexity. 
% %     This provides a PAC-Bayesian explanation of forgetting by interpreting it as generalization degradation on past tasks under later posteriors, and makes cross-task interference explicit through the hyperposterior drift term.

% %     \item We prove a PECL support-reduction theorem for frozen-backbone LoRA adaptation. 
% %     By analyzing the pushforward posterior induced by LoRA parameters on effective adapter updates, we show that the bound-relevant perturbation-smoothed empirical risk and KL complexity are determined by the admissible effective adapter-update geometry. 
% %     This identifies basis-invariant effective adapter-update sharpness, rather than full-parameter-space or raw LoRA factor-space sharpness, as the relevant flatness object under LoRA-constrained updates.

% %     \item We validate the theory-derived predictions across perturbation scopes, LoRA variants, datasets, and architectures. 
% %     Our experiments examine whether adapter-scoped perturbations better match continual-learning behavior than full-space or raw factor-space perturbations, and whether the resulting sharpness diagnostics align with generalization and forgetting trends.
% % \end{itemize}

% % \section{Submission of papers to NeurIPS 2026}


% % Please read the instructions below carefully and follow them faithfully.
% % \begin{figure}[htbp]
% % \centering
% % \begin{minipage}{0.9\linewidth}
% % \centering
% % \includegraphics[width=0.8\linewidth]{figures_TRML/2D_lossLandscape_task0_lora_opt_sam_vs_sgd_3d_compare_withxyz.pdf}
% % \caption{Loss landscape visualization for a ViT model adapted with LoRA, comparing SGD and adapter-only SAM, where sharpness-aware perturbations are applied only to the trainable LoRA parameter. Although perturbation is restricted to the LoRA subspace, full-space probes around the final solution show that adapter-only SAM leads to a flatter local landscape than SGD in this example.}
% % \label{fig:3d_loss_landscape}
% % \end{minipage}
% % \end{figure}

% % {%
% % % {%
% % %
% % % \setlength{\intextsep}{4pt}
% % \setlength{\textfloatsep}{4pt}
% % % \setlength{\abovecaptionskip}{6pt}
% % % \setlength{\belowcaptionskip}{pt}
% % % \captionsetup{skip=10pt}
% % \begin{figure}[ht]
% % \centering
% % \includegraphics[height=4.8cm,keepaspectratio]{figures_TRML/ACC_dataset_pair_Robustness_to_Corruptions_on_ImageNet-C_Robustness_to_Perturbation_on_ImageNet-P2.pdf}
% % % \includegraphics[width=0.9\linewidth]{figures_TRML/hebing_robutness2.pdf}
% % % \vspace{-6pt}
% % \caption{\textbf{Robustness of perturbation types under CL} Classification average accuracy $\operatorname{Acc}_t^{\mathrm{cls}}$ over tasks evaluated under distribution shift on ImageNet-C  and ImageNet-P.
% % Across tasks and LoRA organizations, both $\mathrm{AS}^{(0)}$ and $\mathrm{AS}^{(1)}$ consistently outperform SGD, and the gap widens in later tasks. Here ImageNet-C/P denote the Tiny ImageNet variants; see Experimental Setup.}
% % \label{fig:robustness_imagenet_cp}
% % % \vspace{-10pt}
% % \end{figure}
% % }%

% % \begin{figure}[htbp]
% % \centering
% % \setlength{\textfloatsep}{4pt}
% % \begin{subfigure}[t]{0.3\linewidth}
% %     \centering
% %     \includegraphics[height=5.0cm,keepaspectratio]{figures_TRML/2D_lossLandscape_task0_lora_opt_sam_vs_sgd_3d_compare_withxyz.pdf}
% %     \caption{Loss landscape visualization for a ViT model adapted with LoRA, comparing SGD and adapter-only SAM, where sharpness-aware perturbations are applied only to the trainable LoRA parameter. Although perturbation is restricted to the LoRA subspace, full-space probes around the final solution show that adapter-only SAM leads to a flatter local landscape than SGD in this example.}
% %     \label{fig:3d_loss_landscape}
% % \end{subfigure}
% % \hfill
% % \begin{subfigure}[t]{0.48\linewidth}
% %     \centering
% %     \includegraphics[height=5.0cm,keepaspectratio]{figures_TRML/ACC_dataset_pair_Robustness_to_Corruptions_on_ImageNet-C_Robustness_to_Perturbation_on_ImageNet-P2.pdf}
% %     \caption{\textbf{Robustness of perturbation types under CL} Classification average accuracy $\operatorname{Acc}_t^{\mathrm{cls}}$ over tasks evaluated under distribution shift on ImageNet-C and ImageNet-P. Across tasks and LoRA organizations, both $\mathrm{AS}^{(0)}$ and $\mathrm{AS}^{(1)}$ consistently outperform SGD, and the gap widens in later tasks. Here ImageNet-C/P denote the Tiny ImageNet variants; see Experimental Setup.}
% %     \label{fig:robustness_imagenet_cp}
% % \end{subfigure}
% % \caption{Comparison of local loss geometry and robustness under perturbations.}
% % \label{fig:loss_landscape_and_robustness}
% % \end{figure}

% % \begin{figure}[t]
% % \centering
% % \setlength{\textfloatsep}{4pt}

% % \begin{subfigure}[t]{0.36\linewidth}
% %     \centering
% %     \includegraphics[width=\linewidth,height=4.2cm,keepaspectratio,trim=0.8cm 0.3cm 0.8cm 0.3cm,clip]{figures_TRML/2D_lossLandscape_task0_lora_opt_sam_vs_sgd_3d_compare_withxyz.pdf}
% %     \caption{Local loss landscape of LoRA-adapted ViT under SGD and adapter-only SAM. Although perturbations are applied only in the LoRA subspace, the final solution found by adapter-only SAM exhibits a flatter neighborhood under full-space probing.}
% %     \label{fig:3d_loss_landscape}
% % \end{subfigure}
% % \hfill
% % \begin{subfigure}[t]{0.60\linewidth}
% %     \centering
% %     \includegraphics[width=\linewidth,height=4.2cm,keepaspectratio]{figures_TRML/ACC_dataset_pair_Robustness_to_Corruptions_on_ImageNet-C_Robustness_to_Perturbation_on_ImageNet-P2.pdf}
% %     \caption{Robustness under distribution shift on ImageNet-C and ImageNet-P. Across tasks and LoRA organizations, both $\mathrm{AS}^{(0)}$ and $\mathrm{AS}^{(1)}$ consistently outperform SGD, with larger margins in later tasks.}
% %     \label{fig:robustness_imagenet_cp}
% % \end{subfigure}

% % \vspace{2pt}
% % \caption{Comparison of local loss geometry and robustness under perturbations.}
% % \label{fig:loss_landscape_and_robustness}
% % \vspace{-6pt}
% % \end{figure}

% % Parameter-efficient continual learning (PECL) has emerged as a practical approach for adapting large pretrained models to sequential tasks with limited computational and storage overhead~\cite{pmlr-v97-houlsby19a,Wang_2022_CVPR,dingParameterefficientFinetuningLargescale2023,xu2023parameterefficientfinetuningmethodspretrained}.% ~\cite{pmlr-v97-houlsby19a,Wang_2022_CVPR,dingParameterefficientFinetuningLargescale2023,xu2023parameterefficientfinetuningmethodspretrained}
% % Among PECL methods, Low-Rank Adaptation (LoRA) is particularly attractive because it enables continual adaptation through lightweight low-rank updates while keeping the pretrained backbone fixed~\cite{hu2022lora}. This design naturally turns continual adaptation into a constrained geometric problem: future tasks do not modify the model freely in the full parameter space, but only through task-dependent low-rank update directions induced by the trainable LoRA parameters. Recent LoRA-based continual learning methods have explicitly exploited this structure through adapter modularity, orthogonal constraints, and task-specific subspace allocation to alleviate forgetting and interference~\cite{Gao_2023_ICCV,NEURIPS2020_70d85f35,wang2023orthogonal,wu2025sdlorascalabledecoupledlowrank}. However, these advances are primarily algorithmic. They do not yet provide a principled understanding of how such admissible update geometry shapes continual generalization, forgetting, and the notion of solution stability that should matter in this regime.


% % Another line of work studies continual learning through the geometry of the loss landscape. Flat minima have long been associated with improved generalization and robustness~\cite{FlatMinima1997,wu2017understandinggeneralizationdeeplearning,chaudhariEntropySGDBiasingGradient2019,NEURIPS2020_518a38cc,NEURIPS2021_995f5e03,JMLRAnEmpiricalInvestigatio,yangRevisitingFlatnessAwareOptimization2025}, and sharpness-aware optimization methods such as SAM and its variants operationalize this principle by favoring solutions that remain stable under local perturbations~\cite{foretSharpnessAwareMinimizationEfficiently2021,pmlr-v151-bisla22a,Zhang_GAM,NEURIPS2024_C-Flat}. Prior work has further suggested that flatter solutions can mitigate catastrophic forgetting in continual learning~\cite{NEURIPS2020_518a38cc,JMLRAnEmpiricalInvestigatio,shiOvercomingCatastrophicForgetting2021,liRevisitingCatastrophicForgetting2024}. However, most existing flatness-based arguments are inherited from full-parameter training, where perturbations and curvature are defined over the whole model~\cite{li2024flatlora}. In LoRA-based continual learning, by contrast, adaptation is confined to task-conditioned admissible low-rank update directions. This means that flatness measured over the whole model is not automatically the flatness notion that governs continual adaptation, because flatness is defined relative to the task-conditioned loss, while under frozen-backbone LoRA adaptation, only the admissible low-rank update directions are available to fit the current task data. It therefore remains unclear whether continual generalization should be understood through full-space flatness, or through task-conditioned local flatness along admissible LoRA directions.

% % We sharpen this ambiguity through a motivating empirical observation. As visualized in Figure~\ref{fig:3d_loss_landscape}, applying sharpness-aware minimization only within the low-rank adapter subspace already yields a solution whose surrounding loss landscape appears flatter than that obtained by SGD when probed in the full parameter space. This suggests that geometric control within the trainable LoRA subspace can propagate to measurable differences in full-model sensitivity. 


% % At the same time, such an observation does not resolve which notion of flatness is theoretically relevant: it leaves open whether full-space flatness is itself the governing object for continual generalization, or whether it is only a byproduct of controlling the admissible LoRA subspace geometry.

% % To resolve this question, we develop a sequential hierarchical PAC-Bayesian analysis for LoRA-based continual learning. We derive a generalization bound with an explicit hyperposterior drift term that captures the sequential evolution of task priors across tasks. Under a frozen backbone, we show that both the KL complexity and the sharpness-dependent empirical term reduce to quantities supported on the admissible LoRA update subspace. This identifies task-conditioned local flatness along admissible LoRA directions, rather than arbitrary full-space flatness, as the relevant notion of flatness in this regime. The analysis further yields concrete empirical predictions about perturbation scope and sharpness type, which we test experimentally: sharpness-aware optimization restricted to LoRA updates consistently improves continual performance, whereas perturbations that extend into frozen backbone parameters provide no systematic benefit and can harm stability.

% % \subsection{Style}


% % Papers to be submitted to NeurIPS 2026 must be prepared according to the
% % instructions presented here. Papers may only be up to {\bf nine} pages long, including figures. \textbf{Papers that exceed the page limit will not be reviewed (or in any other way considered) for presentation at the conference.}
% % Additional pages \emph{containing acknowledgments, references, checklist, and optional technical appendices} do not count as content pages. 


% % The margins in 2026 are the same as those in previous years.


% % Authors are required to use the NeurIPS \LaTeX{} style files obtainable at the NeurIPS website as indicated below. Please make sure you use the current files and not previous versions. Tweaking the style files may be grounds for desk rejection.


% % \subsection{Retrieval of style files}


% % The style files for NeurIPS and other conference information are available on the website at
% % \begin{center}
% %   \url{https://neurips.cc}.
% % \end{center}
% % % The file \verb+neurips_2026.pdf+ contains these instructions and illustrates the various formatting requirements your NeurIPS paper must satisfy.


% % The only supported style file for NeurIPS 2026 is \verb+neurips_2026.sty+, rewritten for \LaTeXe{}. \textbf{Previous style files for \LaTeX{} 2.09,  Microsoft Word, and RTF are no longer supported.}


% % The \LaTeX{} style file contains three optional arguments: 
% % \begin{itemize}
% % \item \verb+final+, which creates a camera-ready copy, 
% % \item \verb+preprint+, which creates a preprint for submission to, e.g., arXiv, \item \verb+nonatbib+, which will not load the \verb+natbib+ package for you in case of package clash.
% % \end{itemize}


% % % \paragraph{Preprint option}
% % % If you wish to post a preprint of your work online, e.g., on arXiv, using the NeurIPS style, please use the \verb+preprint+ option. This will create a nonanonymized version of your work with the text ``Preprint. Work in progress.'' in the footer. This version may be distributed as you see fit, as long as you do not say which conference it was submitted to. Please \textbf{do not} use the \verb+final+ option, which should \textbf{only} be used for papers accepted to NeurIPS.


% % % At submission time, please omit the \verb+final+ and \verb+preprint+ options. This will anonymize your submission and add line numbers to aid review. Please do \emph{not} refer to these line numbers in your paper as they will be removed during generation of camera-ready copies.


% % % The file \verb+neurips_2026.tex+ may be used as a ``shell'' for writing your paper. All you have to do is replace the author, title, abstract, and text of the paper with your own.


% % % The formatting instructions contained in these style files are summarized in Sections \ref{gen_inst}, \ref{headings}, and \ref{others} below.


% % % \section{General formatting instructions}
% % % \label{gen_inst}


% % % The text must be confined within a rectangle 5.5~inches (33~picas) wide and
% % % 9~inches (54~picas) long. The left margin is 1.5~inch (9~picas).  Use 10~point
% % % type with a vertical spacing (leading) of 11~points.  Times New Roman is the
% % % preferred typeface throughout, and will be selected for you by default.
% % % Paragraphs are separated by \nicefrac{1}{2}~line space (5.5 points), with no
% % % indentation.


% % % The paper title should be 17~point, initial caps/lower case, bold, centered
% % % between two horizontal rules. The top rule should be 4~points thick and the
% % % bottom rule should be 1~point thick. Allow \nicefrac{1}{4}~inch space above and
% % % below the title to rules. All pages should start at 1~inch (6~picas) from the
% % % top of the page.


% % % For the final version, authors' names are set in boldface, and each name is
% % % centered above the corresponding address. The lead author's name is to be listed
% % % first (left-most), and the co-authors' names (if different address) are set to
% % % follow. If there is only one co-author, list both author and co-author side by
% % % side.


% % % Please pay special attention to the instructions in Section \ref{others}
% % % regarding figures, tables, acknowledgments, and references.

% % % \section{Headings: first level}
% % % \label{headings}


% % % All headings should be lower case (except for first word and proper nouns),
% % % flush left, and bold.


% % % First-level headings should be in 12-point type.


% % % \subsection{Headings: second level}


% % % Second-level headings should be in 10-point type.


% % % \subsubsection{Headings: third level}


% % % Third-level headings should be in 10-point type.


% % % \paragraph{Paragraphs}


% % % There is also a \verb+\paragraph+ command available, which sets the heading in
% % % bold, flush left, and inline with the text, with the heading followed by 1\,em
% % % of space.


% % % \section{Citations, figures, tables, references}
% % % \label{others}


% % % These instructions apply to everyone.


% % % \subsection{Citations within the text}


% % % The \verb+natbib+ package will be loaded for you by default.  Citations may be
% % % author/year or numeric, as long as you maintain internal consistency.  As to the
% % % format of the references themselves, any style is acceptable as long as it is
% % % used consistently.


% % % The documentation for \verb+natbib+ may be found at
% % % \begin{center}
% % %   \url{http://mirrors.ctan.org/macros/latex/contrib/natbib/natnotes.pdf}
% % % \end{center}
% % % Of note is the command \verb+\citet+, which produces citations appropriate for
% % % use in inline text.  For example,
% % % \begin{verbatim}
% % %    \citet{hasselmo} investigated\dots
% % % \end{verbatim}
% % % produces
% % % \begin{quote}
% % %   Hasselmo, et al.\ (1995) investigated\dots
% % % \end{quote}


% % % If you wish to load the \verb+natbib+ package with options, you may add the
% % % following before loading the \verb+neurips_2026+ package:
% % % \begin{verbatim}
% % %    \PassOptionsToPackage{options}{natbib}
% % % \end{verbatim}


% % % If \verb+natbib+ clashes with another package you load, you can add the optional
% % % argument \verb+nonatbib+ when loading the style file:
% % % \begin{verbatim}
% % %    \usepackage[nonatbib]{neurips_2026}
% % % \end{verbatim}


% % % As submission is double blind, refer to your own published work in the third
% % % person. That is, use ``In the previous work of Jones et al.\ [4],'' not ``In our
% % % previous work [4].'' If you cite your other papers that are not widely available
% % % (e.g., a journal paper under review), use anonymous author names in the
% % % citation, e.g., an author of the form ``A.\ Anonymous'' and include a copy of the anonymized paper in the supplementary material.


% % % \subsection{Footnotes}


% % % Footnotes should be used sparingly.  If you do require a footnote, indicate
% % % footnotes with a number\footnote{Sample of the first footnote.} in the
% % % text. Place the footnotes at the bottom of the page on which they appear.
% % % Precede the footnote with a horizontal rule of 2~inches (12~picas).


% % % Note that footnotes are properly typeset \emph{after} punctuation
% % % marks.\footnote{As in this example.}


% % % \subsection{Figures}


% % % \begin{figure}
% % %   \centering
% % %   \fbox{\rule[-.5cm]{0cm}{4cm} \rule[-.5cm]{4cm}{0cm}}
% % %   \caption{Sample figure caption. Explain what the figure shows and add a key take-away message to the caption.}
% % % \end{figure}


% % % All artwork must be neat, clean, and legible. Lines should be dark enough for
% % %  reproduction purposes. The figure number and caption always appear after the
% % % figure. Place one line space before the figure caption and one line space after
% % % the figure. The figure caption should be lower case (except for the first word and proper nouns); figures are numbered consecutively.


% % % You may use color figures.  However, it is best for the figure captions and the
% % % paper body to be legible if the paper is printed in either black/white or in
% % % color.


% % % \subsection{Tables}


% % % All tables must be centered, neat, clean, and legible.  The table number and
% % % title always appear before the table.  See Table~\ref{sample-table}.


% % % Place one line space before the table title, one line space after the
% % % table title, and one line space after the table. The table title must
% % % be lower case (except for the first word and proper nouns); tables are
% % % numbered consecutively.


% % % Note that publication-quality tables \emph{do not contain vertical rules}. We
% % % strongly suggest the use of the \verb+booktabs+ package, which allows for
% % % typesetting high-quality, professional tables:
% % % \begin{center}
% % %   \url{https://www.ctan.org/pkg/booktabs}
% % % \end{center}
% % % This package was used to typeset Table~\ref{sample-table}.


% % % \begin{table}
% % %   \caption{Sample table caption. Explain what the table shows and add a key take-away message to the caption.}
% % %   \label{sample-table}
% % %   \centering
% % %   \begin{tabular}{lll}
% % %     \toprule
% % %     \multicolumn{2}{c}{Part}                   \\
% % %     \cmidrule(r){1-2}
% % %     Name     & Description     & Size ($\mu$m) \\
% % %     \midrule
% % %     Dendrite & Input terminal  & $\approx$100     \\
% % %     Axon     & Output terminal & $\approx$10      \\
% % %     Soma     & Cell body       & up to $10^6$  \\
% % %     \bottomrule
% % %   \end{tabular}
% % % \end{table}

% % % \subsection{Math}
% % % Note that display math in bare TeX commands will not create correct line numbers for submission. Please use LaTeX (or AMSTeX) commands for unnumbered display math. (You really shouldn't be using \$\$ anyway; see \url{https://tex.stackexchange.com/questions/503/why-is-preferable-to} and \url{https://tex.stackexchange.com/questions/40492/what-are-the-differences-between-align-equation-and-displaymath} for more information.)

% % % \subsection{Final instructions}

% % % Do not change any aspects of the formatting parameters in the style files.  In
% % % particular, do not modify the width or length of the rectangle the text should
% % % fit into, and do not change font sizes. Please note that pages should be
% % % numbered.


% % % \section{Preparing PDF files}


% % Please prepare submission files with paper size ``US Letter,'' and not, for
% % example, ``A4.''


% % Fonts were the main cause of problems in the past years. Your PDF file must only
% % contain Type 1 or Embedded TrueType fonts. Here are a few instructions to
% % achieve this.


% % \begin{itemize}


% % \item You should directly generate PDF files using \verb+pdflatex+.


% % \item You can check which fonts a PDF files uses.  In Acrobat Reader, select the
% %   menu Files$>$Document Properties$>$Fonts and select Show All Fonts. You can
% %   also use the program \verb+pdffonts+ which comes with \verb+xpdf+ and is
% %   available out-of-the-box on most Linux machines.


% % \item \verb+xfig+ ``patterned'' shapes are implemented with bitmap fonts.  Use
% %   "solid" shapes instead.


% % \item The \verb+\bbold+ package almost always uses bitmap fonts.  You should use
% %   the equivalent AMS Fonts:
% % \begin{verbatim}
% %    \usepackage{amsfonts}
% % \end{verbatim}
% % followed by, e.g., \verb+\mathbb{R}+, \verb+\mathbb{N}+, or \verb+\mathbb{C}+
% % for $\mathbb{R}$, $\mathbb{N}$ or $\mathbb{C}$.  You can also use the following
% % workaround for reals, natural and complex:
% % \begin{verbatim}
% %    \newcommand{\RR}{I\!\!R} %real numbers
% %    \newcommand{\Nat}{I\!\!N} %natural numbers
% %    \newcommand{\CC}{I\!\!\!\!C} %complex numbers
% % \end{verbatim}
% % Note that \verb+amsfonts+ is automatically loaded by the \verb+amssymb+ package.


% % \end{itemize}


% % If your file contains type 3 fonts or non embedded TrueType fonts, we will ask
% % you to fix it.


% % \subsection{Margins in \LaTeX{}}


% % Most of the margin problems come from figures positioned by hand using
% % \verb+\special+ or other commands. We suggest using the command
% % \verb+\includegraphics+ from the \verb+graphicx+ package. Always specify the
% % figure width as a multiple of the line width as in the example below:
% % \begin{verbatim}
% %    \usepackage[pdftex]{graphicx} ...
% %    \includegraphics[width=0.8\linewidth]{myfile.pdf}
% % \end{verbatim}
% % See Section 4.4 in the graphics bundle documentation
% % (\url{http://mirrors.ctan.org/macros/latex/required/graphics/grfguide.pdf})


% % A number of width problems arise when \LaTeX{} cannot properly hyphenate a
% % line. Please give LaTeX hyphenation hints using the \verb+\-+ command when
% % necessary.

% % \begin{ack}
% % Use unnumbered first level headings for the acknowledgments. All acknowledgments
% % go at the end of the paper before the list of references. Moreover, you are required to declare
% % funding (financial activities supporting the submitted work) and competing interests (related financial activities outside the submitted work).
% % More information about this disclosure can be found at: \url{https://neurips.cc/Conferences/2026/PaperInformation/FundingDisclosure}.


% % Do {\bf not} include this section in the anonymized submission, only in the final paper. You can use the \texttt{ack} environment provided in the style file to automatically hide this section in the anonymized submission.
% % \end{ack}



\section*{References}
\bibliographystyle{unsrtnat}
% \bibliographystyle{plainnat}
\bibliography{NeurIPS_2026/ref}

% References follow the acknowledgments in the camera-ready paper. Use unnumbered first-level heading for
% the references. Any choice of citation style is acceptable as long as you are
% consistent. It is permissible to reduce the font size to \verb+small+ (9 point)
% when listing the references.
% Note that the Reference section does not count towards the page limit.
% \medskip


% {
% \small


% [1] Alexander, J.A.\ \& Mozer, M.C.\ (1995) Template-based algorithms for
% connectionist rule extraction. In G.\ Tesauro, D.S.\ Touretzky and T.K.\ Leen
% (eds.), {\it Advances in Neural Information Processing Systems 7},
% pp.\ 609--616. Cambridge, MA: MIT Press.


% [2] Bower, J.M.\ \& Beeman, D.\ (1995) {\it The Book of GENESIS: Exploring
%   Realistic Neural Models with the GEneral NEural SImulation System.}  New York:
% TELOS/Springer--Verlag.


% [3] Hasselmo, M.E., Schnell, E.\ \& Barkai, E.\ (1995) Dynamics of learning and
% recall at excitatory recurrent synapses and cholinergic modulation in rat
% hippocampal region CA3. {\it Journal of Neuroscience} {\bf 15}(7):5249-5262.
% }


%%%%%%%%%%%%%%%%%%%%%%%%%%%%%%%%%%%%%%%%%%%%%%%%%%%%%%%%%%%%

% \section{Appendix}
\clearpage
\appendix
\section*{Appendix}

\subsection{Scope and Summary of Notation}
\label{app:notation}

This appendix collects the main symbols used in the paper and in the appendices, following (approximately) the order in which they appear.

\begin{itemize}
% \item $t\in\{1,\dots,T\}$: task index in the continual learning sequence. % optional
% % \item $T$: total number of tasks.
% \item $S_t$: training set of task $t$.
% \item $m_t := |S_t|$: number of training samples in $S_t$.
% \item $D_t$: data distribution of task $t$.
% \item $z=(x,y)\sim D_t$: an input label pair drawn from $D_t$.
% % \item $f_W$: predictor (model) parameterized by $W$.
\item $\ell(W;z)$: bounded loss evaluated at parameters $W$ on sample $z$.
\item $L_{D_t}(W)=\mathbb E_{z\sim D_t}[\ell(W;z)]$: population risk on task $t$. 
\item $\widehat L_{S_t}(W)=\frac{1}{m_t}\sum_{z\in S_t}\ell(W;z)$: empirical risk on $S_t$.

% \item $W\in\mathbb R^{d}$: full parameter vector (or a vectorized weight block) of the model.

% \item $d$: dimensionality of the ambient parameter (or increment) space.


\item $\rho$: perturbation radius for adversarial sharpness.
% \item $\epsilon$: adversarial perturbation with $\|\epsilon\|\le \rho$.

% \item $\varepsilon\sim\mathcal N(0,\Sigma)$: Gaussian perturbation variable.
\item $\epsilon$ and $\varepsilon$: adversarial and Gaussian perturbations, respectively, where $\|\epsilon\|\le \rho$ and $\varepsilon\sim\mathcal N(0,\Sigma)$.

% \item $\Sigma\succeq 0$: covariance matrix for Gaussian perturbations.
% \item $\varepsilon\sim\mathcal N(0,\Sigma_t)$: Gaussian perturbation used to define expected (random) sharpness. 
% \item $\nabla_W \hat L_S(W)$: gradient of empirical risk with respect to $W$.
% \item $\mathrm{AS}_S^{(0)}(W,\rho)$: adversarial zeroth-order sharpness.
% \item $\mathrm{AS}_S^{(1)}(W,\rho)$: adversarial first-order sharpness.
\item $\mathrm{AS}^{(0)}_S(W,\rho)$: adversarial zeroth-order sharpness from SAM,
$\mathrm{AS}^{(0)}_S(W,\rho)=\max_{\|\epsilon\|\le\rho}\big[\widehat L_S(W+\epsilon)-\widehat L_S(W)\big]$.

\item $\mathrm{RS}_S^{(0)}(W,\Sigma)$: random zeroth-order sharpness, i.e., the expected empirical loss elevation under $\varepsilon\sim\mathcal N(0,\Sigma)$:
$\mathrm{RS}_S^{(0)}(W,\Sigma)=\mathbb E_{\varepsilon\sim\mathcal N(0,\Sigma)}[\hat L_S(W+\varepsilon)-\hat L_S(W)]$.
\item $\mathrm{AS}^{(1)}_S(W,\rho)$: adversarial first-order sharpness  from GAM,
$\mathrm{AS}^{(1)}_S(W,\rho)=\max_{\|\epsilon\|\le\rho}\|\nabla_W\widehat L_S(W+\epsilon)\|$.

% \item $P$: prior distribution over parameters  in PAC-Bayes analysis.
% \item $Q$: posterior distribution over parameters  in PAC-Bayes analysis.
% \item $\mathrm{KL}(Q\Vert P)$: Kullback--Leibler divergence measuring posterior complexity relative to the prior.
% \item $\delta\in(0,1]$: confidence level in high-probability bounds. 

\item $P_t$: task prior at step $t$ 
% (a distribution over model parameters restricted to the current task). 
\item $Q_t(\cdot \mid P_t,S_t)$: task posterior after learning task $t$, possibly conditioned on $P_t$ and $S_t$.
\item $Q_t^\Delta \ll P_t^\Delta$: absolute continuity assumption
\item $\mathcal P_t$: hyper-posterior (a distribution over task priors), forming a time-varying hyper-process . 
\item $\mathcal F_{t-1}$: filtration ($\sigma$-algebra) capturing the history up to step $t-1$ in the martingale argument.

\item $W_0$: pretrained initialization before continual adaptation.
\item $W_t$: model parameters after training on tasks $1$ to $t$.
% \item $W_t^{\mathrm{froz}}$: frozen component when learning task $t$ (the part not trainable during task $t$).
% \item $\Delta W_t$: trainable update at task $t$ under parameter-efficient adaptation.
\item $W_t = W_t^{\mathrm{froz}} + \Delta W_t$: decomposition of current parameters into frozen and trainable parts. 





% \item $\mathrm{LoRA}$: low-rank adaptation parameterization of a weight block, typically $\Delta W_t = B_tA_t$ with rank $r$.
% \item $A_t\in\mathbb R^{r\times d_{\text{in}}},\; B_t\in\mathbb R^{d_{\text{out}}\times r}$: LoRA factors for task $t$ (shapes may vary by layer/block).
% \item $r$: LoRA rank (adapter dimension).


% \item $\mathcal W_{\Delta,t}\subseteq\mathbb R^d$: the task-admissible linear update subspace at step $t$, containing all feasible increments induced by LoRA while the backbone is frozen.

\item $\mathcal W_{\Delta,t} \subseteq \mathbb R^d$: the task-admissible local linear update support at step $t$, identified with the first-order LoRA-induced tangent space around the learned adapter solution. It contains the feasible infinitesimal increment directions induced by the trainable LoRA coordinates while the backbone is frozen.


\item $\widehat W_t$: the full-model parameter vector after learning task $t$ (a point estimate in the ambient parameter space).
\item $\widehat{\Delta W}_t$: the learned trainable increment at task $t$ under PEFT/LoRA, satisfying $\widehat W_t = W_t^{\mathrm{froz}} + \widehat{\Delta W}_t$.

 


 
% \item $\mathcal B(\cdot)$: Borel sigma-algebra (and its trace on subspaces/affine sets). 
% \item $T_t(\Delta)=W_t^{\mathrm{froz}}+\Delta$: translation map from the update space to the affine model class. 
% \item $T_{t\#}\mu$: pushforward measure induced by $T_t$, defined by $(T_{t\#}\mu)(A)=\mu(T_t^{-1}(A))$. 
\item $P_t^\Delta,\;Q_t^\Delta$: prior/posterior supported on $(\mathcal W_{\Delta,t},\mathcal B(\mathcal W_{\Delta,t}))$.
% \item $P_,\; Q_t$: induced full-model prior/posterior supported on $W_t^{\mathrm{froz}}+\mathcal W_{\Delta,t}$.
% \item $Q_t^\Delta \ll P_t^\Delta$: absolute continuity assumption.
% \item $\frac{dQ}{dP}$: Radon--Nikodym derivative appearing in the KL invariance proof. 

% \item $\Sigma_t$: covariance defining a Gaussian perturbation distribution on $\mathcal W_{\Delta,t}$. 



\item $\mathrm{RS}_t^\Delta(\widehat W_t,\Sigma_t)$: task-conditioned random sharpness with perturbations supported on $\mathcal W_{\Delta,t}$, defined as
$\mathrm{RS}_t^\Delta(\widehat W_t,\Sigma_t)=\mathbb E_{\varepsilon\sim\mathcal N(0,\Sigma_t)}[\hat L_{S_t}(\widehat W_t+\varepsilon)-\hat L_{S_t}(\widehat W_t)]$,
where $\varepsilon$ is restricted to $\mathcal W_{\Delta,t}$, $\Sigma_t$ is defined on $\mathcal{W}_{\Delta,t}$
% (equivalently, $\mathrm{supp}(\varepsilon)\subseteq\mathcal W_{\Delta,t}$).


% \item $\widehat W_t$: posterior mean (or a point estimate) used when specializing $Q_t$ to a Gaussian.
% \item $\mathrm{RS}_t^\Delta(\widehat W_t,\Sigma_t)$:  random sharpness(expected sharpness) restricted to $\mathcal W_{\Delta,t}$ .

% \item $a_{tj}$: accuracy on task $j$ after training up to task $t$.
% \item $\mathrm{Acc}_t=\frac{1}{t}\sum_{j=1}^t a_{tj}$: average accuracy across the first $t$ tasks.
% \item $\mathrm{Acc}_t^{\mathrm{cls}}$: class-incremental classification accuracy at step $t$.
% \item $\mathrm{Acc}_t^{\mathrm{ncm}}$: nearest-class-mean accuracy (prototype-based) reflecting representation quality. 

\item $\alpha_t$ : feature amplification factor

% \item $c_t$: curvature projection ratio 

% \item $\lambda>0$: exponential tilting parameter used in the moment generating function (MGF) step of the proof~\ref{proof:theorem all}. 
% \item $B_T:=\sum_{t=1}^T (m_t-1)$: aggregate sample-size factor in Proof~\ref{proof:theorem all}. 
% \item $\omega_t := (m_t-1)/B_T$: sample-size weight used to form the weighted generalization gap.
% \item $\Delta_t$: step-$t$ generalization gap (defined in ~\ref{proof:theorem all}).
% \item $\overline\Delta := \sum_{t=1}^T \alpha_t \Delta_t$: size-weighted average gap. 
% \item $g_t(W):=\lambda\alpha_t\big(L_{D_t}(W)-\widehat L_{S_t}(W)\big)$: MGF test function at step $t$. 
% \item $h_t(P)$: hyper-level functional used in the Donsker--Varadhan (DV) change-of-measure step. 
% \item $\mathcal F_{t-1}$: filtration ($\sigma$-algebra) capturing the history up to step $t-1$ in the martingale argument. 
% \item $X_t$: exponential increment defined by the conditional MGF bound at step $t$.
% \item $\mathcal M_t:=\prod_{s=1}^t X_s$: nonnegative supermartingale used to obtain a high-probability bound.

% \item $\mathcal W_{\Delta,t}\subseteq\mathbb R^d$: permissible update subspace at task $t$ (task-admissible increment space). 
% \item $\Pi_{\Delta,t}$: orthogonal projector onto $\mathcal W_{\Delta,t}$ when $\mathcal W_{\Delta,t}$ is a linear subspace. 
% \item $J_{\mathrm{LoRA}}$: Jacobian of the LoRA map (used to define the local tangent increment space), with $\mathcal W_{\Delta,t}=\mathrm{Im}(J_{\mathrm{LoRA}})$.


% \item $\Delta W_t=U_{\Delta,t}S_{\Delta,t}V_{\Delta,t}^\top$: rank-$r$ truncated SVD of the learned update . 
% \item $\mathcal S_{\Delta,t}=\{U_{\Delta,t}ZV_{\Delta,t}^\top: Z\in\mathbb R^{r\times r}\}$: bilinear diagnostic subspace induced by the truncated SVD of $\Delta W_t$. 

% % \item $\Pi_{\Delta,t}^{\mathrm{SVD}}(M):=U_{\Delta,t}(U_{\Delta,t}^\top M V_{\Delta,t})V_{\Delta,t}^\top$: Frobenius-orthogonal projection onto $\mathcal S_{\Delta,t}$ (used only for curvature localization). 
% \item $\Pi_{\mathcal S,t}(M):=U_{\Delta,t}(U_{\Delta,t}^\top M V_{\Delta,t})V_{\Delta,t}^\top$: Frobenius-orthogonal projection onto $\mathcal S_{\Delta,t}$ .


% % \item $(U_{W,t},V_{W,t})$: rank-$r$ truncated SVD subspace of the corresponding backbone weight block (Appendix A.4.2). 
% % \item $\mathcal S_{W,t}=\{U_{W,t}ZV_{W,t}^\top: Z\in\mathbb R^{r\times r}\}$: backbone reference bilinear subspace. 
% % \item $\Pi_{W,t}^{\mathrm{SVD}}(M):=U_{W,t}(U_{W,t}^\top M V_{W,t})V_{W,t}^\top$: projection onto $\mathcal S_{W,t}$. 

% \item $(U_{W,t},V_{W,t})$: rank-$r$ truncated SVD bases of the corresponding backbone weight block.
% \item $\mathcal S_{W,t}=\{U_{W,t}ZV_{W,t}^\top: Z\in\mathbb R^{r\times r}\}$: backbone reference bilinear subspace.
% \item $\Pi_{\mathcal S_W,t}(M):=U_{W,t}(U_{W,t}^\top M V_{W,t})V_{W,t}^\top$: Frobenius-orthogonal projection onto $\mathcal S_{W,t}$.

% \item $\{(\lambda_{t,i},v_{t,i})\}_{i=1}^k$: top-$k$ eigenpairs of a PSD curvature operator (e.g., GGN/Fisher) at task $t$. 
% \item $(\delta W^{(0)}_{t,i},\delta A_{t,i},\delta B_{t,i})$: decomposition of eigen-direction $v_{t,i}$ into backbone and LoRA-factor components. 
% \item $\delta W_{t,i}=\delta W^{(0)}_{t,i}+(\delta B_{t,i}A_t+B_t\delta A_{t,i})$: effective weight-space perturbation induced by $(\delta W^{(0)}_{t,i},\delta A_{t,i},\delta B_{t,i})$. 
% \item $\rho_{\Delta,t,i}=\frac{\|\Pi^{\mathrm{SVD}}_{\Delta,t}(\delta W_{t,i})\|_F^2}{\|\delta W_{t,i}\|_F^2}$: per-mode projection ratio onto $\mathcal S_{\Delta,t}$.
% \item $\rho_{W,t,i}=\frac{\|\Pi_{W,t}^{\mathrm{SVD}}(\delta W_{t,i})\|_F^2}{\|\delta W_{t,i}\|_F^2}$: per-mode projection ratio onto $\mathcal S_{W,t}$.
\item $R_{\Delta,t}=\frac{\sum_{i=1}^k\lambda_{t,i}\rho_{\Delta,t,i}}{\sum_{i=1}^k\lambda_{t,i}}$: eigenvalue-weighted localization on $\mathcal S_{\Delta,t}$.
\item $R_{W,t}=\frac{\sum_{i=1}^k\lambda_{t,i}\rho_{W,t,i}}{\sum_{i=1}^k\lambda_{t,i}}$: eigenvalue-weighted localization on $\mathcal S_{W,t}$.
\item $c_t=R_{\Delta,t}/R_{W,t}$: curvature projection ratio comparing $\mathcal S_{\Delta,t}$ against the backbone reference. 
\end{itemize}


\subsection{Proof}

\subsubsection{Proof of Theorem \ref{thm:sequential_pac_bayes}}
\label{proof:theorem all}

\paragraph{Preliminary: a supermartingale tail bound.}
Let $(\mathcal M_t)_{t\ge0}$ be a nonnegative supermartingale adapted to $(\mathcal F_t)$ with $\mathbb E[\mathcal M_0]\le 1$.
Then for any $\delta\in(0,1)$,
\[
\mathbb P\!\left(\mathcal M_T \ge \tfrac{1}{\delta}\right)
\le \delta,
\]
by Markov's inequality and $\mathbb E[\mathcal M_T]\le \mathbb E[\mathcal M_0]\le 1$.
Equivalently, with probability at least $1-\delta$, $\log \mathcal M_T \le \log \tfrac{1}{\delta}$.
This is the fixed-time form of the classical supermartingale maximal inequality  \citet{yangshuxue}.


\begin{theorem}[Pathwise PAC--Bayes with time-varying hyper-posterior]
\label{the: full}
For any $\delta\in(0,1]$, with probability at least $1-\delta$ over $S_{1:T}$,
\begin{equation}
\begin{aligned} 
& \frac{1}{T} \sum_{t=1}^T \underbrace{\mathbb{E}_{P_t \sim \mathcal{P}_t} \mathbb{E}_{W \sim Q_t(\cdot \mid P_t, S_t)} L_{\mathcal{D}_t}(W)}_{\text{test risk at step } t} 
\le \ \frac{1}{T} \sum_{t=1}^T \underbrace{\mathbb{E}_{P_t \sim \mathcal{P}_t} \mathbb{E}_{W \sim Q_t(\cdot \mid P_t, S_t)} \hat{L}_{S_t}(W)}_{\text{empirical risk at step } t} \\ 
& + \sqrt{\frac{ \sum_{t=1}^T \mathbb{E}_{P_t \sim \mathcal{P}_t} \mathrm{KL}\big(Q_t(\cdot \mid P_t, S_t) \,\|\, P_t\big) + \sum_{t=1}^T \mathrm{KL}(\mathcal{P}_t \,\|\, \mathcal{P}_{t-1}) + \log\frac{1}{\delta} }{2 \sum_{t=1}^T (m_t - 1)}}. 
\end{aligned} 
\end{equation}
\end{theorem}

\begin{proof}
    


% \subsection*{A.1\quad Two-level variational inequalities}
% \item $B_T:=\sum_{t=1}^T (m_t-1)$: aggregate sample-size factor in Proof~\ref{proof:theorem all}. 
% \item $\omega_t := (m_t-1)/B_T$: sample-size weight used to form the weighted generalization gap.
Let $B_T=\sum_{t=1}^T (m_t-1)$ denote an aggregate sample-size factor and define $\omega_t=(m_t-1)/B_T$ which serves as a sample-size weight satisfying $\sum_{t=1}^T\omega_t=1$.
% $B_T$ is 

% \paragraph{Goal quantity.} 


Define the step-$t$ generalization gap
\[
\Delta_t
:=\mathbb E_{P_t\sim\mathcal P_{t-1}}\,
\mathbb E_{W\sim Q_t(\cdot\mid P_t,S_t)}
\big(L_{\mathcal D_t}(W)-\hat L_{S_t}(W)\big),
\]
and the size-weighted gap $\bar\Delta:=\sum_{t=1}^T \omega_t \Delta_t$.
Fix $\lambda>0$ and set
\(
g_t(W):=\lambda\omega_t\big(L_{\mathcal D_t}(W)-\hat L_{S_t}(W)\big).
\)

\paragraph{Step 1 Fixed-task PAC-Bayes bound.}
% For any fixed $P$ and $S_t$,
For any fixed prior distribution over the parameters $P$ and posterior distribution $Q$, and only current task dataset $S_t$,   
\begin{equation}
\label{eq:dv-w}
\mathbb E_{W\sim Q_t(\cdot\mid P,S_t)} g_t(W)
\;\le\;
\mathrm{KL}\!\big(Q_t(\cdot\mid P,S_t)\,\|\,P\big)
+\log\mathbb E_{W\sim P}\!\big[e^{g_t(W)}\big].
\end{equation}
Taking $\mathbb E_{S_t}$ and using the standard Hoeffding--MGF bound for bounded losses,
\begin{equation}
\label{eq:mgf-mean-gap}
\mathbb E_{S_t}\log\mathbb E_{W\sim P} e^{g_t(W)}
\;\le\; \frac{(\lambda\omega_t)^2}{8\,(m_t-1)}.
\end{equation}
Therefore,
\begin{equation}
\label{eq:model-step}
\mathbb E_{S_t}\mathbb E_{W\sim Q_t(\cdot\mid P,S_t)} g_t(W)
\;\le\;
\mathbb E_{S_t}\mathrm{KL}\!\big(Q_t(\cdot\mid P,S_t)\,\|\,P\big)
+\frac{(\lambda\omega_t)^2}{8\,(m_t-1)}.
\end{equation}
 % (change of measure in $P$-space)
\paragraph{Step 2 Hyperposterior lifting via the Donsker-Varadhan change-of-measure inequality} 
Define
\[
h_t(P):=\mathbb E_{S_t}\mathbb E_{W\sim Q_t(\cdot\mid P,S_t)} g_t(W)
-\mathbb E_{S_t}\mathrm{KL}\!\big(Q_t(\cdot\mid P,S_t)\,\|\,P\big).
\]
By \eqref{eq:model-step}, $h_t(P)\le \frac{(\lambda\omega_t)^2}{8\,(m_t-1)}$ for all $P$.
Applying Donsker–Varadhan variational representation of the KL divergence~\cite{NEURIPS2021_54a367d6} to change the measure from $\mathcal P_{t-1}$ to $\mathcal P_t$ gives
\begin{equation}
\label{eq:dv-p}
\log\mathbb E_{P\sim\mathcal P_t}\!\big[e^{h_t(P)}\big]
\;\le\; \mathrm{KL}(\mathcal P_t\|\mathcal P_{t-1})
+\log\mathbb E_{P\sim\mathcal P_{t-1}}\!\big[e^{h_t(P)}\big]
\;\le\; \mathrm{KL}(\mathcal P_t\|\mathcal P_{t-1})
+\frac{(\lambda\omega_t)^2}{8\,(m_t-1)}.
\end{equation}
Using $\log\mathbb E[e^X]\ge \mathbb E[X]$ yields
\begin{equation}
\label{eq:p-expect}
\mathbb E_{P_t\sim\mathcal P_t} h_t(P_t)
\;\le\; \mathrm{KL}(\mathcal P_t\|\mathcal P_{t-1})
+\frac{(\lambda\omega_t)^2}{8\,(m_t-1)}.
\end{equation}
Unfolding $h_t(P)$ and interchanging expectations,
\begin{equation}
\label{eq:two-term}
\mathbb E_{S_t}\Big[
\lambda\omega_t\Delta_t
-\omega_t\,\mathbb E_{P_t\sim\mathcal P_{t-1}}\mathrm{KL}\!\big(Q_t(\cdot\mid P_t,S_t)\,\|\,P_t\big)
\Big]
\;\le\;
\mathrm{KL}(\mathcal P_t\|\mathcal P_{t-1})
+\frac{(\lambda\omega_t)^2}{8\,(m_t-1)}.
\end{equation}

% \subsection*{A.2\quad Conditional MGF and supermartingale}
\paragraph{Step 3 Martingale}
% From \eqref{eq:two-term} and the tower property of conditional expectation, 







Using Eq.\eqref{eq:two-term} and the tower property of conditional expectation (iterated expectation) with respect to $\mathcal F_{t-1}$, we obtain the desired conditional exponential bound~\cite{yangshuxue}.
% The conditional exponential bound
\begin{equation}
\label{eq:cond-mgf}
\mathbb E\!\left[
\exp\!\Big(
\lambda\omega_t\Delta_t
-\omega_t\,\mathbb E_{P_t\sim\mathcal P_{t-1}}\mathrm{KL}(Q_t\|P_t)
-\mathrm{KL}(\mathcal P_t\|\mathcal P_{t-1})
-\frac{(\lambda\omega_t)^2}{8\,(m_t-1)}
\Big)\,\Big|\,\mathcal F_{t-1}\right]\;\le\;1.
\end{equation}


From the conditional exponential bound~\eqref{eq:cond-mgf}, define
\[
X_t
:=
\exp\!\Big(
\lambda\omega_t\Delta_t
-\omega_t\,\mathbb E_{P_t\sim\mathcal P_{t-1}}\mathrm{KL}(Q_t\|P_t)
-\mathrm{KL}(\mathcal P_t\|\mathcal P_{t-1})
-\frac{(\lambda\omega_t)^2}{8\,(m_t-1)}
\Big).
\]
Assume $\omega_t$ is deterministic or $\mathcal F_{t-1}$-measurable. Then $X_t\ge 0$ is $\mathcal F_t$-measurable and accoring Eq.~\eqref{eq:cond-mgf}
\[
\mathbb E[X_t\mid \mathcal F_{t-1}] \le 1.
\]
Define the process
\[
\mathcal{M}_0:=1,\qquad
\mathcal{M}_t:=\prod_{s=1}^t X_s.
\]
It follows that $(\mathcal M_t)_{t\ge 0}$ is a nonnegative supermartingale, since
\[
\mathbb E[\mathcal M_t\mid \mathcal F_{t-1}]
=
\mathcal M_{t-1}\,\mathbb E[X_t\mid \mathcal F_{t-1}]
\le \mathcal M_{t-1}.
\]
In particular, $\mathbb E[\mathcal M_T]\le \mathbb E[\mathcal M_0]=1$.
By Markov's inequality,
\[
\mathbb P\!\left(\mathcal M_T \ge \tfrac{1}{\delta}\right)
\le \delta,
\]
so with probability at least $1-\delta$ we have $\log \mathcal M_T \le \log\tfrac{1}{\delta}$.
% Expanding $\log \mathcal M_T = \sum_
Expanding $\log \mathcal M_T = \sum_{t=1}^T \log X_t$ yields~\eqref{eq:hpp}.






% Define $X_t$ to be the exponential increment inside \eqref{eq:cond-mgf} and
% \[
% \mathcal{M}_0:=1,\qquad
% \mathcal{M}_t:=\prod_{s=1}^t X_s.
% \]
% Then $(\mathcal{M}_t)_{t\ge0}$ is a nonnegative supermartingale, so
% $\mathbb E[\mathcal{M}_T]\le 1$.
% By Markov's inequality, with probability at least $1-\delta$,

% so with probability at least $1-\delta$ we have $\log \mathcal M_T \le \log\tfrac{1}{\delta}$

\begin{equation}
\label{eq:hpp}
\sum_{t=1}^T \lambda\omega_t\Delta_t
\;\le\;
\sum_{t=1}^T \omega_t\,\mathbb E_{P_t\sim\mathcal P_{t-1}}\mathrm{KL}(Q_t\|P_t)
\;+\; \sum_{t=1}^T \mathrm{KL}(\mathcal P_t\|\mathcal P_{t-1})
\;+\; \frac{\lambda^2}{8}\sum_{t=1}^T \frac{\omega_t^2}{(m_t-1)}
\;+\; \log\tfrac1\delta.
\end{equation}

% \subsection*{A.3\quad Size-weighted self-normalization and optimization}

By $\omega_t=(m_t-1)/B_T$,
\[
\sum_{t=1}^T \frac{\omega_t^2}{(m_t-1)}
=\sum_{t=1}^T \frac{(m_t-1)^2/B_T^2}{(m_t-1)}
=\frac{1}{B_T}.
\]
Hence, \eqref{eq:hpp} reads, for all $\lambda>0$ and with probability at least $1-\delta$,
\begin{equation}
\label{eq:gap-preopt}
\bar\Delta
\;\le\;
\frac{1}{\lambda}\Big(
\sum_{t=1}^T \omega_t\,\mathbb E_{P_t\sim\mathcal P_{t-1}}\mathrm{KL}(Q_t\|P_t)
+ \sum_{t=1}^T \mathrm{KL}(\mathcal P_t\|\mathcal P_{t-1})
+ \log\tfrac1\delta\Big)
\;+\;
\frac{\lambda}{8B_T}.
\end{equation}
Optimizing the right-hand side over $\lambda>0$ (take $\lambda^\star=\sqrt{8B_T\,A}$ with
$A:=\sum_{t=1}^T \omega_t\,\mathbb E_{P_t\sim\mathcal P_{t-1}}\mathrm{KL}(Q_t\|P_t)
+ \sum_{t=1}^T \mathrm{KL}(\mathcal P_t\|\mathcal P_{t-1}) + \log\tfrac1\delta$) yields
\begin{equation}
\label{eq:weighted-final}
\bar\Delta
\;\le\;
\sqrt{\frac{A}{2B_T}}
\;\le\;
\sqrt{\frac{
\sum_{t=1}^T \mathbb E_{P_t\sim\mathcal P_{t-1}}\mathrm{KL}(Q_t\|P_t)
+ \sum_{t=1}^T \mathrm{KL}(\mathcal P_t\|\mathcal P_{t-1})
+ \log\tfrac1\delta}{2\,\sum_{t=1}^T (m_t-1)}}.
\end{equation}
The last inequality uses $\sum_t \omega_t x_t \le \sum_t x_t$ for nonnegative $(x_t)$.

\paragraph{Step 4 From weighted to the statement of Theorem~\ref{the: full}.}
The theorem controls the unweighted average $\frac1T\sum_{t=1}^T \Delta_t$ and uses the smoothed empirical risk on the right-hand side.
When $m_t$ are homogeneous, $\omega_t\equiv 1/T$ and $\bar\Delta=\frac1T\sum_t\Delta_t$.
For heterogeneous $m_t$, we keep the same right-hand side and upper-bound the numerator $\sum_t \omega_t \mathbb E_{P_t}\mathrm{KL}(Q_t\|P_t)$ by $\sum_t \mathbb E_{P_t}\mathrm{KL}(Q_t\|P_t)$, which yields exactly the denominator $2\sum_t(m_t-1)$ in the theorem.

% \subsection*{A.4\quad Gaussian smoothing (LoRA) corollary}

If $Q_t(\cdot\mid P_t,S_t)=\mathcal N(\hat W_t,\Sigma_t)$ (e.g., supported on an adapter subspace), then
\[
\mathbb E_{W\sim Q_t(\cdot\mid P_t,S_t)}\hat L_{S_t}(W)
=\mathbb E_{\varepsilon\sim\mathcal N(0,\Sigma_t)}\hat L_{S_t}(\hat W_t+\varepsilon),
\]
which converts the empirical term into the smoothed empirical loss (expected sharpness) used in the main text; substituting this identity into \eqref{eq:weighted-final} yields the empirical term in Theorem~\ref{the: full}.

Theorem~A.1 provides a tight complexity penalty of the form.
Using $\sqrt{x+y+z}\le \sqrt{x}+\sqrt{y}+\sqrt{z}$ for $x,y,z\ge 0$ yields the decomposed bound, which is the form stated in Theorem~\ref{thm:sequential_pac_bayes} for interpretability.

\end{proof}

\subsubsection{Proof of Lemma~\ref{lem:kl-decomp-pecl}}
\label{proof:Lemma1}

Let $\mathcal W_{\Delta,t}\subseteq\mathbb R^d$ be a linear subspace and $W_t^\mathrm{froz}+\mathcal W_{\Delta,t}$ be its affine translation. 
We endow both sets with the trace Borel $\sigma$-algebra induced from $\mathbb R^d$. 
Define the translation map
\[
T_t:(\mathcal W_{\Delta,t},\mathcal B(\mathcal W_{\Delta,t}))
\to (W_t^{\mathrm{froz}}+\mathcal W_{\Delta,t},\mathcal B(W_t^{\mathrm{froz}}+\mathcal W_{\Delta,t})),
\qquad
T_t(\Delta)=W_t^{\mathrm{froz}}+\Delta .
\]
For any probability measure $\mu$ on $\mathcal W_{\Delta,t}$, its pushforward
$T_{t\#}\mu$ is the measure on $W_t^{\mathrm{froz}}+\mathcal W_{\Delta,t}$
defined by $(T_{t\#}\mu)(A)=\mu(T_t^{-1}(A))$ for all measurable $A$.
Accordingly, we define the induced full-model prior and posterior supported on
$W_t^{\mathrm{froz}}+\mathcal W_{\Delta,t}$ as
\[
P_t := T_{t\#} P_t^\Delta,
\qquad
Q_t := T_{t\#} Q_t^\Delta,
\]
where $P_t^\Delta$ and $Q_t^\Delta$ are probability measures on
$(\mathcal W_{\Delta,t},\mathcal B(\mathcal W_{\Delta,t}))$.
% $\mathcal B(\mathcal W_{\Delta,t})\ =\{B\cap \mathcal W_{\Delta,t}:\ B\in \mathcal B(\mathbb R^d)\},\quad
% \mathcal B(W^{\mathrm{froz}}_t+\mathcal W_{\Delta,t})\ =\{B\cap (W^{\mathrm{froz}}_t+\mathcal W_{\Delta,t}):\ B\in \mathcal B(\mathbb R^d)\}.$
% These coincide with the Borel $\sigma$-algebras generated by the subspace topologies on $\mathcal W_{\Delta,t}$ and $W^{\mathrm{froz}}_t+\mathcal W_{\Delta,t}$, respectively, so both are Borel spaces.
% To formalize the probabilistic relationship between the LoRA subspace and the full model space, we define the translation map:
% % \begin{align*}
% $T_t: (\mathcal W_{\Delta,t},\mathcal B(\mathcal W_{\Delta,t})) \to (W^{\mathrm{froz}}_t+\mathcal W_{\Delta,t},\mathcal B(W^{\mathrm{froz}}_t+\mathcal W_{\Delta,t})), 
% T_t(\Delta)=W^{\mathrm{froz}}_t+\Delta,$
% % \end{align*}
% which is a Borel isomorphism between the subspace and affine space.  The full-space prior and posterior are defined as pushforwards onto the affine subspace by
% $
% P=(T_t)_{\text{\#}} P^{\Delta}_t,
% Q=(T_t)_{\text{\#}} Q^{\Delta}_t,
% $
% where $P_t^{\Delta}$ and $Q_t^{\Delta}$ are probability measures on $(\mathcal{W}_{\Delta,t}, \mathcal{B}(\mathcal{W}_{\Delta,t}))$.

% \begin{lemma}[KL decomposition with a frozen backbone]
% \label{equation: subspace-full}
% Let $W\in\mathbb R^d$ be fixed and let $P^{\Delta},Q^{\Delta}$ be probability measures on $(\mathcal W_\Delta,\mathcal B(\mathcal W_\Delta))$ with $Q^{\Delta}\ll P^{\Delta}$. Define their pushforwards onto the affine subspace by
% \[
% P:=(T_W)_{\text{\#}} P^{\Delta},\qquad
% Q:=(T_W)_{\text{\#}} Q^{\Delta}.
% \]
% Then $Q\ll P$ and
% \[
% \mathrm{KL}\left(Q\middle|P\right)
% =\mathrm{KL}\left(Q^{\Delta}\middle|P^{\Delta}\right).
% \]
% Equivalently, identifying ${W}\times\mathcal W_\Delta$ with $W+\mathcal W_\Delta$ via $\Phi$, one may write
% $P=\Phi_{\text{\#}}(\delta_W\otimes P^{\Delta})$ and
% $Q=\Phi_{\text{\#}}(\delta_W\otimes Q^{\Delta})$, and the same equality holds.
% \end{lemma}

\begin{lemma}[KL in PECL is adapter-subspace KL]
% Assume $Q_t^\Delta \ll P_t^\Delta$. 
% Assume $Q_t^\Delta \ll P_t^\Delta$, meaning $Q_t^\Delta(A)=0$ for any measurable set $A$ such that $P_t^\Delta(A)=0$.
Assume $Q_t^\Delta \ll P_t^\Delta$ (i.e., $Q_t^\Delta$ is absolutely continuous w.r.t.\ $P_t^\Delta$), meaning $Q_t^\Delta(A)=0$ for any measurable set $A$ such that $P_t^\Delta(A)=0$.
Then $Q_t \ll P_t$ and
$$
\mathrm{KL}\bigl(Q_t\,\Vert\, P_t\bigr)
\;=\;
\mathrm{KL}\bigl(Q_t^\Delta \,\Vert\, P_t^\Delta\bigr).
$$
\end{lemma}

\begin{proof}
  % Absolute continuity is preserved by pushforward under measurable maps, hence $Q_t\ll P_t$. Since $T_t$ is a Borel isomorphism, 
  Since $T_t$ is measurable, absolute continuity is preserved under pushforward, hence
$Q_t=T_{t\#}Q_t^\Delta \ll T_{t\#}P_t^\Delta=P_t$.
Moreover, $T_t$ is a measurable bijection with measurable inverse. 
  By the Radon-Nikodym chain rule
 \[
 \frac{dQ_t}{dP_t}(T_t(\Delta))=\frac{dQ_t^{\Delta}}{dP_t^{\Delta}}(\Delta)\quad\text{for }P^{\Delta}_t\text{-a.e.\ }\Delta.
 \]

\begin{equation*}
     \begin{aligned}
 \mathrm{KL}\left(Q_t\middle|P_t\right)
 &= \int_{W+\mathcal W_{\Delta,t}}
 \log\Big(\frac{dQ_t}{dP_t}(y)\Big)
 Q_t(dy)\\
 &= \int_{\mathcal W_{\Delta,t}}
 \log\Big(\frac{dQ_t}{dP_t}\big(T_t\Delta\big)\Big)
 Q_t^{\Delta}(d\Delta) \quad \text{(} Q_t=(T_t)_{\text{\#}}Q_t^{\Delta}\text{)} 
 \\
 &= \int_{\mathcal W_{\Delta,t}}
 \log\Big(\frac{dQ_t^{\Delta}}{dP_t^{\Delta}}(\Delta)\Big)
 Q_t^{\Delta}(d\Delta) \\
 &= \mathrm{KL}\left(Q_t^{\Delta}\middle|P_t^{\Delta}\right).
 \end{aligned}
\end{equation*} 
\end{proof}

% \section{ Proof of Lemma~\ref{lemma: subspace-full}}
% \begin{lemma}[Sharpness decomposition in PECL]
% \label{lemma: Sharpness in PECL}
% Let $W_0\in\mathbb R^d$ be fixed and let the trainable update lie in the adapter subspace $\mathcal W_\Delta\subseteq\mathbb R^d$. Consider a Gaussian posterior on $(\mathcal W_\Delta,\mathcal B(\mathcal W_\Delta))$,
%  $
%  Q_t^{\Delta}=\mathcal N(\hat{\Delta W}_t,\Sigma_t),
%  $
%  with mean $\hat{\Delta W}_t\in\mathcal W_\Delta$ and covariance operator $\Sigma_t\succeq 0$ acting on $\mathcal W_\Delta$ (viewed in $\mathbb R^d$, this covariance is supported in $\mathcal W_\Delta$). Define the full-space posterior on ${W_0}\times\mathcal W_\Delta$ by
%  $
%  Q_t:=\delta_{W_0}\otimes Q_t^{\Delta},
%  $
% and identify ${W_0}\times\mathcal W_\Delta$ with the affine subspace $W_0+\mathcal W_\Delta$ via $(w,\Delta)\mapsto w+\Delta$. Let $\hat L_{S_t}:\mathbb R^d\to\mathbb R$ be measurable and integrable under $Q_t$. Then
%  \[
%  \mathbb{E}_{W\sim Q_t}\big[\hat{L}_{S_t}(W)\big]
%  = \mathbb{E}_{\varepsilon\sim\mathcal N(0,\Sigma_t)}
%  \big[\hat{L}_{S_t}(W_0+\hat{\Delta W}_t+\varepsilon)\big]
%  = \hat{L}_{S_t}(W_0+\hat{\Delta W}_t) +\mathrm{RS}_t^{\Delta}(\hat{\Delta W}_t,\Sigma_t) \]
%    where the subspace expected sharpness is
%    \[
%    \mathrm{RS}_t^{\Delta}(\hat{\Delta W}_t,\Sigma_t)
%    = \mathbb{E}_{\varepsilon\sim\mathcal N(0,\Sigma_t)}
%    \big[\hat{L}_{S_t}(W_0+\hat{\Delta W}_t+\varepsilon)
%   - \hat{L}_{S_t}(W_0+\hat{\Delta W}_t)\big].
%      \]
% \end{lemma}
% \begin{proof}
%     Sampling $W\sim Q_t$ is equivalent to sampling $\varepsilon\sim\mathcal N(0,\Sigma_t)$ on $\mathcal W_\Delta$ and setting $W=W_0+\hat{\Delta W}_t+\varepsilon$. Taking expectations and adding-subtracting $\hat L_{S_t}(W_0+\hat{\Delta W}_t)$ inside the expectation gives the decomposition. 
% \end{proof}

\subsubsection{Proof of Lemma~\ref{lem:subspace-loss-pecl}}
\label{proof:lemma2}

% \begin{equation*}
% \label{eq:subspace-loss-pecl}
% \begin{aligned}
% &\mathbb E_{W\sim Q_t}\bigl[\hat L_{S_t}(W)\bigr] \\
% &= \mathbb{E}_{\varepsilon\sim\mathcal N(0,\Sigma_t)}
%     \big[\hat{L}_{S_t}(W_\mathrm {frozen }+\hat{\Delta W}_t+\varepsilon)\big],\varepsilon\in\mathcal W_\Delta \\
% & =\hat L_{S_t}(W_{\mathrm{frozen}} + \hat{\Delta W}_t)
% +
% \mathrm{RS}_t^\Delta(\hat{\Delta W}_t,\Sigma_t),\varepsilon\in\mathcal W_\Delta.
% \end{aligned}
% \end{equation*}
We specialize the posterior to a Gaussian supported on $\mathcal W_\Delta$. For task $t$, we consider a Gaussian posterior \( Q_t^{\Delta} = \mathcal{N}(\hat{\Delta W}_t, \Sigma_t) \) on the measurable space \( (\mathcal{W}_\Delta, \mathcal{B}(\mathcal{W}_\Delta)) \), where \( \hat{\Delta W}_t \in \mathcal{W}_\Delta \) is the mean and \( \Sigma_t\) is a covariance on the adapter subspace \( \mathcal{W}_\Delta \). Sampling from \( \Delta W \sim Q_t^{\Delta} \) is equivalent to 
$
\Delta W = \hat{\Delta W}_t + \varepsilon, \quad \varepsilon \sim \mathcal{N}(0, \Sigma_t) \ \text{on } \mathcal{W}_\Delta.
$ 

\begin{lemma}
    The posterior expected empirical loss decomposes as
\begin{equation*}
    \mathbb{E}_{W\sim Q_t}\big[\hat{L}_{S_t}(W)\big]
= \mathbb{E}_{\varepsilon\sim\mathcal N(0,\Sigma_t)}
    \big[\hat{L}_{S_t}(W_\mathrm {frozen }+\hat{\Delta W}_t+\varepsilon)\big]
= \hat{L}_{S_t}(W_\mathrm {frozen }+\hat{\Delta W}_t)
  + \mathrm{RS}_t^{\Delta}(\hat{\Delta W}_t,\Sigma_t),
\end{equation*}
where the random sharpness is
$$
\mathrm{RS}_t^{\Delta}(\hat{\Delta W}_t,\Sigma_t)
:= \mathbb{E}_{\varepsilon\sim\mathcal N(0,\Sigma_t)}
   \big[\hat{L}_{S_t}(W_\mathrm {frozen }+\hat{\Delta W}_t+\varepsilon)
       - \hat{L}_{S_t}(W_\mathrm {frozen }+\hat{\Delta W}_t)\big],\quad
\varepsilon\sim\mathcal N(0,\Sigma_t)\ \text{on } \mathcal W_\Delta$$
\end{lemma}


\begin{proof}
Recall that the adapter posterior is a Gaussian supported on the update subspace,
$
Q_t^\Delta = \mathcal N(\hat{\Delta W}_t,\Sigma_t)
$
on $(\mathcal W_\Delta,\mathcal B(\mathcal W_\Delta))$.
Hence we can represent a draw $\Delta W \sim Q_t^\Delta$ as
$
\Delta W = \hat{\Delta W}_t + \varepsilon
$
with $\varepsilon \sim \mathcal N(0,\Sigma_t)$ on $\mathcal W_\Delta$.

Under the PECL constraint, the full parameter is obtained by translating the adapter update by the frozen part,
$
W = W_{\mathrm{frozen}} + \Delta W.
$
Therefore, for the empirical risk,
\begin{align*}
\mathbb E_{W \sim Q_t}\big[\hat L_{S_t}(W)\big]
&= \mathbb E_{\Delta W \sim Q_t^\Delta}\big[\hat L_{S_t}(W_{\mathrm{frozen}}+\Delta W)\big] \\
&= \mathbb E_{\varepsilon \sim \mathcal N(0,\Sigma_t)}
\big[\hat L_{S_t}(W_{\mathrm{frozen}}+\hat{\Delta W}_t+\varepsilon)\big].
\end{align*}
Finally, add and subtract the mean point loss to obtain
\begin{align*}
\mathbb E_{\varepsilon \sim \mathcal N(0,\Sigma_t)}
\big[\hat L_{S_t}(W_{\mathrm{frozen}}+\hat{\Delta W}_t+\varepsilon)\big]
&=
\hat L_{S_t}(W_{\mathrm{frozen}}+\hat{\Delta W}_t) \\
&\quad +
\mathbb E_{\varepsilon \sim \mathcal N(0,\Sigma_t)}
\Big[
\hat L_{S_t}(W_{\mathrm{frozen}}+\hat{\Delta W}_t+\varepsilon)
-
\hat L_{S_t}(W_{\mathrm{frozen}}+\hat{\Delta W}_t)
\Big].
\end{align*}
The second term is exactly the random Sharpness on the subspace
\[
\mathrm{RS}_t^\Delta(\hat{\Delta W}_t,\Sigma_t)
:=
\mathbb E_{\varepsilon \sim \mathcal N(0,\Sigma_t)}
\Big[
\hat L_{S_t}(W_{\mathrm{frozen}}+\hat{\Delta W}_t+\varepsilon)
-
\hat L_{S_t}(W_{\mathrm{frozen}}+\hat{\Delta W}_t)
\Big],
\]
completes the proof.
\end{proof}



\subsubsection{Proof of Theorem \ref{thm:pecl-bound}}
\label{proof:theorem}

% \subsection{Proof of }





\begin{theorem}[Sharpness-aware PAC--Bayes bound for PECL]
% \label{thm:pecl-bound}

For any $\delta\in(0,1]$, with probability at least $1-\delta$ over $S_1,\ldots,S_T$,
\begin{equation}\label{eq:pecl-main}
\begin{aligned}
&\frac{1}{T}\sum_{t=1}^T
\mathbb{E}_{P_t\sim\mathcal P_{t-1}}\mathbb{E}_{W\sim Q_t}L_{\mathcal{D}_t}(W)
\ \le\
\ \frac{1}{T}\sum_{t=1}^T
\mathbb{E}_{P_t\sim\mathcal P_{t-1}}
\mathbb{E}_{\varepsilon \sim \mathcal{N}(0, \Sigma_t)}[\hat{L}_{S_t}(W_0 + \widehat{\Delta W}_t + \varepsilon)
\\
&\quad +\sqrt{\frac{
\displaystyle \sum_{t=1}^T \mathbb{E}_{P_t^\Delta\sim\mathcal P_{t-1}}\!\Big[\mathrm{KL}\!\big(Q_t^\Delta\|P_t^\Delta\big)\Big]
\;+\;\displaystyle \sum_{t=1}^T \mathrm{KL}\!\big(\mathcal P_t\|\mathcal P_{t-1}\big)
\;+\;\log\!\frac{1}{\delta}
}{2\,\displaystyle\sum_{t=1}^T (m_t-1)}}.
\end{aligned}
\end{equation}
\end{theorem}



\begin{proof}
The bound follows by specializing Theorem \ref{the: full} to the PECL setting through two key decompositions:
By Lemma~\ref{lem:kl-decomp-pecl}, the full-space KL terms reduce to their adapter-space counterparts:
\[
\mathrm{KL}(Q_t \| P_t) = \mathrm{KL}(Q_t^\Delta \| P_t^\Delta).
\]
% \quad
% \mathrm{KL}(\mathcal{P}_t \| \mathcal{P}_{t-1}) = \mathrm{KL}(\mathcal{P}_t^\Delta \| \mathcal{P}_{t-1}^\Delta)
% \textbf{Step 2: Sharpness decomposition via Lemma 6.}  
Under the Gaussian posterior $Q_t^\Delta = \mathcal{N}(\Delta\hat{W}_t, \Sigma_t)$, Lemma \ref{lem:subspace-loss-pecl} gives:
\[
\mathbb{E}_{W \sim Q_t}[\hat{L}_{S_t}(W)] 
= \mathrm{E}_{\varepsilon \sim \mathcal{N}(0, \Sigma_t)}{L}_{S_t}(W^{\mathrm{froz}}_t+\widehat{\Delta W}_t + \varepsilon)
 = \hat{L}_{S_t}(W^{\mathrm{froz}}_t + \widehat{\Delta{W}_t}) + \mathrm{RS}_t^\Delta(\Delta\widehat{W}_t, \Sigma_t),
\]
where the perturbations are confined to the adapter subspace.

% \textbf{Step 3: Substitution into general bound.}
Substituting these decompositions into Lemma \ref{the: full} yields the final PECL-specific bound, where all complexity and sharpness terms are now restricted to the trainable LoRA subspace $\mathcal{W}_\Delta$.
% Substituting these decompositions into Theorem 2 yields the final PECL-specific bound, where all complexity and sharpness terms are now restricted to the trainable LoRA subspace $\mathcal{W}_\Delta$.

\end{proof}



\subsection{Bridging the Permissible Update Subspace and the LoRA Implementation}
\label{app:lora_subspace_bridge}

This appendix clarifies how the abstract permissible update subspace
$\mathcal W_{\Delta,t}$ used in the PAC--Bayesian analysis relates to the concrete LoRA
parameterization implemented in practice.
The key point is that, throughout the paper, $\mathcal W_{\Delta,t}$ denotes a
\emph{linear subspace of the ambient weight-increment space} $\mathbb R^d$ with
$d=\dim(\mathrm{vec}(\Delta W))$, rather than the rank-$r$ factor coordinate space of $(A,B)$.
This distinction is essential because the LoRA mapping $\theta_\Delta=(A,B)\mapsto \Delta W=BA$
is bilinear and therefore does not define a global linear set of reachable increments in weight space.
Our PAC--Bayesian analysis is local and only requires a linear model of permissible \emph{increment directions}
around the learned solution, which is precisely what the Jacobian provides.

% \paragraph{Abstract object in the bound: a linear increment subspace.}
% In the main text, PECL is modeled as restricting task-$t$ updates to a linear subspace
% $\mathcal W_{\Delta,t}\subseteq\mathbb R^d$, leading to the affine model class
% $
% W_{\mathrm{froz}}+\mathcal W_{\Delta,t}.
% $
% All priors, posteriors, and perturbation distributions appearing in
% Theorem~\ref{thm:pathwise_pac} are defined on $(\mathcal W_{\Delta,t},\mathcal B(\mathcal W_{\Delta,t}))$
% and pushed forward to the affine class by translation.
% Concretely, letting $T_{W_{\mathrm{froz}}}(\Delta)=W_{\mathrm{froz}}+\Delta$, we work with
% \[
% P_t := (T_{W_{\mathrm{froz}}})_{\#} P_t^{\Delta},
% \qquad
% Q_t := (T_{W_{\mathrm{froz}}})_{\#} Q_t^{\Delta},
% \]
% where $P_t^\Delta$ and $Q_t^\Delta$ are measures supported on $\mathcal W_{\Delta,t}$.
% Lemma~\ref{lem:kl-decomp-pecl} shows that
% \[
% \mathrm{KL}\!\left(Q_t\Vert P_t\right)
% =
% \mathrm{KL}\!\left(Q_t^{\Delta}\Vert P_t^{\Delta}\right),
% \]
% so under PECL the KL contribution is determined entirely by the distributions supported on
% $\mathcal W_{\Delta,t}$.

% \paragraph{Implemented object: LoRA paramaters.}
For a given adapted weight block (or layer) $\ell$, LoRA introduces trainable coordinates
$\theta_{\Delta,\ell}=(A_\ell,B_\ell)$ and defines the effective increment
\[
\Delta W_\ell(A_\ell,B_\ell)= B_\ell A_\ell,
\qquad
A_\ell\in\mathbb R^{r\times n_\ell},\;
B_\ell\in\mathbb R^{m_\ell\times r}.
\]
Let $\hat\theta_{\Delta,t}=(\hat A_t,\hat B_t)$ denote the learned LoRA parameters after training task $t$.
Consider a small coordinate perturbation
$\delta\theta_\Delta=(\delta A,\delta B)$ around $\hat\theta_{\Delta,t}$, and define
\[
g(\theta_\Delta):=\mathrm{vec}(\Delta W(\theta_\Delta))\in\mathbb R^d,
\qquad d=m_\ell n_\ell .
\]
A first-order expansion yields
\begin{equation}
\label{eq:lora_local_linearization}
g(\hat\theta_{\Delta,t}+\delta\theta_\Delta)
=
g(\hat\theta_{\Delta,t})
+
J_{\mathrm{LoRA}}(\hat\theta_{\Delta,t})\,\delta\theta_\Delta
+
R(\delta\theta_\Delta),
\end{equation}
where $J_{\mathrm{LoRA}}(\hat\theta_{\Delta,t})$ is the Jacobian of $g$ evaluated at $\hat\theta_{\Delta,t}$,
and $\|R(\delta\theta_\Delta)\|_2=O(\|\delta\theta_\Delta\|_2^2)$ as $\|\delta\theta_\Delta\|_2\to 0$.

Crucially, for LoRA the Jacobian image admits an explicit characterization.
Writing $\delta\theta_\Delta=(\delta A,\delta B)$, the induced first-order increment variation in weight space is
\[
\delta(\Delta W_\ell)
=
s_\ell(\delta B\,\hat A_t+\hat B_t\,\delta A),
\]
and hence the Jacobian image at the learned coordinates is the linear set
\begin{equation}
\label{eq:WDelta_Jacobian_image_explicit}
\mathrm{Im}\!\big(J_{\mathrm{LoRA}}(\hat\theta_{\Delta,t})\big)
=
\Big\{
\mathrm{vec}\!\big(s_\ell(\delta B\,\hat A_t+\hat B_t\,\delta A)\big)
:\;
\delta A\in\mathbb R^{r\times n_\ell},\;
\delta B\in\mathbb R^{m_\ell\times r}
\Big\}
\subseteq \mathbb R^d .
\end{equation}
This is a \emph{local subspace in increment space}, not a global reachable set, because it is defined
by the first-order linearization around $\hat\theta_{\Delta,t}$.

% \paragraph{Identification used by the theory: $\mathcal W_{\Delta,t}$ as a local tangent model.}
The abstract subspace $\mathcal W_{\Delta,t}$ in Theorem~\ref{thm:pathwise_pac} is interpreted as a local linear
model of permissible increment \emph{directions} around the task-$t$ solution in the ambient increment space.
Accordingly, we instantiate it via the LoRA-induced tangent space
\begin{equation}
\label{eq:WDelta_t_identification}
\mathcal W_{\Delta,t}
\equiv
\mathrm{Im}\!\big(J_{\mathrm{LoRA}}(\hat\theta_{\Delta,t})\big),
\end{equation}
and define the PAC--Bayesian distributions directly on this linear subspace, as required by
Lemma~\ref{lem:kl-decomp-pecl} and Theorem~\ref{thm:pathwise_pac}.

% \paragraph{How coordinate perturbations induce increment perturbations.}
In implementation, perturbations are applied in the LoRA coordinate space.
Under the local linearization~\eqref{eq:lora_local_linearization}, a coordinate perturbation
$\delta\theta_\Delta=(\delta A,\delta B)$ induces an increment perturbation
\[
\delta\Delta
:=
g(\hat\theta_{\Delta,t}+\delta\theta_\Delta)-g(\hat\theta_{\Delta,t})
\approx
J_{\mathrm{LoRA}}(\hat\theta_{\Delta,t})\,\delta\theta_\Delta,
\]
whose support lies in $\mathrm{Im}(J_{\mathrm{LoRA}}(\hat\theta_{\Delta,t}))=\mathcal W_{\Delta,t}$.
For instance, if $\delta\theta_\Delta\sim\mathcal N(0,\Sigma_{\theta,t})$ and the perturbation radius is chosen
so that the remainder term is negligible with high probability, then the induced increment perturbation is approximately Gaussian in $\mathbb R^d$,
\begin{equation}
\label{eq:Sigma_pushforward}
\delta\Delta \approx \mathcal N(0,\Sigma_t),
\qquad
\Sigma_t := J_{\mathrm{LoRA}}(\hat\theta_{\Delta,t})\,\Sigma_{\theta,t}\,
J_{\mathrm{LoRA}}(\hat\theta_{\Delta,t})^\top,
\end{equation}
and is supported on $\mathcal W_{\Delta,t}$.

\paragraph{Takeaway.}
The bound is formulated on $\mathcal W_{\Delta,t}\subset\mathbb R^d$ because it is a statement about perturbations
and posterior complexity in \emph{weight-increment space} under PECL constraints.
LoRA is a concrete mechanism that realizes such a constraint.
Although the mapping $(A,B)\mapsto BA$ is bilinear and does not define a global linear class of increments,
its local linearization around the learned coordinates induces the tangent increment subspace
$\mathrm{Im}(J_{\mathrm{LoRA}}(\hat\theta_{\Delta,t}))$, which provides an implementation-level realization of
$\mathcal W_{\Delta,t}$ while keeping the KL analysis in Lemma~\ref{lem:kl-decomp-pecl} intact.




\subsection{Feature amplification and curvature projection ratio}
\label{app:weight_projection}

This appendix details two diagnostic measures used to interpret why the perturbation scope matters in LoRA-based continual learning: the feature amplification factor $\alpha_t$ and the curvature projection preference $c_t$. 
This two diagnostics are computed on a single adapted block (layer/module) in experiments rather than the full parameter vector.



\subsubsection{Feature amplification factor}
\label{app:feature_amplification}

Let $\Delta W_t$ denote the learned weight update for a given adapted block at task $t$ (implemented as $\Delta W_t = B_t A_t$).
Let
$
\Delta W_t = U_{\Delta,t} S_{\Delta,t} V_{\Delta,t}^\top
$
be its rank-$r$ truncated SVD. The pair $(U_{\Delta,t},V_{\Delta,t})$ induces the bilinear subspace
$
\mathcal S_{\Delta,t}=\{U_{\Delta,t} Z V_{\Delta,t}^\top : Z\in\mathbb{R}^{r\times r}\}.
$
Note that $\mathcal{S}_{\Delta, t}$ is an empirical diagnostic subspace defined from the learned update $\Delta W_t$ (via its truncated SVD) and is used only for curvature localization. $\mathcal{W}_{\Delta, t}=\operatorname{Im}\left(J_{\text {LoRA }}\right)$ is the local tangent space of LoRA-induced increments in $\operatorname{vec}(\Delta W)$-space, whereas $\mathcal{S}_{\Delta, t}$ captures only the core component that preserves the row and column subspaces of $\Delta W_t$; hence $\mathcal{S}_{\Delta, t}$ is an empirical diagnostic subspace of the current solution and is not assumed to coincide with $\mathcal{W}_{\Delta, t}$.
Under the Frobenius inner product, the orthogonal projection of a matrix $M$ onto $\mathcal S_{\Delta,t}$ is
% $
% \Pi_{\Delta,t}(M)=U_{\Delta,t}(U_{\Delta,t}^\top M V_{\Delta,t})V_{\Delta,t}^\top.
% $
$
\Pi_{\mathcal S,t}(M)=U_{\Delta,t}(U_{\Delta,t}^\top M V_{\Delta,t})V_{\Delta,t}^\top.
$
Equivalently, since $U_{\Delta,t}$ and $V_{\Delta,t}$ have orthonormal columns,
% $
% \|\Pi_{\Delta,t}(M)\|_F = \|U_{\Delta,t}^\top M V_{\Delta,t}\|_F.
% $
$\|\Pi_{\mathcal S,t}(M)\|_F = \|U_{\Delta,t}^\top M V_{\Delta,t}\|_F.
$



Following~\cite{hu2022lora}, we define the feature amplification factor
\[
\alpha_t=\frac{\|\Delta W_t\|_F}{\|U_{\Delta,t}^\top W V_{\Delta,t}\|_F}
=\frac{\|\Delta W_t\|_F}{\|\Pi_{\mathcal S,t}(W)\|_F},
\]
%=\frac{\|\Delta W_t\|_F}{\|\Pi_{\Delta,t}(W)\|_F},
which compares the magnitude of the learned update to the backbone energy restricted to the same bilinear subspace induced by $\Delta W_t$.
Larger $\alpha_t$ indicates that, relative to the backbone component available within $\mathcal S_{\Delta,t}$, the learned update contributes a larger effective change along the same row and column subspaces.

\subsubsection{Curvature projection ratio}
\label{app:curvature_preference}

We next quantify whether dominant task-conditioned curvature modes concentrate more on $\mathcal S_{\Delta,t}$ than on a backbone reference bilinear subspace.

% \paragraph{Curvature eigenpairs in the LoRA-augmented parameter space.}
Let $\{(\lambda_{t,i}, v_{t,i})\}_{i=1}^k$ denote the top-$k$ eigenpairs of a positive semidefinite curvature operator at task $t$ (e.g., GGN/Fisher) computed in the parameter space that includes the backbone block $W$ and the LoRA factors $(A_t,B_t)$ (with $W$ frozen during training).
% \paragraph{Mapping eigenvectors to effective weight-space directions.}
Each eigenvector $v_{t,i}$ is a direction in the LoRA-augmented parameter space.
We decompose it into the corresponding blocks and denote the induced local perturbations by
$(\delta W_{t,i}^{(0)}, \delta A_{t,i}, \delta B_{t,i})$,
where $\delta W_{t,i}^{(0)}$ is reshaped to match the backbone weight block.
Under the first-order pushforward of the LoRA parametrization, this direction induces an effective weight-space perturbation
\[
\delta W_{t,i}
=
\delta W_{t,i}^{(0)} + \big(\delta B_{t,i}A_t + B_t \delta A_{t,i}\big).
\]
% Each eigenvector $v_{t,i}$ is a direction in the full parameter space.
% We map it to an effective direction in the weight-matrix space by aggregating its components on the backbone block and the LoRA factors:
% \[
% dW_{t,i} = dW_{t,i}^{(0)} + s\big(dB_{t,i}A_t + B_t dA_{t,i}\big),
% \]
where $ \delta W_{t,i}^{(0)}$ is the component corresponding to the backbone weight block, $(\delta A_{t,i},\delta B_{t,i})$ are the components corresponding to the LoRA factors.
This mapping matches the implementation, where curvature directions are evaluated after being pushed forward to effective weight-space increments.

% \paragraph{Per-mode projection ratios.}
We define the per-mode projection ratio onto $\mathcal S_{\Delta,t}$ as
% $$
% \rho_{\Delta,t,i}
% =\frac{\|\Pi_{\Delta,t}(\delta W_{t,i})\|_F^2}{\|\delta W_{t,i}\|_F^2}.
% $$
$$
\rho_{\Delta,t,i}
=\frac{\|\Pi_{\mathcal S,t}(\delta W_{t,i})\|_F^2}{\|\delta W_{t,i}\|_F^2}.
$$
For the backbone reference, let $(U_{W,t},V_{W,t})$ be the rank-$r$ truncated SVD subspace of the corresponding backbone weight block (restricted to the same layer/block), inducing
$
\mathcal S_{W,t}=\{U_{W,t} Z V_{W,t}^\top : Z\in\mathbb{R}^{r\times r}\},
\quad
\Pi_{\mathcal S_W,t}(M)=U_{W,t}(U_{W,t}^\top M V_{W,t})V_{W,t}^\top,
$
%\Pi_{W,t}(M)=U_{W,t}(U_{W,t}^\top M V_{W,t})V_{W,t}^\top,
and define
% $$
% \rho_{W,t,i}
% =\frac{\|\Pi_{W,t}(\delta W_{t,i})\|_F^2}{\|\delta W_{t,i}\|_F^2}.
% $$
$$\rho_{W,t,i}
=\frac{\|\Pi_{\mathcal S_W,t}(\delta W_{t,i})\|_F^2}{\|\delta W_{t,i}\|_F^2}.
$$

% \paragraph{Eigenvalue-weighted aggregation and preference ratio.}
We aggregate the localization of the top-$k$ curvature modes by eigenvalue-weighted averages
\[
R_{\Delta,t}=\frac{\sum_{i=1}^k \lambda_{t,i}\rho_{\Delta,t,i}}{\sum_{i=1}^k \lambda_{t,i}},
\qquad
R_{W,t}=\frac{\sum_{i=1}^k \lambda_{t,i}\rho_{W,t,i}}{\sum_{i=1}^k \lambda_{t,i}},
\]
and report the curvature projection ratio
\[
c_t=\frac{R_{\Delta,t}}{R_{W,t}}.
\]
Larger $c_t$ indicates that dominant curvature modes have higher projected Frobenius energy within the bilinear subspace induced by $\Delta W_t$ than within the backbone reference bilinear subspace.








\subsection{Implementation Details}

The code is https://anonymous.4open.science/r/FlatnessICMLNoname9527223/.
% Unless stated otherwise, we use SGD as the base optimizer, train each session for $20$ epochs and SAM radius $\rho = 0.05$. LoRA based variants are trained with learning rate $0.01$ without momentum or weight decay. The same training configuration is used across SeqLoRA, IncLoRA, and OLoRA, so that differences in performance can be attributed to the geometry of the adapter subspace and the choice of flatness optimizer rather than to changes in training budget.

% \subsection{Implementation Details}
\noindent\textbf{Data.}
We evaluate continual learning on ImageNet-R. For robustness under corruptions and perturbations, we use the Tiny ImageNet variants of ImageNet-C and ImageNet-P (200 classes) and construct ImageFolder-style splits that prevent leakage by grouping all corruption/perturbation variants from the same original image (or video) into a single train/test partition. All robustness evaluations use the same $224$-resolution preprocessing pipeline as the main experiments.

\noindent\textbf{Experimental setup.}
Unless stated otherwise, we use SGD as the base optimizer and train each task/session for $20$ epochs with batch size $128$ and a cosine learning-rate schedule. For sharpness-aware optimization, we set the SAM radius to $\rho=0.05$. LoRA-based variants are trained with learning rate $0.01$ and no momentum or weight decay. The same training configuration is used across SeqLoRA, IncLoRA, and OLoRA to isolate the effect of adapter geometry and perturbation design.

\noindent\textbf{LoRA configuration.}
Unless specified otherwise, the LoRA rank is set to $r=16$ for the main ImageNet-R experiments. For the Tiny ImageNet-C/P robustness experiments, we use $r=8$. In rank ablations, we sweep $r\in\{2,8,16\}$.

\noindent\textbf{Evaluation protocol.}
For accuracy reporting, we aggregate results over multiple random seeds. For sharpness and curvature measurements, we fix the random seed to $1993$ to ensure controlled comparisons and reduce estimator variance. Flatness/curvature evaluation is performed on a $10\%$ subset of the evaluation set with batch size $32$. Random sharpness uses $100$ Gaussian samples. Hessian trace and top-eigenvalue estimates are computed with $100$ Lanczos power iterations and $100$ trace samples. Loss-landscape visualizations use radius $1$ with $41$ points along a random basis with norm-filtered directions.

\noindent\textbf{Weight and curvature analysis.}
In the weight/curvature analysis, we include frozen parameters when computing full-parameter Hessian-vector products and use Lanczos with top-$k=16$ eigenvalues. Curvature localization and $\Delta W$ projections are computed on the last transformer block (block 11) and restricted to the $q$ and $v$ subspaces with row ranges $[0,768]$ and $[1536,2304]$, respectively. The localized Rayleigh estimator uses $128$ samples, and the curvature MVP backend is set to GGN. For $\Delta W$ and $\Delta W_{\mathrm{full}}$ projections, we use the QKV weight and its corresponding LoRA factors with rank $16$. The $W$--$\Delta W$ alignment analysis uses rank $16$, seed $42$,, and the seqlora setting.


\subsection{Additional Experimental Results}


\subsubsection{Raw Task-Wise $\operatorname{ACC}^{\mathrm{cls}}_t$ on ImageNet-R and ImageNet-C for Fig.~\ref{fig:robustness_imagenet_cp}}

To complement Fig.~\ref{fig:robustness_imagenet_cp}, we provide the underlying raw task-wise $\operatorname{ACC}^{\mathrm{cls}}_t$ trajectories for ImageNet-C and ImageNet-P in the tables below.
All entries are {means over repeated runs} (averaged across the same set of random seeds used for the main figure).



\begin{table*}[ht]
\centering
\caption{Raw task-wise $\operatorname{ACC}^{\mathrm{cls}}_t$ trajectories ($t=1,\dots,20$) on ImageNet-C under different optimizers (mean over repeated runs).}
\label{tab:taskwise_acc_imagenetr}
\setlength{\tabcolsep}{2.2pt}
\renewcommand{\arraystretch}{0.92}
\scriptsize
\begin{tabular}{@{}l l *{20}{r}@{}}
\toprule
\textbf{Method} & \textbf{Opt.} &
\textbf{t01} & \textbf{t02} & \textbf{t03} & \textbf{t04} & \textbf{t05} &
\textbf{t06} & \textbf{t07} & \textbf{t08} & \textbf{t09} & \textbf{t10} &
\textbf{t11} & \textbf{t12} & \textbf{t13} & \textbf{t14} & \textbf{t15} &
\textbf{t16} & \textbf{t17} & \textbf{t18} & \textbf{t19} & \textbf{t20} \\
\midrule

\multirow{3}{*}{\textbf{SeqLoRA}}
& SGD
& 92.67 & 87.84 & 87.67 & 83.42 & 82.13 & 81.50 & 79.62 & 78.83 & 78.11 & 76.47
& 75.85 & 75.28 & 74.69 & 74.24 & 74.27 & 73.48 & 72.69 & 71.72 & 70.49 & 69.77 \\
& \quad $+\mathrm{AS}^{(0)}_S$
& 92.67 & 87.67 & 87.33 & 84.42 & 83.20 & 82.39 & 80.33 & 79.29 & 79.11 & 78.20
& 77.73 & 77.00 & 76.69 & 77.00 & 77.33 & 77.08 & 76.45 & 75.94 & 75.33 & 73.93 \\
& \quad $+\mathrm{AS}^{(1)}_S$
& 95.67 & 93.67 & 89.78 & 88.84 & 87.67 & 86.61 & 86.67 & 86.33 & 84.85 & 84.14
& 82.61 & 82.17 & 81.87 & 81.60 & 81.42 & 80.44 & 80.10 & 79.96 & 80.21 & 80.03 \\

\addlinespace[2pt]

\multirow{3}{*}{\textbf{IncLoRA}}
& SGD
& 92.33 & 89.84 & 88.22 & 86.58 & 83.40 & 81.22 & 80.38 & 78.46 & 75.89 & 73.77
& 73.21 & 73.61 & 73.02 & 72.95 & 71.62 & 70.17 & 68.82 & 66.30 & 64.86 & 61.77 \\
& \quad $+\mathrm{AS}^{(0)}_S$
& 92.33 & 90.83 & 90.45 & 85.92 & 87.00 & 83.39 & 81.81 & 79.75 & 78.22 & 77.87
& 76.79 & 73.47 & 72.03 & 73.43 & 71.07 & 71.44 & 70.78 & 69.63 & 67.96 & 67.98 \\
& \quad $+\mathrm{AS}^{(1)}_S$
& 96.00 & 92.17 & 88.33 & 86.92 & 85.87 & 85.11 & 86.33 & 85.17 & 84.78 & 84.03
& 82.85 & 82.56 & 81.54 & 81.52 & 80.64 & 79.73 & 79.59 & 80.32 & 79.58 & 79.47 \\

\addlinespace[2pt]

\multirow{3}{*}{\textbf{OLoRA}}
& SGD
& 92.33 & 88.50 & 87.78 & 85.59 & 83.00 & 81.28 & 80.57 & 78.58 & 75.67 & 74.10
& 73.55 & 73.36 & 73.46 & 73.33 & 71.87 & 71.40 & 70.57 & 68.33 & 67.44 & 65.12 \\
& \quad $+\mathrm{AS}^{(0)}_S$
& 92.33 & 88.67 & 90.22 & 85.67 & 86.33 & 83.00 & 81.48 & 79.83 & 78.26 & 78.33
& 77.00 & 73.97 & 71.77 & 72.50 & 70.51 & 70.36 & 69.33 & 68.54 & 66.21 & 66.98 \\
& \quad $+\mathrm{AS}^{(1)}_S$
& 96.00 & 91.17 & 88.00 & 86.42 & 85.47 & 85.83 & 86.33 & 84.96 & 84.19 & 83.07
& 82.30 & 81.53 & 81.46 & 80.69 & 80.27 & 78.75 & 79.14 & 78.98 & 79.28 & 79.35 \\

\bottomrule
\end{tabular}
\end{table*}


\begin{table*}[ht]
\centering
\caption{Raw task-wise $\operatorname{ACC}^{\mathrm{cls}}_t$ trajectories ($t=1,\dots,20$) on ImageNet-P under different optimizers (mean over repeated runs).}
\label{tab:taskwise_acc_new}
\setlength{\tabcolsep}{2.2pt}
\renewcommand{\arraystretch}{0.92}
\scriptsize
\begin{tabular}{@{}l l *{20}{r}@{}}
\toprule
\textbf{Method} & \textbf{Opt.} &
\textbf{t01} & \textbf{t02} & \textbf{t03} & \textbf{t04} & \textbf{t05} &
\textbf{t06} & \textbf{t07} & \textbf{t08} & \textbf{t09} & \textbf{t10} &
\textbf{t11} & \textbf{t12} & \textbf{t13} & \textbf{t14} & \textbf{t15} &
\textbf{t16} & \textbf{t17} & \textbf{t18} & \textbf{t19} & \textbf{t20} \\
\midrule

\multirow{3}{*}{\textbf{SeqLoRA}}
& SGD
& 96.17 & 93.42 & 92.00 & 89.88 & 88.17 & 87.86 & 86.43 & 84.92 & 85.57 & 84.63
& 83.88 & 83.71 & 82.40 & 82.02 & 83.08 & 81.95 & 80.09 & 80.17 & 78.01 & 72.76 \\
& \quad $+\mathrm{AS}^{(0)}_S$
& 95.83 & 92.67 & 92.28 & 89.54 & 88.47 & 88.58 & 87.90 & 87.29 & 87.63 & 86.82
& 85.80 & 85.60 & 84.57 & 84.21 & 84.51 & 83.58 & 83.30 & 83.70 & 82.32 & 80.95 \\
& \quad $+\mathrm{AS}^{(1)}_S$
& 99.17 & 94.92 & 95.00 & 92.71 & 91.60 & 91.92 & 89.55 & 88.69 & 89.29 & 89.10
& 87.70 & 87.71 & 87.68 & 87.62 & 88.19 & 87.51 & 86.89 & 87.02 & 86.93 & 86.08 \\

\addlinespace[2pt]

\multirow{3}{*}{\textbf{IncLoRA}}
& SGD
& 95.67 & 92.50 & 91.78 & 89.71 & 88.80 & 88.08 & 85.69 & 85.29 & 84.59 & 83.25
& 80.94 & 78.58 & 77.72 & 77.79 & 74.13 & 73.24 & 71.70 & 69.29 & 68.85 & 66.92 \\
& \quad $+\mathrm{AS}^{(0)}_S$
& 95.83 & 91.58 & 92.17 & 89.62 & 89.27 & 88.28 & 86.57 & 85.31 & 85.56 & 85.33
& 83.95 & 81.81 & 80.01 & 78.26 & 78.75 & 77.08 & 77.34 & 75.79 & 76.53 & 73.75 \\
& \quad $+\mathrm{AS}^{(1)}_S$
& 98.83 & 96.41 & 95.39 & 93.33 & 91.63 & 90.97 & 89.81 & 88.83 & 88.78 & 88.63
& 87.46 & 86.27 & 86.58 & 86.15 & 86.63 & 86.42 & 86.45 & 85.54 & 85.34 & 85.09 \\

\addlinespace[2pt]

\multirow{3}{*}{\textbf{OLoRA}}
& SGD
& 95.67 & 90.00 & 90.39 & 88.33 & 88.20 & 87.61 & 85.76 & 85.02 & 84.45 & 83.37
& 79.97 & 77.76 & 76.68 & 76.64 & 72.82 & 72.54 & 70.40 & 68.16 & 67.61 & 65.55 \\
& \quad $+\mathrm{AS}^{(0)}_S$
& 95.83 & 90.17 & 91.06 & 88.46 & 88.27 & 87.86 & 86.74 & 84.96 & 85.37 & 84.95
& 83.74 & 82.61 & 81.01 & 79.58 & 78.79 & 77.21 & 76.96 & 75.30 & 75.73 & 73.32 \\
& \quad $+\mathrm{AS}^{(1)}_S$
& 98.83 & 95.33 & 94.39 & 91.46 & 90.17 & 90.14 & 88.60 & 87.48 & 87.02 & 87.23
& 86.23 & 85.47 & 85.99 & 85.79 & 86.19 & 86.23 & 86.01 & 85.29 & 84.61 & 84.47 \\
\bottomrule
\end{tabular}
\end{table*}





\subsubsection{Computational Cost Analysis}


Table~\ref{tab:lora_opt_metric} reports the average per-task training time on ImageNet-R.
Across LoRA variants, adding sharpness regularization increases the computational cost relative to SGD, with $AS^{(1)}_S$ being consistently the most expensive due to the additional gradient-norm maximization, while $AS^{(0)}_S$ shows a moderate overhead.
In contrast, $RS^{(0)}_S$ incurs a smaller and more variable overhead, and can be comparable to SGD in some settings.


\begin{table}[hpb]
\centering
\caption{Average per-task training time (seconds) across methods and optimizers on ImageNet-R.}
\label{tab:lora_opt_metric}
\footnotesize
\setlength{\tabcolsep}{3pt}
\renewcommand{\arraystretch}{0.95}
\begin{tabular}{l cccc}
\toprule
\textbf{Method} & \textbf{SGD} & $+\mathrm{RS}^{(0)}_S$ & $+\mathrm{AS}^{(0)}_S$ & $+\mathrm{AS}^{(1)}_S$ \\
\midrule
SeqLoRA & $282.72 \pm 61.34$  & $244.95 \pm 6.38$  & $472.29 \pm 83.70$  & $670.20 \pm 217.92$ \\
IncLoRA & $310.54 \pm 141.00$ & $658.34 \pm 83.90$ & $777.91 \pm 239.29$ & $1247.56 \pm 5.05$ \\
OLoRA   & $412.99 \pm 146.84$ & $644.19 \pm 55.05$ & $628.44 \pm 181.29$ & $836.99 \pm 327.65$ \\
\bottomrule
\end{tabular}
\end{table}

% \begin{table}[hpb]
% \centering
% \caption{Average per‑task training time (seconds) across methods and optimizers on ImageNet‑R}
% \label{tab:lora_opt_metric}
% \setlength{\tabcolsep}{4pt}
% \renewcommand{\arraystretch}{1}
% \resizebox{\linewidth}{!}{
% \begin{tabular}{l ccccc}
% \toprule
% \textbf{Method} & \textbf{SGD} & $+RS^{(0)}_S$ & $+AS^{(0)}_S$ & $+AS^{(1)}_S$ \\
% \midrule
% SeqLoRA & 282.72 ± 61.34  & 244.95 ± 6.38  & 472.29 ± 83.70  & 670.20 ± 217.92    \\
% IncLoRA & 310.54 ± 141.00 & 658.34 ± 83.90 & 777.91 ± 239.29 & 1247.56 ± 5.05  3 \\
% OLoRA   & 412.99 ± 146.84 & 644.19 ± 55.05 & 628.44 ± 181.29 & 836.99 ± 327.65  \\
% \bottomrule
% \end{tabular}
% }
% \end{table}





% \section{Technical appendices and supplementary material}
% Technical appendices with additional results, figures, graphs, and proofs may be submitted with the paper submission before the full submission deadline (see above). You can upload a ZIP file for videos or code, but do not upload a separate PDF file for the appendix. There is no page limit for the technical appendices. 

% Note: Think of the appendix as ``optional reading'' for reviewers. The paper must be able to stand alone without the appendix; for example, adding critical experiments that support the main claims to an appendix is inappropriate. 

%%%%%%%%%%%%%%%%%%%%%%%%%%%%%%%%%%%%%%%%%%%%%%%%%%%%%%%%%%%%

\newpage
\input{checklist.tex}  % uncomment when checklist is ready


\end{document}
